% concatenated from 1021.tex

\documentclass[12pt,reqno]{amsart}
\usepackage{amsmath}
\usepackage{amssymb}
\usepackage{amstext}
\usepackage{a4wide}
\usepackage{graphicx}
\allowdisplaybreaks \numberwithin{equation}{section}
\usepackage{color}
\usepackage{cases}
%\usepackage{refcheck}
\usepackage{hyperref}
\def\alert#1{\textcolor[rgb]{1,0,0}{#1}}
\def\red#1{\textcolor[rgb]{1,0,0}{#1}}
\def\org#1{\textcolor[rgb]{0,1,0}{#1}}
\def\green#1{\textcolor[rgb]{0,1,0}{#1}}

\numberwithin{equation}{section}

\newtheorem{theorem}{Theorem}[section]
\newtheorem{proposition}[theorem]{Proposition}
\newtheorem{corollary}[theorem]{Corollary}
\newtheorem{lemma}[theorem]{Lemma}

\theoremstyle{definition}


\newtheorem{innercustomthm}{Proposition}
\newenvironment{customthm}[1]
{\renewcommand\theinnercustomthm{#1}\innercustomthm}
{\endinnercustomthm}




\newtheorem{assumption}[theorem]{Assumption}
\newtheorem{definition}[theorem]{Definition}

\theoremstyle{remark}
\newtheorem{remark}[theorem]{Remark}
\newtheorem{remarks}[theorem]{Remarks}
\newtheorem{example}[theorem]{Example}
\newtheorem{examples}[theorem]{Examples}
\newtheorem*{lemmaA}{Lemma A}
\newtheorem*{lemmaB}{Lemma B}

\newcommand{\ep}{\varepsilon}
\newcommand{\Om}{\Omega}
\newcommand{\la}{\lambda}
\newcommand{\ds}{\displaystyle}


\begin{document}
	
	
	
	
	
\title[Nearly parallel helical vortices]
{Co-rotating nearly parallel helical vortices with small cross-section in 3D incompressible Euler equations}
	
	
	
\author{Daomin Cao, Jie Wan}
	
\address{State Key Laboratory of Mathematical Sciences, Academy of Mathematics and Systems Science, Chinese Academy of Sciences, Beijing 100190, and University of Chinese Academy of Sciences, Beijing 100049,  P.R. China}
\email{dmcao@amt.ac.cn}
\address{School of Mathematics and Statistics, Beijing Institute of Technology, Beijing 100081,  P.R. China}
\email{wanjie@bit.edu.cn}

%\footnotetext[0]{Corresponding author}
	
	
	
%\thanks{*  Corresponding author.}
	
	
\begin{abstract}
In this article, we consider clustered solutions to a semilinear elliptic equation in divergence form
\begin{equation*}
\begin{cases}
-\varepsilon^2\text{div}(K(x)\nabla u)=  (u-q|\ln\varepsilon|)^{p}_+,\ \ &x\in \Omega,\\
u=0,\ \ &x\in\partial \Omega
\end{cases}
\end{equation*}
%Using the finite-dimensional reduction method,
for small values of $ \varepsilon $. Using Green’s function of the elliptic operator  $ -\text{div}(K(x)\nabla) $ and finite-dimensional reduction method, we prove that there exist  clustered solutions with cluster point $ 0 $ and cluster distance $ |\ln\varepsilon| ^{-\frac{1}{2}} $ whose small-structure is governed by some functional $ H_N $ determined by $ K $ and $ q $. As an application, we prove the existence of traveling-rotating helical vorticity fields to 3D incompressible Euler equations in infinite cylinders, whose support sets consist  of  helical tubes with small cross-section of radius $ \varepsilon $ and arbitrary circulation $ \kappa $ and concentrates near ``$ 2N $'' and ``$ 2N+1 $'' type of co-rotating helical solutions of nearly parallel vortex filaments model as $ \varepsilon\to0 $, which justifies the result in Klein, Majda and Damodaran [1995, JFM] and generalizes results in Guerra and Musso [arxiv: 2502.01470]. Several kinds of solutions such as ``2 asymmetric'', ``$ 2\times2 $ asymmetric'' and ``$ 2\times2+1 $ asymmetric'' type of co-rotating helical filaments are also considered. 


%In this article, we prove the existence of clustered traveling-rotating helical vortex solutions to 3D Euler equations, whose   support set  consists of $ arbitrary $ many  helical tubes with small cross-section of radius $ \varepsilon $ and concentrates near a single helix as $ \varepsilon\to0 $. %we first consider  solutions to % helical vortices with small cross-section to the 3D incompressible Euler equations in infinite pipes. By considering
%The key is to prove the existence of clustered solutions to a semilinear elliptic problem in divergence form
%\begin{equation*}
%\begin{cases}
%-\varepsilon^2\text{div}(K(x)\nabla u)=  (u-q|\ln\varepsilon|)^{p}_+,\ \ &x\in \Omega,\\
%u=0,\ \ &x\in\partial \Omega
%\end{cases}
%\end{equation*}
%%Using the finite-dimensional reduction method,
%for small values of $ \varepsilon $. By using the variational method, we prove that  there exists a family of clustered solutions which have $ arbitrary $ many bubbles and collapse  into a certain critical point  of  $ q^2\sqrt{\det K} $ as $ \varepsilon\to0. $ %Then as an application, 
\vspace{0.3cm}

\textbf{Keywords:} Incompressible Euler equation;  Helical symmetry; Semilinear elliptic equations; Clustered solution; Lyapunov-Schmidt reduction method.

\vspace{0.3cm}
\textbf{MSC:} Primary 76B47\,\,\,\,\,\,\,\,Secondary 76B03, 35A02, 35Q31
\end{abstract}
	
	
\maketitle

\section{Introduction and main results}
The motion of 3D incompressible ideal flow is governed by   Euler
equations
\begin{equation}\label{Euler eq}
\begin{cases}
\partial_t\mathbf{v}+(\mathbf{v}\cdot \nabla)\mathbf{v}=-\nabla P,\ \ &D\times (0,T),\\
\nabla\cdot \mathbf{v}=0,\ \ &D\times (0,T),\\
\mathbf{v}\cdot \mathbf{n}=0,\ \ &\partial D\times (0,T),
\end{cases}
\end{equation}
where $ D\subseteq \mathbb{R}^3 $, $ \mathbf{v}=(v_1,v_2,v_3) $ is the velocity field, $ P$ is the scalar pressure, $ \mathbf{n} $ is the outward unit normal to $ \partial D $. % The third equation means that the net flux of velocity across the boundary is zero.

For velocity field $\mathbf{v}$, we define the  vorticity field $ \mathbf{w}=\nabla\times \mathbf{v} $.  Then $ \mathbf{w} $ satisfies vorticity equations  
\begin{equation}\label{Euler eq2}
\begin{cases}
\partial_t\mathbf{w}+(\mathbf{v}\cdot \nabla)\mathbf{w}=(\mathbf{w}\cdot \nabla)\mathbf{v},\\
\nabla\times \mathbf{v}=\mathbf{w}.
\end{cases}
\end{equation}

In this paper, we consider the existence of vorticity field of 3D incompressible Euler equations \eqref{Euler eq} concentrating in a small neighborhood of  nearly parallel vortex filaments. The study of the existence   of vorticity field of Euler equations concentrating near some curves  has a long history, traced back to the classical work of Helmholtz \cite{He}. % and Kelvin \cite{KH}.
 Consider the motion of a vortex tube with a small cross-section of radius $ \varepsilon $ and a fixed circulation $ \kappa $, uniformly distributed around one curve $ \gamma(t) $. Da Rios  and Levi-Civita  \cite{DR, LC} formally found that $ \gamma $ asymptotically obeys a law of the form 
\begin{equation}\label{1003}
\partial_t \gamma=\frac{\kappa}{4\pi}|\ln\varepsilon|(\partial_s\gamma\times\partial_{ss}\gamma)=\frac{\kappa\bar{K}}{4\pi}|\ln\varepsilon|\mathbf{b}_{\gamma(t)},
\end{equation}
%Basically, it concerns the motion of vorticity field of \eqref{Euler eq} which cincentrates 
where $ s $ is the parameter of arclength, $ \bar{K} $ is the local curvature and $ \mathbf{b}_{\gamma(t)}  $  is the unit binormal vector. $ \gamma $ and \eqref{1003} are  called the vortex filament and the vortex filament equation respectively.

From mathematical justification, Jerrard and Seis \cite{JS} first %gave a precise form to Da Rios’ computation under some mild conditions on a solution to \eqref{Euler eq2} which remains suitably concentrated around an evolving vortex filament. Their result 
 showed that under some assumptions on   solutions $ \mathbf{w}_\varepsilon $ of \eqref{Euler eq2}, there holds in the sense of distribution 
\begin{equation}\label{1005}
\mathbf{w}_\varepsilon(\cdot,t)-\kappa \delta_{\gamma(t)}\mathbf{t}_{\gamma(t)}\rightharpoonup0,\ \ \text{as}\ \varepsilon\to0,
\end{equation}
where $ \gamma(t) $ satisfies \eqref{1003}, $ \mathbf{t}_{\gamma } $ is the tangent unit vector of $ \gamma $ and $ \delta_{\gamma } $ is the uniform Dirac measure on the curve $ \gamma $. For  any curve $ \gamma $ satisfying \eqref{1003}, whether there are vorticity fields of \eqref{Euler eq} concentrating near $ \gamma $ in sense of \eqref{1005} is   well-known as the vortex filament conjecture, which is unsolved except for   several kinds of special curves: the straight lines, the traveling circles and the traveling-rotating helices, see, e.g., \cite{CLW,CW,CW1,DDMW,DDMW2,DLM,Fr1,FB,GZ,MB} and reference therein.

As for the motion of multiple filaments, Klein, Majda and Damodaran \cite{KMD} have formally derived a simplified asymptotic law
of the evolution of several  nearly parallel vortex filaments. Consider the motion of $ N $ slender vortex tube which are all nearly parallel to the $ x_3 $-axis, whose centerlines are curves parameterized as $ \tilde{X}_{j,\varepsilon}(s,t) $ in the complex plane $ \mathbb{C} $, circulations are $ \kappa_j $ and cross-sections are of radius $ \varepsilon $ for $ j=1,\cdots,N $. The wave-length of each vortex tube is of the order $ O(1) $ and the mutual distance of any two vortex tubes is of the order $ O\left (|\ln\varepsilon|^{-\frac{1}{2}}\right ) $, i.e., $ dist(\tilde{X}_{i,\varepsilon}, \tilde{X}_{j,\varepsilon}) =O\left (|\ln\varepsilon|^{-\frac{1}{2}}\right )$ for $ i\neq j $. Denote $   \tau=|\ln\varepsilon|t $ and $ X_j(s, \tau)=\sqrt{|\ln\varepsilon|}\tilde{X}_{j,\varepsilon}(s, t) $. Then $ X_j(s, \tau) $ asymptotically obeys  (see (2.20) in \cite{KMD})
\begin{equation}\label{NP vor fi} 
\partial_\tau X_j=\frac{1}{4\pi}\left(\mathbf{i}\alpha_j\kappa_j\partial_{ss}X_j+2\mathbf{i}\sum_{k\neq j}\kappa_k\frac{X_j-X_k}{|X_j-X_k|^2} \right),
\end{equation}
where $ \mathbf{i} $ is  the imaginary unit and $ \alpha_j $ are structure constants of the vortex core, which are generally assumed to be 1. 


%It governs the motion of N ≥ 2 filaments and takes into consideration the pairwise interaction between the filaments along with an approximation for motion by self-induction. The model is


%System \eqref{NP vor fi} has been studied for its own (see e.g. \cite{BM,BFM,CGY,KPV,KVG,LM}), in particular its wellposedness and the possibility of colliding filaments. Nevertheless, the justification of the model itself as a limit from a classical fluid mechanics model such as the Euler equation has so far only been obtained formally through matched asymptotic. 
System \eqref{NP vor fi} has been extensively studied  (see, for example, \cite{BM,BFM,CGY,KPV,KVG,LM}), particularly concerning its well-posedness and the potential for filament collisions. A natural question arises as to whether a rigorous justification can be provided to establish the model itself as a limit from  the Euler equations, thereby  giving a  mathematical explanation of the model proposed by \cite{KMD}.

Recently, some progress has been made  for this problem when  filaments $ X_j(s, \tau) $ are some co-rotating helices. In \cite{GM2}, Guerra and Musso considered two special type of  helices $ X_j(s, \tau) $ satisfying \eqref{NP vor fi}: ``$ N $  vortices of equal strength $ 8\pi $" placed at each vertex of a regular $ N $-polygon, and ``$ N + 1 $ vortices   of equal strength $ 8\pi $" $ N $ of them placed at the corners of a regular $ N $-polygon with the $ N + 1 $ vortex   located at the center of vorticity. By considering the existence of solutions to a 2D semilinear elliptic equations of divergence type
\begin{equation*} 
\begin{split}
-\text{div}\left(K_H(x)\nabla v\right)=f_\varepsilon(v) ,\ \ \ \ \text{in}\ \ \mathbb{R}^2,
\end{split}
\end{equation*}
where $ K_H $ is some positive definite matrix (defined by \eqref{coef matrix}) and the profile function $ f_\varepsilon $ is some exponential type function which is chosen properly, the authors constructed vorticity fields of \eqref{Euler eq} with the orthogonality condition whose vortex sets tend to rescaled versions of the   the above two kinds of co-rotating helices as $ \varepsilon\to0. $ Note that by the choice of $ f_\varepsilon $, the associated circulation of each component of vorticity field must be $ 8\pi $ and  support sets of vorticity field are the whole space, rather than a collection of helical vortex tubes with small cross-section. So as $ \varepsilon\to0 $, the vorticity fields tend to filaments of \eqref{NP vor fi} in distributional sense.

Our goal of this paper is to establish the existence of a family of solutions of \eqref{Euler eq}   with vorticities concentrating near co-rotating helical filaments of  \eqref{NP vor fi} with a more general setting. First, as for the two configurations considered in \cite{GM2}, we prove the existence of a family of solutions with vorticities $ \textbf{w}_\varepsilon $ concentrating near helical filaments as $ \varepsilon\to0 $, whose cross-sections are   compact sets with radius $ \varepsilon $ and circulations are arbitrary rather than $8\pi$. So the concentration is in both distributional and topological sense, which coincides with the model deduced in \cite{KMD}. Second, we show the existence of a family of solutions with vorticities $ \textbf{w}_\varepsilon $ concentrating near 2 asymmetric co-rotating helical filaments of \eqref{NP vor fi} with small cross-section as $ \varepsilon\to0 $. Both circulations and  distances between helical tubes and $ x_3 $-axis are different. Some other asymmetric co-rotating helical filaments of \eqref{NP vor fi} are also considered. Following the ansatz in \cite{CW,CW1}, we assume that the fluid domain is an infinitely long cylinder, that is, $ D=B_{R^*}(0)\times\mathbb{R} $ for some $ R^*>0. $


%$ N + 1 $ vortices N of them of equal strength set at the corners of a regular N-polygon with the $ N + 1 $ vortex of arbitrary strength located at the center of vorticity.

We define co-rotating nearly parallel helical filaments as follows. Let $ N\geq 2 $ be an integer. We say $ X_j(s, \tau) $, for $ j=1,\cdots,N $,  co-rotating helical filaments of model \eqref{NP vor fi} with pitch $ h>0 $ and angular velocity $ \alpha\in\mathbb{R} $, if $ X_j(s, \tau)=Z_j(\tau)e^{\mathbf{i}\frac{s}{h}} $ for some $ Z_j(\tau) $ and $   Z_j(\tau)=  Z_j(0)e^{-\mathbf{i}\alpha\tau} $ for $ j=1,\cdots,N $. %Consider $ N $ helical solutions of \eqref{NP vor fi} whose pitch is   $ h>0, $ i.e., $ X_j(s, \tau)=Z_j(\tau)e^{\mathbf{i}\frac{s}{h}} $. Then 
Note that from \eqref{NP vor fi}, $ Z_j $ satisfies for $ j=1,\cdots,N $
\begin{equation}\label{NP heli vor}
\partial_\tau Z_j=\frac{1}{4\pi}\left(-\frac{\mathbf{i}\kappa_j}{h^2}Z_j+2\mathbf{i}\sum_{k\neq j}\kappa_k\frac{Z_j-Z_k}{|Z_j-Z_k|^2} \right).
\end{equation} 
%The co-rotating solution $   Z_j $ of \eqref{NP heli vor} means that there exists $ \alpha\in \mathbb{R} $ such that $   Z_j(\tau)=  Z_j(0)e^{-\mathbf{i}\alpha\tau} $ for $ j=1,\cdots,N. $

Now we give some exact co-rotating helical solutions of \eqref{NP vor fi} and state our main results.

\textbf{Case 1:} ``$ N $  vortices of equal strength $ \kappa $", which was considered in \cite{GM2}. Let $ N\geq 2 $, $ \kappa_j=\kappa $ for   $ \kappa>0 $. Then 
\begin{equation*}
X_{j}(s,\tau)=Z_j(\tau)e^{\mathbf{i}\frac{s}{h}}=r_*e^{-\mathbf{i}\alpha\tau}e^{\mathbf{i}\frac{s}{h}}e^{\mathbf{i}\frac{2\pi (j-1)}{N}},\ \ \ \ \ \ \ \ \  j=1,\cdots,N, 
\end{equation*}
satisfies \eqref{NP vor fi}, where $ \alpha=\frac{\kappa}{4\pi}\left( \frac{1}{h^2}-\frac{N-1}{r_*^2}\right)  $, see also    Appendix \ref{lemA1}. This corresponds to   nearly parallel vortex filaments being of the form
\begin{equation}\label{fila 1}
\tilde{X}_{j,\varepsilon}(s, t)=\frac{r_*}{\sqrt{|\ln\varepsilon|}}e^{-\mathbf{i}\alpha|\ln\varepsilon|t}e^{\mathbf{i}\frac{s}{h}}e^{\mathbf{i}\frac{2\pi (j-1)}{N}},\ \ \ \ j=1,\cdots,N.
\end{equation}
We have 
\begin{theorem}\label{thm01}
	Let $ R^*>0 $, $ h>0 $ and $ N\geq 2 $ be an integer. For $ r_*\in (0,R^*) , \kappa>0 $, let $ \tilde{X}_{j,\varepsilon}$   $(j=1,\cdots,N)$ be $ N $ filaments satisfying \eqref{fila 1}. Then  there exists $ \varepsilon_0>0 $, such that for every $ \varepsilon\in(0,\varepsilon_0] $, \eqref{Euler eq2} has a  family of solutions with vorticities  $ \mathbf{w}_\varepsilon\in C^1(B_{R^*}(0)\times \mathbb{R}\times\mathbb{R}^+) $ satisfying in distributional sense
	\begin{equation*} 
	\mathbf{w}_\varepsilon(x,t)- \kappa\sum_{j=1}^N\delta_{\tilde{X}_{j,\varepsilon}}\textbf{t}_{\tilde{X}_{j,\varepsilon}}\rightharpoonup0\ \ \ \ \text{as}\ \varepsilon\to0.
	\end{equation*}
	Moreover, if we set $ A_{\varepsilon,i}=supp\left (\mathbf{w}_\varepsilon(x_1,x_2,0,0)\right ) \cap B_{\rho_0|\ln\varepsilon|^{-\frac{1}{2}}}\left ( \tilde{X}_{i,\varepsilon}(0,0)\right )  $, where $ \rho_0>0 $ is a small constant, then there exist constants $ R_1,R_2>0 $ independent of $ \varepsilon $  such that
	\begin{equation*}
	R_1\varepsilon \leq diam(A_{\varepsilon,i})\leq R_2\varepsilon .
	\end{equation*}
\end{theorem}

\begin{remark}
In \cite{GM2}, the authors proved the existence of $ \textbf{w}_\varepsilon $ of \eqref{Euler eq} concentrating near $ \eqref{fila 1} $ in sense of distribution, when $ \kappa=8\pi $. In contrast to \cite{GM2}, the circulation of vorticity field in our construction is allowed to be any positive constant, and the support set of each component of $ \textbf{w}_\varepsilon $ is a slender helical tube with cross-section of radius $ \varepsilon $. Our result coincides with formal computational results in \cite{KMD}.
\end{remark}


\textbf{Case 2:} ``$ N+1 $ vortices", $ N $ of them having identical circulation and placed  at each vertex of a regular $ N $-polygon   and the $ (N +1) $th vortex of arbitrary strength
at the center of vorticity. Let  $ N\geq 2 $, $ \kappa_0=\mu $, $ \kappa_j=\kappa $ for $ j=1,\cdots,N $ for   $ \mu, \kappa>0 $. Then 
\begin{equation*} 
X_0(s,\tau)=0,\ \ \ \ X_{j}(s,\tau)=r_*e^{-\mathbf{i}\alpha\tau}e^{\mathbf{i}\frac{s}{h}}e^{\mathbf{i}\frac{2\pi (j-1)}{N}},\ \ \ \ \ \ \  j=1,\cdots,N,
\end{equation*}
satisfies \eqref{NP vor fi}, where $ \alpha=\frac{\kappa}{4\pi}\left( \frac{1}{h^2}-\frac{N-1}{r_*^2}\right)-\frac{\mu}{2\pi r_*^2}  $. This corresponds to  
\begin{equation}\label{fila 2}
\tilde{X}_{0,\varepsilon}(s,t)=0,\ \ \  \tilde{X}_{j,\varepsilon}(s,t)=\frac{r_*}{\sqrt{|\ln\varepsilon|}}e^{-\mathbf{i}\alpha|\ln\varepsilon|t}e^{\mathbf{i}\frac{s}{h}}e^{\mathbf{i}\frac{2\pi (j-1)}{N}},
 \ \ \ \ j=1,\cdots,N.
\end{equation}
We have 
\begin{theorem}\label{thm02}
	Let $ R^*>0 $, $ h>0 $ and $ N\geq 2 $ be an integer. For $ r_*\in (0,R^*) , \kappa>0, \mu>0 $, let $ \tilde{X}_{j,\varepsilon} $ $(j=0,1,\cdots,N)$ be $ N+1 $ filaments satisfying \eqref{fila 2}. Then  there exists $ \varepsilon_0>0 $, such that for every $ \varepsilon\in(0,\varepsilon_0] $, \eqref{Euler eq2} has a family of solutions with vorticities  $ \mathbf{w}_\varepsilon\in C^1(B_{R^*}(0)\times \mathbb{R}\times\mathbb{R}^+) $ satisfying in distributional sense
	\begin{equation*} 
	\mathbf{w}_\varepsilon(x,t)-\mu\delta_{\tilde{X}_{0,\varepsilon}}\textbf{t}_{\tilde{X}_{0,\varepsilon}}- \kappa\sum_{j=1}^N\delta_{\tilde{X}_{j,\varepsilon}}\textbf{t}_{\tilde{X}_{j,\varepsilon}}\rightharpoonup0\ \ \ \ \text{as}\ \varepsilon\to0.
	\end{equation*}
	Moreover, if we set $ A_{\varepsilon,i}=supp\left (\mathbf{w}_\varepsilon(x_1,x_2,0,0)\right ) \cap B_{\rho_0|\ln\varepsilon|^{-\frac{1}{2}}}\left ( \tilde{X}_{i,\varepsilon}(0,0)\right )  $, then there exist constants $ R_1,R_2>0 $ independent of $ \varepsilon $  such that
	\begin{equation*}
	R_1\varepsilon \leq diam(A_{\varepsilon,i})\leq R_2\varepsilon .
	\end{equation*}
\end{theorem}

\begin{remark}
\cite{GM2} considered the above case with the circulation $ \kappa=\mu=8\pi $. In contrast to \cite{GM2}, $ \kappa,\mu $ in our construction can be arbitrary positive constants. 
\end{remark}

We also prove the existence  of a family of solutions of  \eqref{Euler eq} with vorticity fields  concentrating near co-rotating helical solutions of   model \eqref{NP vor fi} when filaments are not symmetric about $ x_3 $-axis. 

\textbf{Case 3:}   ``Co-rotating asymmetric $2 $ vortices" with different circulations and rotation radii. Let $ N= 2 $, $ \kappa_1,\kappa_2>0 $ and $ \lambda_{1,*},\lambda_{2,*}>0 $ be constants satisfying the compatibility condition
\begin{equation}\label{compati cond1} 
\frac{\kappa_2}{\kappa_1}=\frac{2\lambda_{1,*}h^2+\lambda_{1,*}\lambda_{2,*}(\lambda_{1,*}+\lambda_{2,*})}{2\lambda_{2,*}h^2+\lambda_{1,*}\lambda_{2,*}(\lambda_{1,*}+\lambda_{2,*})}.
\end{equation}
Then 
\begin{equation*} 
X_{1}(s,\tau)=\lambda_{1,*}e^{-\mathbf{i}\alpha\tau}e^{\mathbf{i}\frac{s}{h}},\ \ \ \ 	X_{2}(s,\tau)=-\lambda_{2,*}e^{-\mathbf{i}\alpha\tau}e^{\mathbf{i}\frac{s}{h}},
\end{equation*}
satisfies \eqref{NP vor fi}, where $ \alpha=\frac{\kappa_1}{4\pi h^2}  -\frac{\kappa_2}{2\pi \lambda_{1,*}(\lambda_{1,*}+\lambda_{2,*})}=\frac{\kappa_2}{4\pi h^2}  -\frac{\kappa_1}{2\pi \lambda_{2,*}(\lambda_{1,*}+\lambda_{2,*})}  $. This corresponds to nearly parallel vortex filaments  
\begin{equation}\label{fila 3}
\tilde{X}_{j,\varepsilon}(s,t)=\frac{(-1)^{j-1}\lambda_{j,*}}{\sqrt{|\ln\varepsilon|}}e^{-\mathbf{i}\alpha|\ln\varepsilon|t}e^{\mathbf{i}\frac{s}{h}},\ \ \ \ j=1,2. 
\end{equation}
We note that \eqref{compati cond1} is necessary since \eqref{NP vor fi} has solutions being of the form $ X_{1}(s,\tau)=\lambda_{1,*}e^{-\mathbf{i}\alpha\tau}e^{\mathbf{i}\frac{s}{h}}, X_{2}(s,\tau)=-\lambda_{2,*}e^{-\mathbf{i}\alpha\tau}e^{\mathbf{i}\frac{s}{h}}  $ if and only if \eqref{compati cond1} holds. There are also  infinitely many  constant pairs $ (\kappa_1, \kappa_2,\lambda_{1,*},\lambda_{2,*}) $ satisfying \eqref{compati cond1} with $ \kappa_1\neq \kappa_2, \lambda_{1,*}\neq \lambda_{2,*} $. We have 
\begin{theorem}\label{thm03}
	Let $ R^*>0 , h>0 $ and $ N=2 $. For $ \kappa_1,\kappa_2>0 $ and $ \lambda_{1,*},\lambda_{2,*}\in (0,R^*) $   satisfying the compatibility condition \eqref{compati cond1},
%	\begin{equation*} 
%	\frac{\kappa_2}{\kappa_1}=\frac{2\lambda_{1,*}h^2+\lambda_{1,*}\lambda_{2,*}(\lambda_{1,*}+\lambda_{2,*})}{2\lambda_{2,*}h^2+\lambda_{1,*}\lambda_{2,*}(\lambda_{1,*}+\lambda_{2,*})},
%	\end{equation*}
let $ \tilde{X}_{j,\varepsilon} $ $(j=1,2)$ be $ 2 $ filaments satisfying \eqref{fila 3}. Then  there exists $ \varepsilon_0>0 $, such that for every $ \varepsilon\in(0,\varepsilon_0] $, \eqref{Euler eq2} has a family of solutions with vorticities  $ \mathbf{w}_\varepsilon\in C^1(B_{R^*}(0)\times \mathbb{R}\times\mathbb{R}^+) $ satisfying in distributional sense
	\begin{equation*} 
	\mathbf{w}_\varepsilon(x,t)-  \sum_{j=1}^2\kappa_j\delta_{\tilde{X}_{j,\varepsilon}}\textbf{t}_{\tilde{X}_{j,\varepsilon}}\rightharpoonup0\ \ \ \ \text{as}\ \varepsilon\to0.
	\end{equation*}
	Moreover, if we define   $ A_{\varepsilon,j}=supp\left (\mathbf{w}_\varepsilon(x_1,x_2,0,0)\right ) \cap B_{\rho_0|\ln\varepsilon|^{-\frac{1}{2}}}\left ( \tilde{X}_{j,\varepsilon}(0,0)\right )  $, then there exist constants $ R_1,R_2>0 $ independent of $ \varepsilon $  such that
	\begin{equation*}
	R_1\varepsilon \leq diam(A_{\varepsilon,j})\leq R_2\varepsilon .
	\end{equation*}
\end{theorem}

\textbf{Case 4:} ``Co-rotating  $ 2\times2 $ vortices" with different circulations and rotation radii. In this case, $ N= 4 $, $ \kappa_1=\kappa_3=\kappa>0 $, $ \kappa_2=\kappa_4=\mu>0$. $ \kappa,\mu, \lambda_{1,*},\lambda_{2,*}  $ are constants satisfying the compatibility condition
\begin{equation}\label{compa cond2}
\frac{\kappa}{4\pi h^2} -\frac{\kappa}{4\pi \lambda_{1,*}^2} -\frac{\mu}{\pi \left (\lambda_{1,*}^2+\lambda_{2,*}^2\right )}=\frac{\mu}{4\pi h^2} -\frac{\mu}{4\pi \lambda_{2,*}^2} -\frac{\kappa}{\pi \left (\lambda_{1,*}^2+\lambda_{2,*}^2\right )}.
\end{equation}
Then 
\begin{equation*} 
X_{2j-1}(s,\tau)=\lambda_{1,*}e^{-\mathbf{i}\alpha\tau}e^{\mathbf{i}\frac{s}{h}}e^{\mathbf{i}\pi(j-1)},\ \ \ \ 	X_{2j}(s,\tau)=\lambda_{2,*}e^{-\mathbf{i}\alpha\tau}e^{\mathbf{i}\frac{s}{h}}e^{\mathbf{i}\frac{\pi(2j-1)}{2}},\ \ j=1,2,
\end{equation*}
satisfies \eqref{NP vor fi}, where $ \alpha=\frac{\kappa}{4\pi h^2} -\frac{\kappa}{4\pi \lambda_{1,*}^2} -\frac{\mu}{\pi \left (\lambda_{1,*}^2+\lambda_{2,*}^2\right )}=\frac{\mu}{4\pi h^2} -\frac{\mu}{4\pi \lambda_{2,*}^2} -\frac{\kappa}{\pi \left (\lambda_{1,*}^2+\lambda_{2,*}^2\right )}  $. This corresponds to  
\begin{equation}\label{fila 4}
\tilde{X}_{2j-1,\varepsilon}(s,t)=\frac{\lambda_{1,*}}{\sqrt{|\ln\varepsilon|}}e^{-\mathbf{i}\alpha|\ln\varepsilon| t}e^{\mathbf{i}\frac{s}{h}}e^{\mathbf{i}\pi(j-1)},\ \  	\tilde{X}_{2j,\varepsilon}(s,t)=\frac{\lambda_{2,*}}{\sqrt{|\ln\varepsilon|}}e^{-\mathbf{i}\alpha|\ln\varepsilon| t}e^{\mathbf{i}\frac{s}{h}}e^{\mathbf{i}\frac{\pi(2j-1)}{2}},\ \ j=1,2.
\end{equation}
We have
\begin{theorem}\label{thm04}
	Let $ R^*>0, h>0, N=4 $. For $ \kappa,\mu>0 $ and $ \lambda_{1,*},\lambda_{2,*}\in (0,R^*) $ satisfying the compatibility condition \eqref{compa cond2},
%	\begin{equation*} 
%	\frac{\kappa}{4\pi h^2}-\frac{\kappa}{4\pi\lambda_{1,*}^2}-\frac{\mu}{\pi(\lambda_{1,*}^2+\lambda_{2,*}^2)}=\frac{\mu}{4\pi h^2}-\frac{\mu}{4\pi\lambda_{2,*}^2}-\frac{\kappa}{\pi(\lambda_{1,*}^2+\lambda_{2,*}^2)},
%	\end{equation*}  
let $ \tilde{X}_{j,\varepsilon} $ $(j=1,\cdots,4)$ be $ 4 $ filaments satisfying \eqref{fila 4}. Then  there exists $ \varepsilon_0>0 $, such that for every $ \varepsilon\in(0,\varepsilon_0] $, \eqref{Euler eq2} has a family of solutions with vorticities $ \mathbf{w}_\varepsilon\in C^1(B_{R^*}(0)\times \mathbb{R}\times\mathbb{R}^+) $ satisfying in distributional sense
	\begin{equation*} 
	\mathbf{w}_\varepsilon(x,t)-  \kappa\sum_{j=1}^2\delta_{\tilde{X}_{2j-1,\varepsilon}}\textbf{t}_{\tilde{X}_{2j-1,\varepsilon}}-\mu\sum_{j=1}^2\delta_{\tilde{X}_{2j,\varepsilon}}\textbf{t}_{\tilde{X}_{2j,\varepsilon}}\rightharpoonup0\ \ \ \ \text{as}\ \varepsilon\to0.
	\end{equation*}
	Moreover, if we denote   $ A_{\varepsilon,j}=supp\left (\mathbf{w}_\varepsilon(x_1,x_2,0,0)\right ) \cap B_{\rho_0|\ln\varepsilon|^{-\frac{1}{2}}}\left ( \tilde{X}_{j,\varepsilon}(0,0)\right )  $, then there exist constants $ R_1,R_2>0 $ independent of $ \varepsilon $  such that
	\begin{equation*}
	R_1\varepsilon \leq diam(A_{\varepsilon,j})\leq R_2\varepsilon .
	\end{equation*}
\end{theorem}

\textbf{Case 5:}  ``Co-rotating  $ 2\times2+1 $ vortices" with different circulations and rotation radii and the $ 5 $-th vortex of arbitrary strength
at the center of vorticity. Let $ N= 5 $, $ \kappa_1=\kappa_3=\kappa>0 $, $ \kappa_2=\kappa_4=\mu>0$, $ \kappa_0>0 $ and $ \lambda_{1,*},\lambda_{2,*}>0 $. $ \kappa_0,\kappa,\mu, \lambda_{1,*},\lambda_{2,*}  $  satisfy  the compatibility condition
\begin{equation}\label{compa cond3} 
\frac{\kappa}{4\pi h^2}-\frac{\kappa_0}{2\pi \lambda_{1,*}^2} -\frac{\kappa}{4\pi \lambda_{1,*}^2} -\frac{\mu}{\pi \left (\lambda_{1,*}^2+\lambda_{2,*}^2\right )}=\frac{\mu}{4\pi h^2}-\frac{\kappa_0}{2\pi \lambda_{2,*}^2} -\frac{\mu}{4\pi \lambda_{2,*}^2} -\frac{\kappa}{\pi \left (\lambda_{1,*}^2+\lambda_{2,*}^2\right )}.
\end{equation}
Then $ X_0(s,\tau)=0 $,
\begin{equation*}
X_{2j-1}(s,\tau)=\lambda_{1,*}e^{-\mathbf{i}\alpha\tau}e^{\mathbf{i}\frac{s}{h}}e^{\mathbf{i}\pi(j-1)},\ \ \ \ 	X_{2j}(s,\tau)=\lambda_{2,*}e^{-\mathbf{i}\alpha\tau}e^{\mathbf{i}\frac{s}{h}}e^{\mathbf{i}\frac{\pi(2j-1)}{2}},\ \ j=1,2,
\end{equation*}
satisfy \eqref{NP vor fi}, where $ \alpha=\frac{\kappa}{4\pi h^2}-\frac{\kappa_0}{2\pi \lambda_{1,*}^2} -\frac{\kappa}{4\pi \lambda_{1,*}^2} -\frac{\mu}{\pi \left (\lambda_{1,*}^2+\lambda_{2,*}^2\right )}=\frac{\mu}{4\pi h^2}-\frac{\kappa_0}{2\pi \lambda_{2,*}^2} -\frac{\mu}{4\pi \lambda_{2,*}^2} -\frac{\kappa}{\pi \left (\lambda_{1,*}^2+\lambda_{2,*}^2\right )}  $. This corresponds to  $ \tilde{X}_{0,\varepsilon}(s,t)=0 $,
\begin{equation} \label{fila 5}
\tilde{X}_{2j-1,\varepsilon}(s,t)=\frac{\lambda_{1,*}}{\sqrt{|\ln\varepsilon|}}e^{-\mathbf{i}\alpha|\ln\varepsilon|t}e^{\mathbf{i}\frac{s}{h}}e^{\mathbf{i}\pi(j-1)},\ \ 	\tilde{X}_{2j,\varepsilon}(s,t)=\frac{\lambda_{2,*}}{\sqrt{|\ln\varepsilon|}}e^{-\mathbf{i}\alpha|\ln\varepsilon|t}e^{\mathbf{i}\frac{s}{h}}e^{\mathbf{i}\frac{\pi(2j-1)}{2}},\ \ j=1,2.
\end{equation}
We have
\begin{theorem}\label{thm05}
	Let $ R^*>0, h>0, N=5 $. For $ \kappa_0,\kappa,\mu>0 $ and $ \lambda_{1,*},\lambda_{2,*}\in(0,R^*) $ satisfying \eqref{compa cond3}, %the compatibility condition
%	\begin{equation*} 
%	\frac{\kappa}{4\pi h^2}-\frac{\kappa_0}{2\pi\lambda_{1,*}^2}-\frac{\kappa}{4\pi\lambda_{1,*}^2}-\frac{\mu}{\pi(\lambda_{1,*}^2+\lambda_{2,*}^2)}=\frac{\mu}{4\pi h^2}-\frac{\kappa_0}{2\pi\lambda_{2,*}^2}-\frac{\mu}{4\pi\lambda_{2,*}^2}-\frac{\kappa}{\pi(\lambda_{1,*}^2+\lambda_{2,*}^2)},
%	\end{equation*}  
let $ \tilde{X}_{j,\varepsilon} $ $(j=0,1,\cdots,4)$ be $ 5 $ filaments satisfying \eqref{fila 5}. Then  there exists $ \varepsilon_0>0 $, such that for every $ \varepsilon\in(0,\varepsilon_0] $, \eqref{Euler eq2} has a family of solutions with vorticities $ \mathbf{w}_\varepsilon\in C^1(B_{R^*}(0)\times \mathbb{R}\times\mathbb{R}^+) $ satisfying in distributional sense
	\begin{equation*} 
	\mathbf{w}_\varepsilon(x,t)- \kappa_0\delta_{\tilde{X}_{0,\varepsilon}}\textbf{t}_{\tilde{X}_{0,\varepsilon}}- \kappa\sum_{j=1}^2\delta_{\tilde{X}_{2j-1,\varepsilon}}\textbf{t}_{\tilde{X}_{2j-1,\varepsilon}}-\mu\sum_{j=1}^2\delta_{\tilde{X}_{2j,\varepsilon}}\textbf{t}_{\tilde{X}_{2j,\varepsilon}}\rightharpoonup0\ \ \ \ \text{as}\ \varepsilon\to0.
	\end{equation*}
	Moreover, if we define   $ A_{\varepsilon,j}=supp\left (\mathbf{w}_\varepsilon(x_1,x_2,0,0)\right ) \cap B_{\rho_0|\ln\varepsilon|^{-\frac{1}{2}}}\left ( \tilde{X}_{j,\varepsilon}(0,0)\right )  $, then there exist constants $ R_1,R_2>0 $ independent of $ \varepsilon $  such that
	\begin{equation*}
	R_1\varepsilon \leq diam(A_{\varepsilon,j})\leq R_2\varepsilon .
	\end{equation*}
\end{theorem}

\begin{remark}
%Theorems \ref{thm03}, \ref{thm04} and \ref{thm05} consider the cases when filaments of \eqref{NP vor fi} are not symmetric, which are not known in existing results. 
It is possible to extend our results to more general cases, for example, ``$ 2N $ filaments": $ N $ of them locate at vertices of a regular $ N $-polygon with equal circulations $ \kappa_1 $ and rotation radii $ \lambda_1 $, and another $ N $ of them locate at vertices of another regular $ N $-polygon with  equal circulations $ \kappa_2 $ and rotation radii $ \lambda_2 $; and  ``$ 2N+1 $ filaments": $ 2N $ of them are ``$ 2N $ filaments" and the $ 2N+1 $ vortex locates at the center of vorticity with arbitrary positive circulation.

We note that recently, Cao et al. \cite{CFQW2} considered dynamics of nearly parallel helical vortex filaments of 3D incompressible Euler equations. If initial vorticity is
 concentrated in $ \varepsilon- $neighborhoods of $ N $ distinct helical filaments  with 
distances of order $ O(|\ln\varepsilon|^{-\frac{1}{2}}) $, then as time evolves the vorticity field of \eqref{Euler eq} still concentrates near helical vortex filaments evolved by  nearly parallel  vortex filaments model \ref{NP vor fi}, at least for a small time interval. For more results, see \cite{CFQW1}. 


\end{remark}

Our proofs are based on the study of concentrated ``clustered solutions" to some semilinear second-order elliptic equations (see \eqref{key equa} in the following) with cluster point being the origin $ 0 $ and cluster distance being $ |\ln\varepsilon|^{-\frac{1}{2}} $, which means that solutions consist of several bubbles collapsing into $ 0 $ as parameter $ \varepsilon\to0 $  rather than several distinct points, and the distance of any two bubbles is $ O\left ( |\ln\varepsilon|^{-\frac{1}{2}}\right ) $. We call $ 0 $ the cluster point and  $ |\ln\varepsilon|^{-\frac{1}{2}} $  the cluster distance. Following the deduction argument in \cite{CW,CW1,DDMW2, ET,GM}, to solve helical vorticity field of \eqref{Euler eq2} with pitch $ h $ and the orthogonal condition $ \textbf{v}\cdot (x_2,-x_1,h)^T=0 $, it is equivalent to solve the following 2D system















%\subsection{Background}
%In this paper, we study the existence of  concentrated clustered helical vortex solutions to \eqref{Euler eq2}. 
%The research of solutions to 3D Euler equations with helical symmetry has attracted great attention in the past decades, see \cite{AS,CW,DDMW,DDMW2,Du,ET,Jiu} and reference therein. Let us first define helical solutions of \eqref{Euler eq}, see \cite{CW2,ET}.  For fixed $ k>0 $, let $ \mathcal{G}_k=\{H_\rho:\mathbb{R}^3\to\mathbb{R}^3 \} $ be the helical transformation group, where
%\begin{equation*}
%H_\rho(x_1,x_2,x_3)^t=(x_1\cos\rho+x_2\sin\rho, -x_1\sin\rho+x_2\cos\rho, x_3+k\rho)^t.
%\end{equation*}
%Here $ A^t $ is the transposition of a matrix $ A $. 
%Let
%$ R_\rho=\begin{pmatrix}
%\cos\rho & \sin\rho & 0 \\
%-\sin\rho &\cos\rho & 0 \\
%0 & 0 & 1
%\end{pmatrix} $
%be the rotation with respect to $ x_3- $axis. Define the field of tangents of symmetry lines of $ \mathcal{G}_k $
%\begin{equation*}
%\overrightarrow{\zeta}= (x_2, -x_1, k)^t.
%\end{equation*}
%
%Helical solutions must define on helical domains. A domain $ D\in \mathbb{R}^3 $ is called  a  helical~domain, if $ H_\rho(D)=D $ for any $\rho\in\mathbb{R}$.
%Let $ \Omega=D\cap\{x\mid x_3=0\} $ be the section of $ D $ over $ x_1Ox_2 $ plane. Then $ D  $ can be generated by $ \Omega $ by  letting $ D=\cup_{\rho\in\mathbb{R}}H_\rho(\Omega) $. In the following, we always assume that  $ \Omega $ is a simply-connected bounded domain with $ C^\infty $ boundary.
%
%Helical solutions to \eqref{Euler eq} is then defined as follows. We say that  ($\mathbf{v}, P$) is   a  helical solution to \eqref{Euler eq} with pitch $ k $, if ($\mathbf{v}, P$) satisfies \eqref{Euler eq} and the  vector field $ \mathbf{v} $ and scalar function $ P $ satisfies for every $ \rho\in\mathbb{R}, x\in D  $
%\begin{equation}\label{helical func}
%P(H_{\rho}(x))=P(x);\ \ \mathbf{v}(H_{\rho}(x))=R_\rho \mathbf{v}(x).
%\end{equation}
%Moreover, we also impose  $ \mathbf{v} $ to satisfy  the following orthogonality condition:
%\begin{equation}\label{ortho}
%\mathbf{v}\cdot \overrightarrow{\zeta}=0.
%\end{equation}
%Under assumptions \eqref{helical func}  and \eqref{ortho}, it can be proved  that $ \mathbf{w} $ satisfies (see \cite{ET})
%\begin{equation}\label{w formula}
%\mathbf{w}=\frac{\omega}{k}\overrightarrow{\zeta},
%\end{equation}
%where $ \omega:=w_3=\partial_{x_1}v_2-\partial_{x_2}v_1 $, the third component of vorticity field $ \mathbf{w} $, is a helical function. Moreover, $ \omega $ satisfies the 2D vorticity equations
\begin{equation}\label{vor str eq}
\begin{cases}
\partial_t \omega+\nabla \omega\cdot \nabla^\perp\varphi=0, &\Omega\times (0,T),\\
\omega=\mathcal{L}_{K_H}\varphi, &\Omega\times (0,T),\\
\varphi =0,& \partial\Omega,
\end{cases}
\end{equation}
%\begin{equation}\label{vor str eq2}
%\begin{cases}
%\partial_t \omega+\nabla \omega\cdot \nabla^\perp\varphi=0, &\Omega\times (0,T),\\
%\omega=\mathcal{L}_H\varphi, &\Omega\times (0,T),\\
%\nabla^\perp\varphi\cdot \nu|_{\partial \Omega}=v_n|_{\partial \Omega},
%\end{cases}
%\end{equation}
where $ \omega  $ is the third component of $ \mathbf{w} $, $ \varphi $ is the corresponding stream function, $ \perp $ denotes the clockwise rotation through $ \frac{\pi}{2}  $, $ \mathcal{L}_{K_H} =-\text{div}(K_H(x_1,x_2)\nabla ) $ is   a second-order elliptic operator of divergence type with the coefficient matrix
\begin{equation}\label{coef matrix}
K_H(x_1,x_2)=\frac{1}{h^2+x_1^2+x_2^2}
\begin{pmatrix}
h^2+x_2^2 & -x_1x_2 \\
-x_1x_2 &  h^2+x_1^2
\end{pmatrix}.
\end{equation}
%and $ \nu=(n_1,n_2) $ is the two-dimensional vector of the first two component of $ \mathbf{n} $ on $ \partial \Omega $, which is an outward normal to $ \partial \Omega $.
%see \cite{CW2} for more details. 
%Thus for  helical solution pairs ($\mathbf{v}, P$) to \eqref{Euler eq}, it suffices to solve solutions $ \omega $ to \eqref{vor str eq}.
%Note also that when $ v_n\equiv 0 $, which corresponds to the impermeable boundary condition, the third equation of \eqref{vor str eq2} implies that $ \varphi$  is   0 on $ \partial \Omega $. Thus we get the  vorticity equations under the impermeable boundary condition
%Note that since $ \mathcal{L}_{K_H} $ is a uniformly elliptic operator and has the same $ L^p $ estimates as $ -\Delta $, many references get similar well-posedness and stability results of solutions to \eqref{vor str eq} as those to 2D Euler equations. 
%\cite{ET} proved the global well-posedness of $ L^1\cap L^\infty $ weak solutions to  \eqref{vor str eq}, which coincides with the classical Yudovich's result \cite{Y} in 2D Euler flows. \cite{Ben} considered nonlinear stability of stationary smooth Euler flows with helical symmetry  by using the direct method of Lyapunov. For more results of the existence and regularity of helical solutions to Euler equations, see \cite{AS, BLN, Du, Jiu}.
%
%The problem of concentrated helical symmetric solutions to \eqref{Euler eq2}, meanwhile,  has been widely concerned in recent years, see \cite{CW,CW1,CW2,DDMW2,GM,GZ,DLM}. It is also called the $ vortex~filament~conjecture $ to  3D Euler equations with helical symmetry, that is,  constructing ``true'' helical solutions to Euler equations  such that the corresponding vorticity  concentrates near a helix, see \cite{JS2}. The research of this problem  can be traced back to
%Helmholtz \cite{He}
%%, who  first studied the motion of the traveling vortex rings  whose vorticities are supported in toroidal regions with a small  cross-section. Then
%and then many articles proved the existence of vortex solutions to Euler equations concentrating near a straight line and a circle, see \cite{CLW,CPY,DDMW,DV,Fr1,FB,SV} for example. %This problem  is called the vortex desingularization problem to 2D Euler equation and 3D axisymmetric Euler equation.
%For concentrated solutions to Euler equations with helical symmetry,  
%For the case that the filament is a helix, D$\acute{\text{a}}$vila et al.  \cite{DDMW2} first constructed rotational-invariant smooth Euler flows with helical symmetry in the whole space $ \mathbb{R}^3 $. 

For $ \alpha \in\mathbb{R} $, if we consider    solutions to \eqref{vor str eq} with the form
\begin{equation}\label{104}
\omega(x,t)=w\left (\bar{R}_{-\alpha|\ln\varepsilon| t}(x)\right ),\ \ \varphi(x,t)=u\left (\bar{R}_{-\alpha|\ln\varepsilon| t}(x)\right ),
\end{equation}
where   $ \bar{R}_{\alpha t}=\begin{pmatrix}
\cos\alpha t & \sin\alpha t \\
-\sin\alpha t &\cos\alpha t
\end{pmatrix} $, then %Taking \eqref{104} into \eqref{vor str eq}, one has
\begin{equation}\label{rot eq-1}
\begin{cases}
\nabla w\cdot \nabla^\perp \left( u-\frac{\alpha}{2}|x|^2|\ln\varepsilon|\right) =0,\ \ &\Omega,\\
w=\mathcal{L}_{K_H}u,\ \ &\Omega,\\
u=0,\ \ &\partial\Omega.
\end{cases}
\end{equation}
A sufficient condition for $ (w,u) $ satisfying \eqref{rot eq-1} is that $ u $ solves 
\begin{equation}\label{105-1}
\mathcal{L}_{K_H} u=w=f_\varepsilon\left( u-\frac{\alpha}{2}|x|^2|\ln\varepsilon|\right)\ \  \text{in}\ \Omega,\ \ u=0\ \ \text{on}\ \partial\Omega,
\end{equation}
for some profile function $ f_\varepsilon $. Thus the key is to find concentrated clustered solutions to \eqref{105-1} under the proper choice of $ f_\varepsilon $. To make the support set of cross-section of $ \textbf{w}_\varepsilon $ compact, we choose $ f_\varepsilon(t)=\frac{1}{\varepsilon^2}t^p_+ $ with $ p>1 $ in \eqref{105-1}.

%, then \eqref{rot eq-1} automatically holds. Choosing  $ f_\varepsilon(t)=\varepsilon^2e^t $, \cite{DDMW2} proved the existence of solutions to \eqref{105-1} concentrating near several distinct points in the distributional sense as $ \varepsilon\to0 $.
%Note that by the choice of  $ f_\varepsilon $, the support set of vorticity is still the whole plane. Recently, the authors in \cite{CW} considered
%rotational-invariant  concentrated helical solutions with small cross-section to \eqref{Euler eq2} in infinite cylinders $ B_{R^*}(0)\times\mathbb{R} $. From \cite{CW}, it suffices to solve 
%a semilinear elliptic equations in divergence form
%\begin{equation}\label{rot eq2}
%\begin{split}
%-\text{div}(K(x)\nabla u)=f_\varepsilon\left(u-q |\ln\varepsilon|\right)\ \ \text{in}\  \Omega;\ \
%u(x)=0\ \ \text{on}\  \partial \Omega
%\end{split}
%\end{equation}
%for some profile function $ f_\varepsilon  $, where  $ K $ is a positive-definite matrix and the function $ q $ is positive. Denote   $ \det (K) $ the determinant of $ K $. By taking $ f_\varepsilon(t)=\frac{1}{\varepsilon^2}t^p_+ $ for $ p>1 $ and using the asymptotic expansion of Green's function $ G_K $ of the elliptic operator $ -\text{div}\cdot(K(x)\nabla  ) $ being of the form
%\begin{equation*}
%G_K(x, y)=\frac{\sqrt{\det K(x)}^{-1}+\sqrt{\det K(y)}^{-1}}{2}\Gamma\left (\frac{T_x+T_y}{2}(x-y)\right )+S_K(x,y),
%\end{equation*}
%where $ \Gamma(x)=-\frac{1}{2\pi}\ln|x| $, $ (T_x)^{-1}(T_x)^{-t}=K(x) $ and $ S_K(x,y)\in C^{0,\gamma}(\Omega\times\Omega) $ for $ \gamma\in(0,1) $, \cite{CW} proved the following results:
%\begin{customthm}{A}[\cite{CW}]\label{thm1A}
%For any given $ l $ distinct strict local minimum (maximum) points $ x_{0,j}(j=1,\cdots, l)$ of $ q^2\sqrt{\det(K)} $ in $ \Omega $, there exists $ \varepsilon_0>0 $  such that for every $ \varepsilon\in(0,\varepsilon_0] $, \eqref{rot eq2} has a solution $ u_\varepsilon $ satisfying
%\begin{enumerate}
%	\item  Define $ \bar{A}_{\varepsilon,i}=\left\{u_\varepsilon>q\ln\frac{1}{\varepsilon}\right\}\cap B_{\bar{\rho}}(x_{0,i})  $, where $ \bar{\rho}>0 $ is small. Then there exist $ (z_{1,\varepsilon}, \cdots, z_{l,\varepsilon}) $ and $ R_1,R_2>0 $ independent of $ \varepsilon $ satisfying
%	\begin{equation*}
%	\lim_{\varepsilon\to 0}(z_{1,\varepsilon}, \cdots, z_{l,\varepsilon})=(x_{0,1}, \cdots, x_{0,l});\ \ 	 B_{R_1\varepsilon}(z_{i,\varepsilon})\subseteq \bar{A}_{\varepsilon,i}\subseteq B_{R_2\varepsilon}(z_{i,\varepsilon}).
%	\end{equation*}	
%	\item
%	\begin{equation*}
%	\lim_{\varepsilon\to 0}\frac{1}{\varepsilon^2}\int_{B_{\bar{\rho}}(x_{0,i})}\left( u_\varepsilon-q|\ln\varepsilon|\right)^{p}_+dx=2\pi q\sqrt{\det(K)}(x_{0,i}).
%	\end{equation*}
%\end{enumerate}
%\end{customthm}
%By choosing $ K=K_H $, $ q=\frac{\alpha}{2}|x|^2+\beta $ for proper constants $ \alpha,\beta $ in Theorem \ref{thm1A}, \cite{CW} constructed  multiple traveling-rotating helical vortices with small cross-section in $ B_{R^*}(0)\times \mathbb{R} $. The concentrating locations $ x_{0,j} $ are $l $  distinct  points, which constitute the vertices of a regular polygon. For more results, see  \cite{CW1, CW2}.
%
%
%Results in \cite{CW,CW1,CW2} indicate that there exist concentrated helical  vortex solutions concentrating near several $ distinct $ helices in $D$. A natural question is, are there  helical vortex solutions to \eqref{Euler eq2}, whose support sets consist of several helical tubes and collapse into a single helix as parameter $ \varepsilon\to 0 $? We call this $ clustered~solutions$ of \eqref{Euler eq2}. From the deduction of \eqref{rot eq2}, the question becomes whether there exists a family of solutions to \eqref{rot eq2}, such that solutions consist of several bubbles which collapse into a single point as $ \varepsilon\to 0 $. Note that when $ K(x)\equiv Id $,  \eqref{rot eq2} becomes vorticity equations of  2D steady Euler equations. In this case, classical results (see \cite{CGPY,CPY}) indicate that limiting locations of concentrated solutions must be critical points of the Kirchhoff-Routh function, which are $ l $ distinct points in $\Omega$. So clustered solutions to 2D steady Euler equations seem to not exist. As for vortex rings to 3D Euler equations, \cite{ALW} constructed smooth clustered solutions shrinking to a circle. For helical symmetric solutions to \eqref{Euler eq2}, very recently   \cite{GM}   constructed  smooth clustered  solutions  to \eqref{105-1} by choosing proper $ f_\varepsilon $, whose support sets consist of several bubbles and collapse into a single point in $ \mathbb{R}^2 $ as $ \varepsilon\to0 $. This corresponds to clustered helical vorticity fields of \eqref{Euler eq2} in $ \mathbb{R}^3 $. Because of the choice of the vortex profile $ f(t)=e^t $, support sets of  vorticity fields constructed in \cite{ALW,GM} are the whole space $ \mathbb{R}^3 $, rather than helical tubes with small cross-section.
%
%Our goal in this paper is to construct clustered helical vortex solutions to Euler equations \eqref{Euler eq2} with small cross-section in helical domains, such that the support of vortices consists of several helical tubes   collapsing  into a single helix as $ \varepsilon\to 0 $. From the deduction of \eqref{rot eq2}, it suffices to construct clustered solutions to a semilinear elliptic equations in divergence form \eqref{rot eq2}.  We prove that, for any $ x_0 $ being a $ strict~local~maximizer $ of $ q^2\sqrt{\det (K)} $ in $\Omega$ and any positive integer $ m $, there exists a family of clustered solutions to \eqref{rot eq2} concentrating near $ m $ points $ (z_{1,\varepsilon},z_{2,\varepsilon},\cdots,z_{m,\varepsilon}) $, which satisfy  $ \lim_{\varepsilon\to 0}z_{i,\varepsilon} =x_0 $ for $ i=1,\cdots,m $. 
%%The key of proof is to get $ C^1-$asymptotic expansion of Green's function $ G_K $ of the operator $ -\text{div}(K(x)\nabla  ) $ (see Lemma \ref{Green expansion2}), the $ C^1- $dependence of the error term $ \omega_{\delta,Z} $ with respect to $ Z $ (see Proposition \ref{differ of w}) and the existence of critical points of energy $ K_\delta(Z) $ (see Proposition \ref{order of main term}).
%Therefore in our construction, solutions consist of several bubbles  concentrating near  a single point rather than $ m $ distinct points, which is quite different from known results in \cite{CW,CW1,DDMW}.
%
%\subsection{Main results}
%Now we begin to show our main results. Our first result is the existence of clustered solutions to a semilinear elliptic equation in divergence form. Consider \eqref{rot eq2} with $ f_\varepsilon(t)=\frac{1}{\varepsilon^2}t^p_+ $ for $ p>1 $, that is, 

Let us directly consider   clustered solutions to the following general semilinear second-order elliptic equations 
\begin{equation}\label{key equa}
\begin{cases}
-\varepsilon^2\text{div}\left (K(x)\nabla v\right )= \left (v-q|\ln\varepsilon|\right )^{p}_+,\ \ &x\in \Omega,\\
v=0,\ \ &x\in\partial \Omega,
\end{cases}
\end{equation}
with cluster point being the origin $ 0 $ and cluster distance being $ |\ln\varepsilon|^{-\frac{1}{2}} $, where $ \Omega\subset \mathbb{R}^2 $ is a smooth bounded domain containing $ 0 $, $ \varepsilon\in(0,1) $. $ K=(K_{i,j})_{2\times2} $  and $ q $ satisfy the following assumptions
\begin{enumerate}
	\item[(K1).] %$ -\text{div}(K(x)\nabla \cdot) $ is a uniformly elliptic operator, that is, 
	There exist  $ \Lambda_1,\Lambda_2>0 $ such that $$ \Lambda_1|\zeta|^2\le (K(x)\zeta|\zeta) \le \Lambda_2|\zeta|^2,\ \ \ \ \forall\ x\in \Omega, \ \zeta\in \mathbb{R}^2.$$
%\end{enumerate}
%%$ q  $ is a function defined in $ \overline{\Omega} $ satisfying
%\begin{enumerate}
	\item[(Q1).]  $ q \in C^{\infty}(\overline{\Omega}) $ and $ q(x)>0 $ for any $ x\in\overline{\Omega}. $
%\end{enumerate}
%$ K $ and $ q $ also need to satisfy another assumption
%\begin{enumerate}
	\item[(KQ).]   $ \nabla \left (q^2\sqrt{\det K}\right )(0)=0 $.
\end{enumerate}

Let us remark that the assumption $ (KQ) $ means that $ 0 $ is a critical point of $ q^2\sqrt{\det K} $, which seems to be artificial. However, when we prove Theorems \ref{thm01}-\ref{thm05} by the special form of $q$ and $K$, this assumption is satisfied automatically.

For any given integer $ N\geq 2 $, we define an auxiliary function from $ \Omega^{(N)} $ to $ \mathbb{R} $
\begin{equation}\label{def of H_N}
\begin{split}
H_N(\tilde{z}_1,\cdots, \tilde{z}_N)=&\frac{1}{2}\sum_{j=1}^N\tilde{z}_j\cdot \nabla^2\left(q^2\sqrt{\det K} \right)(0 )\cdot\tilde{z}_j^T\\
&+\sum_{1\leq i\neq j\leq N}\left( q^2  \sqrt{\det K}\right)(0)\ln| K(0)^{-\frac{1}{2}}\left( \tilde{z}_i - \tilde{z}_j\right)  |.
\end{split}
\end{equation}
Here $ \tilde{z}_j^T $ is the transpose of $ \tilde{z}_j $.
Under  assumptions $ (K1),(Q1)$  and  $(KQ) $, we can obtain the existence of clustered solutions to \eqref{key equa} with cluster point $ 0 $ and cluster distance $ |\ln\varepsilon|^{-\frac{1}{2}} $, whose bubbles concentrate near strict local maximizers (or minimizers) of $ H_N $ after $ |\ln\varepsilon|^{\frac{1}{2}} $ scaling. Indeed, we have the following result.%Our first result is as follows.
\begin{theorem}\label{thm0}
	Let $(K1)$, $ (Q1) $ and $ (KQ) $ hold.  %Assume that  $ \nabla \left (q^2\sqrt{\det K}\right )(0)=0 $. 
	Let $ \left( \tilde{z}_1^*,\cdots, \tilde{z}_N^*\right)  $ be a strict local maximizer or minimizer of $ H_N $ defined by \eqref{def of H_N}, i.e., there exists a $ \bar{\rho}- $neighborhood $ B_{\bar{\rho}} $ of $ \left( \tilde{z}_1^*,\cdots, \tilde{z}_N^*\right) $ for some $ \bar{\rho}>0 $ small  such that
	\begin{equation*}
	H_N\left( \tilde{z}_1,\cdots, \tilde{z}_N\right)<(\text{or} >) H_N\left( \tilde{z}_1^*,\cdots, \tilde{z}_N^*\right)\ \ \ \ \forall \left( \tilde{z}_1,\cdots, \tilde{z}_N\right)\in B_{\bar{\rho}}\backslash\{\left( \tilde{z}_1^*,\cdots, \tilde{z}_N^*\right) \}.
	\end{equation*} 
	Then  there exists $ \varepsilon_0>0 $, such that for every $ \varepsilon\in(0,\varepsilon_0] $, \eqref{key equa} has a  solution  $ v_\varepsilon $. Moreover,  
	\begin{enumerate}
		\item There exist $ \left (\tilde{z}_{1,\varepsilon}, \cdots, \tilde{z}_{N,\varepsilon}\right ) $ such that $ \tilde{z}_{i,\varepsilon}=\tilde{z}_{i}^*+o(1) $ for $ i=1,\cdots,N $ and 
		\begin{equation*}
		\left \{v_\varepsilon>q|\ln\varepsilon|\right \}\subseteq \cup_{i=1}^N B_{|\ln\varepsilon|^{-1}}\left (\frac{\tilde{z}_{i,\varepsilon}}{\sqrt{|\ln\varepsilon|}}\right ).
		\end{equation*}
		
		\item Define  sets $ A_{\varepsilon,i}=\left\{v_\varepsilon>q|\ln\varepsilon|\right\}\cap B_{|\ln\varepsilon|^{-1}}\left (\frac{\tilde{z}_{i,\varepsilon}}{\sqrt{|\ln\varepsilon|}}\right )  $. Then there exist constants $ R_1,R_2>0 $ independent of $ \varepsilon $  such that
		\begin{equation*}
		B_{R_1\varepsilon}\left (\frac{\tilde{z}_{i,\varepsilon}}{\sqrt{|\ln\varepsilon|}}\right )\subseteq A_{\varepsilon,i}\subseteq B_{R_2\varepsilon}\left (\frac{\tilde{z}_{i,\varepsilon}}{\sqrt{|\ln\varepsilon|}}\right ).
		\end{equation*}
		\item 
		\begin{equation*}
		\frac{1}{\varepsilon^2}\int_{A_{\varepsilon,i}}\left (v_\varepsilon-q|\ln \varepsilon|\right )^{p}_+dx=2\pi \left( q\sqrt{\det K}\right) (0)+O\left (\frac{1}{\sqrt{|\ln\varepsilon|}}\right ).
		\end{equation*}
	\end{enumerate}
\end{theorem}	



%\begin{theorem}\label{thm1}
%Suppose that $(\mathcal{K}1)$ and $ (Q1) $ hold. Let $ x_0 $ be a strict local maximum point of $ q^2\sqrt{\det (K)} $ in $ \Omega $, i.e., there exists $ \bar{\rho}>0 $ small such that
%\begin{equation*}
%q^2\sqrt{\det (K)}(y)< q^2\sqrt{\det (K)}(x_{0})\ \ \ \ \forall y\in B_{\bar{\rho}}(x_{0})\backslash\{x_{0}\}.
%\end{equation*} Then, for any $ m\in \mathbb{N}^* $  there exists $ \varepsilon_0>0 $, such that for every $ \varepsilon\in(0,\varepsilon_0] $, \eqref{eq1-1} has a family of clustered solutions $ u_\varepsilon $ with
%\begin{equation*}
%\frac{1}{\varepsilon^2}\int_{\Omega}\left( u_\varepsilon-q|\ln\varepsilon|\right)^{p}_+dx\to 2\pi m q\sqrt{\det (K)}(x_{0})  \ \ \text{as}\ \varepsilon\to0.
%\end{equation*}
%Moreover, there exist $ (z_{1,\varepsilon}, \cdots, z_{m,\varepsilon})\in \Omega^{(m)} $ such that
%\begin{equation*}
%|z_{i,\varepsilon}-z_{j,\varepsilon}|\geq |\ln\varepsilon|^{-m^2-1},\ \ \forall i\neq j; \ \ \left \{u_\varepsilon>q|\ln\varepsilon|\right \}\subseteq \cup_{i=1}^m B_{|\ln\varepsilon|^{-m^2-2}}(z_{i,\varepsilon})
%\end{equation*}
%and
%\begin{equation*}
%\lim_{\varepsilon\to 0}(z_{1,\varepsilon}, \cdots, z_{m,\varepsilon})=(x_{0}, \cdots, x_{0}).
%\end{equation*}	
%Define the  set $ A_{\varepsilon,i}=\left\{u_\varepsilon>q|\ln\varepsilon|\right\}\cap B_{|\ln\varepsilon|^{-m^2-2}}(z_{i,\varepsilon})  $. Then there exist constants $ R_1,R_2>0 $ independent of $ \varepsilon $  such that
%\begin{equation*}
%B_{R_1\varepsilon}(z_{i,\varepsilon})\subseteq A_{\varepsilon,i}\subseteq B_{R_2\varepsilon}(z_{i,\varepsilon}).
%\end{equation*}
%\end{theorem}

Indeed, by choosing properly $ K $ and $ q $ in Theorem \ref{thm0}, we can directly prove Theorems \ref{thm01}, \ref{thm02}, \ref{thm03}, \ref{thm04} and \ref{thm05}. 


\begin{remark}
%The assumption $ (KQ) $ in Theorem \ref{thm0} means that $ 0 $ is a critical point of $ q^2\sqrt{\det K} $, which seems to be artificial. However, when using Theorem \ref{thm0} to prove Theorems \ref{thm01}-\ref{thm05}, this assumption automatically holds. 
Indeed, if the assumption $ (KQ) $ fails, using our idea of proof of Theorem \ref{thm0} it is still possible to construct clustered solutions to \eqref{key equa} with cluster point $ x_0\neq 0  $ and cluster radius $ |\ln\varepsilon|^{-1} $. In \cite{GM}, the authors used gluing method to construct clustered solutions to \eqref{key equa} with cluster point $ x_0\neq 0  $ and cluster radius $ |\ln\varepsilon|^{-1} $ when $ \Omega=\mathbb{R}^2 $, $ K=K_H $ and profile functions $ f_\varepsilon  $ are exponential type functions. Note that $ x_0 $ is not a critical point of $ q^2\sqrt{\det K_H} $ in their results. Using methods in our paper, it is possible to generalize results in \cite{GM}, that is, constructing clustered solutions to \eqref{key equa} in case that  $ \Omega $ is  any bounded domain and $ K $ is any positive-definite matrix.
%It is quite surprising that accumulation of bubbles can occur for elliptic equations \eqref{rot eq2}. When $ K= Id $, such phenomenon may not exist, see \cite{CGPY} for example. %However when $ K= a(x)Id $ for some function $ a $,  clustered solutions of \eqref{rot eq2} with proper $ f_\varepsilon $ exist, see \cite{ALW,WYZ}. 
%In case that $ K= a(x)Id $ for some positive function $ a $, Wei-Ye-Zhou \cite{WYZ} considered an anisotropic Emden-Fowler equation
%\begin{equation*}
%\begin{cases}
%\text{div}(a(x)\nabla u)+\varepsilon^2a(x)e^u=0& \text{in}\ \Omega,\\
%u=0& \text{on}\ \partial\Omega.
%\end{cases}
%\end{equation*}
%%By using the expansion of Green's function of $ \nabla(a(x)\nabla ) $, the authors prove the existence of clustering solutions.where $ a $ is a smooth positive function in $ \Omega $. 
%For any  maximum points $ x^* $ of $ a $,  they constructed clustered 
%solutions concentrating near $ x^*  $. See also \cite{ALW}.
%For general $ K $ satisfying $(\mathcal{K}1)$, the only known results for such phenomena is given by \cite{GM, GM2}. The construction of clustered solutions of \eqref{rot eq2} in \cite{GM, GM2} depends on the choice of $ f_\varepsilon $ and the accurate expression of $ K_H $ in \eqref{coef matrix}. In contrast to \cite{GM}, Theorem \ref{thm1} shows the existence of clustered solutions to equations \eqref{eq1-1} with any positive-definite matrix $ K $. % by using $ C^1 $-asymptotic estimates of Green's function $ G_K $. 
%Another interesting phenomenon  is the multiplicity of solutions to \eqref{eq1-1}. Indeed, Theorem \ref{thm1} shows that there exists solutions of \eqref{eq1-1} with arbitrarily
%many bubbles   at given local maximum points of $ q^2\sqrt{\det (K)} $. So the quantity
%\begin{equation*}
%\frac{1}{\varepsilon^2}\int_{\Omega}\left( u_\varepsilon-q|\ln\varepsilon|\right)^{p}_+dx
%\end{equation*}
%can tend to $ +\infty $ as $ \varepsilon\to0 $. This result shows a striking difference with Theorem A in \cite{CW} and related results in \cite{CW1,DDMW2}.
\end{remark}


%\begin{remark}
%In \cite{WYZ}, the authors considered an anisotropic Emden-Fowler equation
%\begin{equation*}
%\begin{cases}
%\text{div}(a(x)\nabla u)+\varepsilon^2a(x)e^u=0& \text{in}\ \Omega,\\
% u=0& \text{on}\ \partial\Omega,
% \end{cases}
%\end{equation*}
%%By using the expansion of Green's function of $ \nabla(a(x)\nabla ) $, the authors prove the existence of clustering solutions.
%where $ a $ is a smooth positive function in $ \Omega $. For any given maximum points $ x^* $ of $ a $,  they constructed clustered 
%solutions concentrating near $ x^*  $. Note that when choosing $ K(x)=a(x)Id $ in \eqref{eq1-1},   results  in Theorem \ref{thm1} 
%coincides with those in \cite{WYZ}. %Thus to some extent, Theorem \ref{thm1} extends results in \cite{WYZ}.
%\end{remark}
Our proof of Theorem \ref{thm0} is based on proper decomposition of the Green's function of elliptic operators in divergence form and the finite-dimensional reduction method. % is based on the Lyapunov-Schmidt reduction method. %Set $ \delta=\varepsilon|\ln\varepsilon|^{-\frac{p-1}{2}} $ and $ u=|\ln\varepsilon|v $, then \eqref{eq1-1} becomes
%\begin{equation}\label{111}
%\begin{cases}
%-\delta^2\text{div}(K(x)\nabla v)=  (v-q)^{p}_+,\ \ &x\in \Omega,\\
%v=0,\ \ &x\in\partial \Omega.
%\end{cases}
%\end{equation}
To get solutions to \eqref{key equa}, we  first  give the $ C^1- $expansion  of Green's function $ G_K $ of the operator $ -\text{div}(K(x)\nabla) $. %, especially the $ C^1- $expansion of $ G_K $, see Lemma \ref{Green expansion2}. 
 Then, we   give ansatz of approximate solutions of \eqref{key equa} of the form $ \sum_{j=1}^N\left (V_{\varepsilon, z_j, \hat{q}_j}+H_{\varepsilon, z_j, \hat{q}_j}\right )+\omega_{\varepsilon,Z} $,  where $ V_{\varepsilon, z_j, \hat{q}_j} $, $ H_{\varepsilon, z_j, \hat{q}_j} $ and $ \omega_{\varepsilon,Z}$ are the main term,  projection term and  error term respectively. The key ingredient is that the configuration space  $ \Lambda_{\varepsilon, N} $ and   parameters $ \hat{q}_j $ must be chosen properly   %Under the choice of $ H_{\delta, z_j, \hat{q}_j} $, we get the  equation \eqref{3-03} for $ \omega_\delta $. 
to ensure that %the distance between $ z_i $ and $ z_j $ is of the $ |\ln \varepsilon|^{-m^2-1} $ order and that the $ L^{\infty} $ norm of
 the difference between $ V_{\varepsilon,Z}-q|\ln\varepsilon| $ and $V_{\varepsilon, z_i, \hat{q}_{i}}-\hat{q}_{i}|\ln\varepsilon| $ is small in some neighborgood of $ z_i $, see \eqref{203}. Then we give the linear and nonlinear theory of the linearized operator of \eqref{key equa} at $ \sum_{j=1}^N\left (V_{\varepsilon, z_j, \hat{q}_j}+H_{\varepsilon, z_j, \hat{q}_j}\right ) $.  %Using the non-degeneracy of some limiting equations, 
We get the existence, uniqueness  and the $ C^1 $ differentiability of $ \omega_{\varepsilon,Z}  $ with respect to $ Z $, see  Proposition \ref{exist and uniq of w} and  
%Note that to prove that the energy functional is $ C^1 $ with respect to the variable $ Z $, we need  the $ C^1 $ differentiability of $ \omega_{\varepsilon,Z}  $ about $ Z $, 
 Proposition \ref{differ of w}. Finally, by calculating   coefficients of each order of energy $ K_\varepsilon(Z) $, that is, the coefficients of $ O(|\ln\varepsilon|), O(\ln|\ln\varepsilon|) $ and $ O(1) $ terms, % and choosing test functions as  the vertices of a $ m $–sided regular polygon,  
we prove the existence of   clustered solutions to \eqref{key equa} with cluster point $ 0 $ and cluster radius $  |\ln\varepsilon|^{-\frac{1}{2}} $.

We give a remark that model \eqref{NP vor fi} has also appeared in many other PDE models, such as Gross-Pitaevskii equations and Ginzburg-Landau equations. \cite{CJ1} established connection between the Ginzburg–Landau functional and the energy of nearly parallel filaments, and constructed  solutions of the Ginzburg–Landau equations  whose small-scale structure is governed by the Ginzburg–Landau functional. \cite{JS2} rigorously
derived the corresponding asymptotic motion law in the context of the Gross–Pitaevskii
equation.  As for the helix,  \cite{DDMR,DDMR2} considered interacting helical traveling waves
for the Gross-Pitaevskii equations and the Ginzburg-Landau equations. For more results, see \cite{Cho, WY}.

%As an application of Theorem \ref{thm1}, for any given $R^*>0$, we get 
%As an direct application of Theorem \ref{thm1}, we can show the existence of clustered helical  rotational-invariant solutions with small cross-section to Euler equations \eqref{Euler eq} in an infinite cylinder $ B_{R^*}(0)\times\mathbb{R} $ for $R^*>0$.    According to    \cite{DDMW2,ET}, if we look for rotating-invariant solution pairs $ (\omega,\varphi) $ to \eqref{vor str eq} with angular velocity $ \alpha|\ln\varepsilon| $ for some $ \alpha\in\mathbb{R} $, %that is, $ (\omega,\varphi) $ satisfies  \eqref{104}. 
%then it suffices to solve solutions $ u $ of a second-order elliptic equations
%\begin{equation*}\label{rot eq3}
%\begin{cases}
%-\text{div}\cdot(K_H(x)\nabla u)=f_\varepsilon\left(u-\frac{\alpha}{2}|x|^2 |\ln\varepsilon|\right),\  &x\in B_{R^*}(0),\\
%u(x)=0,\ &x\in \partial B_{R^*}(0)
%\end{cases}
%\end{equation*}
%for some function $ f_\varepsilon  $, %then $ \omega(x,t)=w\left (\bar{R}_{-\alpha|\ln\varepsilon| t}(x)\right )$ and  
%where $ \varphi(x,t)=u\left (\bar{R}_{-\alpha|\ln\varepsilon| t}(x)\right )  $ and $ f_\varepsilon $ can be chozen arbitrary. % satisfy vorticity equations \eqref{vor str eq}, which corresponds helical solutions to Euler equations \eqref{Euler eq}. 
%By choosing $ f_\varepsilon(t)=\frac{1}{\varepsilon^2}(t-\beta|\ln\varepsilon|)^p_+ $ for   $ \beta\in\mathbb{R} $, we have
%\begin{theorem}\label{thm2}
%Let $ R^*>0$ and $k>0$. %be two given positive numbers. 
%Suppose that $\alpha, \beta\in \mathbb{R}  $  are two numbers such that  $ \min_{x\in B_{R^*}(0)}\left( \frac{\alpha|x|^2}{2}+\beta\right)>0  $ and that $ \left( \frac{\alpha|x|^2}{2}+\beta\right)^2\sqrt{\det (K_H)} $ has a strict local maximum point $ x_0\in B_{R^*}(0)$ up to a rotation, i.e.,   $ |x_0| $ is a strict local maximum  point of $ \left( \frac{\alpha|x|^2}{2}+\beta\right)^2\sqrt{\det (K_H)} $ in $ [0,{R^*}) $.   Then for any $ m\in \mathbb{N}^* $  there exists $ \varepsilon_0>0 $ such that for $ \varepsilon\in(0,\varepsilon_0] $, \eqref{Euler eq} has a family of clustered helical Euler flows $ (\mathbf{v}_\varepsilon, P_\varepsilon)(x,t)\in C^1(B_{R^*}(0)\times \mathbb{R}\times \mathbb{R}^+) $. Moreover, the associated vorticity-stream function pair $ (\omega_\varepsilon,\varphi_\varepsilon) $ is a  rotational-invariant solution to \eqref{vor str eq} with the following properties:
%	\begin{enumerate}
%		\item The angular velocity is $ \alpha|\ln\varepsilon| $ and the circulation satisfies 
%		\begin{equation*}
%	\int_{B_{R^*}(0)} \omega_\varepsilon dx\to \frac{\pi km(\alpha|x_0|^2+2\beta
%		)}{\sqrt{k^2+|x_0|^2}}\ \ \ \ \text{as}\ \  \varepsilon\to 0.
%		\end{equation*}
%		\item There exist $ (z_{1,\varepsilon}, \cdots, z_{m,\varepsilon})\in B_{R^*}(0)^{(m)} $ such that
%		\begin{equation*}
%		|z_{i,\varepsilon}-z_{j,\varepsilon}|\geq |\ln\varepsilon|^{-m^2-1},\ \ \forall i\neq j; \ \ \text{supp}(\omega_\varepsilon)\subseteq \cup_{i=1}^m B_{|\ln\varepsilon|^{-m^2-2}}(z_{i,\varepsilon})
%		\end{equation*}
%		and
%		\begin{equation*}
%		\lim_{\varepsilon\to 0}(z_{1,\varepsilon}, \cdots, z_{m,\varepsilon})=(x_{0}, \cdots, x_{0}).
%		\end{equation*}	
%		\item  There exist constants $ R_1,R_2>0 $ independent of $ \varepsilon $  such that
%		\begin{equation*}
%		B_{R_1\varepsilon}(z_{i,\varepsilon})\subseteq \text{supp}(\omega_\varepsilon)\cap B_{|\ln\varepsilon|^{-m^2-2}}(z_{i,\varepsilon})\subseteq B_{R_2\varepsilon}(z_{i,\varepsilon}).
%		\end{equation*}
%	\end{enumerate}
%\end{theorem}
%
%Let us give an example to show that $ \alpha $ and $ \beta $ satisfying the assumptions of Theorem \ref{thm2}  do exist and correspond to a class of clustered helical vorticity field of Euler equations \eqref{Euler eq}. Actually, by choosing proper $ \alpha $ and $ \beta $, we can show the existence of rotational-invariant clustered helical vorticity fields to \eqref{Euler eq} in infinite cylinders, whose support sets consist of several helical tubes and collapse into the $ x_3- $axis as parameter $ \varepsilon\to 0 $.
%
%%A direct consequence
%%%, by choosing proper constants $ \alpha$ and  $\beta$ in 
%%of Theorem \ref{thm2} is the existence of rotational-invariant clustered helical vorticity solutions to 3D incompressible Euler equation in infinite cylinders, whose support sets consist of several helical tubes and collapse into  $ x_3- $axis as parameter $ \varepsilon\to 0 $.
%\begin{corollary}\label{cor3}
%Let $ R^*>0$ and $k>0$. Suppose that  $ \alpha$ and  $\beta$ are constants such that $ \alpha\leq 0 $ and $ \min_{x\in B_{R^*}(0)}\left( \frac{\alpha|x|^2}{2}+\beta\right)>0 $. Then for any $ m\in \mathbb{N}^* $  there exists $ \varepsilon_0>0 $ such that for $ \varepsilon\in(0,\varepsilon_0] $, \eqref{Euler eq} has a family of clustered helical Euler flows $ (\mathbf{v}_\varepsilon, P_\varepsilon)(x,t)\in C^1(B_{R^*}(0)\times \mathbb{R}\times \mathbb{R}^+) $. Moreover, we have
%\begin{enumerate}
%	\item The associated vorticity field $ \mathbf{w}_\varepsilon=\nabla\times \mathbf{v}_\varepsilon  $ is a  rotational-invariant solution to \eqref{Euler eq2} whose support set consists of $ m $ helical tubes with pitch $ k $ and collapses into the $ x_3- $axis as  $ \varepsilon\to 0 $. 
%	\item The circulation of $ \mathbf{w}_\varepsilon $ tends to $ 2\pi m\beta $ as $ \varepsilon\to0 $.  
%	\item the associated vorticity-stream function pair $ (\omega_\varepsilon,\varphi_\varepsilon) $ is a  rotational-invariant solution to \eqref{vor str eq} with the angular velocity  $ \alpha|\ln\varepsilon| $.
%\end{enumerate}
%
%\end{corollary}
%
%
%%We now give some examples of the choice of $ \alpha$ and  $\beta$ so that the assumptions of Theorem \ref{thm2} are satisfied. 
%%The idea of  proof is as follows. We choose constants $ \alpha$ and  $\beta$ so that $ \alpha<0 $ and $ \min_{x\in B_{R^*}(0)}\left( \frac{\alpha|x|^2}{2}+\beta\right)>0 $ in Theorem \ref{thm2}. Direct computations show that $ (0,0) $ is a strict local maximum point of $ \left( \frac{\alpha|x|^2}{2}+\beta\right)^2\sqrt{\det K_H} $ up to a rotation. From Theorem \ref{thm2}, there exist clustered helical solutions concentrating near $ x_3- $axis. This phenomenon is not found in any existing literatures.
%\begin{remark}
%%Indeed, it is also possible to construct clustered steady helical solutions to Euler equations \eqref{Euler eq} in general  helical domains, see \cite{Ben,CW2}. Moreover, it is interesting whether there exist clustered solutions to \eqref{vor str eq} with different vortex profiles, such as vortex patch solutions.
%In \cite{GM2},  Guerra-Musso considered the existence of clustered helical vorticity fields of \eqref{Euler eq} concentrating near  some helices collapsing into 
%the $ x_3- $axis as  $ \varepsilon\to 0 $. The motion of these helices satisfies a model deduced by Damodaran-Klein-Majda in \cite{DKM}, which is called  ``nearly parallel  vortex filaments'' and the distance of two different helices is of the $ \frac{1}{\sqrt{|\ln\varepsilon|}} $ order. It is interesting whether solutions constructed in \cite{GM2} and that in Corollary \ref{cor3} are the same.
%\end{remark}
%
%\begin{remark}
%In \cite{MR}, Martel and Rapha$\ddot{\text{e}}$l considered the  existence of clustered solutions for the mass critical two dimensional nonlinear Schr\"odinger equation
%\begin{equation}\label{NLS}
%i\partial_t u+\Delta u+|u|^2u=0,\ \ t\in\mathbb{R},\ \ x\in\mathbb{R}^2.
%\end{equation}
%Given any integer $ K\geq 2 $, the authors constructed a global (for $ t > 0 $) $ K $-solitary wave solution $ u(t) $ of \eqref{NLS} that decomposes asymptotically into a sum of
%solitary waves centered at the vertices of a $ K $-sided regular polygon and concentrating at a logarithmic rate as $ t\to\infty$ so that the solution blows up in infinite time with the rate $ ||\nabla u||_{L^2}\sim |\ln t| $ as $ t\to \infty. $ Moreover, such solution concentrates $ K $ bubbles at a point $ x_0\in\mathbb{R}^2 $. In contrast to \cite{MR}, Theorem \ref{thm2} constructed clustered helical rotational-invariant solutions to 3D incompressible Euler equations that decomposes asymptotically into a sum of
%bubbles collapsing to a point $ x_0  $ as $ \varepsilon\to0 $, rather than $ t\to \infty. $ It is interesting whether one can construct clustered helical solutions to 3D incompressible Euler equations which blow up in infinite time and finite time. %To our knowledge, it is also unknown.
%\end{remark}


This paper is organized as follows.  %To construct clustered solutions to \eqref{111},
In section 2, we give the $ C^1- $expansion of Green's function $ G_K $, the definition of   $ \Lambda_{\varepsilon, N} $ and ansatz of approximate solutions. In section 3, we study the linear theory of \eqref{key equa}. Using the non-degeneracy of solutions to limiting equations \eqref{limit eq}, we get coercive estimates of the linearized operator $ Q_\varepsilon L_\varepsilon $. In section 4, the nonlinear theory of \eqref{key equa} is considered. The existence and  the differentiability of the error term $ \omega_{\varepsilon,Z} $   are proved.
In sections 5 and 6, we give the main terms and remaining higher order term of the energy functional $ K_\varepsilon  $ and   complete the proof of Theorem \ref{thm0}. By proper choice of $ K $ and $ q $ in Theorem \ref{thm0}, we give proofs of  Theorems \ref{thm01},\ref{thm02},\ref{thm03},\ref{thm04} and \ref{thm05} in section 7. Some calculations such as  exact helical solutions of \eqref{NP vor fi} and $ C^1-$expansion of Green’s function are given in Appendix.









\section{Green's function and approximate solutions}

To  prove Theorem \ref{thm0}, the key ingredient is to show the existence of clustered solutions to \eqref{key equa}:
\begin{equation*} 
\begin{cases}
-\varepsilon^2\text{div}(K(x)\nabla v)=  \left (v-q(x)|\ln\varepsilon| \right )^{p}_+,\ \ &\text{in}\  \Omega,\\
v=0,\ \ &\text{on}\ \partial \Omega.
\end{cases}
\end{equation*}
Here $ 0<\varepsilon<1, $ $ p>1 $, $ K $ and $ q $ satisfy $(K1)$, $ (Q1) $ and $ (KQ) $. %$ \delta>0 $ sufficiently small, $ \textbf{1}_{D} $ is Heaviside function of some set $ D $, $ C_{\frac{r_*}{\sqrt{|\ln\varepsilon|}}}=\left \{x\mid |x|=\frac{r_*}{\sqrt{|\ln\varepsilon|}}\right \} $ is the circle centered at the origin of radius $ \frac{r_*}{\sqrt{|\ln\varepsilon|}} $, and $ D_{\frac{\delta}{\sqrt{|\ln\varepsilon|}}}\left (C_{\frac{r_*}{\sqrt{|\ln\varepsilon|}}}\right )=\left \{x\mid |x|\in \left (\frac{r_*-\delta}{\sqrt{|\ln\varepsilon|}},\frac{r_*+\delta}{\sqrt{|\ln\varepsilon|}} \right )\right \} $ is the $ \frac{\delta}{\sqrt{|\ln\varepsilon|}} $-neighborhood of $ C_{\frac{r_*}{\sqrt{|\ln\varepsilon|}}} $.  The constants $ \alpha $ and $ \beta $ in \eqref{key equa} are chosen by
%\begin{equation}\label{choice of a,b} 
%\alpha=\frac{\kappa}{4\pi}\left( \frac{1}{h^2}-\frac{N-1}{r_*^2}\right),\ \ \ \ \beta=\frac{\kappa}{2\pi}.
%\end{equation}
The purpose of this section is to give the $ C^1-$expansion of Green's function $ G_{K} $ of the elliptic operator $  -\text{div}(K(x)\nabla \cdot) $ and give the ansatz of  approximate solutions to \eqref{key equa}. %For the sake of simplicity in the notation of the formula, we always denote $ B_{R^*}(0) $ as $ \Omega $, $ K_H $ as $ K $ and  $ \frac{\alpha}{2}|x|^2+\beta $ as $ q(x) $ in the subsequent sections. %the following elliptic equations in divergence form
%\begin{equation}\label{111}
%\begin{cases}
%-\delta^2\text{div}(K(x)\nabla w)=  (w-q)^{p}_+,\ \ &\text{in}\  \Omega,\\
%w=0,\ \ &\text{on}\ \partial \Omega.
%\end{cases}
%\end{equation}
%and give some basic estimates for the approximate  solutions. %Note that $ -\text{div}(K(x)\nabla \cdot) $ is a uniformly elliptic operator.
%Throughout this paper, we denote $ C,C_1,C_2\cdots $ positive constants independent of $ \varepsilon $, whose values may change from line to line.



\subsection{$ C^1-$Expansion of Green's function $ G_K $}
The expansion of Green's function $ G_K $ of the elliptic operator $ -\text{div}(K(x)\nabla \cdot) $ with  Dirichlet condition  plays an essential role, since it appears in the construction of approximate solutions of \eqref{key equa} and the $ C^1-$regularity of the finite-dimensional functional.
%We first give some properties of the Green's function $ G_K $ of the operator $ -\text{div}(K(x)\nabla \cdot) $ with Dirichlet condition.
Let $ G_K(x,y) $ be the Green's function of $ -\text{div}(K(x)\nabla \cdot) $ with zero Dirichlet condition on  the boundary of $ \Omega $, that is, solutions of the following linear elliptic problem:
\begin{equation}\label{diver form}
\begin{cases}
-\text{div}(K(x)\nabla u)= f,\ \ & \Omega,\\
u=0,\ \ & \partial \Omega
\end{cases}
\end{equation}
can be expressed by $ u(x)=\int_{\Omega}G_K(x,y)f(y)dy $ for $ x\in \Omega $, see \cite{GT}.
%Since $ K $ is a $ C^\infty $ positive definite matrix with all eigenvalues having uniformly positive lower and upper bounds, by the Cholesky decomposition \textcolor{red}{(see \cite{G})} one can find the unique $ T\in C^\infty(\overline{\Omega}) $ such that for any $ x\in \Omega $, $ T(x) $ is invertible and
%Let $ T $ be the unique positive-definite matrix-valued function satisfying
%\begin{equation}\label{T_z choice}
%(T(x)^{-1})(T(x)^{-1})^t=K(x),\ \ \ \ x\in \Omega.
%\end{equation}
%For simplicity, we denote $ T_{x}=T(x) $. Define $ \Gamma(x)=-\frac{1}{2\pi}\ln|x|$ the fundamental solutions of $ -\Delta $ in $ \mathbb{R}^2 $.

%From lemma 3.1 in \cite{CW2}  (see also Lemma \ref{A1} in Appendix), we know that there exists a function $ S_K\in C^{0,\gamma}_{loc} (\Omega\times \Omega) $ for any $ \gamma\in(0,1) $, such that for every $f \in L^{p}(\Omega)$  with $ p>2 $ and every $ x\in \Omega $, the solution of \eqref{diver form} can be expressed by
%\begin{equation*}
%\begin{split}
%u(x)=\int_{\Omega}\left (\frac{\sqrt{\det K(x)}^{-1}+\sqrt{\det K(y)}^{-1}}{2}\Gamma\left (\frac{T_x+T_y}{2}(x-y)\right )+S_K(x,y)\right)f(y)dy.
%\end{split}
%\end{equation*}
%Moreover, $ S_K(x,y)=S_K(y,x) $  for any  $ x,y \in \Omega. $ In other words, the Green's function $ G_K $ of the operator $ -\text{div}(K(x)\nabla  ) $ with Dirichlet condition has the decomposition
%\begin{equation}\label{2-11}
%G_K(x,y)=\frac{\sqrt{\det K(x)}^{-1}+\sqrt{\det K(y)}^{-1}}{2}\Gamma\left (\frac{T_x+T_y}{2}(x-y)\right )+S_K(x,y).
%\end{equation}
  Note that \cite{CW,CW1}  obtained the following $ C^0-$asymptotic expansion of $ G_K(x,y)$. %We need to use $ C^1 $-asymptotic expansion of $ G_K(x,y) $ in the present paper.
\begin{lemma}[lemma 3.1, \cite{CW}]\label{Green expansion}
For $ y\in\Omega $, let $ T_y=T(y) $ be the unique positive-definite matrix satisfying
$$ T_y^{-1}(T_y^{-1})^{T}=K(y). $$ 
Then for any $ \gamma\in(0,1) $, there exists a function $\bar{S}_K\in C^{0,\gamma}_{loc} (\Omega\times \Omega) $ such that
\begin{equation*}
G_K(x,y)=\sqrt{\det K(y)}^{-1}\Gamma\left (T_y(x-y)\right )+\bar{S}_K(x,y),\ \ \forall\ x,y\in\Omega.
\end{equation*}
\end{lemma}
For $ i,j=1,2 $, we denote $ T_{ij}=(T_y)_{ij}=(T(y))_{ij} $ the component of row $ i $, column $ j $ of the matrix $ T_y $. Based on Lemma \ref{Green expansion} and some more detailed calculations, we give $ C^1- $asymptotic expansion of $ G_K(x,y) $ as follows, the proof of which will be given in the Appendix (see Appendix A.2). %, which improves results in Lemma \ref{Green expansion}.
\begin{lemma}\label{Green expansion2}
%There exists a function $\bar{S}_K\in C^{0,\gamma}_{loc} (\Omega\times \Omega) $ for any $ \gamma\in(0,1) $, such that
There holds
\begin{equation*}
\bar{S}_K(x,y)=-F_{1,y}(x)-F_{2,y}(x)+\bar{H}_1(x,y) \ \ \ \ \forall\ x,y\in\Omega,
\end{equation*}
where
\begin{equation}\label{exp of F_{1,y}}
\begin{split}
F_{1,y}(x)=-\frac{1}{4\pi}\sqrt{\det K(y)}^{-1}\sum_{i,j,m=1}^2T_{mj}\partial_{x_i}K_{ij}(y)\left( T_y(x-y)\right)_m\ln|T_y(x-y)|,
\end{split}
\end{equation}
\begin{equation*}
\begin{split}
F_{2,y}(x)=&\frac{1}{\pi}\sqrt{\det K(y)}^{-1}\sum_{i,j,\alpha=1}^2\partial_{x_\alpha}K_{ij}(y)\cdot\\
&\bigg\{T^{-1}_{\alpha1}T_{1j}T_{1i}\left( -\frac{1}{8}\frac{\left( T_y\left (x-y\right )\right) _1^3}{|T_y\left (x-y\right )|^2}+\frac{1}{8}\left( T_y\left (x-y\right )\right)_1\ln|T_y\left (x-y\right )|\right) \\
+&T^{-1}_{\alpha1}T_{1j}T_{2i}\left( -\frac{1}{8}\frac{\left( T_y\left (x-y\right )\right) _1^2\left( T_y\left (x-y\right )\right) _2}{|T_y\left (x-y\right )|^2}+\frac{1}{8}\left( T_y\left (x-y\right )\right)_2\ln|T_y\left (x-y\right )|\right) \\
+&T^{-1}_{\alpha1}T_{2j}T_{1i}\left( -\frac{1}{8}\frac{\left( T_y\left (x-y\right )\right) _1^2\left( T_y\left (x-y\right )\right) _2}{|T_y\left (x-y\right )|^2}+\frac{1}{8}\left( T_y\left (x-y\right )\right)_2\ln|T_y\left (x-y\right )|\right)\\
+&T^{-1}_{\alpha1}T_{2j}T_{2i}\left( -\frac{1}{8}\frac{\left( T_y\left (x-y\right )\right) _1\left( T_y\left (x-y\right )\right) _2^2}{|T_y\left (x-y\right )|^2}-\frac{1}{8}\left( T_y\left (x-y\right )\right)_1\ln|T_y\left (x-y\right )|\right)
\end{split}
\end{equation*}
\begin{equation}\label{exp of F_{2,y}}
\begin{split}
+&T^{-1}_{\alpha2}T_{1j}T_{1i}\left( -\frac{1}{8}\frac{\left( T_y\left (x-y\right )\right) _1^2\left( T_y\left (x-y\right )\right) _2}{|T_y\left (x-y\right )|^2}-\frac{1}{8}\left( T_y\left (x-y\right )\right)_2\ln|T_y\left (x-y\right )|\right)\\
+&T^{-1}_{\alpha2}T_{1j}T_{2i}\left( -\frac{1}{8}\frac{\left( T_y\left (x-y\right )\right) _1\left( T_y\left (x-y\right )\right) _2^2}{|T_y\left (x-y\right )|^2}+\frac{1}{8}\left( T_y\left (x-y\right )\right)_1\ln|T_y\left (x-y\right )|\right)\\
+&T^{-1}_{\alpha2}T_{2j}T_{1i}\left( -\frac{1}{8}\frac{\left( T_y\left (x-y\right )\right) _1\left( T_y\left (x-y\right )\right) _2^2}{|T_y\left (x-y\right )|^2}+\frac{1}{8}\left( T_y\left (x-y\right )\right)_1\ln|T_y\left (x-y\right )|\right)\\
+&T^{-1}_{\alpha2}T_{2j}T_{2i}\left( -\frac{1}{8}\frac{\left( T_y\left (x-y\right )\right) _2^3 }{|T_y\left (x-y\right )|^2}+\frac{1}{8}\left( T_y\left (x-y\right )\right)_2\ln|T_y\left (x-y\right )|\right)\bigg\},
\end{split}
\end{equation}
 and $ x\to \bar{H}_1(x,y)\in C^{1,\gamma}(\overline{\Omega})$ for all $ y\in \Omega $, $ \gamma\in (0, 1) $.  Moreover, the function $(x, y)\to \bar{H}_1(x,y)\in C^1(\Omega\times\Omega),$ and in particular the corresponding Robin function $x\to \bar{S}_K(x,x)\in C^1(\Omega)$.
\end{lemma}


\begin{remark}
From Lemma \ref{Green expansion2}, Green's function $ G_K $ has an expansion
\begin{equation*} 
G_K(x,y)=\sqrt{\det K(y)}^{-1}\Gamma\left (T_y(x-y)\right )-F_{1,y}(x)-F_{2,y}(x)+\bar{H}_1(x,y),\ \ \forall\ x,y\in\Omega.
\end{equation*}
%Indeed, Lemma \ref{Green expansion2} holds not only for $ \Omega=B_{R^*}(0), K=K_H $, but also for $ \Omega $ being any bounded domian and $ K $ being any positive-definite matrix. %We give some examples to understand this result. to explain results in Lemma \ref{Green expansion2}.\\

\iffalse
\noindent$ Example $ 1. If $ K(x)= Id $, then \eqref{diver form} is the standard Laplacian problem. In this case, one computes directly that $ F_{1,y}=F_{2,y}\equiv 0 $. From \eqref{exp of G_K},  Green's function has an expansion
\begin{equation*}
G_1(x,y)=\Gamma(x-y)+S_1(x,y),\ \ \ \ \forall\ x,y\in \Omega.
\end{equation*}
Thus we have $ S_1(x,y)=-H(x,y) $, where $ H(x,y) $ is the regular part of Green's function of $ -\Delta $ in $ \Omega $ with zero-Dirichlet data, which coincides with classical results in \cite{GT}.

\noindent$ Example $ 2. If $ K(x)=\frac{1}{b(x)}Id $, where $ b\in C^2(\overline{\Omega}) $ and $ \inf_{\Omega}b>0 $, then $ \det K=\frac{1}{b^2} $ and $ T=\sqrt{b}Id $. By \eqref{exp of F_{1,y}} and \eqref{exp of F_{2,y}} it is not hard to get that $ F_{1,y}(x)= \frac{\nabla b(y)\cdot (x-y)}{4\pi}\ln|x-y|+F^*_1(x,y) $ and $ F_{2,y}(x)=F^*_2(x,y) $ for $ x,y\in\Omega $, where $ F^*_1,F^*_2$ are two functions belonging to $ C^{1}(\Omega\times\Omega) $. From  \eqref{exp of G_K},  Green's function has the expansion
\begin{equation*}
G_b(x,y)=b(y)\Gamma\left (x-y\right )-\frac{\nabla b(y)\cdot (x-y)}{4\pi}\ln|x-y|+S_b(x,y) \ \ \ \ \forall\ x,y\in \Omega,
\end{equation*}
where $ S_b\in C^1(\Omega\times\Omega) $, which coincides with results in \cite{De2,WYZ}.

\fi

It is also noted that recently, \cite{CW3} constructed the asymptotic expansion of Green's function and the regularity of Robin's function of an elliptic operator in divergence form in bounded domains in all dimensions. Especially when the dimension is 2 and the coefficient matrix is $ K_H $, they proved that $ \bar{S}_K $ is smooth in $ \Omega. $ In contrast to \cite{CW3}, we  calculate  in detail the specific expansion of  Green's function.
\end{remark}
%\begin{proof}
%%Since the subsequent construction of concentrated  solutions only involves this property, 
%We place  the proof of this result in Proposition \ref{lemA2} in Appendix. 
%\end{proof}


\subsection{Ansatz for approximate solutions}

The argument that  $ \bar{S}_K(x,x)$ is a $ C^1 $  function  plays an important role in our argument to get the $ C^1 $ regularity
of the error term and the finite-dimensional variational functional with respect to parameter $ Z $, see sections 4 and 5.


We give approximate solutions of \eqref{key equa} and   the configuration space $ \Lambda_{\varepsilon, N} $ for the parameter $ Z=(z_1,\cdots,z_N) $  as follows. Let $ N> 1 $ be an integer.  For any $ \hat{x}\in \Omega, \hat{q}>0 $, we define
\begin{equation}\label{eq2-2}
V_{\varepsilon, \hat{x}, \hat{q}}(x)=\begin{cases}
\hat{q}\ln \frac{1}{\varepsilon}+\varepsilon^{\frac{2}{p-1}}s_\varepsilon^{-\frac{2}{p-1}}\phi\left(\frac{|T_{\hat{x}}(x-\hat{x})|}{s_\varepsilon}\right),\ \ &|T_{\hat{x}}(x-\hat{x})|\le s_\varepsilon,\\
 \hat{q} \ln \frac{1}{\varepsilon}\frac{\ln |T_{\hat{x}}(x-\hat{x})|}{\ln s_\varepsilon},\ \  &  |T_{\hat{x}}(x-\hat{x})|>s_\varepsilon,
\end{cases}
\end{equation}
where $ T_{\hat{x}}^{-1}\left (T_{\hat{x}}^{-1}\right )^{T}=K\left (\hat{x}\right ) $, $ \phi\in H^1_0(B_1(0)) $ satisfies (see, e.g., \cite{CLW})
\begin{equation*}
-\Delta\phi=\phi^p, \ \ \phi>0\ \ \text{in}\ B_1(0),
\end{equation*}
and $ s_\varepsilon $ satisfies
\begin{equation}\label{201}
\varepsilon^{\frac{2}{p-1}}s_\varepsilon^{-\frac{2}{p-1}}\phi'(1)=\hat{q}\frac{ \ln\frac{1}{\varepsilon}}{\ln s_\varepsilon}.
\end{equation}
Clearly, $ V_{\varepsilon, \hat{x}, \hat{q}}\in C^1 $ satisfies
\begin{equation}\label{eq2}
\begin{cases}
-\varepsilon^2\text{div}\left (K(\hat{x})\nabla V_{\varepsilon, \hat{x}, \hat{q}}\right )=  \left (V_{\varepsilon, \hat{x}, \hat{q}}-\hat{q}\ln\frac{1}{\varepsilon}\right )^{p}_+,\ \ & \text{in}\  \mathbb{R}^2,\\
V_{\varepsilon, \hat{x}, \hat{q}}=\hat{q}\ln\frac{1}{\varepsilon},\ \ &\text{on} \ \{x\mid |T_{\hat{x}}(x-\hat{x})|=s_\varepsilon\},
\end{cases}
\end{equation}
and  for $ \varepsilon $ sufficiently small, \eqref{201} is uniquely solvable with
\begin{equation*}
\frac{s_\varepsilon}{\varepsilon}\to \left( \frac{|\phi'(1)|}{\hat{q}}\right) ^{\frac{p-1}{2}}\ \ \ \ \text{as}\ \varepsilon\to0.
\end{equation*}
The Pohozaev identity implies
\begin{equation}\label{PI}
\int_{B_1(0)}\phi^{p+1}=\frac{\pi(p+1)}{2}|\phi'(1)|^2,\ \ \int_{B_1(0)}\phi^{p}= 2\pi|\phi'(1)|.
\end{equation}


Since $ V_{\varepsilon, \hat{x}, \hat{q}} $ is not 0 on $ \partial \Omega $, we need to make a projection on $ H^1_0(\Omega). $ %Note that the operator $ \text{div}(K(\hat{x})\nabla \cdot) $ in \eqref{eq2} is different from $ \text{div}(K(x)\nabla \cdot) $ appeared in \eqref{111}, we introduce a projection term 
Let $ H_{\varepsilon, \hat{x}, \hat{q}} $ be a solution of the following equations
\begin{equation}\label{eq3}
\begin{cases}
-\text{div}\left (K(x)\nabla H_{\varepsilon, \hat{x}, \hat{q}}\right )=\text{div}\left (\left (K(x)-K(\hat{x})\right )\nabla V_{\varepsilon, \hat{x}, \hat{q}}\right ),\ \ &\Omega,\\
H_{\varepsilon, \hat{x}, \hat{q}}=-V_{\varepsilon, \hat{x}, \hat{q}},\ \ &\partial\Omega.
\end{cases}
\end{equation}
Then  $ H_{\varepsilon, \hat{x}, \hat{q}}\in W^{2,p}(\Omega)\subset  C^{1,\alpha}(\overline{\Omega}) $ for any $ p>1, \alpha\in(0,1) $. Similar to lemma 3.2 in \cite{CW}, we give estimates of the difference between $ H_{\varepsilon, \hat{x}, \hat{q}} $ and  $ -\frac{2\pi\hat{q}\sqrt{\det K(\hat{x})}\ln\frac{1}{\varepsilon}}{\ln s_\varepsilon}\bar{S}_K(\cdot,\hat{x}) $.

\begin{lemma}\label{H estimate}
Define $ \zeta_{\varepsilon, \hat{x}, \hat{q}}(x)=H_{\varepsilon, \hat{x}, \hat{q}}(x)+\frac{2\pi\hat{q}\sqrt{\det K(\hat{x})}\ln\frac{1}{\varepsilon}}{\ln s_\varepsilon}\bar{S}_K(x,\hat{x}) $ for $ x\in\Omega $. Then for any $ p\in(1,2) $, there exists a constant $ C>0 $ independent of $ \varepsilon $ such that
\begin{equation*}
||\zeta_{\varepsilon, \hat{x}, \hat{q}}||_{C^{0,2-\frac{2}{p}}(\Omega)}\leq Cs_\varepsilon^{\frac{2}{p}-1}.
\end{equation*}
\end{lemma}
%\begin{proof}
Since the idea of proof is similar to that in lemma 3.2 in \cite{CW}, we just give a sketch of   proof for the above lemma   in the Appendix (see Appendix A.3).
%\end{proof}


Using \eqref{eq2-2}, the definition of $ H_{\varepsilon, \hat{x}, \hat{q}} $ in \eqref{eq3} and the classical $ L^p $-theory of elliptic equations, it is also not hard to get 
\begin{equation}\label{3-001}
||H_{\varepsilon, \hat{x}, \hat{q}}||_{W^{2,p}(\Omega)}\leq
\begin{cases}
\frac{C}{\varepsilon^{1-\frac{2}{p}}},\ \ &p>2,\\
C|\ln\varepsilon|,\ \ &p=2,\\
C,\ \ &1\leq p<2.
\end{cases}
\end{equation}
%\begin{equation}\label{3-002}
%||\frac{\partial H_{\delta, \hat{x}, \hat{q}}}{\partial \hat{x}_h}||_{W^{1,p}(\Omega)}\leq
%\begin{cases}
%\frac{C}{\varepsilon^{1-\frac{2}{p}}|\ln\varepsilon|},\ \ &p>2,\\
%C,\ \ &p=2,\\
%\frac{C}{|\ln\varepsilon|},\ \ &1\leq p<2.
%\end{cases}
%\end{equation}


%Let $ Z\in\Omega^{(N)} $. Since $ x_{0} $ is a strict local maximum point  of $ q^2\sqrt{\det(K)} $ in $ \Omega $,   
Let $ \left( \tilde{z}_1^*,\cdots, \tilde{z}_N^*\right)  $ be a local strict maximizer or minimizer of $ H_N $. The configuration space $ \Lambda_{\varepsilon, N} $ for $ Z=(z_1,\cdots,z_N) $ is defined as follows:%. . Hence we can choose $ \bar{\rho}>0 $ sufficiently small such that $ B_{\bar{\rho}}(x_{0})\Subset\Omega   $ and
%\begin{equation*}
%\max_{\partial B_{\bar{\rho}}(x_{0})}q^2\sqrt{\det(K)}< \max_{  B_{\bar{\rho}}(x_{0})}q^2\sqrt{\det(K)}=q^2\sqrt{\det(K)}(x_0).
%\end{equation*}
%The admissible class  we choose is as follows
\begin{equation}\label{admis set}
\begin{split}
\Lambda_{\varepsilon, N}=\bigg\{&Z=(z_1,\cdots, z_N)\in\Omega^{(N)}\mid z_i\in B_{\frac{\rho_0}{2\sqrt{|\ln\varepsilon|}}}\left ( \frac{\tilde{z}_i^*}{\sqrt{|\ln\varepsilon|}}\right )\bigg\}.%\\
%& \min_{i\neq j}|z_i-z_j|\geq |\ln\varepsilon|^{-1}, \ \ \  \forall\  1\leq i\neq j\leq N.
\end{split}
\end{equation}
%where $ M=m^2+1. $ 
Here $ \rho_0>0 $ is small. Note that by \eqref{admis set},
\begin{equation}\label{3-003}
G_K(z_i,z_j)\leq C|\ln|z_i-z_j||\leq C\ln|\ln\varepsilon|,\ \ \ \ \forall\  Z\in \Lambda_{\varepsilon, N}.
\end{equation}



In the following, we will construct solutions of \eqref{key equa} being of the form
\begin{equation}\label{solu config}
v_\varepsilon= V_{\varepsilon, Z}+\omega_\varepsilon=\sum_{j=1}^NV_{\varepsilon, Z,j}+\omega_{\varepsilon,Z}=\sum_{j=1}^N\left (V_{\varepsilon, z_j, \hat{q}_j}+H_{\varepsilon, z_j, \hat{q}_j}\right )+\omega_{\varepsilon,Z},
\end{equation}
where $ Z\in \Lambda_{\varepsilon, N} $, $ \sum_{j=1}^NV_{\varepsilon, Z,j} $ is the main term and $ \omega_{\varepsilon,Z}$ is an error term. The choice of $ \hat{q}_j $ will be determined later on. First from \eqref{key equa}, one computes directly that
\begin{equation*}
\begin{split}
0=&\sum_{j=1}^N-\varepsilon^2\text{div}(K(x)\nabla (V_{\varepsilon, z_j, \hat{q}_j}+H_{\varepsilon, z_j, \hat{q}_j}))-\varepsilon^2\text{div}(K(x)\nabla \omega_{\varepsilon,Z})-\left (V_{\varepsilon, Z}+\omega_{\varepsilon,Z}-q|\ln\varepsilon|\right )^{p}_+\\
=&-\sum_{j=1}^N\varepsilon^2\text{div}(K(z_j)\nabla V_{\varepsilon, z_j, \hat{q}_j})-\sum_{j=1}^N\varepsilon^2\text{div}((K(x)-K(z_j))\nabla V_{\varepsilon, z_j, \hat{q}_j})\\
&-\sum_{j=1}^N\varepsilon^2\text{div}(K(x)\nabla H_{\varepsilon, z_j, \hat{q}_j})+\left (-\varepsilon^2\text{div}(K(x)\nabla \omega_{\varepsilon,Z})-p\left ( V_{\varepsilon, Z}-q|\ln\varepsilon|\right )^{p-1}_+\omega_{\varepsilon,Z} \right )\\
&-\left (\left (V_{\varepsilon, Z}+\omega_{\varepsilon,Z}-q|\ln\varepsilon|\right )^{p}_+-p\left (V_{\varepsilon, Z}-q|\ln\varepsilon|\right )^{p-1}_+\omega_{\varepsilon,Z}\right ).
%=&\sum_{j=1}^m(V_{\delta, z_j, \hat{q}_j}-\hat{q}_j)^{p}_+-\left (\sum_{j=1}^mV_{\delta, Z,j}-q\right )^{p}_+\\
%&+\left (-\delta^2\text{div}(K(x)\nabla \omega_\delta)-p\left (\sum_{j=1}^mV_{\delta, Z,j}-q\right )^{p-1}_+\omega_\delta \right )\\
%&-\left (\left (\sum_{j=1}^mV_{\delta, Z,j}+\omega_\delta-q\right )^{p}_+-\left (\sum_{j=1}^mV_{\delta, Z,j}-q\right )^{p}_+-p\left (\sum_{j=1}^mV_{\delta, Z,j}-q\right )^{p-1}_+\omega_\delta\right )\\
%=&-l_\varepsilon+L_\varepsilon \omega_{\varepsilon,Z}- R_\varepsilon(\omega_{\varepsilon,Z}),
\end{split}
\end{equation*}
Thus it suffices to find solutions  $ \omega=\omega_{\varepsilon,Z} $ of the following equation 
\begin{equation}\label{3-03}
L_\varepsilon \omega=l_\varepsilon+ R_\varepsilon(\omega),
\end{equation}
where
\begin{equation*}
l_\varepsilon:=\left (\sum_{j=1}^NV_{\varepsilon, Z,j}-q|\ln\varepsilon|\right )^{p}_+-\sum_{j=1}^N\left (V_{\varepsilon, z_j, \hat{q}_j}-\hat{q}_j|\ln\varepsilon|\right )^{p}_+,
\end{equation*}
$ L_\varepsilon $ is the linearized operator of \eqref{key equa} at $  V_{\varepsilon, Z} $ defined by
\begin{equation*}
L_\varepsilon \omega:=-\varepsilon^2\text{div}(K(x)\nabla \omega)- p\left (\sum_{j=1}^NV_{\varepsilon, Z,j}-q|\ln\varepsilon|\right )^{p-1}_+\omega,
\end{equation*}
and $ R_\varepsilon(\omega_{\varepsilon,Z}) $ is the high-order error term defined by
\begin{equation*}
\begin{split}
R_\varepsilon(\omega):=&\left (V_{\varepsilon, Z}+\omega-q|\ln\varepsilon|\right )^{p}_+-\left (V_{\varepsilon, Z}-q|\ln\varepsilon|\right )^{p}_+-p\left (V_{\varepsilon, Z}-q|\ln\varepsilon|\right )^{p-1}_+\omega.
\end{split}
\end{equation*}



%To make the  norm of $ \omega_\delta $ as small as possible, we choose  $ \hat{q}_j $ suitably close to $ q(z_j) $.
Next we will choose parameters $ \hat{q}_j $ suitably to make the term $ l_\varepsilon $ as small as possible in some norm. For any $ Z\in \Lambda_{\varepsilon, N}, $ let $ \hat{q}_i=\hat{q}_{\varepsilon,i}(Z) $   satisfy for $ i=1,\cdots,N $
\begin{equation}\label{q_i choice}
\hat{q}_i=q(z_i)+\frac{2\pi\hat{q}_i\sqrt{\det K(z_i)}}{\ln s_{\varepsilon,i}}\bar{S}_K(z_i, z_i)+\sum_{j\neq i}\frac{2\pi\hat{q}_j\sqrt{\det K(z_j)}}{\ln s_{\varepsilon,j}}G_K(z_i, z_j),
\end{equation}
where $ s_{\varepsilon,i} $ satisfies for $ i=1,\cdots,N $
\begin{equation*}
\varepsilon^{\frac{2}{p-1}}s_{\varepsilon,i}^{-\frac{2}{p-1}}\phi'(1)=\hat{q}_i\frac{ \ln\frac{1}{\varepsilon}}{\ln s_{\varepsilon,i}}.
\end{equation*}

Then from \cite{Ku}, %the Poincar$\acute{\text{e}}$ -- Miranda Theorem (see \cite{Ku}), 
for any $ \varepsilon $ sufficiently small  and $ Z\in \Lambda_{\varepsilon, N} $, there exists $ \hat{q}_{\varepsilon,i}(Z) $ satisfying \eqref{q_i choice}. Moreover, $ \hat{q}_i>0 $ and %one computes directly that for $ i=1,\cdots,m $
%\begin{equation}\label{1999}
%\hat{q}_i=q(z_i)+O\left (\frac{\ln|\ln\varepsilon|}{|\ln\varepsilon|}\right ),
%\end{equation}
%\begin{equation}\label{1999-1}
%\frac{\partial \hat{q}_i}{\partial z_{j,h}}=\frac{\partial q}{\partial x_{h}}(z_i)\delta_{i,j}+O\left (|\ln\varepsilon|^{M}\right ),
%\end{equation}
\begin{equation}\label{2000}
\hat{q}_i=q(z_i)+O\left (\frac{\ln|\ln\varepsilon|}{|\ln\varepsilon|}\right );\ \ \frac{1}{\ln\frac{1}{s_{\varepsilon,i}}}=\frac{1}{\ln\frac{1}{\varepsilon}}+O\left( \frac{\ln|\ln\varepsilon|}{|\ln\varepsilon|^2}\right).
\end{equation}


By the choice of $ \Lambda_{\varepsilon, N} $ in \eqref{admis set} and  $ \hat{q}_{\varepsilon,j} $ in \eqref{q_i choice}, we claim that for any $ Z\in \Lambda_{\varepsilon, N} $, $ \gamma\in(0,1), L>1 $ and $ x\in B_{Ls_{\varepsilon,i}}(z_i) $
\begin{equation}\label{203}
\begin{split}
\sum_{j=1}^NV_{\varepsilon, Z,j}(x)-q(x)|\ln\varepsilon|=V_{\varepsilon, z_i, \hat{q}_{\varepsilon,i}}(x)-\hat{q}_{\varepsilon,i}|\ln\varepsilon|+O\left( \varepsilon^\gamma\right).
\end{split}
\end{equation}
Indeed, for $ x\in B_{Ls_{\varepsilon,i}}(z_i) $
\begin{equation*}
\begin{split}
V_{\varepsilon,Z,i}(x)&-q(x)|\ln\varepsilon|\\
=&V_{\varepsilon, z_i, \hat{q}_{\varepsilon,i}}(x)+H_{\varepsilon, z_i, \hat{q}_{\varepsilon,i}}(x)-q(x)|\ln\varepsilon|\\
=&V_{\varepsilon, z_i, \hat{q}_{\varepsilon,i}}(x)-q(z_i)|\ln\varepsilon|-\frac{2\pi\hat{q}_{\varepsilon,i}\sqrt{\det K(z_i)}|\ln\varepsilon|}{\ln s_{\varepsilon,i}}\bar{S}_K(x, z_i)+O(s_{\varepsilon,i}|\ln\varepsilon|)+O\left( s_{\varepsilon,i}^{\gamma}\right) \\
=&V_{\varepsilon, z_i, \hat{q}_{\varepsilon,i}}(x)-q(z_i)|\ln\varepsilon|-\frac{2\pi\hat{q}_{\varepsilon,i}\sqrt{\det K(z_i)}|\ln\varepsilon|}{\ln s_{\varepsilon,i}}\bar{S}_K(z_i, z_i)+ O\left( \varepsilon^\gamma\right).
\end{split}
\end{equation*}
For any $ j\neq i $, since $ |\ln\varepsilon|^{-1}>2Ls_{\varepsilon,i} $, for $ \varepsilon $ sufficiently small  one has
\begin{equation*}
\begin{split}
V_{\varepsilon,Z,j}(x)=&V_{\varepsilon, z_j, \hat{q}_{\varepsilon,j}}(x)+H_{\varepsilon, z_j, \hat{q}_{\varepsilon,j}}(x)\\
=&\frac{\hat{q}_{\varepsilon,j}|\ln\varepsilon|}{\ln s_{\varepsilon,j}}\ln |T_{z_j}(x-z_j)|-\frac{2\pi\hat{q}_{\varepsilon,j}\sqrt{\det K(z_j)}|\ln\varepsilon|}{\ln s_{\varepsilon,j}}\bar{S}_K(x, z_j)+O\left(  s_{\varepsilon,j}^{\gamma}\right)\\
=&-\frac{2\pi\hat{q}_{\varepsilon,j}\sqrt{\det K(z_j)}|\ln\varepsilon|}{\ln s_{\varepsilon,j}}G_K(x, z_j)+O\left(  s_{\varepsilon,j}^{\gamma}\right) \\
=&-\frac{2\pi\hat{q}_{\varepsilon,j}\sqrt{\det K(z_j)}|\ln\varepsilon|}{\ln s_{\varepsilon,j}}G_K(z_i, z_j)+O\left( \varepsilon^\gamma\right),
\end{split}
\end{equation*}
where we have used  Lemma \ref{H estimate} and the fact that  for $ x\in B_{Ls_{\varepsilon,i}}(z_i) $
\begin{equation*}
G_K(x, z_j)=G_K(z_i, z_j)+O(|\nabla_{z_i}G_K(z_i, z_j)(x-z_j)|)=G_K(z_i, z_j)+O(\varepsilon|\ln\varepsilon|).
\end{equation*}
Adding up the above inequalities and using \eqref{q_i choice}, we get \eqref{203}.

Using the definition of $ V_{\varepsilon, z_i, \hat{q}_{\varepsilon,i}} $, we obtain

\begin{equation}\label{200}
\begin{split}
\frac{\partial V_{\varepsilon, z_i, \hat{q}_{\varepsilon,i}}(x)}{\partial x_\hbar}=\begin{cases}
\frac{1}{s_{\varepsilon,i}}(\frac{\varepsilon}{s_{\varepsilon,i}})^{\frac{2}{p-1}}\phi'(\frac{|T_{z_i}(x-z_i)|}{s_{\varepsilon,i}})\frac{(T_{z_i})_\hbar^T\cdot T_{z_i}(x-z_i)}{|T_{z_i}(x-z_i)|},\ \ &|T_{z_i}(x-z_i)|\leq s_{\varepsilon,i},\\
\frac{\hat{q}_{\varepsilon,i}|\ln\varepsilon|}{\ln s_{\varepsilon,i}}\frac{(T_{z_i})_\hbar^T\cdot T_{z_i}(x-z_i)}{|T_{z_i}(x-z_i)|^2},\ \ &|T_{z_i}(x-z_i)|> s_{\varepsilon,i},
\end{cases}
\end{split}
\end{equation}
for $ \hbar=1,2, $ where $ (T_{z_i})_\hbar^T $ is the $\hbar$-th row of $ (T_{z_i})^T $.

We end this section by giving %by using \eqref{203},
 some estimates of approximate solutions $ V_{\varepsilon,Z} $, which will be frequently used in the subsequent sections.
\begin{lemma}\label{lemA-5}
	Let $\gamma\in(0,1)$ and $ Z\in \Lambda_{\varepsilon, N} $. Then there exists a constant $ L > 1 $ such that for $ \varepsilon $ small
	\begin{equation*}
	V_{\varepsilon,Z}-q|\ln\varepsilon|>0,\ \ \  \ \text{in}\ \   \cup_{j=1}^N\left( T_{z_j}^{-1}B_{\left( 1-L \varepsilon^\gamma \right) s_{\varepsilon,j}}(0)+z_j\right),
	\end{equation*}
	\begin{equation*}
	V_{\varepsilon,Z}-q|\ln\varepsilon|<0,\ \ \  \ \text{in}\ \  \Omega\backslash\cup_{j=1}^N\left( T_{z_j}^{-1}B_{Ls_{\varepsilon,j}}(0)+z_j\right).
	\end{equation*}
	
\end{lemma}

\begin{proof}
If $ |T_{z_j}(x-z_j)|\leq \left( 1-L \varepsilon^\gamma \right) s_{\varepsilon,j}   $, then by \eqref{203} and $ \phi'(1)<0 $ we have
	\begin{equation*}
	\begin{split}
	V_{\varepsilon,Z}(x)-q(x)|\ln\varepsilon|=&V_{\varepsilon,z_j,\hat{q}_{\varepsilon,j}}(x)-\hat{q}_{\varepsilon,j}|\ln\varepsilon|+O\left(  \varepsilon^\gamma\right) \\
	= &\frac{\hat{q}_{\varepsilon,j}|\ln\varepsilon|}{|\phi'(1)|\ln\frac{1}{s_{\varepsilon,j}}}\phi\left( \frac{|T_{z_j}(x-z_j)|}{s_{\varepsilon,j}}\right) +O\left( \varepsilon^\gamma\right) >0,
	\end{split}
	\end{equation*}
	if $ L $ is sufficiently large.
	
	
	On the other hand, if $ \tau>0 $ small, $ x\in\Omega$  %\textbf{1}_{D_{\frac{\delta}{2\sqrt{|\ln\varepsilon|}}}\left (C_{\frac{r_*}{\sqrt{|\ln\varepsilon|}}}\right )} $,
	and $ |T_{z_j}(x-z_j)|\geq s_{\varepsilon,j}^\tau $ for any $ j=1,\cdots,N $, then by the definition of $ V_{\varepsilon,z_j,\hat{q}_{\varepsilon,j}} $ and Lemma \ref{H estimate},
	\begin{equation*}
	\begin{split}
	V_{\varepsilon,Z}(x)-q(x)|\ln\varepsilon|=&\sum_{j=1}^N\left( V_{\varepsilon,z_j,\hat{q}_{\varepsilon,j}}(x)+H_{\varepsilon,z_j,\hat{q}_{\varepsilon,j}}(x)\right) -q(x)|\ln\varepsilon|\\
	\leq &\left( \sum_{j=1}^N\frac{\hat{q}_{\varepsilon,j}\ln s_{\varepsilon,j}^\tau}{\ln s_{\varepsilon,j}}-C\right)|\ln\varepsilon| \\
	\leq &\left( \tau \sum_{j=1}^N \hat{q}_{\varepsilon,j}-C\right) |\ln\varepsilon|<0.
	\end{split}
	\end{equation*}
	If $ Ls_{\varepsilon,j}\leq |T_{z_j}(x-z_j)|\leq s_{\varepsilon,j}^\tau $, then by \eqref{q_i choice} for $ L $ sufficiently large
	\begin{equation*}
	\begin{split}
	V_{\varepsilon,Z}&(x)-q(x)|\ln\varepsilon|\\
	=&V_{\varepsilon,z_j,\hat{q}_{\varepsilon,j}}(x)+H_{\varepsilon,z_j,\hat{q}_{\varepsilon,j}}(x)-q(x)|\ln\varepsilon|+\sum_{i\neq j}\left(V_{\varepsilon, z_i, \hat{q}_{\varepsilon,i}}(x)+H_{\varepsilon, z_i, \hat{q}_{\varepsilon,i}}(x)\right)\\
	=&V_{\varepsilon, z_j, \hat{q}_{\varepsilon,j}}(x)-q(z_j)|\ln\varepsilon|-\frac{2\pi\hat{q}_{\varepsilon,j}\sqrt{\det K(z_j)}|\ln\varepsilon|}{\ln s_{\varepsilon,j}}\bar{S}_K(x, z_j)\\
	&-\sum_{i\neq j}\frac{2\pi\hat{q}_{\varepsilon,i}\sqrt{\det K(z_i)}|\ln\varepsilon|}{\ln s_{\varepsilon,i}}G_K(x, z_i)+O(s_{\varepsilon,j}^\tau)
	\end{split}
\end{equation*}	
	\begin{equation*}
\begin{split}	
	=&V_{\varepsilon,z_j,\hat{q}_{\varepsilon,j}}(x)-q(z_j)|\ln\varepsilon|-\frac{2\pi\hat{q}_{\varepsilon,j}\sqrt{\det K(z_j)}|\ln\varepsilon|}{\ln s_{\varepsilon,j}}\bar{S}_K(z_j, z_j)\\
	&-\sum_{i\neq j}\frac{2\pi\hat{q}_{\varepsilon,i}\sqrt{\det K(z_i)}|\ln\varepsilon|}{\ln s_{\varepsilon,i}}G_K(z_j, z_i)
+O\left( \varepsilon^{\tau\gamma}\right) \\
	=&V_{\varepsilon,z_j,\hat{q}_{\varepsilon,j}}(x)-\hat{q}_{\varepsilon,j}|\ln\varepsilon|+O\left( \varepsilon^{\tau\gamma}\right) \\
	\leq &-\frac{\hat{q}_{\varepsilon,j}\ln L|\ln\varepsilon|}{\ln\frac{1}{s_{\varepsilon,j}}}+O\left( \varepsilon^{\tau\gamma}\right)<0.
	\end{split}
	\end{equation*}

	
	
	
	
\end{proof}
\section{The linear theory: solvability of  a linear problem}
In this section,  we prove the solvability of  a linear equation related to the linearized operator $ L_\varepsilon $ of \eqref{key equa} at the approximate solution $ \sum_{j=1}^NV_{\varepsilon, Z,j} $. For given $ \xi\in F_{\varepsilon, Z} $, following the road map used in \cite{CPY}, we obtain the existence and uniqueness of $ u\in E_{\varepsilon,Z} $ satisfying equations
\begin{equation*} 
 Q_\varepsilon L_\varepsilon u=\xi.
\end{equation*}
Definitions of $ E_{\varepsilon,Z}, F_{\varepsilon,Z} $ and $ Q_\varepsilon $ can be seen in \eqref{204}, \eqref{205} and \eqref{206} below.
%First we prove that for any $ Z $ satisfying \eqref{admis set}, there exists $ \omega_{\delta,Z} $ such that $ V_{\delta, Z}+\omega_{\delta,Z} $ solves \eqref{eq1} in a co-dimensional $ 2m $ subspace of $ H^1_0 $. In the next section we choose proper $ Z=Z(\delta) $ such that $ V_{\delta, Z}+\omega_\delta $ is a solution.

We first note that the equation
\begin{equation}\label{eq5}
-\Delta w=w^p_+,\ \ \text{in}\ \mathbb{R}^2 
\end{equation}
has the unique $ C^1 $ solution 
\begin{equation*}
w(x)=\begin{cases}
\phi(x),\ \ &|x|\leq 1,\\
\phi'(1)\ln|x|,\ \ &|x|> 1.
\end{cases}
\end{equation*}
By the classical theory for elliptic equations, $ w\in C^{2,\alpha}(\mathbb{R}^2) $ for any $ \alpha\in(0,1) $. The linearized equation of \eqref{eq5} at $ w $ is
\begin{equation}\label{limit eq}
-\Delta v-pw^{p-1}_+v=0, \ \ v\in L^{\infty}(\mathbb{R}^2).
\end{equation}
Clearly, $ \frac{\partial w}{\partial x_\hbar} $ $(\hbar=1,2) $  are solutions of \eqref{limit eq}. Indeed, from \cite{DY} (see also \cite{CLW}), we have the non-degeneracy of $ w $ as follows.
\begin{proposition}[Non-degeneracy]\label{Non-degenerate}
$ w $ is non-degenerate, i.e., the kernel of the linearized equation \eqref{limit eq} is $$ span\left \{\frac{\partial w}{\partial x_1}, \frac{\partial w}{\partial x_2}\right \}. $$
\end{proposition}

Let $ \eta $ be a smooth  truncation function satisfying
\begin{equation*}
supp(\eta)\subseteq B_1(0),\ \ 0\leq \eta\leq 1\ \text{in}\ B_{1}(0),\   \  \eta\equiv1\ \text{in}\ B_{\frac{1}{2}}(0).
\end{equation*}
Define the scaled cut off function 
\begin{equation}\label{cut off fun}
\eta_i(x)=\eta\left ((x-z_i)|\ln\varepsilon|^2\right ).
\end{equation}
 Then $ supp(\eta_i)\subseteq B_{|\ln\varepsilon|^{-2}}(z_i) $ and $ supp(\eta_i)\cap supp(\eta_j)=\varnothing $ for $ i\neq j $ and $ \varepsilon $ sufficiently small. Moreover, $ ||\nabla\eta_i||_{L^\infty}\leq C|\ln\varepsilon|^{2} $ and $ ||\nabla^2\eta_i||_{L^\infty}\leq C|\ln\varepsilon|^{4} $.

Denote
\begin{equation}\label{204}
F_{\varepsilon,Z}=\left \{u\in L^p(\Omega)\mid \int_{\Omega}u\left( \eta_j\frac{\partial V_{\varepsilon, Z,j}}{\partial x_{\tilde{h}}}\right) =0,\ \ \forall j=1,\cdots,N,\ \tilde{h}=1,2\right \},
\end{equation}
and
\begin{equation}\label{205}
E_{\varepsilon,Z}=\left \{u\in W^{2,p}\cap H^1_0(\Omega)\mid \int_{\Omega}\left (K(x)\nabla u| \nabla \left( \eta_j\frac{\partial V_{\varepsilon, Z,j}}{\partial x_{\tilde{h}}}\right)\right )=0,\ \ \forall j=1,\cdots,N,\ \tilde{h}=1,2\right \}.
\end{equation}
Then $ F_{\varepsilon,Z}$ and $ E_{\varepsilon,Z} $ are co-dimensional $ 2N $ subspaces of $ L^p $ and $ W^{2,p}\cap H^1_0(\Omega) $, respectively.

For any $ u\in L^p(\Omega) $, we  define the projection operator $ Q_\varepsilon: L^p \to  F_{\varepsilon,Z} $
\begin{equation}\label{206}
Q_\varepsilon u:=u-\sum_{j=1}^N\sum_{\tilde{h}=1}^2b_{j,\tilde{h}}\left( -\varepsilon^2\text{div}\left (K(z_j)\nabla   \frac{\partial V_{\varepsilon, z_j, \hat{q}_{\varepsilon,j}}}{\partial x_{\tilde{h}}}\right )\right) ,
\end{equation}
where $ b_{j,\tilde{h}} (j=1,\cdots,N,\ \tilde{h}=1,2) $ satisfies a linear system as following
\begin{equation*}%\label{207}
\sum_{j=1}^N\sum_{\tilde{h}=1}^2b_{j,\tilde{h}}\int_{\Omega}\left( -\varepsilon^2\text{div}\left (K(z_j)\nabla   \frac{\partial V_{\varepsilon, z_j, \hat{q}_{\varepsilon,j}}}{\partial x_{\tilde{h}}}\right )\right)\left( \eta_i\frac{\partial V_{\varepsilon, Z,i}}{\partial x_\hbar}\right)=\int_{\Omega}u \left( \eta_i\frac{\partial V_{\varepsilon, Z,i}}{\partial x_\hbar}\right)
\end{equation*}
for $ i=1,\cdots,N,\ \hbar=1,2. $

We claim that $ Q_\varepsilon $ is a well-defined  linear projection  from $ L^p $ to $ F_{\varepsilon,Z}$. Indeed, using  \eqref{eq2}, \eqref{3-001} and \eqref{200},  for  $ Z\in \Lambda_{\varepsilon, N} $ the entry of coefficient matrix of the linear system satisfied by $b_{j,\tilde{h}}$ is
\begin{equation}\label{coef of C}
\begin{split}
\int_{\Omega}&\left( -\varepsilon^2\text{div}\left (K(z_j)\nabla   \frac{\partial V_{\varepsilon, z_j, \hat{q}_{\varepsilon,j}}}{\partial x_{\tilde{h}}}\right )\right)\left( \eta_i\frac{\partial V_{\varepsilon, Z,i}}{\partial x_\hbar}\right)\\
&=p\int_{\Omega}\eta_i\left (V_{\varepsilon, z_j, \hat{q}_{\varepsilon,j}}-\hat{q}_{\varepsilon,j}|\ln\varepsilon|\right )^{p-1}_+ \frac{\partial V_{\varepsilon, z_j, \hat{q}_{\varepsilon,j}}}{\partial x_{\tilde{h}}}\frac{\partial V_{\varepsilon, z_i, \hat{q}_{\varepsilon,i}}}{\partial x_\hbar}+O\left( \varepsilon^\gamma\right)\\
&=\delta_{i,j}(M_{i})_{\tilde{h},\hbar}+O\left( \varepsilon^\gamma\right) ,
\end{split}
\end{equation}
where $ \delta_{i,j}= 1 $ if $ i = j $ and $ \delta_{i,j}= 0 $ otherwise.  $ M_{i} $ are $ N $ positive definite matrices %and there exist positive constants $ \bar{c}_1,\bar{c}_2 $ independent of $ \delta, Z $
such that all eigenvalues of $ M_{i} $ belong to $ (\bar{c}_1,\bar{c}_2) $ for constants $ \bar{c}_1,\bar{c}_2>0 $.
This implies the existence and uniqueness of $ b_{j,\tilde{h}} $.
Note that for $ u\in L^p $, $ Q_{\varepsilon} u \equiv u $ in $ \Omega \backslash\cup_{i=1}^NB_{Ls_{\varepsilon,i}}(z_{i}) $ for some $ L>1 $. Moreover, one can easily get that there exists a constant $ C>0 $ independent of $ \varepsilon $, such that for any $ q\in[1,+\infty) $, $ u\in L^q(\Omega) $ with $ supp(u)\subset  \cup_{j=1}^NB_{Ls_{\varepsilon,j}}(z_j)$,
\begin{equation*}
||Q_\varepsilon u||_{L^q(\Omega)}\leq C||u||_{L^q(\Omega)}.
\end{equation*}

The linearized operator of \eqref{key equa} at $ V_{\varepsilon, Z} $ is
\begin{equation*}
L_\varepsilon \omega =-\varepsilon^2\text{div}(K(x)\nabla \omega)- p\left (V_{\varepsilon, Z}-q|\ln\varepsilon|\right )^{p-1}_+ \omega.
\end{equation*}
We give  coercive estimates of the linear operator $ Q_\varepsilon L_\varepsilon $.
\begin{proposition}\label{coercive esti}
There exist  $ \rho_0>0, \varepsilon_0>0 $ such that for any $ \varepsilon\in(0,\varepsilon_0],  Z\in \Lambda_{\varepsilon, N} $, if $ u\in E_{\varepsilon,Z} $ satisfying $ Q_\varepsilon L_\varepsilon u=0 $ in $ \Omega\backslash\cup_{j=1}^NB_{Ls_{\varepsilon,j}}(z_j) $ for some $ L>1 $ large, then
\begin{equation*}
||Q_\varepsilon L_\varepsilon u||_{L^p}\geq \rho_0\varepsilon^{\frac{2}{p}}||u||_{L^\infty}. 
\end{equation*}
\end{proposition}

\begin{proof}
We argue by contradiction. Suppose that there are $ \varepsilon_m \to 0 $, $ Z_m=(z_{m,1},\cdots,z_{m,N})\to (z_{1},\cdots,z_{N})\in \Omega^{(N)} $  and $  u_m \in E_{\varepsilon_m, Z_m} $ with $ Q_{\varepsilon_m} L_{\varepsilon_m} u_m = 0 $ in $ \Omega \backslash\cup_{j=1}^NB_{Ls_{\varepsilon_m,j}}(z_{m,j}) $ for some $ L $ large and $ ||u_m||_{L^\infty}=1 $ such that
\begin{equation*}
||Q_{\varepsilon_m} L_{\varepsilon_m} u_m||_{L^p}\leq \frac{1}{m}\varepsilon_m^{\frac{2}{p}}.
\end{equation*}
By definition of $ Q_{\varepsilon}  $,
\begin{equation}\label{208}
Q_{\varepsilon_m} L_{\varepsilon_m} u_m=  L_{\varepsilon_m} u_m-\sum_{j=1}^N\sum_{\tilde{h}=1}^2b_{j,\tilde{h},m}\left( -\varepsilon_m^2\text{div}\left (K(z_{m,j})\nabla   \frac{\partial V_{\varepsilon_m, z_{m,j},\hat{q}_{\varepsilon_m,j}}}{\partial x_{\tilde{h}}}\right )\right) .
\end{equation}
We now estimate $ b_{j,\tilde{h},m} $. For fixed $ i=1,\cdots,N, \hbar=1,2 $, multiplying \eqref{208} by $ \eta_i\frac{\partial V_{\varepsilon_m, Z_{m},i}}{\partial x_\hbar} $ and integrating on $ \Omega $ we get
\begin{equation*}
\begin{split}
\int_{\Omega}u_m&L_{\varepsilon_m}\left( \eta_i\frac{\partial V_{\varepsilon_m, Z_{m},i}}{\partial x_\hbar}\right) =\int_{\Omega}L_{\varepsilon_m}u_m\left( \eta_i\frac{\partial V_{\varepsilon_m, Z_{m},i}}{\partial x_\hbar}\right) \\
=&\sum_{j=1}^N\sum_{\tilde{h}=1}^2b_{j,\tilde{h},m}\int_{\Omega}-\varepsilon_m^2\text{div}\left (K(z_{m,j})\nabla   \frac{\partial V_{\varepsilon_m, z_{m,j},\hat{q}_{\varepsilon_m,j}}}{\partial x_{\tilde{h}}}\right )\left( \eta_i\frac{\partial V_{\varepsilon_m, Z_{m},i}}{\partial x_\hbar}\right) .
\end{split}
\end{equation*}

%We estimate $ \int_{\Omega}u_NL_{\delta_N}\left( \eta_i\frac{\partial V_{\delta_N, Z_{N},i}}{\partial x_\hbar}\right). $ 
For the left hand side of the above equality, we have
\begin{equation*}
\begin{split}
&\int_{\Omega}u_mL_{\varepsilon_m}\left( \eta_i\frac{\partial V_{\varepsilon_m, Z_{m},i}}{\partial x_\hbar}\right) \\
=&-\int_{\Omega}u_m\varepsilon_m^2\text{div}\left (K(x)\nabla \left( \eta_i\frac{\partial V_{\varepsilon_m, Z_{m},i}}{\partial x_\hbar}\right) \right )- p\int_{\Omega}u_m\left (V_{\varepsilon_m, Z_m}-q|\ln\varepsilon_m|\right )^{p-1}_+\left( \eta_i\frac{\partial V_{\varepsilon_m, Z_{m},i}}{\partial x_\hbar}\right)  \\
=&-\int_{\Omega}\eta_iu_m\varepsilon_m^2\text{div}\left (K(x)\nabla \frac{\partial V_{\varepsilon_m, Z_{m},i}}{\partial x_\hbar} \right )-2\int_{\Omega}u_m\varepsilon_m^2\left (K(x)\nabla\eta_i|\nabla   \frac{\partial V_{\varepsilon_m, Z_{m},i}}{\partial x_\hbar}  \right )\\
&-\int_{\Omega}u_m\varepsilon_m^2\text{div}\left (K(x)\nabla   \eta_i  \right )\frac{\partial V_{\varepsilon_m, Z_{m},i}}{\partial x_\hbar}- p\int_{\Omega}u_m\left (V_{\varepsilon_m, Z_m}-q|\ln\varepsilon_m|\right )^{p-1}_+\left( \eta_i\frac{\partial V_{\varepsilon_m, Z_{m},i}}{\partial x_\hbar}\right)
\end{split}
\end{equation*}
\begin{equation}\label{301}
\begin{split}
=&\int_{\Omega}\eta_iu_mp\left( V_{\varepsilon_m, z_{m,i}, \hat{q}_{\varepsilon_m,i}}-\hat{q}_{\varepsilon_m,i}|\ln\varepsilon_m|\right)^{p-1}_+\frac{\partial V_{\varepsilon_m, z_{m,i}, \hat{q}_{\varepsilon_m,i}}}{\partial x_\hbar}+\int_{\Omega}\eta_iu_m\varepsilon_m^2\text{div}\left (\frac{\partial K(x)}{\partial x_\hbar}\nabla V_{\varepsilon_m, Z_{m},i}  \right )\\
&-2\int_{\Omega}u_m\varepsilon_m^2\left (K(x)\nabla\eta_i|\nabla   \frac{\partial V_{\varepsilon_m, Z_{m},i}}{\partial x_\hbar}  \right ) -\int_{\Omega}u_m\varepsilon_m^2\text{div}\left (K(x)\nabla   \eta_i  \right )\frac{\partial V_{\varepsilon_m, Z_{m},i}}{\partial x_\hbar}\\
&- p\int_{\Omega}u_m\left (V_{\varepsilon_m, Z_m}-q|\ln\varepsilon_m|\right )^{p-1}_+\left( \eta_i\frac{\partial V_{\varepsilon_m, Z_{m},i}}{\partial x_\hbar}\right).
%=:&I_1+I_2+I_3+I_4+I_5.
\end{split}
\end{equation}

By \eqref{3-001}, \eqref{203} and Lemma \ref{lemA-5},  
\begin{equation}\label{302}
\begin{split}
\int_{\Omega}\eta_i&u_mp\left( V_{\varepsilon_m, z_{m,i}, \hat{q}_{\varepsilon_m,i}}-\hat{q}_{\varepsilon_m,i}|\ln\varepsilon_m|\right)^{p-1}_+\frac{\partial V_{\varepsilon_m, z_{m,i}, \hat{q}_{\varepsilon_m,i}}}{\partial x_\hbar}-  p\int_{\Omega}u_m\left (V_{\varepsilon_m, Z_m}-q|\ln\varepsilon_m|\right )^{p-1}_+\left( \eta_i\frac{\partial V_{\varepsilon_m, Z_{m},i}}{\partial x_\hbar}\right)\\
=&p\int_{\Omega}u_m\left( V_{\varepsilon_m, z_{m,i}, \hat{q}_{\varepsilon_m,i}}-\hat{q}_{\varepsilon_m,i}|\ln\varepsilon_m|\right)^{p-1}_+\frac{\partial V_{\varepsilon_m, z_{m,i}, \hat{q}_{\varepsilon_m,i}}}{\partial x_\hbar}\\
&-p\int_{\Omega}u_m\left( V_{\varepsilon_m, z_{m,i}, \hat{q}_{\varepsilon_m,i}}-\hat{q}_{\varepsilon_m,i}|\ln\varepsilon_m|+O\left( \varepsilon_m^\gamma\right) \right)^{p-1}_+\frac{\partial V_{\varepsilon_m, z_{m,i}, \hat{q}_{\varepsilon_m,i}}}{\partial x_\hbar}+O\left( \varepsilon_m^{1+\gamma}\right)\\
=&O\left( \varepsilon_m^{1+\gamma}\right).
\end{split}
\end{equation}
As for the remaining terms of \eqref{301}, %from the choice of $ \eta_i $, we have 
using $ ||\nabla\eta_i||_{L^\infty}\leq C|\ln\varepsilon_m|^{2}$ and $ ||\nabla^2\eta_i||_{L^\infty}\leq C|\ln\varepsilon_m|^{4} $ we have
\begin{equation}\label{304}
\begin{array}{ll}
\int_{\Omega}\eta_iu_m\varepsilon_m^2\text{div}\left (\frac{\partial K(x)}{\partial x_\hbar}\nabla V_{\varepsilon_m, Z_{m},i}  \right )&\\
\qquad=\int\limits_{|T_{z_{m,i}}(x-z_{m,i})|\leq s_{\varepsilon_m,i}}\eta_iu_m \left( \varepsilon_m^2\text{div}\left( \frac{\partial K(x)}{\partial x_\hbar}\nabla V_{\varepsilon_m, z_{m,i}, \hat{q}_{\varepsilon_m,i}}\right)\right)&\\
\qquad\,\,\,\,\,+\int\limits_{ s_{\varepsilon_m,i}<|T_{z_{m,i}}(x-z_{m,i})|\leq |\ln\varepsilon_m|^{2}}
\eta_iu_m \left( \varepsilon_m^2\text{div}\left( \frac{\partial K(x)}{\partial x_\hbar}\nabla V_{\varepsilon_m, z_{m,i}, \hat{q}_{\varepsilon_m,i}}\right)\right)+O\left( \varepsilon_m^2\right)  &\\
\qquad=O\left(\varepsilon_m^2\right) +O\left (\varepsilon_m^2|\ln\varepsilon_m|\right )+O\left( \varepsilon_m^2\right)&\\
\qquad=O\left( \varepsilon_m^2|\ln\varepsilon_m|\right),
\end{array}
\end{equation}
\begin{equation}\label{305}
\begin{split}
-2\int_{\Omega}&u_m\varepsilon_m^2\left (K(x)\nabla\eta_i|\nabla   \frac{\partial V_{\varepsilon_m, Z_{m},i}}{\partial x_\hbar}  \right )\\
=&-2\int_{B_{|\ln\varepsilon_m|^{2}}(z_{m,i})\backslash B_{\frac{|\ln\varepsilon_m|^{2}}{2}}(z_{m,i})}u_m\varepsilon_m^2\left (K(x)\nabla\eta_i|\nabla   \frac{\partial V_{\varepsilon_m, z_{m,i},\hat{q}_{\varepsilon_m,i}}}{\partial x_\hbar}  \right )+O\left( \varepsilon_m^2|\ln\varepsilon_m|^{2}\right)\\
=&O\left( \varepsilon_m^2|\ln\varepsilon_m|^{2}\right),
\end{split}
\end{equation}
and
\begin{equation}\label{306}
\begin{split}
-\int_{\Omega}&u_m\varepsilon_m^2\text{div}\left (K(x)\nabla   \eta_i  \right )\frac{\partial V_{\varepsilon_m, Z_{m},i}}{\partial x_\hbar}\\
=&-\int_{B_{|\ln\varepsilon_m|^{2}}(z_{m,i})\backslash B_{\frac{|\ln\varepsilon_m|^{2}}{2}}(z_{m,i})}u_m\varepsilon_m^2\text{div}\left (K(x)\nabla   \eta_i  \right )\frac{\partial V_{\varepsilon_m, z_{m,i},\hat{q}_{\varepsilon_m,i}}}{\partial x_\hbar}+O\left( \varepsilon_m^2|\ln\varepsilon_m|^{4}\right)\\
=&O\left( \varepsilon_m^2|\ln\varepsilon_m|^{4}\right),
\end{split}
\end{equation}
where we have used \eqref{200}. Inserting \eqref{302}, \eqref{304}, \eqref{305} and \eqref{306} into \eqref{301}, we finally get
\begin{equation*}
\begin{split}
&\int_{\Omega}u_mL_{\varepsilon_m}\left( \eta_i\frac{\partial V_{\varepsilon_m, Z_{m},i}}{\partial x_\hbar}\right) =O\left( \varepsilon_m^{1+\gamma}\right),
\end{split}
\end{equation*}
which combined with \eqref{coef of C} deduces
\begin{equation*}
b_{j,\tilde{h},m}=O\left (\varepsilon_m^{1+\gamma}\right ).
\end{equation*}
Consequently
\begin{equation*}
\begin{split}
\sum_{j=1}^N\sum_{\tilde{h}=1}^2&b_{j,\tilde{h},m}\left( -\varepsilon_m^2\text{div}\left (K(z_{m,j})\nabla   \frac{\partial V_{\varepsilon_m, z_{m,j},\hat{q}_{\varepsilon_m,j}}}{\partial x_{\tilde{h}}}\right )\right)\\
%=&p\sum_{j=1}^m\sum_{h=1}^2b_{j,h,N}(V_{\delta_N, z_{N,j}, \hat{q}_{\delta_N,j}}-\hat{q}_{\delta_N,j})^{p-1}_+ \frac{\partial V_{\delta_N, z_{N,j},\hat{q}_{\delta_N,j}}}{\partial x_h} \\
&=O\left( \sum_{j=1}^N\sum_{\tilde{h}=1}^2\varepsilon_m^{\frac{2}{p}-1}|b_{j,\tilde{h},m}|\right)
=O\left( \varepsilon_m^{\frac{2}{p}+\gamma}\right),\ \ \ \ \text{in}\ L^p(\Omega).
\end{split}
\end{equation*}
Taking this into \eqref{208}, we have
\begin{equation}\label{307}
\begin{split}
L_{\varepsilon_m} u_m=&Q_{\varepsilon_m} L_{\varepsilon_m} u_m+\sum_{j=1}^N\sum_{\tilde{h}=1}^2b_{j,\tilde{h},m}\left( -\varepsilon_m^2\text{div}\left (K(z_{m,j})\nabla   \frac{\partial V_{\varepsilon_m, z_{m,j},\hat{q}_{\varepsilon_m,j}}}{\partial x_{\tilde{h}}}\right )\right)\\
=&O\left( \frac{1}{m}\varepsilon_m^{\frac{2}{p}}\right) +O\left(  \varepsilon_m^{\frac{2}{p}+\gamma}\right)
=o\left( \varepsilon_m^{\frac{2}{p}}\right),\ \ \ \ \text{in}\ L^p(\Omega).
\end{split}
\end{equation}

For fixed $ i, $ we define the scaled function $ \tilde{u}_{m,i}(y)=u_m(s_{\varepsilon_m, i}y+z_{m,i}) $ for $ y\in \Omega_{m,i}:=\{y\in\mathbb{R}^2\mid s_{\varepsilon_m, i}y+z_{m,i}\in \Omega\} $.
Define
\begin{equation*}
\begin{split}
\tilde{L}_{m,i}u=&-\text{div}\left (K(s_{\varepsilon_m, i}y+z_{m,i})\nabla u\right )\\
&-p\frac{s_{\varepsilon_m,i}^2}{\varepsilon_m^2}\left (V_{\varepsilon_m, Z_m}(s_{\varepsilon_m, i}y+z_{m,i})-q(s_{\varepsilon_m, i}y+z_{m,i})|\ln\varepsilon_m|\right )^{p-1}_+ u.
\end{split}
\end{equation*}
Then
\begin{equation*}
\begin{split}
||\tilde{L}_{m,i}\tilde{u}_{m,i}||_{L^p(\Omega_{m,i})}
%=&\left( \int_{\Omega_{N,i}}\left( \tilde{L}_{N,i}\tilde{u}_{N,i} \right)^pdy\right)^{\frac{1}{p}}  \\
%=&\left( \frac{1}{s_{\delta_N, i}^2}\int_{\Omega}\left( -s_{\delta_N, i}^2\text{div}(K(x)\nabla u_N)-p\frac{s_{\delta_N, i}^2}{\delta_N^2}(V_{\delta_N,Z_N}-q)^{p-1}_+u_N \right)^pdx\right)^{\frac{1}{p}}\\
=\frac{s_{\varepsilon_m, i}^2}{s_{\varepsilon_m, i}^{\frac{2}{p}}\varepsilon_m^2}||L_{\varepsilon_m}u_m||_{L^p(\Omega)}.
\end{split}
\end{equation*}
%i.e., $ s_{\delta_N, i}^{\frac{2}{p}}\cdot\frac{\delta_N^2}{s_{\delta_N, i}^2}||\tilde{L}_{N,i}\tilde{u}_{N,i}||_{L^p(\Omega_{N,i})}=||L_{\delta_N}u_N||_{L^p(\Omega)}. $
Since  $ s_{\varepsilon_m, i}=O(\varepsilon_m) $, by \eqref{307} we have
\begin{equation*}
\tilde{L}_{m,i}\tilde{u}_{m,i}=o(1)\ \ \ \  \text{in}\ \ L^p(\Omega_{m,i}).
\end{equation*}
Since $ ||\tilde{u}_{m,i}||_{L^\infty(\Omega_{m,i})}=1 $, by the classical regularity theory of  elliptic equations, $ \tilde{u}_{m,i} $ is uniformly bounded in $ W^{2,p}_{loc}(\mathbb{R}^2) $, which implies that there is $u_i$ such that
\begin{equation*}
\tilde{u}_{m,i}\to u_i\ \ \ \ \text{in}\ \ C^1_{loc}(\mathbb{R}^2).
\end{equation*}

We claim that $ u_i\equiv 0. $ On the one hand, in our ansatz for $ Z\in \Lambda_{\varepsilon, N} $, $ |z_i-z_j|\geq |\ln\varepsilon|^{-1} $. So by \eqref{203}, $ z_{m,i}\to z_i $ as $ m\to\infty $ and the fact that $ \lim_{\varepsilon\to 0}\varepsilon|\ln\varepsilon|=0 $, we get
\begin{equation*}
\begin{split}
\frac{s_{\varepsilon_m,i}^2}{\varepsilon_m^2}&\left (V_{\varepsilon_m, Z_m}(s_{\varepsilon_m, i}y+z_{m,i})-q\left (s_{\varepsilon_m, i}y+z_{m,i}\right )|\ln\varepsilon_m|\right )^{p-1}_+\\
=&\frac{s_{\varepsilon_m,i}^2}{\varepsilon_m^2}\left( V_{\varepsilon_m, z_{m,i}, \hat{q}_{\varepsilon_m,i}}(s_{\varepsilon_m, i}y+z_{m,i})-\hat{q}_{\varepsilon_m,i}|\ln\varepsilon_m|+O\left (\varepsilon_m^{\gamma}\right )\right)^{p-1}_+ \\
\to&\phi(T_{z_i}y)^{p-1}_+\ \  \  \text{in}\ C^0_{loc}(\mathbb{R}^2)\ \ \ \ \ \ \  \text{as}\ m\to\infty.
\end{split}
\end{equation*}
So $ u_i $ satisfies
\begin{equation*}
-\text{div}(K(z_i)\nabla u_i(x))-p\phi(T_{z_i}x)^{p-1}_+u_i(x)=0,\ \ x\in\mathbb{R}^2.
\end{equation*}
Let $ \hat{u}_i(x)=u_i(T_{z_i}^{-1}x) $. Since $ T_{z_i}^{-1}(T_{z_i}^{-1})^T=K(z_i) $, we have
\begin{equation*}
-\Delta \hat{u}_i(x)=-\text{div}(K(z_i)\nabla u_i)(T_{z_i}^{-1}x)=p\phi(x)^{p-1}_+\hat{u}_i(x),\ \ \ \   \ x\in \mathbb{R}^2.
\end{equation*}
By Proposition \ref{Non-degenerate}, there exist $ c_1,c_2 $ such that
\begin{equation}\label{210}
\hat{u}_i=c_1\frac{\partial \phi}{\partial x_1}+c_2\frac{\partial \phi}{\partial x_2}.
\end{equation}

On the other hand, since $ u_m\in E_{\varepsilon_m,Z_m} $, we get
\begin{equation*}
\int_{\Omega}-\varepsilon_m^2\text{div}\left( K(x)\nabla \left( \eta_i\frac{\partial V_{\varepsilon_m, Z_{m}, i}}{\partial x_\hbar}\right)\right)  u_m=0,\ \ \forall\ i=1,\cdots,N,\  \hbar=1,2,
\end{equation*}
which implies that
\begin{equation}\label{209}
\begin{split}
0%=&-\int_{\Omega}u_N\eta_i\delta_N^2\text{div}\left( K(x)\nabla \frac{\partial V_{\delta_N, z_{N,i}, \hat{q}_{\delta_N,i}}}{\partial x_\hbar}\right) -2\int_{\Omega}u_N\delta_N^2\left (K(x)\nabla\eta_i|\nabla   \frac{\partial V_{\delta_N, z_{N,i},\hat{q}_{\delta_N,i}}}{\partial x_\hbar}  \right )\\
%& -\int_{\Omega}u_N\delta_N^2\text{div}\left (K(x)\nabla   \eta_i  \right )\frac{\partial V_{\delta_N, z_{N,i},\hat{q}_{\delta_N,i}}}{\partial x_\hbar}\\
=&p\int_{\Omega}u_m\left( V_{\varepsilon_m, z_{m,i}, \hat{q}_{\varepsilon_m,i}}-\hat{q}_{\varepsilon_m,i}|\ln\varepsilon_m|\right)^{p-1}_+\frac{\partial V_{\varepsilon_m, z_{m,i}, \hat{q}_{\varepsilon_m,i}}}{\partial x_\hbar}+\int_{\Omega}\eta_iu_m\varepsilon_m^2\text{div}\left (\frac{\partial K(x)}{\partial x_\hbar}\nabla V_{\varepsilon_m, Z_{m},i}  \right )\\
&-2\int_{\Omega}u_m\varepsilon_m^2\left (K(x)\nabla\eta_i|\nabla   \frac{\partial V_{\varepsilon_m, Z_{m},i}}{\partial x_\hbar}  \right ) -\int_{\Omega}u_m\varepsilon_m^2\text{div}\left (K(x)\nabla   \eta_i  \right )\frac{\partial V_{\varepsilon_m, Z_{m},i}}{\partial x_\hbar}.
\end{split}
\end{equation}
Using again \eqref{304}, \eqref{305} and \eqref{306}, we have
\begin{equation}\label{209-1}
\begin{split}
\int_{\Omega}&\eta_iu_m\varepsilon_m^2\text{div}\left (\frac{\partial K(x)}{\partial x_\hbar}\nabla V_{\varepsilon_m, Z_{m},i}  \right )-2\int_{\Omega}u_m\varepsilon_m^2\left (K(x)\nabla\eta_i|\nabla   \frac{\partial V_{\varepsilon_m, Z_{m},i}}{\partial x_\hbar}  \right ) \\
&-\int_{\Omega}u_m\varepsilon_m^2\text{div}\left (K(x)\nabla   \eta_i  \right )\frac{\partial V_{\varepsilon_m, Z_{m},i}}{\partial x_\hbar}
=O\left( \varepsilon_m^2|\ln\varepsilon_m|^{4}\right).
\end{split}
\end{equation}
%and
%\begin{equation}\label{209-2}
%-2\int_{\Omega}u_N\delta_N^2\left (K(x)\nabla\eta_i|\nabla   \frac{\partial V_{\delta_N, z_{N,i},\hat{q}_{\delta_N,i}}}{\partial x_\hbar}  \right ) -\int_{\Omega}u_N\delta_N^2\text{div}\left (K(x)\nabla   \eta_i  \right )\frac{\partial V_{\delta_N, z_{N,i},\hat{q}_{\delta_N,i}}}{\partial x_\hbar}=O\left( \frac{\varepsilon_N^2}{|\ln\varepsilon_N|^{p}}\right).
%\end{equation}
From  \eqref{200},
\begin{equation}\label{209-3}
\begin{split}
&p\int_{\Omega} u_m\left (V_{\varepsilon_m, z_{m,i}, \hat{q}_{\varepsilon_m,i}}-\hat{q}_{\varepsilon_m,i}|\ln\varepsilon_m|\right )^{p-1}_+ \frac{\partial V_{\varepsilon_m, z_{m,i}, \hat{q}_{\varepsilon_m,i}}}{\partial x_\hbar}\\
=&p\int_{\Omega}\frac{1}{s_{\varepsilon_m,i}}\left( \frac{\varepsilon_m}{s_{\varepsilon_m,i}}\right)^{\frac{2p}{p-1}} \phi\left( \frac{T_{z_{m,i}}(x-z_{m,i})}{s_{\varepsilon_m,i}}\right)^{p-1}_+\phi'\left( \frac{T_{z_{m,i}}(x-z_{m,i})}{s_{\varepsilon_m,i}}\right)\frac{(T_{z_{m,i}})_\hbar^T\cdot T_{z_{m,i}}(x-z_{m,i})}{|T_{z_{m,i}}(x-z_{m,i})|}u_m\\
=&ps_{\varepsilon_m,i}\left( \frac{\varepsilon_m}{s_{\varepsilon_m,i}}\right)^{\frac{2p}{p-1}}\int_{\mathbb{R}^2}\phi(T_{z_{m,i}}y)^{p-1}_+\phi'(T_{z_{m,i}}y)\frac{(T_{z_{m,i}})_\hbar^T\cdot T_{z_{m,i}}y}{|T_{z_{m,i}}y|}\tilde{u}_{m,i}(y)dy.
\end{split}
\end{equation}
Taking \eqref{209-1}  and \eqref{209-3} into \eqref{209}
%, we have
%\begin{equation}\label{209-0}
%\begin{split}
%0=&ps_{\delta_N,i}\left( \frac{\delta_N}{s_{\delta_N,i}}\right)^{\frac{2p}{p-1}}\int_{\mathbb{R}^2}\phi(T_{z_{N,i}}y)^{p-1}_+\phi'(T_{z_{N,i}}y)\frac{(T_{z_{N,i}})_\hbar^t\cdot T_{z_{N,i}}y}{|T_{z_{N,i}}y|}\tilde{u}_{N,i}(y)dy+O\left( \delta_N^2|\ln\varepsilon_N|^{2M+1}\right).
%\end{split}
%\end{equation}
, dividing both sides of \eqref{209} into $ ps_{\varepsilon_m,i}(\frac{\varepsilon_m}{s_{\varepsilon_m,i}})^{\frac{2p}{p-1}} $ and passing $ m $ to the limit, we get
\begin{equation*}
\begin{split}
0%=&\int_{\mathbb{R}^2}\phi(T_{z_{i}}y)^{p-1}_+\phi'(T_{z_{i}}y)\frac{(T_{z_{i}})_\hbar^t\cdot T_{z_{i}}y}{|T_{z_{i}}y|}u_{i}(y)dy\\
=\int_{\mathbb{R}^2}\phi(x)^{p-1}_+\phi'(x)\frac{(T_{z_{i}})_\hbar^T\cdot x}{|x|}\hat{u}_{i}(x)\sqrt{\det(K(z_i))}dx,\ \  \hbar=1,2,
\end{split}
\end{equation*}
which implies that
\begin{equation}\label{211}
0=\int_{B_1(0)}\phi^{p-1}_+\frac{\partial \phi}{\partial x_\hbar}\hat{u}_i.
\end{equation}
So $ c_1=c_2=0. $ That is, $ u_i\equiv 0. $ We conclude that $ \tilde{u}_{N,i}\to 0 $ in  $ C^1(B_L(0)) $, which implies that
\begin{equation}\label{212}
||u_m||_{L^\infty(B_{Ls_{\varepsilon_m,i}}(z_{m,i}))}=o(1).
\end{equation}

Since $ Q_{\varepsilon_m} L_{\varepsilon_m} u_m = 0 $ in $ \Omega \backslash\cup_{i=1}^NB_{Ls_{\varepsilon_m,i}}(z_{m,i}) $, by the definition of $  Q_{\varepsilon_m} $ we have for $ L $ large
\begin{equation*}
 L_{\varepsilon_m} u_m = 0\  \ \text{in} \ \Omega \backslash\cup_{i=1}^NB_{Ls_{\varepsilon_m,i}}(z_{m,i}).
\end{equation*}
From Lemma \ref{lemA-5},  $ (V_{\varepsilon_m,Z_m}-q|\ln\varepsilon_m|)_+=0\  \ \text{in} \ \Omega \backslash\cup_{i=1}^NB_{Ls_{\varepsilon_m,i}}(z_{m,i}). $
So $ -\text{div}(K(x)\nabla u_m)=0 $ in $ \Omega \backslash\cup_{i=1}^NB_{Ls_{\varepsilon_m,i}}(z_{m,i}). $
%Since $ u_N=o(1) $ on $ \cup_{i=1}^m\partial B_{Ls_{\delta_N,i}}(z_{N,i}) $ and $ u_N=0 $ on $ \partial \Omega $,
Thus by the maximum principle, 
\begin{equation*}
||u_m||_{L^\infty(\Omega\backslash \cup_{i=1}^NB_{Ls_{\varepsilon_m,i}}(z_{m,i}))}=o(1),
\end{equation*}
which combined with \eqref{212} leads to
\begin{equation*}
||u_m||_{L^\infty(\Omega)}=o(1).
\end{equation*}
This is a contradiction since $ ||u_m||_{L^\infty(\Omega)}=1. $



\end{proof}


As a  consequence of Proposition \ref{coercive esti}, we have that $ Q_\varepsilon L_\varepsilon $ is indeed a one to one and onto map from $ E_{\varepsilon,Z} $ to $ F_{\varepsilon, Z}. $
\begin{proposition}\label{one to one and onto}
$ Q_\varepsilon L_\varepsilon $ is  a one to one and onto map from $ E_{\varepsilon,Z} $ to $ F_{\varepsilon, Z}. $
\end{proposition}

\begin{proof}
If $ Q_\varepsilon L_\varepsilon u\equiv0 $, by Lemma \ref{coercive esti}, $ u\equiv0 $. So $ Q_\varepsilon L_\varepsilon $ is  one to one.

% Denote
%\begin{equation*}
%\hat{E}=\left \{u\in H^1_0(\Omega)\mid \int_{\Omega}\left( K(x)\nabla u|\nabla\left(\eta_i \frac{\partial V_{\varepsilon, Z,i}}{\partial x_h}\right)\right)  =0,\ \ \ \ i=1,\cdots,m,\ h=1,2\right \}.
%\end{equation*}
%Then $ E_{\varepsilon,Z}=\hat{E}\cap W^{2,p}(\Omega) $. 
For any $ \hat{h}\in F_{\varepsilon,Z}  $, by the Riesz representation theorem there is a unique $ u\in H^1_0(\Omega) $ such that
\begin{equation}\label{213}
\varepsilon^2\int_{\Omega}\left( K(x)\nabla u|\nabla\varphi\right) =\int_{\Omega}\hat{h}\varphi,\ \ \   \ \forall \varphi\in H^1_0(\Omega).
\end{equation}

Since $ \hat{h}\in F_{\varepsilon,Z} $, using the classical $ L^p $ theory of elliptic equations, we conclude that $ u\in W^{2,p}(\Omega) $, which implies that $ u\in E_{\varepsilon,Z}. $ Thus $ -\varepsilon^2\text{div}(K(x)\nabla)=Q_{\varepsilon}(-\varepsilon^2\text{div}(K(x)\nabla)) $ is a one to one and onto map from $ E_{\varepsilon,Z} $ to $ F_{\varepsilon, Z}. $

For any $ \xi\in F_{\varepsilon,Z} $,  $ Q_\varepsilon L_\varepsilon u=\xi $ is equivalent to
\begin{equation}\label{214}
\begin{split}
u=&\left (Q_{\varepsilon}\left (-\varepsilon^2\text{div}(K(x)\nabla)\right )\right )^{-1}pQ_\varepsilon\left (V_{\varepsilon,Z}-q|\ln\varepsilon|\right )^{p-1}_+u\\
&+\left (Q_{\varepsilon}\left (-\varepsilon^2\text{div}(K(x)\nabla)\right )\right )^{-1}\xi.%,\ \ u\in E_{\varepsilon,Z}.
\end{split}
\end{equation}
Note that $  (Q_{\varepsilon}(-\varepsilon^2\text{div}(K(x)\nabla)))^{-1}pQ_\varepsilon\left (V_{\varepsilon,Z}-q|\ln\varepsilon|\right )^{p-1}_+u $ is a compact operator in $ E_{\varepsilon,Z}, $ by the Fredholm alternative, \eqref{214} is solvable if and only if
\begin{equation*}
u=\left (Q_{\varepsilon}\left (-\varepsilon^2\text{div}(K(x)\nabla)\right )\right )^{-1}pQ_\varepsilon\left (V_{\varepsilon,Z}-q|\ln\varepsilon|\right )^{p-1}_+u
\end{equation*}
has only trivial solution, which is true since $ Q_\varepsilon L_\varepsilon  $ is one to one. %So $ Q_\varepsilon L_\varepsilon $ is  an onto  map from $ E_{\varepsilon,Z} $ to $ F_{\varepsilon, Z}  $ and the proof is complete.
\end{proof}






\section{The nonlinear theory: solvability of a nonlinear equation}
%Now we  find solution of \eqref{111} being of the form $$ V_{\delta, Z}+\omega_\delta. $$ First we prove that for any $ Z $ satisfying \eqref{admis set}, there exists $ \omega_{\delta,Z} $ such that $ V_{\delta, Z}+\omega_{\delta,Z} $ solves \eqref{eq1} in a co-dimensional $ 2m $ subspace of $ H^1_0 $. In the next section we choose proper $ Z=Z(\delta) $ such that $ V_{\delta, Z}+\omega_\delta $ is a solution.


%From \eqref{3-03}, to construct solutions of \eqref{111} being the form  $ V_{\delta,Z}+\omega_\delta $, it suffices to solve
%\begin{equation}\label{215}
%\begin{split}
%L_\delta \omega_\delta=l_\delta+R_\delta(\omega_\delta),
%\end{split}
%\end{equation}
%where $ l_\delta=\left (\sum_{j=1}^mV_{\delta,Z,j}-q\right )^{p}_+-\sum_{j=1}^m(V_{\delta, z_j, \hat{q}_j}-\hat{q}_j)^{p}_+  $
%and $ R_\delta(\omega_\delta):=\left (\sum_{j=1}^mV_{\delta,Z,j}+\omega_\delta-q\right )^{p}_+-\left (\sum_{j=1}^mV_{\delta,Z,j}-q\right )^{p}_+-p\left (\sum_{j=1}^mV_{\delta,Z,j}-q\right )^{p-1}_+\omega_\delta. $

In this section, we look for  solutions $ \omega\in E_{\varepsilon,Z}  $ of the following nonlinear equation
\begin{equation}\label{216}
Q_\varepsilon L_\varepsilon \omega=Q_\varepsilon l_\varepsilon+Q_\varepsilon R_\varepsilon(\omega),
\end{equation}
or equivalently,
\begin{equation*}
\begin{split}
\omega=T_\varepsilon(\omega):=(Q_\varepsilon L_\varepsilon)^{-1}Q_\varepsilon l_\varepsilon+(Q_\varepsilon L_\varepsilon)^{-1}Q_\varepsilon R_\varepsilon(\omega).
\end{split}
\end{equation*}
%And then we can reduce \eqref{215} to a finite dimensional problem.
We have
\begin{proposition}\label{exist and uniq of w}
There exists $ \varepsilon_0>0, $ such that for any $ \gamma\in (0,1) $, $ 0<\varepsilon<\varepsilon_0 $ and $ Z\in \Lambda_{\varepsilon,N} $, \eqref{216} has the unique solution $ \omega_{\varepsilon,Z}\in E_{\varepsilon,Z} $ with
\begin{equation*}
||\omega_{\varepsilon,Z}||_{L^\infty(\Omega)}=O\left( \varepsilon^\gamma \right).
\end{equation*}
\end{proposition}

\begin{proof}
It follows from  Lemma \ref{lemA-5} that for $ L $ sufficiently large and $ \varepsilon $ small,
\begin{equation*}
\left (V_{\varepsilon,Z}-q|\ln\varepsilon|\right )_+=0,\ \ \ \ \text{in}\ \Omega \backslash\cup_{i=1}^NB_{Ls_{\varepsilon,i}}(z_{i}).
\end{equation*}
Let $ \mathcal{N}= E_{\varepsilon,Z} \cap\{\omega\mid ||\omega||_{L^\infty(\Omega)}\leq \frac{1}{|\ln\varepsilon|^{1-\theta_0}}\}$ for some $ \theta_0\in(0,1). $ Then $ \mathcal{N} $ is complete under $ L^\infty $ norm and $ T_\varepsilon $ is a map from $ E_{\varepsilon,Z} $ to $ E_{\varepsilon,Z} $. We prove that $ T_\varepsilon $ is a contraction map from $ \mathcal{N} $ to $ \mathcal{N} $.

First, we claim that $ T_\varepsilon $ is a map from $ \mathcal{N} $ to $ \mathcal{N} $. For any $ \omega\in \mathcal{N} $, by Lemma \ref{lemA-5} we get that for $ L>1 $ large and  $ \varepsilon $ small,
\begin{equation*}
\left (V_{\varepsilon,Z}+\omega-q|\ln\varepsilon|\right )_+=0,\ \ \ \ \text{in}\ \Omega \backslash\cup_{i=1}^NB_{Ls_{\varepsilon,i}}(z_{i}).
\end{equation*}
So $ l_\varepsilon=R_\varepsilon(\omega)=0 $ in $ \Omega \backslash\cup_{i=1}^NB_{Ls_{\varepsilon,i}}(z_{i}). $ By the definition of $ Q_\varepsilon $,
\begin{equation*}
Q_\varepsilon l_\varepsilon+Q_\varepsilon R_\varepsilon(\omega)=0,\ \ \ \ \text{in}\ \Omega \backslash\cup_{i=1}^NB_{Ls_{\varepsilon,i}}(z_{i}).
\end{equation*}
Using Proposition \ref{coercive esti}, we obtain
\begin{equation*}
||T_\varepsilon(\omega)||_{L^\infty}=||(Q_\varepsilon L_\varepsilon)^{-1}(Q_\varepsilon l_\varepsilon+Q_\varepsilon R_\varepsilon(\omega))||_{L^\infty}\leq C\varepsilon^{-\frac{2}{p}}||Q_\varepsilon l_\varepsilon+Q_\varepsilon R_\varepsilon(\omega)||_{L^p}.
\end{equation*}
Note that $ ||Q_\varepsilon l_\varepsilon+Q_\varepsilon R_\varepsilon(\omega)||_{L^p}\leq C(|| l_\varepsilon||_{L^p}+ ||R_\varepsilon(\omega)||_{L^p}). $
It follows from \eqref{203}, the definition of $ l_\varepsilon, R_\varepsilon(\omega) $ and Lemma \ref{lemA-5} that
\begin{equation*}
\begin{split}
||l_\varepsilon||_{L^p}=&||\left (V_{\varepsilon,Z}-q|\ln\varepsilon|\right )^p_+-\sum_{j=1}^N\left (V_{\varepsilon,z_j,\hat{q}_{\varepsilon,j}}-\hat{q}_{\varepsilon,j}|\ln\varepsilon|\right )^p_+||_{L^p}\leq C\varepsilon^{\frac{2}{p}+\gamma},
\end{split}
\end{equation*}
and
\begin{equation*}
\begin{split}
||R_\varepsilon(\omega)||_{L^p}&=||\left (V_{\varepsilon,Z}+\omega-q|\ln\varepsilon|\right )^p_+-\left (V_{\varepsilon,Z}-q|\ln\varepsilon|\right )^p_+-p\left (V_{\varepsilon,Z}-q|\ln\varepsilon|\right )^{p-1}_+\omega||_{L^p}\\
&\leq C\varepsilon^{\frac{2}{p}}||\omega||_{L^\infty}^2.
\end{split}
\end{equation*}
Hence we get
\begin{equation}\label{217}
\begin{split}
||T_\varepsilon(\omega)||_{L^\infty}\leq C\varepsilon^{-\frac{2}{p}}\left( \varepsilon^{\frac{2}{p}+\gamma}+\varepsilon^{\frac{2}{p}}||\omega||_{L^\infty}^2\right)
\leq \frac{1}{|\ln\varepsilon|^{1-\theta_0}}.
\end{split}
\end{equation}
So $ T_\varepsilon $ is a map from $ \mathcal{N}$ to $ \mathcal{N} $.

Then we prove that  $ T_\varepsilon $ is a contraction map. For any $ \omega_1,\omega_2\in \mathcal{N} $,
\begin{equation*}
T_\varepsilon(\omega_1)-T_\varepsilon(\omega_2)=(Q_\varepsilon L_\varepsilon)^{-1}Q_\varepsilon(R_\varepsilon(\omega_1)-R_\varepsilon(\omega_2)).
\end{equation*}
Note that $ R_\varepsilon(\omega_1)=R_\varepsilon(\omega_2)=0 $ in $ \Omega \backslash\cup_{i=1}^NB_{Ls_{\varepsilon,i}}(z_{i}). $ By Proposition \ref{coercive esti} and the definition of $ \mathcal{N} $, for $ \varepsilon $ sufficiently small
\begin{equation*}
\begin{split}
||T_\varepsilon(\omega_1)-T_\varepsilon(\omega_2)||_{L^\infty}\leq &C\varepsilon^{-\frac{2}{p}}||R_\varepsilon(\omega_1)-R_\varepsilon(\omega_2)||_{L^p}\\
\leq &C\varepsilon^{-\frac{2}{p}}\varepsilon^{\frac{2}{p}}\left(||\omega_1||_{L^\infty}+||\omega_2||_{L^\infty} \right) ||\omega_1-\omega_2||_{L^\infty}\\
\leq& \frac{1}{2}||\omega_1-\omega_2||_{L^\infty}.
\end{split}
\end{equation*}
So $ T_\varepsilon $ is a contraction map.

To conclude,  %$ T_\varepsilon $ is a contraction map from $ \mathcal{N} $ to $ \mathcal{N} $ and thus 
there is a unique $ \omega_{\varepsilon,Z}\in \mathcal{N} $ such that $ \omega_{\varepsilon,Z}=T_\varepsilon(\omega_{\varepsilon,Z}) $. Moreover, by \eqref{217}  we have $ ||\omega_{\varepsilon,Z}||_{L^\infty(\Omega)}=O\left( \varepsilon^\gamma\right).  $


\end{proof}

%\begin{remark}\label{rk1}
%Indeed, since $ K,q \in C^\infty$ and $ p>1 $,  we can also check that $ \omega_{\delta,Z} $ is a $ C^1 $ map about $ Z $, see \cite{CLW,CPY} for example.
%\end{remark}

The result of Proposition \ref{exist and uniq of w} implies that there exists a unique solution $ \omega_{\varepsilon,Z}\in E_{\varepsilon,Z} $ to \eqref{216}. % satisfying $ ||\omega_{\delta,Z}||_{L^\infty(\Omega)}=O\left(\frac{\varepsilon^\gamma}{|\ln\varepsilon|}\right).  $
This implies that for some $ b_{j,\tilde{h}}=b_{j,\tilde{h}}(Z) $
\begin{equation}\label{5-000}
L_\varepsilon \omega_{\varepsilon,Z}= l_\varepsilon+  R_\varepsilon(\omega_{\varepsilon,Z})+\sum_{j=1}^N\sum_{\tilde{h}=1}^2b_{j,\tilde{h}}\left( -\varepsilon^2\text{div}\left (K(z_j)\nabla   \frac{\partial V_{\varepsilon, z_j, \hat{q}_{\varepsilon,j}}}{\partial x_{\tilde{h}}}\right )\right),
\end{equation}
or equivalently
\begin{equation}\label{5-00}
\begin{split}
-\varepsilon^2\text{div}&\left (K(x)\nabla (V_{\varepsilon,Z}+\omega_{\varepsilon,Z})\right )-\left (V_{\varepsilon,Z}+\omega_{\varepsilon,Z}-q|\ln\varepsilon|\right )^p_+\\
=&\sum_{j=1}^N\sum_{\tilde{h}=1}^2b_{j,\tilde{h}}\left( -\varepsilon^2\text{div}\left (K(z_j)\nabla   \frac{\partial V_{\varepsilon, z_j, \hat{q}_{\varepsilon,j}}}{\partial x_{\tilde{h}}}\right )\right).
\end{split}
\end{equation}

In the following result, we prove  the differentiability of $ \omega_{\varepsilon,Z}  $ with respect to the variable $ Z $, which will be used to show the $ C^1-$regularity of the finite-dimensional functional $ K_\varepsilon(Z) $. Similar to results in \cite{CPY,WYZ}, we give the $ L^\infty $ estimate of  $ \frac{\partial \omega_{\varepsilon,Z}}{\partial z_{i,\hbar}} $ and show that $ \omega_{\varepsilon,Z} $ is a $ C^1 $ map of $ Z $ in $ H^1_0 (\Omega) $.

\begin{proposition}\label{differ of w}
Let $ \omega_{\varepsilon,Z} $ be the function obtained in Proposition \ref{exist and uniq of w}. Then $ \omega_{\varepsilon,Z} $ is a $ C^1 $ map of $ Z $ in the
norm of $ H^1_0 (\Omega) $, and for any $ \gamma\in (0,1) $, $ l=1,\cdots,N $, $ \bar{h}=1,2 $
\begin{equation*}
\bigg|\bigg|\frac{\partial \omega_{\varepsilon,Z}}{\partial z_{l,\bar{h}}}\bigg|\bigg|_{L^\infty(\Omega)}=O\left (\frac{1}{\varepsilon^{1-\gamma}}\right ).
\end{equation*}

\end{proposition}

\begin{proof}
Note that from Lemma \ref{Green expansion2}, the regular part of Green's function $ \bar{S}_K(x,x)\in C^1(\Omega) $. Thus taking $ \frac{\partial}{\partial z_{l,\bar{h}}} $ in  \eqref{q_i choice}, we get
\begin{equation}\label{5-01}
\begin{split}
\frac{\partial \hat{q}_{\varepsilon,i}}{\partial z_{l,\bar{h}}}=&\frac{\partial q}{\partial x_{\bar{h}}}(z_i)\delta_{i,l}+\sum_{j\neq l}\frac{2\pi\hat{q}_{\varepsilon,j}\sqrt{\det K(z_j)}}{\ln s_{\varepsilon,j}}\frac{\partial G_K(z_l,z_j)}{\partial z_{l,\bar{h}}}+o(1)\sum_{j=1}^N\bigg|\frac{\partial \hat{q}_{\varepsilon,i}}{\partial z_{l,\bar{h}}}\bigg|+o(1)\\
=&O\left (|\ln\varepsilon|\right ),
\end{split}
\end{equation}
where we have used $ |\nabla_{z_l} G(z_l, z_j)|\leq\frac{ C}{|z_l-z_j|}\leq C|\ln\varepsilon| $ for $ Z\in \Lambda_{\varepsilon, N}. $ By the definition of $ V_{\varepsilon, z_j, \hat{q}_{\varepsilon,j}} $ and \eqref{5-01}, we get
\begin{equation}\label{5-02}
\bigg|\bigg|\frac{\partial V_{\varepsilon, z_j, \hat{q}_{\varepsilon,j}}}{\partial z_{l,\bar{h}}}\bigg|\bigg|_{L^\infty(\Omega)}=O\left (\frac{1}{\varepsilon}\right )+O\left (|\ln\varepsilon|^2\right )=O\left (\frac{1}{\varepsilon}\right ).
\end{equation}
Using the definition of $ H_{\varepsilon, z_j, \hat{q}_{\varepsilon,j}} $ in \eqref{eq3} and the $ L^p $-theory of elliptic equations, one has
\begin{equation}\label{5-03}
\bigg|\bigg|\frac{\partial H_{\varepsilon, z_j, \hat{q}_{\varepsilon,j}}}{\partial z_{l,\bar{h}}}\bigg|\bigg|_{W^{1,p}(\Omega)}\leq
\begin{cases}
\frac{C}{\varepsilon^{1-\frac{2}{p}}},\ \ &p>2,\\
C|\ln\varepsilon|,\ \ &p=2,\\
C,\ \ &1\leq p<2.
\end{cases}
\end{equation}
Combining \eqref{5-02} and \eqref{5-03}, we have
\begin{equation}\label{5-04}
\bigg|\bigg|\frac{\partial V_{\varepsilon, Z, j}}{\partial z_{l,\bar{h}}}\bigg|\bigg|_{L^\infty(\Omega)}=O\left (\frac{1}{\varepsilon}\right ).
\end{equation}

Now we calculate the $ L^\infty $ norm of $ \frac{\partial \omega_{\varepsilon,Z}}{\partial z_{l,\bar{h}}}. $ We first estimate $ b_{j,\tilde{h}} $ in \eqref{5-00}. Recall the definition of $ \eta_i $ in \eqref{cut off fun}. Using \eqref{5-00}%and the fact that $ supp(\eta_i)\subset D_{\frac{\delta}{2\sqrt{|\ln\varepsilon|}}}\left (C_{\frac{r_*}{\sqrt{|\ln\varepsilon|}}}\right ) $
, $ b_{j,h} $ is determined by
\begin{equation}\label{5-05}
\begin{split}
\sum_{j=1}^N&\sum_{\tilde{h}=1}^2b_{j,\tilde{h}}\int_\Omega\left( -\varepsilon^2\text{div}\left (K(z_j)\nabla   \frac{\partial V_{\varepsilon, z_j, \hat{q}_{\varepsilon,j}}}{\partial x_{\tilde{h}}}\right )\right)\left( \eta_i\frac{\partial V_{\varepsilon, Z,i}}{\partial x_\hbar}\right)\\
=&\int_\Omega\varepsilon^2\left (K(x)\nabla(V_{\varepsilon,Z}+\omega_{\varepsilon,Z})|\nabla\left( \eta_i\frac{\partial V_{\varepsilon, Z,i}}{\partial x_\hbar}\right)\right )-\int_\Omega\left (V_{\varepsilon,Z}+\omega_{\varepsilon,Z}-q|\ln\varepsilon|\right )^p_+\left( \eta_i\frac{\partial V_{\varepsilon, Z,i}}{\partial x_\hbar}\right)\\
=&\int_\Omega\varepsilon^2\left (K(x)\nabla V_{\varepsilon,Z}|\nabla\left( \eta_i\frac{\partial V_{\varepsilon, Z,i}}{\partial x_\hbar}\right)\right )-\int_\Omega\left (V_{\varepsilon,Z}+\omega_{\varepsilon,Z}-q|\ln\varepsilon|\right )^p_+\left( \eta_i\frac{\partial V_{\varepsilon, Z,i}}{\partial x_\hbar}\right)\\
=&\int_\Omega\left( \sum_{j=1}^N\left (V_{\varepsilon, z_j, \hat{q}_{\varepsilon,j}}-\hat{q}_{\varepsilon,j}|\ln\varepsilon|\right )^{p}_+-\left (V_{\varepsilon,Z}+\omega_{\varepsilon,Z}-q|\ln\varepsilon|\right )^{p}_+\right) \left( \eta_i\frac{\partial V_{\varepsilon, Z,i}}{\partial x_\hbar}\right),
\end{split}
\end{equation}
where we have used $ \omega_{\varepsilon,Z}\in E_{\varepsilon,Z} $. By Lemma \ref{lemA-5},
\begin{equation*}
\begin{split}
\int_\Omega&\left( \sum_{j=1}^N\left (V_{\varepsilon, z_j, \hat{q}_{\varepsilon,j}}-\hat{q}_{\varepsilon,j}|\ln\varepsilon|\right )^{p}_+-\left (V_{\varepsilon,Z}+\omega_{\varepsilon,Z}-q|\ln\varepsilon|\right )^{p}_+\right) \left( \eta_i\frac{\partial V_{\varepsilon, Z,i}}{\partial x_\hbar}\right)\\
&=O\left (\left (\varepsilon^\gamma+||\omega_{\varepsilon,Z}||_{L^\infty}\right )\int_{\cup_{j=1}^NB_{L\varepsilon}(z_j)}\bigg|\frac{\partial V_{\varepsilon, Z,i}}{\partial x_\hbar} \bigg|\right )=O\left (\varepsilon^{1+\gamma}\right ).
\end{split}
\end{equation*}
Thus combining this with \eqref{coef of C} and \eqref{5-05}, we get
\begin{equation}\label{5-06}
b_{j,\tilde{h}}=O\left ( \varepsilon^{1+\gamma} \right ).
\end{equation}

We now estimate  $  \frac{\partial b_{j,\tilde{h}}}{\partial z_{l,\bar{h}}} $. On the one hand, taking $ \frac{\partial}{\partial z_{l,\bar{h}}} $ in both sides of \eqref{5-05}, we obtain
\begin{equation}\label{5-07}
\begin{split}
\sum_{j=1}^N&\sum_{\tilde{h}=1}^2\frac{\partial b_{j,\tilde{h}}}{\partial z_{l,\bar{h}}}\int_\Omega\left( -\varepsilon^2\text{div}\left (K(z_j)\nabla   \frac{\partial V_{\varepsilon, z_j, \hat{q}_{\varepsilon,j}}}{\partial x_{\tilde{h}}}\right )\right)\left( \eta_i\frac{\partial V_{\varepsilon, Z,i}}{\partial x_\hbar}\right)\\
=&-\sum_{j=1}^N\sum_{\tilde{h}=1}^2b_{j,\tilde{h}} \frac{\partial }{\partial z_{l,\bar{h}}}\left\{  \int_\Omega\left( -\varepsilon^2\text{div}\left (K(z_j)\nabla   \frac{\partial V_{\varepsilon, z_j, \hat{q}_{\varepsilon,j}}}{\partial x_{\tilde{h}}}\right )\right)\left( \eta_i\frac{\partial V_{\varepsilon, Z,i}}{\partial x_\hbar}\right)\right\} \\
&+\frac{\partial }{\partial z_{l,\bar{h}}}\left\{\int_\Omega\left( \sum_{j=1}^N\left (V_{\varepsilon, z_j, \hat{q}_{\varepsilon,j}}-\hat{q}_{\varepsilon,j}|\ln\varepsilon|\right )^{p}_+-\left (V_{\varepsilon,Z}+\omega_{\varepsilon,Z}-q|\ln\varepsilon|\right )^{p}_+\right) \left( \eta_i\frac{\partial V_{\varepsilon, Z,i}}{\partial x_\hbar}\right)\right\}.
\end{split}
\end{equation}
Note that from \eqref{5-02}, \eqref{5-03} and \eqref{5-04},
\begin{equation*}
 \frac{\partial }{\partial z_{l,\bar{h}}}\left\{  \int_\Omega\left( -\varepsilon^2\text{div}\left (K(z_j)\nabla   \frac{\partial V_{\varepsilon, z_j, \hat{q}_{\varepsilon,j}}}{\partial x_{\tilde{h}}}\right )\right)\left( \eta_i\frac{\partial V_{\varepsilon, Z,i}}{\partial x_\hbar}\right)\right\}=O\left( \frac{1}{\varepsilon}\right)
\end{equation*}
and
\begin{equation*}
\begin{split}
 \int_\Omega&\left( \sum_{j=1}^N\left (V_{\varepsilon, z_j, \hat{q}_{\varepsilon,j}}-\hat{q}_{\varepsilon,j}|\ln\varepsilon|\right )^{p}_+-\left (V_{\varepsilon,Z}+\omega_{\varepsilon,Z}-q|\ln\varepsilon|\right )^{p}_+\right)  \frac{\partial }{\partial z_{l,\bar{h}}}\left( \eta_i\frac{\partial V_{\varepsilon, Z,i}}{\partial x_\hbar}\right) \\
 &=O\left (\left ( \varepsilon^\gamma+||\omega_{\varepsilon,Z}||_{L^\infty}\right )\int_{\cup_{j=1}^NB_{L\varepsilon}(z_j)}\bigg|\frac{\partial^2 V_{\varepsilon, Z,i}}{\partial z_{l,\bar{h}}\partial x_\hbar} \bigg|+|\ln\varepsilon|^2\bigg|\frac{\partial V_{\varepsilon, Z,i}}{\partial x_\hbar} \bigg|\right )=O\left (\varepsilon^{\gamma}\right ).
\end{split}
\end{equation*}
Taking these into \eqref{5-07}, we obtain
\begin{equation}\label{5-08}
\begin{split}
 \frac{\partial b_{j,\tilde{h}}}{\partial z_{l,\bar{h}}}=&O\left (\frac{ |b_{j,\tilde{h}}|}{\varepsilon}+\varepsilon^{\gamma}\right )\\
 &+O\left(  \int_\Omega\frac{\partial }{\partial z_{l,\bar{h}}}\left( \sum_{j=1}^N\left (V_{\varepsilon, z_j, \hat{q}_{\varepsilon,j}}-\hat{q}_{\varepsilon,j}|\ln\varepsilon|\right )^{p}_+-\left (V_{\varepsilon,Z}+\omega_{\varepsilon,Z}-q|\ln\varepsilon|\right )^{p}_+\right) \left( \eta_i\frac{\partial V_{\varepsilon, Z,i}}{\partial x_\hbar}\right) \right) .
\end{split}
\end{equation}
Using \eqref{203}, \eqref{5-01}, \eqref{5-02} and \eqref{5-03}, we get
\begin{equation*}\label{5-09}
\begin{split}
\frac{\partial }{\partial z_{l,\bar{h}}}&\left( \sum_{j=1}^N\left (V_{\varepsilon, z_j, \hat{q}_{\varepsilon,j}}-\hat{q}_{\varepsilon,j}|\ln\varepsilon|\right )^{p}_+-\left (V_{\varepsilon,Z}+\omega_{\varepsilon,Z}-q|\ln\varepsilon|\right )^{p}_+\right)\\
=& p\sum_{j=1}^N\left (V_{\varepsilon, z_j, \hat{q}_{\varepsilon,j}}-\hat{q}_{\varepsilon,j}|\ln\varepsilon|\right )^{p-1}_+\frac{\partial }{\partial z_{l,\bar{h}}}\left (V_{\varepsilon, z_j, \hat{q}_{\varepsilon,j}}-\hat{q}_{\varepsilon,j}|\ln\varepsilon|\right )\\
&-p\left (V_{\varepsilon,Z}+\omega_{\varepsilon,Z}-q|\ln\varepsilon|\right )^{p-1}_+\frac{\partial }{\partial z_{l,\bar{h}}}\left(V_{\varepsilon,Z}+\omega_{\varepsilon,Z} \right)  \\
=&-p\left (V_{\varepsilon,Z}+\omega_{\varepsilon,Z}-q|\ln\varepsilon|\right )^{p-1}_+\frac{\partial \omega_{\varepsilon,Z} }{\partial z_{l,\bar{h}}}\\
&+p\sum_{j=1}^N\left( \left (V_{\varepsilon, z_j, \hat{q}_{\varepsilon,j}}-\hat{q}_{\varepsilon,j}|\ln\varepsilon|\right )^{p-1}_+-\left (V_{\varepsilon,Z}+\omega_{\varepsilon,Z}-q|\ln\varepsilon|\right )^{p-1}_+\right)\frac{\partial V_{\varepsilon, z_j, \hat{q}_{\varepsilon,j}}}{\partial z_{l,\bar{h}}} \\
&-p\sum_{j=1}^N\left (V_{\varepsilon,Z}+\omega_{\varepsilon,Z}-q|\ln\varepsilon|\right )^{p-1}_+\frac{\partial H_{\varepsilon, z_j, \hat{q}_{\varepsilon,j}}}{\partial z_{l,\bar{h}}}-p|\ln\varepsilon|\sum_{j=1}^N\left (V_{\varepsilon, z_j, \hat{q}_{\varepsilon,j}}-\hat{q}_{\varepsilon,j}|\ln\varepsilon|\right )^{p-1}_+\frac{\partial \hat{q}_{\varepsilon,j} }{\partial z_{l,\bar{h}}} \\
=&-p\left (V_{\varepsilon,Z}+\omega_{\varepsilon,Z}-q|\ln\varepsilon|\right )^{p-1}_+\frac{\partial \omega_{\varepsilon,Z} }{\partial z_{l,\bar{h}}}+O\left( \frac{1}{\varepsilon^{1-\gamma}}\right),
\end{split}
\end{equation*}
from which we deduce,
\begin{equation}\label{5-10}
\begin{split}
\int_\Omega&\frac{\partial }{\partial z_{l,\bar{h}}}\left( \sum_{j=1}^N\left (V_{\varepsilon, z_j, \hat{q}_{\varepsilon,j}}-\hat{q}_{\varepsilon,j}|\ln\varepsilon|\right )^{p}_+-\left (V_{\varepsilon,Z}+\omega_{\varepsilon,Z}-q|\ln\varepsilon|\right )^{p}_+\right) \left( \eta_i\frac{\partial V_{\varepsilon, Z,i}}{\partial x_\hbar}\right)\\
&=-p\int_\Omega \left (V_{\varepsilon,Z}+\omega_{\varepsilon,Z}-q|\ln\varepsilon|\right )^{p-1}_+\frac{\partial \omega_{\varepsilon,Z} }{\partial z_{l,\bar{h}}}\left( \eta_i\frac{\partial V_{\varepsilon, Z,i}}{\partial x_\hbar}\right)+O\left( \varepsilon^\gamma\right).
\end{split}
\end{equation}
Inserting \eqref{5-10} into \eqref{5-08}, we obtain
\begin{equation}\label{5-11}
\begin{split}
\frac{\partial b_{j,\tilde{h}}}{\partial z_{l,\bar{h}}}=&O\left( \int_\Omega \left (V_{\varepsilon,Z}+\omega_{\varepsilon,Z}-q|\ln\varepsilon|\right )^{p-1}_+\frac{\partial \omega_{\varepsilon,Z} }{\partial z_{l,\bar{h}}}\left( \eta_i\frac{\partial V_{\varepsilon, Z,i}}{\partial x_\hbar}\right)\right) +O\left( \varepsilon^\gamma\right).
\end{split}
\end{equation}

On the other hand, we note that   $ \int_{\Omega}\varepsilon^2\left (K(x) \nabla\omega_{\varepsilon,Z}| \nabla \left( \eta_i\frac{\partial V_{\varepsilon, Z,i}}{\partial x_\hbar}\right)\right )=0$. Taking $ \frac{\partial}{\partial z_{l,\bar{h}}} $ into both sides of this equality, we have
\begin{equation}\label{5-12}
\int_{\Omega}\varepsilon^2 \left (K(x) \nabla \frac{\partial \omega_{\varepsilon,Z} }{\partial z_{l,\bar{h}}}| \nabla \left( \eta_i\frac{\partial V_{\varepsilon, Z,i}}{\partial x_\hbar}\right)\right )=-\int_{\Omega} \varepsilon^2 \left (K(x) \nabla\omega_{\varepsilon,Z}| \nabla \frac{\partial  }{\partial z_{l,\bar{h}}}\left( \eta_i\frac{\partial V_{\varepsilon, Z,i}}{\partial x_\hbar}\right)\right ) .
\end{equation}
One computes directly  that
\begin{equation}\label{5-13}
\begin{split}
-\int_{\Omega} \varepsilon^2 \left (K(x) \nabla\omega_{\varepsilon,Z}| \nabla \frac{\partial  }{\partial z_{l,\bar{h}}}\left( \eta_i\frac{\partial V_{\varepsilon, Z,i}}{\partial x_\hbar}\right)\right )=\int_{\Omega} \varepsilon^2 \text{div} \left (K(x) \nabla\omega_{\varepsilon,Z}  \right )\frac{\partial  }{\partial z_{l,\bar{h}}}\left( \eta_i\frac{\partial V_{\varepsilon, Z,i}}{\partial x_\hbar}\right).
\end{split}
\end{equation}
Since $ ||l_\varepsilon+R_\varepsilon(\omega_{\varepsilon,Z})||_{L^p}\leq C\varepsilon^{\frac{2}{p}+\gamma} $, by \eqref{5-000}  and \eqref{5-06}, we get
\begin{equation*}
L_\varepsilon \omega_{\varepsilon,Z}= l_\varepsilon+  R_\varepsilon(\omega_{\varepsilon,Z})+\sum_{j=1}^N\sum_{\tilde{h}=1}^2b_{j,\tilde{h}}\left( -\varepsilon^2\text{div}\left (K(z_j)\nabla   \frac{\partial V_{\varepsilon, z_j, \hat{q}_{\varepsilon,j}}}{\partial x_{\tilde{h}}}\right )\right)=O\left( \varepsilon^{\frac{2}{p}+\gamma}\right) \  \ \ \text{in}\ L^p(\Omega),
\end{equation*}
which implies that
\begin{equation*}
\begin{split}
-\varepsilon^2 \text{div}\left (K(x) \nabla\omega_{\varepsilon,Z}  \right )=&L_\varepsilon \omega_{\varepsilon,Z}
+p\left( V_{\varepsilon,Z}-q|\ln\varepsilon|\right)^{p-1}_+\omega_{\varepsilon,Z}
= O\left( \varepsilon^{\frac{2}{p}+\gamma}\right) \  \ \ \text{in}\ L^p(\Omega).
\end{split}
\end{equation*}
Taking this into \eqref{5-13}, using the definition of $ V_{\varepsilon, z_j, \hat{q}_{\varepsilon,j}} $ and \eqref{5-03}, one has
\begin{equation}\label{5-14}
\begin{split}
-\int_{\Omega} \varepsilon^2 \left (K(x) \nabla\omega_{\varepsilon,Z}| \nabla \frac{\partial  }{\partial z_{l,\bar{h}}}\left( \eta_i\frac{\partial V_{\varepsilon, Z,i}}{\partial x_\hbar}\right)\right )=O\left( \varepsilon^\gamma\right).
\end{split}
\end{equation}
We also note that
\begin{equation}\label{5-15}
\begin{split}
\int_{\Omega}\varepsilon^2& \left (K(x) \nabla \frac{\partial \omega_{\varepsilon,Z} }{\partial z_{l,\bar{h}}}| \nabla \left( \eta_i\frac{\partial V_{\varepsilon, Z,i}}{\partial x_\hbar}\right)\right )=\int_{\Omega}-\varepsilon^2  \text{div}\left (K(x) \nabla \left( \eta_i\frac{\partial V_{\varepsilon, Z,i}}{\partial x_\hbar}\right)\right )\frac{\partial \omega_{\varepsilon,Z} }{\partial z_{l,\bar{h}}}\\
=&\int_{\Omega}\bigg\{\eta_ip\left (V_{\varepsilon, z_i, \hat{q}_{\varepsilon,i}}-\hat{q}_{\varepsilon,i}|\ln\varepsilon|\right )^{p-1}_+\frac{\partial V_{\varepsilon, z_i, \hat{q}_{\varepsilon,i}}}{\partial x_\hbar}+\eta_i\varepsilon^2\text{div}\left (\frac{\partial K(x)}{\partial x_\hbar}\nabla V_{\varepsilon, Z,i}  \right ) \\
& -2 \varepsilon^2\left (K(x)\nabla\eta_i|\nabla   \frac{\partial V_{\varepsilon, Z,i}}{\partial x_\hbar}  \right ) -\varepsilon^2\text{div}\left (K(x)\nabla   \eta_i  \right )\frac{\partial V_{\varepsilon, Z,i}}{\partial x_\hbar} \bigg\}\frac{\partial \omega_{\varepsilon,Z} }{\partial z_{l,\bar{h}}}\\
=& p\int_{\Omega}\eta_i\frac{\partial \omega_{\varepsilon,Z} }{\partial z_{l,\bar{h}}} \left (V_{\varepsilon, z_i, \hat{q}_{\varepsilon,i}}-\hat{q}_{\varepsilon,i}|\ln\varepsilon|\right )^{p-1}_+\frac{\partial V_{\varepsilon, z_i, \hat{q}_{\varepsilon,i}}}{\partial x_\hbar}+O\left( \varepsilon^2|\ln\varepsilon|^{4}\bigg|\bigg|\frac{\partial \omega_{\varepsilon,Z} }{\partial z_{l,\bar{h}}}\bigg|\bigg|_{L^\infty}\right).
\end{split}
\end{equation}
Combining \eqref{5-12} with \eqref{5-14} and \eqref{5-15}, we obtain
\begin{equation}\label{5-16}
\begin{split}
\int_{\Omega}\eta_i\frac{\partial \omega_{\varepsilon,Z} }{\partial z_{l,\bar{h}}} \left (V_{\varepsilon, z_i, \hat{q}_{\varepsilon,i}}-\hat{q}_{\varepsilon,i}|\ln\varepsilon|\right )^{p-1}_+\frac{\partial V_{\varepsilon, z_i, \hat{q}_{\varepsilon,i}}}{\partial x_\hbar}=O\left(  \varepsilon^\gamma +\varepsilon^2|\ln\varepsilon|^{4}\bigg|\bigg|\frac{\partial \omega_{\varepsilon,Z} }{\partial z_{l,\bar{h}}}\bigg|\bigg|_{L^\infty}\right).
\end{split}
\end{equation}
%\begin{equation*}
%\begin{split}
%-\int_{\Omega}&\omega_{\delta,Z}\frac{\partial  }{\partial z_{l,\bar{h}}}\left( -\delta^2\nabla\cdot\left (K(x)  \nabla \left( \eta_i\frac{\partial V_{\delta, Z,i}}{\partial x_\hbar}\right)\right ) \right)\\
%=&-\int_{\Omega} \omega_{\delta,Z}\frac{\partial  }{\partial z_{l,\bar{h}}}\bigg\{ \eta_ip\left (V_{\delta, z_j, \hat{q}_{\delta,j}}-\hat{q}_{\delta,j}\right )^{p-1}_+\frac{\partial V_{\delta, z_j, \hat{q}_{\delta,j}}}{\partial x_h}+\eta_i\delta^2\nabla\cdot\left (\frac{\partial K(x)}{\partial x_\hbar}  \nabla   V_{\delta, Z,i} \right )\\
%& -\delta^2\left( K(x)\nabla\eta_i|\nabla \frac{\partial V_{\delta, Z,i}}{\partial x_\hbar}\right) \bigg\}
%\end{split}
%\end{equation*}

Taking \eqref{5-16} into \eqref{5-11} and using \eqref{203} and Proposition \ref{exist and uniq of w}, we conclude that
\begin{equation}\label{5-17}
\begin{split}
\frac{\partial b_{j,\tilde{h}}}{\partial z_{l,\bar{h}}}=& O\left( \varepsilon^\gamma+ \varepsilon^{1+\gamma}\bigg|\bigg|\frac{\partial \omega_{\varepsilon,Z} }{\partial z_{l,\bar{h}}}\bigg|\bigg|_{L^\infty}\right).
\end{split}
\end{equation}

Now we calculate the $ L^\infty $ norm of $ \frac{\partial \omega_{\varepsilon,Z} }{\partial z_{l,\bar{h}}} $. Taking $ \frac{\partial}{\partial z_{l,\bar{h}}} $ in both sides of \eqref{5-00}, we get
%\begin{equation}
%\begin{split}
%L_\delta \omega_{\delta,Z}= l_\delta+  R_\delta(\omega_{\delta,Z})+\sum_{j=1}^m\sum_{h=1}^2b_{j,h}\left( -\delta^2\text{div}\left (K(z_j)\nabla   \frac{\partial V_{\delta, z_j, \hat{q}_{\delta,j}}}{\partial x_h}\right )\right),
%\end{split}
%\end{equation}
\begin{equation}\label{5-18}
\begin{split}
-\varepsilon^2\text{div}&\left( K(x)\nabla \frac{\partial\omega_{\varepsilon,Z}}{\partial z_{l,\bar{h}}}\right)-p\left( V_{\varepsilon,Z}+\omega_{\varepsilon,Z}-q|\ln\varepsilon|\right)^{p-1}_+\frac{\partial\omega_{\varepsilon,Z}}{\partial z_{l,\bar{h}}}\\
=&\varepsilon^2\text{div}\left( K(x)\nabla \frac{\partial V_{\varepsilon,Z}}{\partial z_{l,\bar{h}}}\right)+p\left( V_{\varepsilon,Z}+\omega_{\varepsilon,Z}-q|\ln\varepsilon|\right)^{p-1}_+\frac{\partial V_{\varepsilon,Z}}{\partial z_{l,\bar{h}}}\\
% L_\delta \frac{\partial\omega_{\delta,Z}}{\partial z_{l,\bar{h}}}
%=& p(p-1)\left( V_{\delta,Z}-q\right)^{p-2}_+\frac{\partial V_{\delta,Z}}{\partial z_{l,\bar{h}}}\omega_{\delta,Z}+\frac{\partial l_\delta}{\partial z_{l,\bar{h}}}+  \frac{\partial R_\delta(\omega_{\delta,Z})}{\partial z_{l,\bar{h}}}\\
&+\sum_{j=1}^N\sum_{\tilde{h}=1}^2\frac{\partial b_{j,\tilde{h}}}{\partial z_{l,\bar{h}}}p\left(   V_{\varepsilon, z_j, \hat{q}_{\varepsilon,j}}-\hat{q}_{\varepsilon,j}|\ln\varepsilon| \right )^{p-1}_+\frac{\partial V_{\varepsilon, z_j, \hat{q}_{\varepsilon,j}}}{\partial x_{\tilde{h}}}\\
&+\sum_{j=1}^N\sum_{\tilde{h}=1}^2b_{j,\tilde{h}}\frac{\partial }{\partial z_{l,\bar{h}}}\left( p\left(   V_{\varepsilon, z_j, \hat{q}_{\varepsilon,j}}-\hat{q}_{\varepsilon,j}|\ln\varepsilon| \right )^{p-1}_+\frac{\partial V_{\varepsilon, z_j, \hat{q}_{\varepsilon,j}}}{\partial x_{\tilde{h}}}\right) .
\end{split}
\end{equation}
Note that the function $ \frac{\partial\omega_{\varepsilon,Z}}{\partial z_{l,\bar{h}}} $ may not be in $ E_{\varepsilon,Z} $. We make the following decomposition:
\begin{equation}\label{5-19}
\frac{\partial\omega_{\varepsilon,Z}}{\partial z_{l,\bar{h}}}=\omega_{\varepsilon}^*+\sum_{j=1}^N\sum_{\tilde{h}=1}^2C_{j,\tilde{h}}\zeta_j\frac{\partial V_{\varepsilon, z_j, \hat{q}_{\varepsilon,j}}}{\partial x_{\tilde{h}}},
\end{equation}
where $ \omega_{\varepsilon}^*\in E_{\varepsilon,Z} $,  $\zeta_j(x):=\eta\left (\frac{|T_{z_j}\left( x-z_j\right) |}{s_{\varepsilon,j}}\right )$ and $ C_{j,\tilde{h}} $ is determined by
\begin{equation*}
\begin{split}
\sum_{j=1}^N\sum_{\tilde{h}=1}^2&C_{j,\tilde{h}}\int_{\Omega}\zeta_j\frac{\partial V_{\varepsilon, z_j, \hat{q}_{\varepsilon,j}}}{\partial x_{\tilde{h}}}\cdot\left( -\varepsilon^2\text{div}\left( K(x)\nabla\left(\eta_i\frac{\partial V_{\varepsilon,Z,i}}{\partial x_\hbar} \right) \right)\right) \\
=& \int_{\Omega}\frac{\partial\omega_{\varepsilon,Z}}{\partial z_{l,\bar{h}}}\cdot\left( -\varepsilon^2\text{div}\left( K(x)\nabla\left(\eta_i\frac{\partial V_{\varepsilon,Z,i}}{\partial x_\hbar} \right)\right)\right)  \ \ \ \ i=1,\cdots,N,\ \hbar=1,2.
\end{split}
\end{equation*}
Note that
\begin{equation*}
\int_{\Omega}\zeta_j\frac{\partial V_{\varepsilon, z_j, \hat{q}_{\varepsilon,j}}}{\partial x_{\tilde{h}}}\cdot\left( -\varepsilon^2\text{div}\left( K(x)\nabla\left(\eta_i\frac{\partial V_{\varepsilon,Z,i}}{\partial x_\hbar} \right) \right)\right) = (\tilde{M}_i)_{\tilde{h},\hbar}\delta_{i,j}+o(1),
\end{equation*}
where $ \tilde{M}_{i} $ are $ m $ positive definite matrices. Combining this with \eqref{5-12} and \eqref{5-14}, we obtain
\begin{equation}\label{5-20}
C_{j,\tilde{h}}=O\left( \varepsilon^\gamma \right).
\end{equation}
Taking \eqref{5-19} into \eqref{5-18}, we get
\begin{equation}\label{5-21}
\begin{split}
-\varepsilon^2&\text{div}\left( K(x)\nabla  \omega_{\varepsilon}^*\right)-p\left( V_{\varepsilon,Z} -q|\ln\varepsilon|\right)^{p-1}_+\omega_{\varepsilon}^*\\
=&\sum_{j=1}^N\sum_{\tilde{h}=1}^2C_{j,\tilde{h}} \varepsilon^2\text{div}\left( K(x)\nabla \left( \zeta_j\frac{\partial V_{\varepsilon, z_j, \hat{q}_{\varepsilon,j}}}{\partial x_{\tilde{h}}}\right) \right)+ p\sum_{j=1}^N\sum_{\tilde{h}=1}^2C_{j,\tilde{h}}\left( V_{\varepsilon,Z} -q|\ln\varepsilon|\right)^{p-1}_+\zeta_j\frac{\partial V_{\varepsilon, z_j, \hat{q}_{\varepsilon,j}}}{\partial x_{\tilde{h}}}\\
&+p\left( \left( V_{\varepsilon,Z}+\omega_{\varepsilon,Z}-q|\ln\varepsilon|\right)^{p-1}_+-\left( V_{\varepsilon,Z} -q|\ln\varepsilon|\right)^{p-1}_+\right) \frac{\partial\omega_{\varepsilon,Z}}{\partial z_{l,\bar{h}}}\\
&-p\sum_{j=1}^N\left (V_{\varepsilon, z_j, \hat{q}_{\varepsilon,j}}-\hat{q}_{\varepsilon,j}|\ln\varepsilon|\right )^{p-1}_+\frac{\partial }{\partial z_{l,\bar{h}}}\left (V_{\varepsilon, z_j, \hat{q}_{\varepsilon,j}}-\hat{q}_{\varepsilon,j}|\ln\varepsilon|\right )   \\
% L_\delta \frac{\partial\omega_{\delta,Z}}{\partial z_{l,\bar{h}}}
%=& p(p-1)\left( V_{\delta,Z}-q\right)^{p-2}_+\frac{\partial V_{\delta,Z}}{\partial z_{l,\bar{h}}}\omega_{\delta,Z}+\frac{\partial l_\delta}{\partial z_{l,\bar{h}}}+  \frac{\partial R_\delta(\omega_{\delta,Z})}{\partial z_{l,\bar{h}}}\\
&+p\left (V_{\varepsilon,Z}+\omega_{\varepsilon,Z}-q|\ln\varepsilon|\right )^{p-1}_+\frac{\partial V_{\varepsilon,Z}}{\partial z_{l,\bar{h}}}+\sum_{j=1}^N\sum_{\tilde{h}=1}^2\frac{\partial b_{j,\tilde{h}}}{\partial z_{l,\bar{h}}}p\left(   V_{\varepsilon, z_j, \hat{q}_{\varepsilon,j}}-\hat{q}_{\varepsilon,j}|\ln\varepsilon| \right )^{p-1}_+\frac{\partial V_{\varepsilon, z_j, \hat{q}_{\varepsilon,j}}}{\partial x_{\tilde{h}}}\\
&+\sum_{j=1}^N\sum_{\tilde{h}=1}^2b_{j,\tilde{h}}\frac{\partial }{\partial z_{l,\bar{h}}}\left( p\left(   V_{\varepsilon, z_j, \hat{q}_{\varepsilon,j}}-\hat{q}_{\varepsilon,j}|\ln\varepsilon| \right )^{p-1}_+\frac{\partial V_{\varepsilon, z_j, \hat{q}_{\varepsilon,j}}}{\partial x_{\tilde{h}}}\right) \\
=&A_1+A_2+A_3+A_4+A_5+A_6+A_7.
\end{split}
\end{equation}
By \eqref{5-01}, \eqref{5-06}, \eqref{5-17} and \eqref{5-20}, one computes directly that
\begin{equation*}
\begin{split}
%\sum_{j=1}^m&\sum_{h=1}^2C_{j,h} \delta^2\text{div}\left( K(x)\nabla \left( \zeta_j\frac{\partial V_{\delta, z_j, \hat{q}_{\delta,j}}}{\partial x_h}\right) \right)+ \sum_{j=1}^m\sum_{h=1}^2C_{j,h}p\left( V_{\delta,Z} -q\right)^{p-1}_+\zeta_j\frac{\partial V_{\delta, z_j, \hat{q}_{\delta,j}}}{\partial x_h} \\
A_1+A_2=&\left( |C_{j,\tilde{h}}|\varepsilon^{\frac{2}{p}-1}\right)=O\left(  \varepsilon^{\gamma+\frac{2}{p}-1}\right)\ \ \text{in}\ L^p(\Omega),
\end{split}
\end{equation*}
\begin{equation*}
\begin{split}
%p\left( \left( V_{\delta,Z}+\omega_{\delta,Z}-q\right)^{p-1}_+-\left( V_{\delta,Z} -q\right)^{p-1}_+\right) \frac{\partial\omega_{\delta,Z}}{\partial z_{l,\bar{h}}}
A_3=O\left(  \varepsilon^{\gamma+\frac{2}{p}}\bigg|\bigg|\frac{\partial \omega_{\varepsilon,Z} }{\partial z_{l,\bar{h}}}\bigg|\bigg|_{L^\infty}\right)\ \ \text{in}\ L^p(\Omega),
\end{split}
\end{equation*}
\begin{equation*}
\begin{split}
%-p&\sum_{j=1}^m\left (V_{\delta, z_j, \hat{q}_{\delta,j}}-\hat{q}_{\delta,j}\right )^{p-1}_+\frac{\partial }{\partial z_{l,\bar{h}}}\left (V_{\delta, z_j, \hat{q}_{\delta,j}}-\hat{q}_{\delta,j}\right )+p\left (V_{\delta,Z}+\omega_{\delta,Z}-q\right )^{p-1}_+\frac{\partial V_{\delta,Z}}{\partial z_{l,\bar{h}}}\\
A_4+A_5=&O\left( \varepsilon^{\frac{2}{p}}|\ln\varepsilon|^{4}+  \varepsilon^{\gamma+\frac{2}{p}-1}\right) =O\left(  \varepsilon^{\gamma+\frac{2}{p}-1}\right)\ \ \text{in}\ L^p(\Omega),
\end{split}
\end{equation*}
\begin{equation*}
\begin{split}
%&\sum_{j=1}^m\sum_{h=1}^2\frac{\partial b_{j,h}}{\partial z_{l,\bar{h}}}p\left(   V_{\delta, z_j, \hat{q}_{\delta,j}}-\hat{q}_{\delta,j} \right )^{p-1}_+\frac{\partial V_{\delta, z_j, \hat{q}_{\delta,j}}}{\partial x_h}\\
A_6&=O\left(   \frac{\varepsilon^{ \frac{2}{p}}}{\varepsilon}\bigg|\frac{\partial b_{j,\tilde{h}}}{\partial z_{l,\bar{h}}}\bigg|\right)=O\left(   \varepsilon^{ \frac{2}{p}-1}\left(   \varepsilon^\gamma +  \varepsilon^{1+\gamma}   \bigg|\bigg|\frac{\partial \omega_{\varepsilon,Z} }{\partial z_{l,\bar{h}}}\bigg|\bigg|_{L^\infty}\right) \right)\ \ \text{in}\ L^p(\Omega),
\end{split}
\end{equation*}
\begin{equation*}
\begin{split}
%&\sum_{j=1}^m\sum_{h=1}^2b_{j,h}\frac{\partial }{\partial z_{l,\bar{h}}}\left( p\left(   V_{\delta, z_j, \hat{q}_{\delta,j}}-\hat{q}_{\delta,j} \right )^{p-1}_+\frac{\partial V_{\delta, z_j, \hat{q}_{\delta,j}}}{\partial x_h}\right)
A_7=O\left(  \varepsilon^{\gamma+\frac{2}{p}-1}\right)\ \ \text{in}\ L^p(\Omega).
\end{split}
\end{equation*}
Combining these with Proposition \ref{coercive esti} and \eqref{5-21}, we are led to
\begin{equation}\label{5-22}
||\omega_{\varepsilon}^*||_{L^\infty }\leq C\varepsilon^{-\frac{2}{p}}||L_\varepsilon\omega_{\varepsilon}^*||_{L^p }\leq C\varepsilon^{-\frac{2}{p}} \left(\varepsilon^{\gamma+\frac{2}{p}-1}+ \varepsilon^{\gamma+\frac{2}{p}}\bigg|\bigg|\frac{\partial \omega_{\varepsilon,Z} }{\partial z_{l,\bar{h}}}\bigg|\bigg|_{L^\infty} \right).
\end{equation}
From the decomposition \eqref{5-19}, we have
\begin{equation*}
\bigg|\bigg|\frac{\partial \omega_{\varepsilon,Z} }{\partial z_{l,\bar{h}}}\bigg|\bigg|_{L^\infty}\leq ||\omega_{\varepsilon}^*||_{L^\infty }+O\left( \frac{1}{\varepsilon^{1-\gamma}}\right) .
\end{equation*}
Taking this into \eqref{5-22}, we obtain
\begin{equation*}
||\omega_{\varepsilon}^*||_{L^\infty }\leq C  \frac{1}{\varepsilon^{1-\gamma}},
\end{equation*}
from which we deduce, $ \big|\big|\frac{\partial \omega_{\varepsilon,Z} }{\partial z_{l,\bar{h}}}\big|\big|_{L^\infty(\Omega)}=O\left( \frac{1}{\varepsilon^{1-\gamma}} \right)  $.

Finally, we prove that $ \omega_{\varepsilon,Z} $ is a $ C^1 $ map of $ Z\in \Lambda_{\varepsilon, N} $ in $ H^1(\Omega) $. To prove the continuity of $ \omega_{\varepsilon,Z} $ of $ Z $, let $ Z_j\to Z_0 $. By Proposition \ref{exist and uniq of w}, $ \omega_{\varepsilon,Z_j}  $ is uniformly bounded in $ L^\infty(\Omega) $. Thus  using \eqref{5-000} and \eqref{5-06}, we conclude that $ ||\omega_{\varepsilon,Z_j}||_{H^1_0(\Omega)} $ is bounded by a constant $ C $  which is independent of $ j $. Then there is a subsequence (still denoted by $ Z_j $) such that
\begin{equation*}
\omega_{\varepsilon,Z_j}\rightarrow \omega^{**}\ \ \ \  \text{weakly in } H^1_0(\Omega)
\end{equation*}
and
\begin{equation*}
\omega_{\varepsilon,Z_j}\rightarrow \omega^{**}\ \ \ \  \text{strongly in } L^2(\Omega).
\end{equation*}
Using the equation again, we can get that
\begin{equation*}
\omega_{\varepsilon,Z_j}\rightarrow \omega^{**}\ \ \ \  \text{strongly in } H^1_0(\Omega),
\end{equation*}
from which  we deduce that $ \omega^{**}\in E_{\varepsilon,Z_0} $ and $ \omega^{**} $ satisfies \eqref{216} with $ Z_j $ replaced by $ Z_0 $.  By the uniqueness, we get $ \omega^{**}=\omega_{\varepsilon,Z_0} $ and hence  $ \omega_{\varepsilon,Z} $ is continuous with respect  to  $Z$ in the norm of $ H^1_0(\Omega) $. Moreover, using similar method as  Proposition 3.7 in  \cite{CPY}, we can show that  $ \frac{\partial \omega_{\varepsilon,Z}}{\partial z_{l,\tilde{h}}} $ is continuous with respect  to $ Z $ in $ H^1(\Omega) $. The proof is thus complete.
\end{proof}






\section{The finite-dimensional energy expansion}

Our aim in the rest part of this paper is to find $Z$ properly so that $V_{\varepsilon, Z}+\omega_{\varepsilon, Z}$ is a solution of \eqref{key equa}.  In this section we give expansion of the main term $I(V_{\varepsilon, Z})$ in term of $Z$. From Proposition \ref{exist and uniq of w}, given any $ \varepsilon $ small and $ Z\in \Lambda_{\varepsilon, N} $, there exists a unique $ \omega_{\varepsilon,Z}\in E_{\varepsilon,Z} $ satisfying $ Q_\varepsilon L_\varepsilon \omega_{\varepsilon,Z}=Q_\varepsilon l_\varepsilon+Q_\varepsilon R_\varepsilon(\omega_{\varepsilon,Z}), $
i.e., for some $ b_{j,\tilde{h}}=b_{j,\tilde{h}}(Z) $
\begin{equation*}
 L_\varepsilon \omega_{\varepsilon,Z}= l_\varepsilon+  R_\varepsilon(\omega_{\varepsilon,Z})+\sum_{j=1}^N\sum_{\tilde{h}=1}^2b_{j,\tilde{h}}\left( -\varepsilon^2\text{div}\left (K(z_j)\nabla   \frac{\partial V_{\varepsilon, z_j, \hat{q}_{\varepsilon,j}}}{\partial x_{\tilde{h}}}\right )\right).
\end{equation*}
Thus, it suffices to find $Z$ properly so that
\begin{equation*}
 b_{j,\tilde{h}}(Z)=0,\ \ \forall\ j=1,\cdots,N,\ \tilde{h}=1,2.
\end{equation*}

%In the following we find proper $ Z=Z(\delta) $ such that $ V_{\delta,Z}+\omega_{\delta,Z} $ is a solution of \eqref{111}. Note that
Define
\begin{equation}\label{functional I}
I_\varepsilon(u)=\frac{\varepsilon^2}{2}\int_{\Omega}\left( K(x)\nabla u|\nabla u\right)-\frac{1}{p+1}\int_{\Omega}\left (u-q|\ln\varepsilon|\right )^{p+1}_+,%\textbf{1}_{D_{\frac{\delta}{\sqrt{|\ln\varepsilon|}}}\left (C_{\frac{r_*}{\sqrt{|\ln\varepsilon|}}}\right )},
\end{equation}
and
\begin{equation*}
K_\varepsilon(Z)=I_\varepsilon(V_{\varepsilon,Z}+\omega_{\varepsilon,Z}).
\end{equation*}
%Then by Remark \ref{rk1}, we know that $ P_\delta(Z) $ is a $ C^1 $ function.
%By the regularity of $ K $ and $ q $, it is not hard to check that if $ Z $ is a critical point of $ P_\delta(Z)  $, then $ V_{\delta,Z}+\omega_{\delta,Z} $ is a critical point of $ I_\delta $, i.e., a solution of \eqref{111}. We now give estimates of  $ P_\delta(Z) $.
From Proposition \ref{differ of w}, $ K_\varepsilon(\cdot) $ is a $ C^1 $ function.
The following lemma shows that, to find solutions of \eqref{key equa}, it suffices to find critical points of  $ K_\varepsilon(Z) $.
\begin{lemma}\label{choice of Z}
If $ Z\in \Lambda_{\varepsilon, N}$ is a critical point of $ K_\varepsilon(Z) $, then
$ V_{\varepsilon,Z}+\omega_{\varepsilon,Z} $  is a solution to \eqref{key equa}.
\end{lemma}
\begin{proof}
It follows from Proposition \ref{exist and uniq of w} that
\begin{equation}\label{611}
\begin{split}
\langle I'(V_{\varepsilon,Z}+\omega_{\varepsilon,Z}), \phi\rangle=\sum_{j=1}^N\sum_{\tilde{h}=1}^2b_{j,\tilde{h}}\int_{\Omega} -\varepsilon^2\text{div}\left (K(z_j)\nabla   \frac{\partial V_{\varepsilon, z_j, \hat{q}_{\varepsilon,j}}}{\partial x_{\tilde{h}}}\right )\phi,\ \ \forall \phi\in H^1_0(\Omega).
\end{split}
\end{equation}
If we can choose $ Z $ such that the corresponding constants $ b_{j,\tilde{h}} $  are all zero, then $ V_{\varepsilon,Z}+\omega_{\varepsilon,Z} $  is a solution to \eqref{key equa}.
Suppose that $ Z $ is a critical point of $ K_\varepsilon(Z) $. Then from \eqref{611} and Proposition \ref{differ of w}, for $ i=1,\cdots,N, \hbar=1,2 $
\begin{equation}\label{603}
\begin{split}
0=&\frac{\partial K_\varepsilon(Z)}{\partial z_{i,\hbar}}=\left \langle I'(V_{\varepsilon,Z}+\omega_{\varepsilon,Z}), \frac{\partial (V_{\varepsilon,Z}+\omega_{\varepsilon,Z})}{\partial z_{i,\hbar}}\right \rangle\\
=&\sum_{j=1}^N\sum_{\tilde{h}=1}^2b_{j,\tilde{h}}\int_{\Omega} -\varepsilon^2\text{div}\left (K(z_j)\nabla   \frac{\partial V_{\varepsilon, z_j, \hat{q}_{\varepsilon,j}}}{\partial x_{\tilde{h}}}\right )\frac{\partial (V_{\varepsilon,Z}+\omega_{\varepsilon,Z})}{\partial z_{i,\hbar}}\\
=&\sum_{j=1}^N\sum_{\tilde{h}=1}^2b_{j,\tilde{h}}((M_i)_{\tilde{h},\hbar}\delta_{i,j}+o(\varepsilon^\gamma))+O\left (\varepsilon\sum_{j=1}^N\sum_{\tilde{h}=1}^2|b_{j,\tilde{h}}|\bigg|\bigg|\frac{\partial \omega_{\varepsilon,Z}}{\partial z_{j,\tilde{h}}}\bigg|\bigg|_{L^\infty}\right )\\
=&\sum_{j=1}^N\sum_{\tilde{h}=1}^2b_{j,\tilde{h}}((M_i)_{\tilde{h},\hbar}\delta_{i,j}+o(\varepsilon^\gamma))+O\left (\varepsilon^\gamma\sum_{j=1}^N\sum_{\tilde{h}=1}^2|b_{j,\tilde{h}}|\right ),
\end{split}
\end{equation}
from which we deduce that  $ b_{j,\tilde{h}}(Z)=0. $

\end{proof}


Now we give the  expansion of the energy functional $ K_\varepsilon(Z) $. First we obtain that $I_\varepsilon(V_{\varepsilon,Z})$ is its main part.
\begin{proposition}\label{pro401}
There holds
\begin{equation*}
K_{\varepsilon}(Z)=I_\varepsilon(V_{\varepsilon,Z})+O\left( \varepsilon^{2+\gamma}\right).
\end{equation*}
\end{proposition}

\begin{proof}
Note  that
\begin{equation*}
\begin{split}
K_{\varepsilon}(Z)=&I_\varepsilon(V_{\varepsilon,Z})+\varepsilon^2\int_{\Omega}\left( K(x)\nabla V_{\varepsilon,Z}|\nabla \omega_{\varepsilon,Z}\right) +\frac{\varepsilon^2}{2}\int_{\Omega}\left( K(x)\nabla \omega_{\varepsilon,Z}|\nabla \omega_{\varepsilon,Z}\right)\\
&-\frac{1}{p+1}\bigg( \int_{\Omega}\left (V_{\varepsilon,Z}+\omega_{\varepsilon,Z}-q|\ln\varepsilon|\right )^{p+1}_+-\int_{\Omega}\left (V_{\varepsilon,Z}-q|\ln\varepsilon|\right )^{p+1}_+\bigg).
\end{split}
\end{equation*}
It follows from Proposition \ref{exist and uniq of w} that
\begin{equation*}
\begin{split}
\int_{\Omega}&\left (V_{\varepsilon,Z}+\omega_{\varepsilon,Z}-q|\ln\varepsilon|\right )^{p+1}_+-\int_{\Omega}\left (V_{\varepsilon,Z}-q|\ln\varepsilon|\right )^{p+1}_+\\
&=(p+1)\sum_{j=1}^N\int_{B_{Ls_{\varepsilon,j}}(z_j)}\left (V_{\varepsilon,Z}-q|\ln\varepsilon|\right )^{p}_+\omega_{\varepsilon,Z}
+O\left( \sum_{j=1}^N\int_{B_{Ls_{\varepsilon,j}}(z_j)}\left (V_{\varepsilon,Z}-q|\ln\varepsilon|\right )^{p-1}_+\omega_{\varepsilon,Z}^2\right)\\
%&=O\left( \sum_{j=1}^m\frac{s_{\delta,j}^2||\omega_{\delta,Z}||_{L^\infty}}{|\ln\varepsilon|^p}\right) +O\left( \sum_{j=1}^m\frac{s_{\delta,j}^2||\omega_{\delta,Z}||_{L^\infty}^2}{|\ln\varepsilon|^{p-1}}\right) \\
&=O\left( \varepsilon^{2+\gamma}\right).
\end{split}
\end{equation*}
Since $  -\varepsilon^2\text{div}(K(x)\nabla V_{\varepsilon, Z,j})=\left (V_{\varepsilon,z_j,\hat{q}_{\varepsilon,j}}-\hat{q}_{\varepsilon,j}|\ln\varepsilon|\right )^p_+, $
we get
\begin{equation*}
\begin{split}
\varepsilon^2\int_{\Omega}\left( K(x)\nabla V_{\varepsilon,Z}|\nabla \omega_{\varepsilon,Z}\right)
=&\sum_{j=1}^N\int_{B_{Ls_{\varepsilon,j}}(z_j)}\left (V_{\varepsilon,z_j,\hat{q}_{\varepsilon,j}}-\hat{q}_{\varepsilon,j}|\ln\varepsilon|\right )^p_+\omega_{\varepsilon,Z}
%=&O\left( \sum_{j=1}^m\frac{s_{\delta,j}^2||\omega_{\delta,Z}||_{L^\infty}}{|\ln\varepsilon|^p}\right)\\
=O\left(  \varepsilon^{2+\gamma}\right).
\end{split}
\end{equation*}
As for the term $ \frac{\varepsilon^2}{2}\int_{\Omega}\left( K(x)\nabla \omega_{\varepsilon,Z}|\nabla \omega_{\varepsilon,Z}\right), $ since  $ \omega_{\varepsilon,Z}\in E_{\varepsilon,Z} $,  $ -\varepsilon^2\text{div}(K(x)\nabla \omega_{\varepsilon,Z})\in F_{\varepsilon,Z} $. So
\begin{equation*}
\begin{split}
Q_\varepsilon L_\varepsilon \omega_{\varepsilon,Z}%=&Q_\delta(-\delta^2\text{div}(K(x)\nabla \omega_{\delta,Z}))-Q_\delta(p(V_{\delta,Z}-q)^{p-1}_+\omega_{\delta,Z})\\
=&-\varepsilon^2\text{div}(K(x)\nabla \omega_{\varepsilon,Z})-Q_\varepsilon\left (p\left (V_{\varepsilon,Z}-q|\ln\varepsilon|\right )^{p-1}_+\omega_{\varepsilon,Z}\right ),
\end{split}
\end{equation*}
which combined with $ Q_\varepsilon L_\varepsilon \omega_{\varepsilon,Z}=Q_\varepsilon l_\varepsilon+Q_\varepsilon R_\varepsilon(\omega_{\varepsilon,Z})  $ yields
\begin{equation*}
-\varepsilon^2\text{div}(K(x)\nabla \omega_{\varepsilon,Z})=Q_\varepsilon\left (p\left (V_{\varepsilon,Z}-q|\ln\varepsilon|\right )^{p-1}_+\omega_{\varepsilon,Z}\right )+Q_\varepsilon l_\varepsilon+Q_\varepsilon R_\varepsilon(\omega_{\varepsilon,Z}).
\end{equation*}
Hence by Lemma \ref{lemA-5} and Proposition \ref{exist and uniq of w}, 
\begin{equation*}
\begin{split}
\varepsilon^2\int_{\Omega}&\left( K(x)\nabla \omega_{\varepsilon,Z}|\nabla \omega_{\varepsilon,Z}\right)\\
=&\int_{\Omega}Q_\varepsilon\left (p\left (V_{\varepsilon,Z}-q|\ln\varepsilon|\right )^{p-1}_+ \omega_{\varepsilon,Z}\right )\omega_{\varepsilon,Z}+\int_{\Omega}Q_\varepsilon l_\varepsilon\omega_{\varepsilon,Z}+\int_{\Omega}Q_\varepsilon R_\varepsilon(\omega_{\varepsilon,Z})\omega_{\varepsilon,Z}\\
%\leq &||Q_\delta(p(V_{\delta,Z}-q)^{p-1}_+\omega_{\delta,Z})||_{L^1}||\omega_{\delta,Z}||_{L^\infty}+||Q_\delta l_\delta||_{L^1}||\omega_{\delta,Z}||_{L^\infty}+||Q_\delta R_\delta(\omega_{\delta,Z})||_{L^1}||\omega_{\delta,Z}||_{L^\infty}\\
=&O\bigg(|| \left (V_{\varepsilon,Z}-q|\ln\varepsilon|\right )^{p-1}_+\omega_{\varepsilon,Z}||_{L^1}||\omega_{\varepsilon,Z}||_{L^\infty}+|| l_\varepsilon||_{L^1}||\omega_{\varepsilon,Z}||_{L^\infty}+|| R_\varepsilon(\omega_{\varepsilon,Z})||_{L^1} ||\omega_{\varepsilon,Z}||_{L^\infty}\bigg )\\
%= &O\left( \frac{\varepsilon^2||\omega_{\delta,Z}||_{L^\infty}^2}{|\ln\varepsilon|^{p-1}}\right) +O\left( \frac{\varepsilon^{2+\gamma}||\omega_{\delta,Z}||_{L^\infty}}{|\ln\varepsilon|^{p}}\right) +O\left( \frac{\varepsilon^2||\omega_{\delta,Z}||_{L^\infty}^3}{|\ln\varepsilon|^{p-2}}\right) \\
=&O\left( \varepsilon^{2+\gamma}\right).
\end{split}
\end{equation*}
To conclude, we have $ K_{\varepsilon}(Z)=I_\varepsilon(V_{\varepsilon,Z})+O\left( \varepsilon^{2+\gamma}\right). $




\end{proof}
For $ I_\varepsilon(V_{\varepsilon,Z}) $,  the energy expansion is as follows.
\begin{proposition}\label{order of main term}
There holds
\begin{equation}\label{344}
\begin{split}
I_\varepsilon(V_{\varepsilon,Z})=&\sum_{j=1}^N \pi\varepsilon^2 |\ln \varepsilon|q(z_j)^2\sqrt{\det K(z_j)}+\sum_{j=1}^N\frac{(p-1)}{4}\pi\varepsilon^2q(z_j)^2\sqrt{\det K(z_j)}\\
&-\sum_{j=1}^N2\pi^2\varepsilon^2 q(z_j)^2  \det K(z_j)\bar{S}_K(z_j,z_j)\\
&-\sum_{1\leq i\neq j\leq N} 2\pi^2\varepsilon^2q(z_i)q(z_j)\sqrt{\det K(z_i)}\sqrt{\det K(z_j)} G_K(z_i,z_j)+O\left( \frac{\varepsilon^2 \ln|\ln\varepsilon| }{|\ln\varepsilon|}\right).
\end{split}
\end{equation}
As a consequence, 
\begin{equation}\label{345}
\begin{split}
K_\varepsilon(Z)=&\sum_{j=1}^N \pi\varepsilon^2 |\ln \varepsilon|q(z_j)^2\sqrt{\det K(z_j)}+\sum_{j=1}^N\frac{(p-1)}{4}\pi\varepsilon^2q(z_j)^2\sqrt{\det K(z_j)}\\
&-\sum_{j=1}^N2\pi^2\varepsilon^2 q(z_j)^2  \det K(z_j)\bar{S}_K(z_j,z_j)\\
&-\sum_{1\leq i\neq j\leq N} 2\pi^2\varepsilon^2q(z_i)q(z_j)\sqrt{\det K(z_i)}\sqrt{\det K(z_j)} G_K(z_i,z_j)+O\left( \frac{\varepsilon^2 \ln|\ln\varepsilon| }{|\ln\varepsilon|}\right).
\end{split}
\end{equation}
\end{proposition}
\begin{proof}
Note that
\begin{equation}\label{218}
\begin{split}
I_\varepsilon(V_{\varepsilon,Z})=&\frac{1}{2}\int_{\Omega}-\varepsilon^2\text{div}\left( K(x)\nabla V_{\varepsilon,Z}  \right)V_{\varepsilon,Z}-\frac{1}{p+1}\int_{\Omega}\left (V_{\varepsilon,Z}-q|\ln\varepsilon|\right )^{p+1}_+ \\
%=&\frac{1}{2}\sum_{j=1}^m\int_{\Omega}(V_{\delta,z_j,\hat{q}_{\delta,j}}-\hat{q}_{\delta,j})^p_+V_{\delta,Z}-\frac{1}{p+1}\int_{\Omega}(V_{\delta,Z}-q)^{p+1}_+\\
=&\frac{1}{2}\sum_{j=1}^N\int_{\Omega}\left (V_{\varepsilon,z_j,\hat{q}_{\varepsilon,j}}-\hat{q}_{\varepsilon,j}|\ln\varepsilon|\right )^p_+V_{\varepsilon,Z,j} +\frac{1}{2}\sum_{1\leq i\neq j\leq N}\int_{\Omega}\left (V_{\varepsilon,z_j,\hat{q}_{\varepsilon,j}}-\hat{q}_{\varepsilon,j}|\ln\varepsilon|\right )^p_+V_{\varepsilon,Z,i}\\
&-\frac{1}{p+1}\int_{\Omega}\left (V_{\varepsilon,Z}-q|\ln\varepsilon|\right )^{p+1}_+.
\end{split}
\end{equation}
By   the definition of $ V_{\varepsilon,Z,j} $, 
\begin{equation*}
\begin{split}
\int_{\Omega}\left (V_{\varepsilon,z_j,\hat{q}_{\varepsilon,j}}-\hat{q}_{\varepsilon,j}|\ln\varepsilon|\right )^p_+V_{\varepsilon,Z,j}
%=&\int_{\Omega}(V_{\delta,z_j,\hat{q}_{\delta,j}}-\hat{q}_{\delta,j})^p_+ (V_{\delta,z_j,\hat{q}_{\delta,j}}+H_{\delta,z_j,\hat{q}_{\delta,j}})\\
=&\hat{q}_{\varepsilon,j}|\ln\varepsilon|\int_{\Omega}\left (V_{\varepsilon,z_j,\hat{q}_{\varepsilon,j}}-\hat{q}_{\varepsilon,j}|\ln\varepsilon|\right )^p_+
+\int_{\Omega}\left (V_{\varepsilon,z_j,\hat{q}_{\varepsilon,j}}-\hat{q}_{\varepsilon,j}|\ln\varepsilon|\right )^{p+1}_+
\\&+\int_{\Omega}\left (V_{\varepsilon,z_j,\hat{q}_{\varepsilon,j}}-\hat{q}_{\varepsilon,j}|\ln\varepsilon|\right )^p_+H_{\varepsilon,z_j,\hat{q}_{\varepsilon,j}}.
\end{split}
\end{equation*}
%By the definition of $ V_{\delta,z_j,\hat{q}_{\delta,j}} $, the fact that 
By $ T_{z_j}^{-1}(T_{z_j}^{-1})^T=K(z_j) $ and \eqref{PI}, we get
\begin{equation*}
\begin{split}
\hat{q}_{\varepsilon,j}|\ln\varepsilon|\int_{\Omega}\left (V_{\varepsilon,z_j,\hat{q}_{\varepsilon,j}}-\hat{q}_{\varepsilon,j}|\ln\varepsilon|\right )^p_+
=&\hat{q}_{\varepsilon,j}|\ln\varepsilon| s_{\varepsilon,j}^2\left (\frac{\varepsilon}{s_{\varepsilon,j}}\right )^{\frac{2p}{p-1}}\int_{|T_{z_j}x|\leq 1}\phi\left (T_{z_j}x\right )^pdx\\
%=&\hat{q}_{\delta,j} s_{\delta,j}^2(\frac{\delta}{s_{\delta,j}})^{\frac{2p}{p-1}}\sqrt{\det(K(z_j))}\cdot 2\pi|\phi'(1)|\\
%=&\hat{q}_{\varepsilon,j}|\ln\varepsilon|\varepsilon^2|\phi'(1)|^{p-1}\left( \frac{\ln\frac{1}{s_{\delta,j}}}{\hat{q}_{\delta,j}}\right)^{p-1}|\phi'(1)|^{-p}\left( \frac{\ln\frac{1}{s_{\delta,j}}}{\hat{q}_{\delta,j}}\right)^{-p}\sqrt{\det(K(z_j))}\cdot 2\pi|\phi'(1)|\\
=&\frac{2\pi\varepsilon^2|\ln\varepsilon|^2}{|\ln s_{\varepsilon,j}|}\hat{q}_{\varepsilon,j}^2\sqrt{\det(K(z_j))},
\end{split}
\end{equation*}
and
\begin{equation*}
\begin{split}
\int_{\Omega}\left (V_{\varepsilon,z_j,\hat{q}_{\varepsilon,j}}-\hat{q}_{\varepsilon,j}|\ln\varepsilon|\right )^{p+1}_+
=& s_{\varepsilon,j}^2\left( \frac{\varepsilon}{s_{\varepsilon,j}}\right)^{\frac{2(p+1)}{p-1}}\sqrt{\det(K(z_j))}\cdot \frac{(p+1)\pi}{2}|\phi'(1)|^2\\
%=&\delta^2|\phi'(1)|^{p-1}\left( \frac{\ln\frac{1}{s_{\delta,j}}}{\hat{q}_{\delta,j}}\right)^{p-1}|\phi'(1)|^{-(p+1)}\left( \frac{\ln\frac{1}{s_{\delta,j}}}{\hat{q}_{\delta,j}}\right)^{-(p+1)}\sqrt{det(K(z_j))}\cdot \frac{(p+1)\pi}{2}|\phi'(1)|^2\\
=&\frac{(p+1)\pi\varepsilon^2|\ln\varepsilon|^2}{2 |\ln s_{\varepsilon,j}|^2}\hat{q}_{\varepsilon,j}^2\sqrt{\det(K(z_j))}.
\end{split}
\end{equation*}
By Lemma \ref{H estimate}, we have
\begin{equation*}
\begin{split}
\int_{\Omega}&\left (V_{\varepsilon,z_j,\hat{q}_{\varepsilon,j}}-\hat{q}_{\varepsilon,j}|\ln\varepsilon|\right )^p_+H_{\varepsilon,z_j,\hat{q}_{\varepsilon,j}}\\
&=\frac{2\pi\hat{q}_{\varepsilon,j}\sqrt{\det K(z_j)}|\ln\varepsilon|}{|\ln  s_{\varepsilon,j}|}\int_{\Omega}\left (V_{\varepsilon,z_j,\hat{q}_{\varepsilon,j}}-\hat{q}_{\varepsilon,j}|\ln\varepsilon|\right )^p_+\bar{S}_K(x,z_j)dx+O\left( \varepsilon^{2+\gamma}\right)  \\
&=\frac{4\pi^2\varepsilon^2 \bar{S}_K(z_j,z_j)|\ln\varepsilon|^2}{|\ln s_{\varepsilon,j}|^2}\hat{q}_{\varepsilon,j}^2\cdot \det(K(z_j))+O\left( \varepsilon^{2+\gamma} \right),
\end{split}
\end{equation*}
from which we deduce that, 
\begin{equation}\label{219}
\begin{split}
\int_{\Omega}\left (V_{\varepsilon,z_j,\hat{q}_{\varepsilon,j}}-\hat{q}_{\varepsilon,j}|\ln\varepsilon|\right )^p_+V_{\varepsilon,Z,j}
=&\frac{2\pi\varepsilon^2|\ln\varepsilon|^2}{|\ln s_{\varepsilon,j}|}\hat{q}_{\varepsilon,j}^2\sqrt{\det(K(z_j))}+\frac{(p+1)\pi\varepsilon^2|\ln\varepsilon|^2}{2 |\ln s_{\varepsilon,j}|^2}\hat{q}_{\varepsilon,j}^2\sqrt{\det(K(z_j))}\\
&+\frac{4\pi^2\varepsilon^2 \bar{S}_K(z_j,z_j)|\ln\varepsilon|^2}{|\ln s_{\varepsilon,j}|^2}\hat{q}_{\varepsilon,j}^2\cdot \det(K(z_j))+O\left( \varepsilon^{2+\gamma} \right).
\end{split}
\end{equation}
By Lemmas \ref{Green expansion}, \ref{H estimate}, the definition of $ \Lambda_{\varepsilon, N} $ and the fact that $ \lim_{\varepsilon\to 0}\varepsilon|\ln\varepsilon|=0 $, for $ 1\leq i\neq j\leq N $
\begin{equation}\label{220}
\begin{split}
\int_{\Omega}&\left (V_{\varepsilon,z_j,\hat{q}_{\varepsilon,j}}-\hat{q}_{\varepsilon,j}|\ln\varepsilon|\right )^p_+V_{\varepsilon,Z,i}\\
%=&\int_{\Omega}(V_{\delta,z_j,\hat{q}_{\delta,j}}-\hat{q}_{\delta,j})^p_+\left( \frac{\hat{q}_{\delta,i}}{\ln s_{\delta,i}}\ln |T_{z_i}(x-z_i)|-\frac{2\pi\hat{q}_{\delta,i}\sqrt{\det K(z_i)}}{\ln  s_{\delta,i}} \bar{S}_K(x,z_i) \right)dx+O\left(\frac{\varepsilon^{2+\gamma}}{|\ln\varepsilon|^{p+1}}\right) \\
&=\frac{2\pi\hat{q}_{\varepsilon,i}\sqrt{\det K(z_i)}|\ln\varepsilon|}{\ln|s_{\varepsilon,i}|}\int_{\Omega}\left (V_{\varepsilon,z_j,\hat{q}_{\varepsilon,j}}-\hat{q}_{\varepsilon,j}|\ln\varepsilon|\right )^p_+G_K(x,z_i)dx+O\left( \varepsilon^{2+\gamma} \right)\\
&=\frac{4\pi^2\varepsilon^2G_K(z_j,z_i)|\ln\varepsilon|^2}{|\ln s_{\varepsilon,i}||\ln s_{\varepsilon,j}|}\hat{q}_{\varepsilon,i}\hat{q}_{\varepsilon,j}\sqrt{\det(K(z_i))}\sqrt{\det(K(z_j))}+O\left( \varepsilon^{2+\gamma} \right).
\end{split}
\end{equation}

Finally by \eqref{203},
\begin{equation}\label{223}
\begin{split}
\int_{\Omega}&\left (V_{\varepsilon,Z}-q|\ln\varepsilon|\right )^{p+1}_+\\
=&\sum_{j=1}^N\int_{B_{Ls_{\varepsilon,j}}(z_j)}\left( V_{\varepsilon, z_j, \hat{q}_{\varepsilon,j}}-\hat{q}_{\varepsilon,j}|\ln\varepsilon|+O\left(  \varepsilon^\gamma \right) \right)^{p+1}_+\\
=&\sum_{j=1}^N\int_{\Omega}\left (V_{\varepsilon, z_j, \hat{q}_{\varepsilon,j}}-\hat{q}_{\varepsilon,j}|\ln\varepsilon|\right )^{p+1}_++O\left(  \varepsilon^\gamma \sum_{j=1}^N\int_{\Omega}\left (V_{\varepsilon, z_j, \hat{q}_{\varepsilon,j}}-\hat{q}_{\varepsilon,j}|\ln\varepsilon|\right )^{p}_+\right) \\
=&\sum_{j=1}^N\frac{(p+1)\pi\varepsilon^2|\ln\varepsilon|^2}{2 |\ln s_{\varepsilon,j}|^2}\hat{q}_{\varepsilon,j}^2\sqrt{\det(K(z_j))}+O\left( \varepsilon^{2+\gamma} \right).
\end{split}
\end{equation}
Taking \eqref{219}, \eqref{220} and \eqref{223} into \eqref{218}, one has
\begin{equation}\label{225}
\begin{split}
I_\varepsilon(V_{\varepsilon,Z})=&\sum_{j=1}^N\frac{\pi\varepsilon^2|\ln\varepsilon|^2}{|\ln s_{\varepsilon,j}|}\hat{q}_{\varepsilon,j}^2\sqrt{\det(K(z_j))}+\sum_{j=1}^N\frac{(p+1)\pi\varepsilon^2|\ln\varepsilon|^2}{4 |\ln s_{\varepsilon,j}|^2}\hat{q}_{\varepsilon,j}^2\sqrt{\det(K(z_j))}\\
&+\sum_{j=1}^N\frac{2\pi^2\varepsilon^2 \bar{S}_K(z_j,z_j)|\ln\varepsilon|^2}{|\ln s_{\varepsilon,j}|^2}\hat{q}_{\varepsilon,j}^2\cdot \det(K(z_j))\\
&+\sum_{1\leq i\neq j\leq N}\frac{2\pi^2\varepsilon^2G_K(z_j,z_i)|\ln\varepsilon|^2}{|\ln s_{\varepsilon,i}||\ln s_{\varepsilon,j}|}\hat{q}_{\varepsilon,i}\hat{q}_{\varepsilon,j}\sqrt{\det(K(z_i))}\sqrt{\det(K(z_j))}\\
&-\sum_{j=1}^N\frac{\pi\varepsilon^2|\ln\varepsilon|^2}{2 |\ln s_{\varepsilon,j}|^2}\hat{q}_{\varepsilon,j}^2\sqrt{\det(K(z_j))}+O\left( \varepsilon^{2+\gamma} \right).
\end{split}
\end{equation}
%\begin{equation}\label{225-1}
%\begin{split}
%I_\delta(V_{\delta,Z})=&\sum_{j=1}^m\frac{\pi\delta^2}{\ln\frac{1}{s_{\delta,j}}}\hat{q}_{\delta,j}^2\sqrt{\det(K(z_j))}+\sum_{j=1}^m\frac{(p+1)\pi\delta^2}{4\left(  \ln\frac{1}{s_{\delta,j}}\right)^2}\hat{q}_{\delta,j}^2\sqrt{\det(K(z_j))}\\
%&+\sum_{j=1}^m\frac{2\pi^2\delta^2 \bar{S}_K(z_j,z_j)}{\left( \ln\frac{1}{s_{\delta,j}}\right)^2}\hat{q}_{\delta,j}^2  \det(K(z_j))\\
%&+\sum_{1\leq i\neq j\leq m}\frac{2\pi^2\delta^2G_K(z_j,z_i)}{\ln\frac{1}{s_{\delta,i}}\ln\frac{1}{s_{\delta,j}}}\hat{q}_{\delta,i}\hat{q}_{\delta,j}\sqrt{\det(K(z_i))}\sqrt{\det(K(z_j))}\\
%&-\sum_{j=1}^m\frac{\pi\delta^2}{2\left(  \ln\frac{1}{s_{\delta,j}}\right)^2}\hat{q}_{\delta,j}^2\sqrt{\det(K(z_j))}+O\left(\frac{\varepsilon^{2+\gamma}}{|\ln\varepsilon|^{p+1}}\right).
%\end{split}
%\end{equation}
Taking \eqref{3-003}, \eqref{q_i choice} and \eqref{2000} into \eqref{225}, we get
%\begin{equation*}
%\begin{split}
%I_\delta(V_{\delta,Z})=&\sum_{j=1}^m\frac{\pi\delta^2}{\ln\frac{1}{\varepsilon}}q(z_j)^2\sqrt{\det K(z_j)}+\sum_{j=1}^m\frac{(p-1)\pi\delta^2}{4\left(  \ln\frac{1}{\varepsilon}\right)^2}q(z_j)^2\sqrt{\det K(z_j)}\\
%&-\sum_{j=1}^m\frac{2\pi^2\delta^2 q(z_j)^2  \det K(z_j)}{\left( \ln\frac{1}{\varepsilon}\right)^2}\bar{S}_K(z_j,z_j)\\
%&-\sum_{1\leq i\neq j\leq m}\frac{2\pi^2\delta^2q(z_i)q(z_j)\sqrt{\det K(z_i)}\sqrt{\det K(z_j)}}{\left( \ln\frac{1}{\varepsilon}\right) ^2}G_K(z_i,z_j)+O\left( \frac{\left( \ln|\ln\varepsilon|\right) ^2}{|\ln\varepsilon|^3}\right).
%\end{split}
%\end{equation*}
\eqref{344}. By \eqref{344} and Proposition \ref{pro401}, we have \eqref{345}.%The proof is thus finished.

%\textbf{Proof of Theorem \ref{thm0}:}





\end{proof}
\section{Proof of Theorem \ref{thm0}}

Having made all preparations, we are now ready to prove Theorem \ref{thm0} in this section. Assume that  $ \nabla \left (q^2\sqrt{\det K}\right )(0)=0 $ and $ \left( \tilde{z}_1^*,\cdots, \tilde{z}_N^*\right)  $ is a local maximizer or minimizer of $ H_N $ defined by \eqref{def of H_N}. We will find critical points of $ K_\varepsilon(Z) $  of the form $ Z=\left (\frac{\tilde{z}_1}{\sqrt{|\ln\varepsilon|}},\frac{\tilde{z}_2}{\sqrt{|\ln\varepsilon|}},\cdots,\frac{\tilde{z}_N}{\sqrt{|\ln\varepsilon|}}\right ) $ with $ \tilde{z}_i=\tilde{z}_i^*+o(1) $. Using Proposition \ref{order of main term}, Lemma \ref{Green expansion} and Taylor's expansion, we have
\begin{equation*}\label{601}
\begin{split}
K_\varepsilon&\left (\frac{\tilde{z}_1}{\sqrt{|\ln\varepsilon|}},\frac{\tilde{z}_2}{\sqrt{|\ln\varepsilon|}},\cdots,\frac{\tilde{z}_N}{\sqrt{|\ln\varepsilon|}}\right )\\
=&\sum_{j=1}^N\pi\varepsilon^2|\ln\varepsilon|\left(\left( q^2\sqrt{\det K}\right)(0)+ \frac{1}{2|\ln\varepsilon|}\tilde{z}_j\cdot \nabla^2\left(q^2\sqrt{\det K} \right)(0 )\cdot\tilde{z}_j^T+O\left( \frac{1}{|\ln\varepsilon|^{\frac{3}{2}}}\right)    \right) \\
&+\sum_{j=1}^N\frac{(p-1)}{4}\pi\varepsilon^2\left( \left( q^2\sqrt{\det K}\right) (0)+O\left( \frac{1}{|\ln\varepsilon|^{\frac{1}{2}}}\right) \right) \\
&-\sum_{j=1}^N2\pi^2\varepsilon^2 \left( \left( q^2  \det K\right)(0) \bar{S}_K(0,0)+O\left( \frac{1}{|\ln\varepsilon|^{\frac{1}{2}}}\right)\right) \\
&-\sum_{1\leq i\neq j\leq N}2\pi^2\varepsilon^2\bigg( \left( q^2  \sqrt{\det K}\right)(0)\cdot-\frac{1}{2\pi}\ln\bigg|\frac{K(0)^{-\frac{1}{2}}\left( \tilde{z}_i-\tilde{z}_j\right) }{\sqrt{|\ln\varepsilon|}}\bigg|+\left( q^2  \det K\right)(0) \bar{S}_K(0,0)\\
&+O\left( \frac{\ln|\ln\varepsilon|}{|\ln\varepsilon|^{\frac{1}{2}}}\right)\bigg) +O\left( \frac{\varepsilon^2\ln|\ln\varepsilon|}{|\ln\varepsilon|}\right),
\end{split}
\end{equation*}
from which we deduce that
\begin{equation}\label{602}
\begin{split}
K_\varepsilon&\left (\frac{\tilde{z}_1}{\sqrt{|\ln\varepsilon|}},\frac{\tilde{z}_2}{\sqrt{|\ln\varepsilon|}},\cdots,\frac{\tilde{z}_N}{\sqrt{|\ln\varepsilon|}}\right )\\
=&N\pi\varepsilon^2|\ln\varepsilon|\left( q^2\sqrt{\det K}\right)(0)-\frac{N(N-1)\pi\varepsilon^2}{2}\left( q^2 \sqrt{\det K} \right)(0)\ln|\ln\varepsilon|\\
&+\pi\varepsilon^2\bigg(  \frac{1}{2}\sum_{j=1}^N\tilde{z}_j\cdot \nabla^2\left(q^2\sqrt{\det K} \right)(0 )\cdot\tilde{z}_j^T+\sum_{1\leq i\neq j\leq N}\left( q^2  \sqrt{\det K}\right)(0)\ln| K(0)^{-\frac{1}{2}}\left( \tilde{z}_i - \tilde{z}_j\right)  |     \\
&-2N^2\pi   \left( q^2  \det K\right)(0) \bar{S}_K(0,0)+ \frac{N(p-1)}{4}   \left( q^2\sqrt{\det K}\right) (0)   \bigg)
+O\left( \frac{\varepsilon^2\ln|\ln\varepsilon|}{|\ln\varepsilon|^{\frac{1}{2}}}\right).
\end{split}
\end{equation}

Note that %Define an auxiliary function 
\begin{equation*} 
\begin{split}
H_N(\tilde{z}_1,\cdots, \tilde{z}_N)=&\frac{1}{2}\sum_{j=1}^N\tilde{z}_j\cdot \nabla^2\left(q^2\sqrt{\det K} \right)(0 )\cdot\tilde{z}_j^T\\
&+\sum_{1\leq i\neq j\leq N}\left( q^2  \sqrt{\det K}\right)(0)\ln| K(0)^{-\frac{1}{2}}\left( \tilde{z}_i - \tilde{z}_j\right)  |.
\end{split}
\end{equation*}

From assumptions, if $ \left( \tilde{z}_1^*,\cdots, \tilde{z}_N^*\right)  $ is a local maximizer or minimizer of $ H_N $, then using perturbation theory we have 
\begin{proposition}\label{prop6-1}
	There exists $$ Z_\varepsilon=\left( \tilde{z}_{1,\varepsilon},\cdots, \tilde{z}_{N,\varepsilon}\right)=\left( \tilde{z}_1^*,\cdots, \tilde{z}_N^*\right)+o(1)$$ such that $ \left( \frac{\tilde{z}_{1,\varepsilon}}{\sqrt{|\ln\varepsilon|}},\cdots, \frac{\tilde{z}_{N,\varepsilon}}{\sqrt{|\ln\varepsilon|}}\right) $ is a critical point of $ K_\varepsilon(Z).$ As a consequence,  \eqref{key equa} has a solution $ v_\varepsilon= V_{\varepsilon,Z_\varepsilon}+\omega_{\varepsilon,Z_\varepsilon} $, where $ V_{\varepsilon,Z_\varepsilon} $ satisfies \eqref{solu config} and $ \omega_{\varepsilon,Z_\varepsilon} $ is determined by Proposition \ref{exist and uniq of w}.
%O\left( \frac{\left( \ln|\ln\varepsilon|\right)^{\frac{1}{2}}}{|\ln\varepsilon|^{\frac{1}{4}}}\right)  $
\end{proposition}
\begin{proof}
Using \eqref{602} and the former deduction, we can easily get the result.
\end{proof}

We can also estimate     $ L^1 $ norm of $ v_\varepsilon $, which corresponds to the circulation of each component of vorticity fields in our construction in subsequent sections.
\begin{lemma}\label{circulation}
For $ i=1,\cdots,N $, we have
	\begin{equation*}
	\lim_{\varepsilon\to 0}\frac{1}{\varepsilon^2}\int_{B_{|\ln\varepsilon|^{-1}}\left (\frac{\tilde{z}_{i,\varepsilon}}{\sqrt{|\ln\varepsilon|}}\right )}\left (v_\varepsilon-q|\ln \varepsilon|\right )^{p}_+dx=2\pi \left( q\sqrt{\det K}\right) (0).
	\end{equation*}
%	As a consequence,
%	\begin{equation*}
%	\lim_{\varepsilon\to 0}\frac{1}{\varepsilon^2}\int_{\Omega}\left (u_\varepsilon-q\ln\frac{1}{\varepsilon}\right )^{p}_+dx=2\pi m q\sqrt{\det K}(x_{0}).
%	\end{equation*}
\end{lemma}
\begin{proof}
	
	It follows from  \eqref{q_i choice}, \eqref{2000}, \eqref{203} and Proposition \ref{exist and uniq of w} that
	\begin{equation*}
	\begin{split}
	\frac{1}{\varepsilon^2}&\int_{B_{|\ln\varepsilon|^{-1}}\left (\frac{\tilde{z}_{i,\varepsilon}}{\sqrt{|\ln\varepsilon|}}\right )}\left (v_\varepsilon-q|\ln \varepsilon|\right )^{p}_+dx
	\\
	&=\frac{1}{\varepsilon^2}\int_{B_{Ls_{\varepsilon,i}}\left (\frac{\tilde{z}_{i,\varepsilon}}{\sqrt{|\ln\varepsilon|}}\right )}\left(V_{\varepsilon, z_{i,\varepsilon}, \hat{q}_{\varepsilon,i}}(x)-\hat{q}_{\varepsilon,i}|\ln\varepsilon|+O\left( \varepsilon^\gamma\right)\right)^{p}_+dx\\
	%&=\frac{|\ln\varepsilon|^p}{\varepsilon^2}s_{\delta,i}^2\left( \frac{\delta}{s_{\delta,i}}\right)^{\frac{2p}{p-1}}\int_{|T_{z_{i,\delta}}x|\leq 1}\phi(T_{z_{i,\delta}}x)^pdx+o(1)\\
	&=\frac{s_\varepsilon^2}{\varepsilon^2}\left( \frac{\hat{q}_{\varepsilon,i}}{|\phi'(1)|}\right) ^{p}\left( \frac{|\ln\varepsilon|}{|\ln s_{\varepsilon,i}|}\right)^{p}\sqrt{\det K}\left (\frac{\tilde{z}_{i,\varepsilon}}{\sqrt{|\ln\varepsilon|}}\right )\cdot 2\pi|\phi'(1)|+O(\varepsilon^\gamma)\\
	&= \frac{2\pi\hat{q}_{\varepsilon,i} |\ln\varepsilon|}{|\ln s_{\varepsilon,i}|}\cdot \sqrt{\det K }\left (\frac{\tilde{z}_{i,\varepsilon}}{\sqrt{|\ln\varepsilon|}}\right )+O(\varepsilon^\gamma) \\
	&= 2\pi \left( q \sqrt{\det K }\right) (0)+O\left (\frac{1}{\sqrt{|\ln\varepsilon|}}\right ).
	\end{split}
	\end{equation*}
\end{proof}




\textbf{Proof of Theorem \ref{thm0}:}

The rest of properties of $ v_\varepsilon $ can be easily deduced from the decomposition of $ v_\varepsilon $ in \eqref{solu config}, Proposition \ref{prop6-1} and Lemma \ref{circulation}. The proof of Theorem \ref{thm0} is thus finished.



	
	
	











\section{Proof of Theorems \ref{thm01}, \ref{thm02}, \ref{thm03}, \ref{thm04}, \ref{thm05}}

The proof of Theorems \ref{thm01}-\ref{thm05} are similar to that of Theorem \ref{thm0} and we therefore just show the difference to keep this paper not too long.

\textbf{Proof of Theorem \ref{thm01}:}
 Let $ N\geq 2, R^*>0, r_*>0, h>0, \kappa>0 $. Consider  the following equations
\begin{equation}\label{equa 1}
\begin{cases}
-\varepsilon^2\text{div}(K_H(x)\nabla v)=  \left (v-\left(\frac{\alpha}{2}|x|^2+\beta \right)|\ln\varepsilon| \right )^{p}_+\textbf{1}_{D_{\frac{\rho_0}{\sqrt{|\ln\varepsilon|}}}\left (C_{\frac{r_*}{\sqrt{|\ln\varepsilon|}}}\right )},\ \ &\text{in}\  B_{R^*}(0),\\
v=0,\ \ &\text{on}\ \partial B_{R^*}(0).
\end{cases}
\end{equation}
Here  $ \rho_0>0 $ sufficiently small, $ \textbf{1}_{D} $ is Heaviside function of some set $ D $, $ C_{\frac{r_*}{\sqrt{|\ln\varepsilon|}}}=\left \{x\mid |x|=\frac{r_*}{\sqrt{|\ln\varepsilon|}}\right \} $ is the circle centered at the origin of radius $ \frac{r_*}{\sqrt{|\ln\varepsilon|}} $, and $ D_{\frac{\rho_0}{\sqrt{|\ln\varepsilon|}}}\left (C_{\frac{r_*}{\sqrt{|\ln\varepsilon|}}}\right )=\left \{x\mid |x|\in \left (\frac{r_*-\rho_0}{\sqrt{|\ln\varepsilon|}},\frac{r_*+\rho_0}{\sqrt{|\ln\varepsilon|}} \right )\right \} $ is the $ \frac{\rho_0}{\sqrt{|\ln\varepsilon|}} $-neighborhood of $ C_{\frac{r_*}{\sqrt{|\ln\varepsilon|}}} $.  The constants $ \alpha $ and $ \beta $ in \eqref{equa 1} are chosen by
\begin{equation*} 
\alpha=\frac{\kappa}{4\pi}\left( \frac{1}{h^2}-\frac{N-1}{r_*^2}\right),\ \ \ \ \beta=\frac{\kappa}{2\pi}.
\end{equation*}

So if we take  $ K(x)=K_H(x) $  and $ q(x)=\frac{\alpha}{2}|x|^2+\beta $ and $ \Omega=B_{R^*}(0) $ in \eqref{key equa}, we get \eqref{equa 1}. 
In this case we choose 
\begin{equation*} 
\begin{split}
\Lambda_{\varepsilon, N,1}=\bigg\{&Z=(z_1,\cdots, z_N)\in\left(  B_{R^*}(0)\right) ^{(N)}\mid z_j\in D_{\frac{\rho_0}{2\sqrt{|\ln\varepsilon|}}}\left (C_{\frac{r_*}{\sqrt{|\ln\varepsilon|}}}\right ),\ \  z_j=e^{\textbf{i}\frac{2\pi(j-1)}{N}}z_1.\bigg\}
\end{split}
\end{equation*}
Our goal is to construct solutions of \eqref{equa 1} being of the form $ V_{\varepsilon,Z}+\omega_{\varepsilon,Z} $, where $ Z\in \Lambda_{\varepsilon, N,1} $.
Following the proof of Theorem \ref{thm0}, we only need to calculate critical points of $ K_\varepsilon(Z) $ in \eqref{functional I}. Note that $  K_H$ and  $ q $ are rotational-invariant. Since $ z_j=e^{\textbf{i}\frac{2\pi(j-1)}{N}}z_1 $, one can prove that $ K_\varepsilon(Z) $ is a   function of $ z_1 $ %, i,e, $ K_\varepsilon(Z)=  \hat{K}_\varepsilon(z_1)   $ for some $ \hat{K}_\varepsilon $. Moreover, $ \hat{K}_\varepsilon $ 
and is rotational-invariant, i.e., $ K_\varepsilon(Z)=\tilde{K}_\varepsilon(|z_1|) $ for some $ \tilde{K}_\varepsilon $. Let $ z_1=\left( \frac{r}{\sqrt{|\ln\varepsilon|}},0\right)  $ for $ r\in \left( r_*-\frac{\rho_0}{2},r_*+\frac{\rho_0}{2} \right)  $. %Using  Proposition \ref{order of main term} and 
Since   $ \nabla \left (q^2\sqrt{\det K_H}\right )(0)=0$, using \eqref{602}, one computes that
\begin{equation} \label{701}
\begin{split}
\tilde{K}_\varepsilon\left (\frac{r}{\sqrt{|\ln\varepsilon|}}\right )
=&N\pi\varepsilon^2|\ln\varepsilon|\left( q^2\sqrt{\det K_H}\right)(0)-\frac{N(N-1)\pi\varepsilon^2}{2}\left( q^2\det K_H\right)(0)\ln|\ln\varepsilon|\\
&+\pi\varepsilon^2\bigg(  \frac{N}{2} \left(q^2\sqrt{\det K_H} \right)^{''}(0 )r^2+ \frac{N(p-1)}{4}   \left( q^2\sqrt{\det K_H}\right) (0)      \\
&-2N\pi   \left( q^2  \det K_H\right)(0) \bar{S}_K(0,0)+N(N-1)\left( q^2  \det K_H\right)(0)\ln r\\
&+\sum_{1\leq j\neq k\leq N}\left( q^2  \det K_H\right)(0)\ln| e^{\textbf{i}\frac{2\pi(j-1)}{N}} - e^{\textbf{i}\frac{2\pi(k-1)}{N}} |\bigg)  
+O\left( \frac{\varepsilon^2\ln|\ln\varepsilon|}{|\ln\varepsilon|^{\frac{1}{2}}}\right).
\end{split}
\end{equation}
 So the existence of critical points of  $ \tilde{K}_\varepsilon $ is guaranteed by the existence of critical points of the function 
$$H_1(r)=\frac{N}{2} \left(q^2\sqrt{\det K_H} \right)^{''}(0 )r^2+N(N-1)\left( q^2  \det K_H\right)(0)\ln r.$$
We have
\begin{lemma}\label{lm701}
The function 
$ H_1(r) $
has the unqiue critical point $ r=r_* $ in $  \left( r_*-\frac{\rho_0}{2},r_*+\frac{\rho_0}{2} \right) $, which is a maximizer.
\end{lemma}  
\begin{proof}
We can prove that $  \left(q^2\sqrt{\det K_H} \right)^{''}(0 )=\frac{\beta(2\alpha h^2-\beta)}{h^2} $, $ \left( q^2  \det K_H\right)(0)=\beta^2 $ and $ H'_1(r)=N\left(\frac{\beta(2\alpha h^2-\beta)}{h^2} r+\frac{(N-1)\beta^2}{r} \right)  $. Taking $ \alpha,\beta $ into $ H'_1 $, we show that $ r_* $ is the unique maximizer.
\end{proof}

Note that from \eqref{701},
\begin{equation*}  
\begin{split}
K_\varepsilon\left (Z\right )
=&N\pi\varepsilon^2|\ln\varepsilon|\left( q^2\sqrt{\det K_H}\right)(0)-\frac{N(N-1)\pi\varepsilon^2}{2}\left( q^2\det K_H\right)(0)\ln|\ln\varepsilon|\\
&+\pi\varepsilon^2\bigg(  H_1(r)+ \frac{N(p-1)}{4}   \left( q^2\sqrt{\det K_H}\right) (0)      -2N\pi   \left( q^2  \det K_H\right)(0) \bar{S}_K(0,0)\\
&+\sum_{1\leq j\neq k\leq N}\left( q^2  \det K_H\right)(0)\ln| e^{\textbf{i}\frac{2\pi(j-1)}{N}} - e^{\textbf{i}\frac{2\pi(k-1)}{N}} |\bigg)  
+\tilde{N}_\varepsilon(z_1),
\end{split}
\end{equation*}
where $ \tilde{N}_\varepsilon(z_1) $ is a $O\left( \frac{\varepsilon^2\ln|\ln\varepsilon|}{|\ln\varepsilon|^{\frac{1}{2}}}\right) -$perturbation term which is invariant under rotation. %Since $ r_* $ is a strict local maximizer of $ H_1 $, 
By Lemma \ref{lm701}, there exists $ z_{1,\varepsilon}=\left( \frac{r_\varepsilon}{\sqrt{|\ln\varepsilon|}},0\right)  $ near $ \left( \frac{r_*}{\sqrt{|\ln\varepsilon|}},0\right) $ being a maximizer of $ \tilde{K}_\varepsilon  $.  Let $ z_{j,\varepsilon}=e^{\textbf{i}\frac{2\pi(j-1)}{N}}z_{1,\varepsilon} $ for $ j=1,\cdots,N. $ Then $ \left( z_{1,\varepsilon},\cdots,z_{N,\varepsilon} \right) $ is a maximizer of $ K_\varepsilon(Z) $, which yields a solution $ v_\varepsilon $ of \eqref{equa 1}. Similar to Theorem \ref{thm0}, %there exist $ \left (\tilde{z}_{1,\varepsilon}, \cdots, \tilde{z}_{N,\varepsilon}\right ) $ such that $ \tilde{z}_{i,\varepsilon}=\tilde{z}_{i}^*+o(1) $ for $ i=1,\cdots,N $ and 
\begin{equation*}
\left \{v_\varepsilon>q|\ln\varepsilon|\right \}\subseteq \cup_{i=1}^N B_{|\ln\varepsilon|^{-1}}\left (z_{i,\varepsilon}\right ),
\end{equation*}
and if we define $ A_{\varepsilon,i}=\left\{v_\varepsilon>q|\ln\varepsilon|\right\}\cap B_{|\ln\varepsilon|^{-1}}\left (z_{i,\varepsilon}\right )  $, then there are $ R_1,R_2>0 $ independent of $ \varepsilon $  such that $ B_{R_1\varepsilon}\left ( z_{i,\varepsilon}\right )\subseteq A_{\varepsilon,i}\subseteq B_{R_2\varepsilon}\left (z_{i,\varepsilon}\right ) $ and 
\begin{equation*}
\frac{1}{\varepsilon^2}\int_{A_{\varepsilon,i}}\left( v_\varepsilon-q|\ln\varepsilon|\right)^{p}_+dx=2\pi \left( q\sqrt{\det K_H }\right) (0)+O\left(|\ln\varepsilon|^{-\frac{1}{2}} \right) = \kappa+O\left(|\ln\varepsilon|^{-\frac{1}{2}} \right).
\end{equation*}

%Hence by Theorem \ref{thmA}, for any $ \varepsilon $ small there exists a solution $ u_\varepsilon $ of \eqref{eq01} concentrating near $ (r_*,0) $ with $ \alpha=\frac{c}{4\pi k\sqrt{k^2+r_*^2}} $ and $\frac{1}{\varepsilon^2} \int_{A_\varepsilon}\left( u_\varepsilon-\left( \frac{\alpha|x|^2}{2}+\beta\right)\ln\frac{1}{\varepsilon}\right)^p_+dx $ tending to $ c $. 
Define for any $ (x_1,x_2,x_3)\in B_{R^*}(0)\times\mathbb{R}, t\in\mathbb{R}  $
\begin{equation*}
\mathbf{w}_\varepsilon(x_1,x_2,x_3,t)= \frac{w_\varepsilon(x_1,x_2,x_3,t)}{h}(x_2,-x_1,h),
\end{equation*}
where $ w_\varepsilon(x_1,x_2,x_3,t) $ is a helical function satisfying
\begin{equation*} 
w_\varepsilon(x_1,x_2,0,t)=\frac{1}{\varepsilon^2}\left( v_\varepsilon(\bar{R}_{-\alpha|\ln\varepsilon| t}(x_1,x_2))-\left( \frac{\alpha|(x_1,x_2)|^2}{2}+\beta\right)|\ln \varepsilon|\right)^p_+\textbf{1}_{D_{\frac{\rho_0}{\sqrt{|\ln\varepsilon|}}}\left (C_{\frac{r_*}{\sqrt{|\ln\varepsilon|}}}\right )}.
\end{equation*}
Then $ w_\varepsilon(x_1,x_2,0,t) $ satisfies \eqref{vor str eq} and $ \mathbf{w}_\varepsilon $ is a  helical vorticity field of \eqref{Euler eq2}.  %Moreover, $ w_\varepsilon(x_1,x_2,0,t) $ rotates clockwise around the origin with angular velocity $ \alpha|\ln\varepsilon| $. 
The circulation of $ \mathbf{w}_\varepsilon $ is
\begin{equation*}
\iint_{A_{\varepsilon,i}}\mathbf{w}_\varepsilon\cdot\mathbf{n}d\sigma=\frac{1}{\varepsilon^2} \int_{A_{\varepsilon,i}}\left( v_\varepsilon-\left( \frac{\alpha|x|^2}{2}+\beta\right)|\ln \varepsilon|\right)^p_+dx=\kappa+O\left(|\ln\varepsilon|^{-\frac{1}{2}} \right).
\end{equation*}

Recall $ \tilde{X}_{j,\varepsilon}(s, t)=\frac{r_*}{\sqrt{|\ln\varepsilon|}}e^{-\mathbf{i}\alpha|\ln\varepsilon|t}e^{\mathbf{i}\frac{s}{h}}e^{\mathbf{i}\frac{2\pi (j-1)}{N}} $ for $ j=1,\cdots,N $. Based on the above derivation,   $ \mathbf{w}_\varepsilon $ tends asymptotically to $ \cup_{j=1}^N\tilde{X}_{j,\varepsilon} $ in sense of both topology and distribution. %, i.e.,
%\begin{equation*} 
%\mathbf{w}_\varepsilon(x,t)- \kappa\sum_{j=1}^N\delta_{\tilde{X}_{j,\varepsilon}}\textbf{t}_{\tilde{X}_{j,\varepsilon}}\rightharpoonup0\ \ \ \ \text{as}\ \varepsilon\to0.
%\end{equation*}
The proof of Theorem \ref{thm01} is complete.



\textbf{Proof of Theorem \ref{thm02}:}
 Let $ N\geq 2, R^*>0, r_*>0, h>0, \kappa>0, \mu>0 $. Consider  the following problem
\begin{equation}\label{equa 2}
\begin{cases}
-\varepsilon^2\text{div}(K_H(x)\nabla v)=&\left (v-\left(\frac{\alpha}{2}|x|^2+\beta_1 \right)|\ln\varepsilon| \right )^{p}_+\textbf{1}_{D_{\frac{\rho_0}{\sqrt{|\ln\varepsilon|}}}\left (C_{\frac{r_*}{\sqrt{|\ln\varepsilon|}}}\right )}\\
&+\left (v-\left(\frac{\alpha}{2}|x|^2+\beta_2 \right)|\ln\varepsilon| \right )^{p}_+\textbf{1}_{B_{\frac{\rho_0}{\sqrt{|\ln\varepsilon|}}}\left (0\right )},\ \ \text{in}\  B_{R^*}(0),\\
v=0,&\ \ \ \ \ \ \ \ \ \ \ \ \ \ \ \ \ \ \ \ \ \ \ \ \ \ \ \ \ \ \ \ \ \ \ \ \ \ \ \ \ \ \ \ \ \ \ \ \ \ \ \ \ \text{on}\ \partial B_{R^*}(0),
\end{cases}
\end{equation}
where  $ \rho_0>0 $ sufficiently small, %, $ \textbf{1}_{D} $ is Heaviside function of some set $ D $, $ C_{\frac{r_*}{\sqrt{|\ln\varepsilon|}}}=\left \{x\mid |x|=\frac{r_*}{\sqrt{|\ln\varepsilon|}}\right \} $ is the circle centered at the origin of radius $ \frac{r_*}{\sqrt{|\ln\varepsilon|}} $, and $ D_{\frac{\delta}{\sqrt{|\ln\varepsilon|}}}\left (C_{\frac{r_*}{\sqrt{|\ln\varepsilon|}}}\right )=\left \{x\mid |x|\in \left (\frac{r_*-\delta}{\sqrt{|\ln\varepsilon|}},\frac{r_*+\delta}{\sqrt{|\ln\varepsilon|}} \right )\right \} $ is the $ \frac{\delta}{\sqrt{|\ln\varepsilon|}} $-neighborhood of $ C_{\frac{r_*}{\sqrt{|\ln\varepsilon|}}} $.  
 $ \alpha $ and $ \beta_i $ in \eqref{equa 2} are chosen by
\begin{equation*}
\alpha=\frac{\kappa}{4\pi}\left( \frac{1}{h^2}-\frac{N-1}{r_*^2}\right)-\frac{\mu}{2\pi r_*^2},\ \ \ \ \beta_1=\frac{\kappa}{2\pi}, \ \ \ \ \beta_2=\frac{\mu}{2\pi}.
\end{equation*}
Denote $ q_1(x)=\frac{\alpha}{2}|x|^2+\beta_1 $ and $ q_2(x)=\frac{\alpha}{2}|x|^2+\beta_2. $

We choose 
\begin{equation*} 
\begin{split}
\Lambda_{\varepsilon, N+1}=\bigg\{&Z=(z_0, z_1,\cdots, z_N) \mid z_0\in B_{\frac{\rho_0}{2\sqrt{|\ln\varepsilon|}}}\left (0\right ), z_j\in D_{\frac{\rho_0}{2\sqrt{|\ln\varepsilon|}}}\left (C_{\frac{r_*}{\sqrt{|\ln\varepsilon|}}}\right ),\   z_j=e^{\textbf{i}\frac{2\pi(j-1)}{N}}z_1\bigg\}.
\end{split}
\end{equation*}
%Our goal is to construct solutions of \eqref{equa 1} being of the form $ V_{\varepsilon,Z}+\omega_{\varepsilon,Z} $, where $ Z\in \Lambda_{\varepsilon, N,1} $.
Similar to the proof of Theorem \ref{thm0}, we find critical points of $ K_\varepsilon(Z) $ over $ \Lambda_{\varepsilon, N+1} $. Since $ K_\varepsilon\left (z_0, z_1,\cdots, z_N\right )=K_\varepsilon\left (e^{\textbf{i}\frac{2\pi(j-1)}{N}}z_0, z_1,\cdots, z_N\right ) $ for $ j=1,\cdots,N $, we have $\nabla K_\varepsilon(0, z_1,\cdots, z_N)=0$ for any $ (z_1,\cdots,z_N) $.  Since $ z_j=e^{\textbf{i}\frac{2\pi(j-1)}{N}}z_1 $,   $ K_\varepsilon(0,z_1,\cdots,z_N) $ is a   function of $ z_1 $ %, i,e, $ K_\varepsilon(Z)=  \hat{K}_\varepsilon(z_1)   $ for some $ \hat{K}_\varepsilon $. Moreover, $ \hat{K}_\varepsilon $ 
and is rotational-invariant, i.e., $ K_\varepsilon(0,z_1,\cdots,z_N)=\tilde{K}_\varepsilon(|z_1|) $ for some $ \tilde{K}_\varepsilon $. Let $ z_1=\left( \frac{r}{\sqrt{|\ln\varepsilon|}},0\right)  $ for $ r\in \left (r_*-\frac{\rho_0}{2},r_*+\frac{\rho_0}{2} \right ) $. %Using  Proposition \ref{order of main term} and 
Using \eqref{602}, one computes that
\begin{equation} \label{702}
\begin{split}
\tilde{K}_\varepsilon&\left (\frac{r}{\sqrt{|\ln\varepsilon|}}\right )\\
=&\pi\varepsilon^2|\ln\varepsilon|\left( N\left( q_1^2\sqrt{\det K_H}\right)(0)+\left( q_2^2\sqrt{\det K_H}\right)(0)\right) \\
&-\frac{N(N-1)\pi\varepsilon^2}{2}\left( q_1^2 \det K_H \right)(0)\ln|\ln\varepsilon|-N\pi\varepsilon^2\left( q_1q_2 \det K_H \right)(0)\ln|\ln\varepsilon|\\
&+\pi\varepsilon^2\bigg(  \frac{N}{2} \left(q_1^2\sqrt{\det K_H} \right)^{''}(0 )r^2+ \frac{N(p-1)}{4}   \left( q_1^2\sqrt{\det K_H}\right) (0)+ \frac{(p-1)}{4}   \left( q_2^2\sqrt{\det K_H}\right) (0)     \\
&-2N\pi   \left( q_1^2  \det K_H\right)(0) \bar{S}_K(0,0)-2\pi   \left( q_2^2  \det K_H\right)(0) \bar{S}_K(0,0)+N(N-1)\left( q_1^2  \det K_H\right)(0)\ln r\\
&+2N\left( q_1q_2 \det K_H\right)(0)\ln r+\sum_{1\leq j\neq k\leq N}\left( q_1^2  \det K_H\right)(0)\ln| e^{\textbf{i}\frac{2\pi(j-1)}{N}} - e^{\textbf{i}\frac{2\pi(k-1)}{N}} |\bigg)  
+O\left( \frac{\varepsilon^2\ln|\ln\varepsilon|}{|\ln\varepsilon|^{\frac{1}{2}}}\right).
\end{split}
\end{equation}
Define 
$$H_2(r)=\frac{N}{2} \left(q_1^2\sqrt{\det K_H} \right)^{''}(0 )r^2+N(N-1)\left( q_1^2  \det K_H\right)(0)\ln r+2N\left( q_1q_2 \det K_H\right)(0)\ln r.$$
We have
\begin{lemma}\label{lm702}
	The function 
	$ H_2(r) $
	has the unqiue critical point $ r=r_* $ in $  \left (r_*-\frac{\rho_0}{2},r_*+\frac{\rho_0}{2} \right ) $, which is a maximizer.
\end{lemma}  
\begin{proof}
Note that $  \left(q_1^2\sqrt{\det K_H} \right)^{''}(0 )=\frac{\beta_1(2\alpha h^2-\beta_1)}{h^2} $, $ \left( q_1q_i \det K_H\right)(0)=\beta_1\beta_i $ and $$ H'_2(r)=N\left(\frac{\beta_1(2\alpha h^2-\beta_1)}{h^2} r+\frac{(N-1)\beta_1^2+2\beta_1\beta_2}{r} \right).  $$ Taking $ \alpha,\beta_i $ into $ H'_2 $, we show that $ r_* $ is the unique maximizer.
\end{proof}
%
%Note that from \eqref{701},
%\begin{equation*}  
%\begin{split}
%K_\varepsilon\left (Z\right )
%=&N\pi\varepsilon^2|\ln\varepsilon|\left( q^2\sqrt{\det K_H}\right)(0)-\frac{N(N-1)\pi\varepsilon^2}{2}\left( q^2\sqrt{\det K_H}\right)(0)\ln|\ln\varepsilon|\\
%&+\pi\varepsilon^2\bigg(  H_1(r)+ \frac{N(p-1)}{4}   \left( q^2\sqrt{\det K_H}\right) (0)      -2N\pi   \left( q^2  \det K_H\right)(0) \bar{S}_K(0,0)\\
%&+\sum_{1\leq j\neq k\leq N}\left( q^2  \det K_H\right)(0)\ln| e^{\textbf{i}\frac{2\pi(j-1)}{N}} - e^{\textbf{i}\frac{2\pi(k-1)}{N}} |\bigg)  
%+\tilde{N}_\varepsilon(z_1).
%\end{split}
%\end{equation*}
%where $ \tilde{N}_\varepsilon(z_1) $ is a $O\left( \frac{\varepsilon^2\ln|\ln\varepsilon|}{|\ln\varepsilon|^{\frac{1}{2}}}\right) -$perturbation term which is invariant under rotation. %Since $ r_* $ is a strict local maximizer of $ H_1 $, 
By Lemma \ref{lm702}, there exists $ z_{1,\varepsilon}=\left( \frac{r_\varepsilon}{\sqrt{|\ln\varepsilon|}},0\right)  $ near $ \left( \frac{r_*}{\sqrt{|\ln\varepsilon|}},0\right) $ being  a maximizer of $ \tilde{K}_\varepsilon  $. Let $ z_{0,\varepsilon}=0 $ and $ z_{j,\varepsilon}=e^{\textbf{i}\frac{2\pi(j-1)}{N}}z_{1,\varepsilon} $ for $ j=1,\cdots,N. $ Then $ \left( z_{0,\varepsilon},z_{1,\varepsilon},\cdots,z_{N,\varepsilon} \right)  $ is a critical point  of $ K_\varepsilon  $, which yields a solution $ v_\varepsilon $ of \eqref{equa 1}. Define 
\begin{equation*} 
\begin{split}
w_\varepsilon(x_1,x_2,0,t)=&\frac{1}{\varepsilon^2}\left (v_\varepsilon\left(\bar{R}_{-\alpha|\ln\varepsilon| t}(x_1,x_2) \right) -q_1(x_1,x_2)|\ln\varepsilon| \right )^{p}_+\textbf{1}_{D_{\frac{\rho_0}{\sqrt{|\ln\varepsilon|}}}\left (C_{\frac{r_*}{\sqrt{|\ln\varepsilon|}}}\right )}\\
&+\frac{1}{\varepsilon^2}\left (v_\varepsilon\left(\bar{R}_{-\alpha|\ln\varepsilon| t}(x_1,x_2) \right)-q_2(x_1,x_2)|\ln\varepsilon| \right )^{p}_+\textbf{1}_{B_{\frac{\rho_0}{\sqrt{|\ln\varepsilon|}}}\left (0\right )}.
\end{split}
\end{equation*}
 %there exist $ \left (\tilde{z}_{1,\varepsilon}, \cdots, \tilde{z}_{N,\varepsilon}\right ) $ such that $ \tilde{z}_{i,\varepsilon}=\tilde{z}_{i}^*+o(1) $ for $ i=1,\cdots,N $ and 
Then $ supp(w_\varepsilon(\cdot,0,0))\subseteq \cup_{i=0}^N B_{|\ln\varepsilon|^{-1}}\left (z_{i,\varepsilon}\right ) $
and if we define $ A_{\varepsilon,i}=supp(w_\varepsilon(\cdot,0,0))\cap B_{|\ln\varepsilon|^{-1}}\left (z_{i,\varepsilon}\right )  $ for $ i=0,1,\cdots,N $, then there are $ R_1,R_2>0 $ independent of $ \varepsilon $  such that $ B_{R_1\varepsilon}\left ( z_{i,\varepsilon}\right )\subseteq A_{\varepsilon,i}\subseteq B_{R_2\varepsilon}\left (z_{i,\varepsilon}\right ) $ and  
\begin{equation*}
\int_{A_{\varepsilon,i}}w_\varepsilon dx=\begin{cases}
&2\pi \left( q_1\sqrt{\det K_H }\right) (0)+O\left(|\ln\varepsilon|^{-\frac{1}{2}} \right)=\kappa+O\left(|\ln\varepsilon|^{-\frac{1}{2}} \right)\ \ \ \ i=1,\cdots,N,\\
 &2\pi \left( q_2\sqrt{\det K_H }\right) (0)+O\left(|\ln\varepsilon|^{-\frac{1}{2}} \right)=\mu+O\left(|\ln\varepsilon|^{-\frac{1}{2}} \right)\ \ \ \ i=0.
\end{cases}
\end{equation*}

Recall $ 	 
\tilde{X}_{0,\varepsilon}(s,t)=0$ and $ \tilde{X}_{j,\varepsilon}(s,t)=\frac{r_*}{\sqrt{|\ln\varepsilon|}}e^{-\mathbf{i}\alpha|\ln\varepsilon|t}e^{\mathbf{i}\frac{s}{h}}e^{\mathbf{i}\frac{2\pi (j-1)}{N}} 
  $ for $ j=1,\cdots,N $. Based on the above derivation, we get Theorem \ref{thm02}.




\textbf{Proof of Theorem \ref{thm03}:}
Let $   R^*>0, h>0, $ $ \kappa_1,\kappa_2>0 $ and $ \lambda_{1,*},\lambda_{2,*}>0 $ are constants satisfying the compatibility condition
\begin{equation*} 
\frac{\kappa_2}{\kappa_1}=\frac{2\lambda_{1,*}h^2+\lambda_{1,*}\lambda_{2,*}(\lambda_{1,*}+\lambda_{2,*})}{2\lambda_{2,*}h^2+\lambda_{1,*}\lambda_{2,*}(\lambda_{1,*}+\lambda_{2,*})}.
\end{equation*}
Consider  the following problem
\begin{equation}\label{equa 3}
\begin{cases}
-\varepsilon^2\text{div}(K_H(x)\nabla v)=&\left (v-\left(\frac{\alpha}{2}|x|^2+\beta_1 \right)|\ln\varepsilon| \right )^{p}_+\textbf{1}_{B_{\frac{\rho_0}{\sqrt{|\ln\varepsilon|}}}\left (\frac{(\lambda_{1,*},0)}{\sqrt{|\ln\varepsilon|}}\right )}\\
&+\left (v-\left(\frac{\alpha}{2}|x|^2+\beta_2 \right)|\ln\varepsilon| \right )^{p}_+\textbf{1}_{B_{\frac{\rho_0}{\sqrt{|\ln\varepsilon|}}}\left (\frac{(-\lambda_{2,*},0)}{\sqrt{|\ln\varepsilon|}}\right )},\ \ \text{in}\  B_{R^*}(0),\\
v=0,&\ \ \ \ \ \ \ \ \ \ \ \ \ \ \ \ \ \ \ \ \ \ \ \ \ \ \ \ \ \ \ \ \ \ \ \ \ \ \ \ \ \ \ \ \ \ \ \ \ \ \ \ \ \ \ \ \ \ \ \ \text{on}\ \partial B_{R^*}(0),
\end{cases}
\end{equation}
where  $ \rho_0>0 $ sufficiently small, %, $ \textbf{1}_{D} $ is Heaviside function of some set $ D $, $ C_{\frac{r_*}{\sqrt{|\ln\varepsilon|}}}=\left \{x\mid |x|=\frac{r_*}{\sqrt{|\ln\varepsilon|}}\right \} $ is the circle centered at the origin of radius $ \frac{r_*}{\sqrt{|\ln\varepsilon|}} $, and $ D_{\frac{\delta}{\sqrt{|\ln\varepsilon|}}}\left (C_{\frac{r_*}{\sqrt{|\ln\varepsilon|}}}\right )=\left \{x\mid |x|\in \left (\frac{r_*-\delta}{\sqrt{|\ln\varepsilon|}},\frac{r_*+\delta}{\sqrt{|\ln\varepsilon|}} \right )\right \} $ is the $ \frac{\delta}{\sqrt{|\ln\varepsilon|}} $-neighborhood of $ C_{\frac{r_*}{\sqrt{|\ln\varepsilon|}}} $.  
$ \alpha $ and $ \beta_i $ in \eqref{equa 3} are chosen by
\begin{equation*}
\alpha=\frac{\kappa_1}{4\pi h^2}  -\frac{\kappa_2}{2\pi \lambda_{1,*}(\lambda_{1,*}+\lambda_{2,*})}=\frac{\kappa_2}{4\pi h^2}  -\frac{\kappa_1}{2\pi \lambda_{2,*}(\lambda_{1,*}+\lambda_{2,*})},\ \ \ \ \beta_i=\frac{\kappa_i}{2\pi}.
\end{equation*}
Denote $ q_1(x)=\frac{\alpha}{2}|x|^2+\beta_1 $ and $ q_2(x)=\frac{\alpha}{2}|x|^2+\beta_2. $

We choose 
\begin{equation*} 
\begin{split}
\Lambda_{\varepsilon, 2}=\bigg\{Z=(  z_1,  z_2) \mid &z_1\in B_{\frac{\rho_0}{2\sqrt{|\ln\varepsilon|}}}\left (\frac{(\lambda_{1,*},0)}{\sqrt{|\ln\varepsilon|}}\right ), z_2\in B_{\frac{\rho_0}{2\sqrt{|\ln\varepsilon|}}}\left (\frac{(-\lambda_{2,*},0)}{\sqrt{|\ln\varepsilon|}}\right )\bigg\}.
\end{split}
\end{equation*}
We will find critical points of $ K_\varepsilon(Z) $ over $ \Lambda_{\varepsilon, 2} $. By rotational-invariance of $ K_\varepsilon(z_1,z_2) $, we assume that $ z_1 $ is on $ x_1 $-axis. Define $ \overline{z_2} $ as the symmetry point of $ z_2 $ with respect to the $ x_1 $-axis. Since $ K_\varepsilon(z_1,z_2)=K_\varepsilon(z_1,\overline{z_2}) $, by  principle of symmetric criticality (see \cite{W}), we  can assume that $ z_2 $ is on $ x_1 $-axis. %Our goal is to construct solutions of \eqref{equa 1} being of the form $ V_{\varepsilon,Z}+\omega_{\varepsilon,Z} $, where $ Z\in \Lambda_{\varepsilon, N,1} $.
%Similar to the proof of Theorem \ref{thm0}, we find critical points of $ K_\varepsilon(Z) $ over $ \Lambda_{\varepsilon, N+1} $. Since $ K_\varepsilon(z_0, z_1,\cdots, z_N)=K_\varepsilon(-z_0, z_1,\cdots, z_N) $, $\nabla K_\varepsilon(0, z_1,\cdots, z_N)=0$ for any $ (z_1,\cdots,z_N) $.  Since $ z_j=e^{\textbf{i}\frac{2\pi(j-1)}{N}}z_1 $,   $ K_\varepsilon(0,z_1,\cdots,z_N) $ is a   function of $ z_1 $ %, i,e, $ K_\varepsilon(Z)=  \hat{K}_\varepsilon(z_1)   $ for some $ \hat{K}_\varepsilon $. Moreover, $ \hat{K}_\varepsilon $ 
%and is rotational-invariant, i.e., $ K_\varepsilon(0,z_1,\cdots,z_N)=\tilde{K}_\varepsilon(|z_1|) $ for some $ \tilde{K}_\varepsilon $. 
Let $ z_1=\left( \frac{\lambda_1}{\sqrt{|\ln\varepsilon|}},0\right)  $ and $ z_2=\left( \frac{-\lambda_2}{\sqrt{|\ln\varepsilon|}},0\right)  $ for $\lambda_i\in \left( \lambda_{i,*}-\frac{\rho_0}{2}, \lambda_{i,*}+\frac{\rho_0}{2}\right)   $. %Using  Proposition \ref{order of main term} and 
Using \eqref{602}, one computes that
\begin{equation} \label{703}
\begin{split}
K_\varepsilon&\left (z_1,z_2\right )\\
=&\pi\varepsilon^2|\ln\varepsilon|\sum_{i=1}^2\left( q_i^2\sqrt{\det K_H}\right) (0) -\pi\varepsilon^2\left( q_1q_2 \det K_H \right)(0)\ln|\ln\varepsilon|\\
&+\pi\varepsilon^2\bigg(  \frac{1}{2}\sum_{i=1}^2 \left(q_i^2\sqrt{\det K_H} \right)^{''}(0 )\lambda_i^2+ \frac{(p-1)}{4}   \sum_{i=1}^2\left( q_i^2\sqrt{\det K_H}\right) (0)     \\
&-2\pi  \sum_{i=1}^2 \left( q_i^2  \det K_H\right)(0) \bar{S}_K(0,0)+2\left( q_1q_2  \det K_H\right)(0)\ln (\lambda_1+\lambda_2)\bigg)  
+O\left( \frac{\varepsilon^2\ln|\ln\varepsilon|}{|\ln\varepsilon|^{\frac{1}{2}}}\right).
\end{split}
\end{equation}
Define 
$$H_3(\lambda_1,\lambda_2)=\frac{1}{2}\sum_{i=1}^2 \left(q_i^2\sqrt{\det K_H} \right)^{''}(0 )\lambda_i^2+2\left( q_1q_2 \det K_H\right)(0)\ln (\lambda_1+\lambda_2).$$
We have
\begin{lemma}\label{lm703}
	The function 
	$ H_3(\lambda_1,\lambda_2) $
	has the unqiue critical point $ \lambda_1=\lambda_{1,*}, \lambda_2=\lambda_{2,*} $ in $  \left( \lambda_{1,*}-\frac{\rho_0}{2}, \lambda_{1,*}+\frac{\rho_0}{2}\right)\times \left( \lambda_{2,*}-\frac{\rho_0}{2}, \lambda_{2,*}+\frac{\rho_0}{2}\right) $, which is a maximizer.
\end{lemma}  
\begin{proof}
	Note that $  \left(q_i^2\sqrt{\det K_H} \right)^{''}(0 )=\frac{\beta_i(2\alpha h^2-\beta_i)}{h^2}=-\frac{\kappa_1,\kappa_2}{2\pi^2\lambda_{i,*}(\lambda_{1,*}+\lambda_{2,*})}<0 $, $ \left( q_1q_2 \det K_H\right)(0)=\beta_1\beta_2=\frac{\kappa_1\kappa_2}{4\pi^2} $ and $$ \nabla H_3(\lambda_1,\lambda_2)=\begin{pmatrix}
   \frac{\beta_1(2\alpha h^2-\beta_1)}{h^2} \lambda_1+\frac{2\beta_1\beta_2}{\lambda_1+\lambda_2}    \\
 \frac{\beta_2(2\alpha h^2-\beta_2)}{h^2} \lambda_2+\frac{2\beta_1\beta_2}{\lambda_1+\lambda_2} 
	\end{pmatrix}.  $$ 
So $ \nabla H_3=0 $ implies that $ \lambda_1=\lambda_{1,*}, \lambda_2=\lambda_{2,*} $. Since
$$\nabla^2 H_3(\lambda_{1,*},\lambda_{2,*})=\begin{pmatrix}
\frac{\beta_1(2\alpha h^2-\beta_1)}{h^2}  -\frac{2\beta_1\beta_2}{(\lambda_{1,*}+\lambda_{2,*})^2} &  -\frac{2\beta_1\beta_2}{(\lambda_{1,*}+\lambda_{2,*})^2}   \\
-\frac{2\beta_1\beta_2}{(\lambda_{1,*}+\lambda_{2,*})^2} & \frac{\beta_2(2\alpha h^2-\beta_2)}{h^2}  -\frac{2\beta_1\beta_2}{(\lambda_{1,*}+\lambda_{2,*})^2}
\end{pmatrix}$$
is negative definite, we conclude that $  (\lambda_{1,*},\lambda_{2,*}) $ is the unique maximizer.
\end{proof}
%
%Note that from \eqref{701},
%\begin{equation*}  
%\begin{split}
%K_\varepsilon\left (Z\right )
%=&N\pi\varepsilon^2|\ln\varepsilon|\left( q^2\sqrt{\det K_H}\right)(0)-\frac{N(N-1)\pi\varepsilon^2}{2}\left( q^2\sqrt{\det K_H}\right)(0)\ln|\ln\varepsilon|\\
%&+\pi\varepsilon^2\bigg(  H_1(r)+ \frac{N(p-1)}{4}   \left( q^2\sqrt{\det K_H}\right) (0)      -2N\pi   \left( q^2  \det K_H\right)(0) \bar{S}_K(0,0)\\
%&+\sum_{1\leq j\neq k\leq N}\left( q^2  \det K_H\right)(0)\ln| e^{\textbf{i}\frac{2\pi(j-1)}{N}} - e^{\textbf{i}\frac{2\pi(k-1)}{N}} |\bigg)  
%+\tilde{N}_\varepsilon(z_1).
%\end{split}
%\end{equation*}
%where $ \tilde{N}_\varepsilon(z_1) $ is a $O\left( \frac{\varepsilon^2\ln|\ln\varepsilon|}{|\ln\varepsilon|^{\frac{1}{2}}}\right) -$perturbation term which is invariant under rotation. %Since $ r_* $ is a strict local maximizer of $ H_1 $, 
By Lemma \ref{lm703}, there exists $ z_{i,\varepsilon}=\left( \frac{(-1)^{i-1}\lambda_{i,\varepsilon}}{\sqrt{|\ln\varepsilon|}},0\right)  $ near $ \left( \frac{(-1)^{i-1}\lambda_{i,*}}{\sqrt{|\ln\varepsilon|}},0\right) $  for $ i=1,2 $ being a maximizer of $ K_\varepsilon  $, which yields a solution $ v_\varepsilon $ of \eqref{equa 3}.  Define 
\begin{equation*} 
\begin{split}
w(x_1,x_2,0,0)=&\frac{1}{\varepsilon^2}\sum_{i=1}^2\left (v_\varepsilon\left(x_1,x_2 \right) -q_i(x_1,x_2)|\ln\varepsilon| \right )^{p}_+\textbf{1}_{B_{\frac{\rho_0}{\sqrt{|\ln\varepsilon|}}}\left (\left( \frac{(-1)^{i-1}\lambda_{i,*}}{\sqrt{|\ln\varepsilon|}},0\right)\right )} 
\end{split}
\end{equation*}
and $ w(x_1,x_2,0,t)=w(\bar{R}_{-\alpha|\ln\varepsilon| t}(x_1,x_2),0,0)  $.
%there exist $ \left (\tilde{z}_{1,\varepsilon}, \cdots, \tilde{z}_{N,\varepsilon}\right ) $ such that $ \tilde{z}_{i,\varepsilon}=\tilde{z}_{i}^*+o(1) $ for $ i=1,\cdots,N $ and 
Then $ supp(w_\varepsilon(\cdot,0,0))\subseteq \cup_{i=1}^2 B_{|\ln\varepsilon|^{-1}}\left (z_{i,\varepsilon}\right ) $
and if we define $ A_{\varepsilon,i}=supp(w_\varepsilon(\cdot,0,0))\cap B_{|\ln\varepsilon|^{-1}}\left (z_{i,\varepsilon}\right )  $ for $ i=1,2 $, then %there are $ R_1,R_2>0 $ independent of $ \varepsilon $  such that $ B_{R_1\varepsilon}\left ( z_{i,\varepsilon}\right )\subseteq A_{\varepsilon,i}\subseteq B_{R_2\varepsilon}\left (z_{i,\varepsilon}\right ) $ and  
\begin{equation*}
\int_{A_{\varepsilon,i}}w_\varepsilon dx=\begin{cases}
&\kappa_1+O\left(|\ln\varepsilon|^{-\frac{1}{2}} \right),\ \ \ \ i=1,\\
&\kappa_2+O\left(|\ln\varepsilon|^{-\frac{1}{2}} \right),\ \ \ \ i=2.
\end{cases}
\end{equation*}

Recall  $ \tilde{X}_{j,\varepsilon}(s,t)=\frac{(-1)^{j-1}\lambda_{j,*}}{\sqrt{|\ln\varepsilon|}}e^{-\mathbf{i}\alpha|\ln\varepsilon|t}e^{\mathbf{i}\frac{s}{h}} 
$ for $ j=1,2 $. Based on the above derivation, we get Theorem \ref{thm03}.







\textbf{Proof of Theorem \ref{thm04}:}
Let $   R^*>0, h>0, N=4, $ $ \kappa,\mu>0 $ and $ \lambda_{1,*},\lambda_{2,*}>0 $ are constants satisfying the compatibility condition
\begin{equation*} 
\frac{\kappa}{4\pi h^2}-\frac{\kappa}{4\pi\lambda_{1,*}^2}-\frac{\mu}{\pi(\lambda_{1,*}^2+\lambda_{2,*}^2)}=\frac{\mu}{4\pi h^2}-\frac{\mu}{4\pi\lambda_{2,*}^2}-\frac{\kappa}{\pi(\lambda_{1,*}^2+\lambda_{2,*}^2)}.
\end{equation*}. 
Consider  the following problem
\begin{equation}\label{equa 4}
	\begin{cases}
		-\varepsilon^2\text{div}(K_H(x)\nabla v)=&\sum_{i=1}^2\left (v-\left(\frac{\alpha}{2}|x|^2+\beta_1 \right)|\ln\varepsilon| \right )^{p}_+\textbf{1}_{B_{\frac{\rho_0}{\sqrt{|\ln\varepsilon|}}}\left (\frac{((-1)^{i-1}\lambda_{1,*},0)}{\sqrt{|\ln\varepsilon|}}\right )}\\
		&+\sum_{i=1}^2\left (v-\left(\frac{\alpha}{2}|x|^2+\beta_2 \right)|\ln\varepsilon| \right )^{p}_+\textbf{1}_{B_{\frac{\rho_0}{\sqrt{|\ln\varepsilon|}}}\left (\frac{(0,(-1)^{i-1}\lambda_{2,*})}{\sqrt{|\ln\varepsilon|}}\right )},\ \ \text{in}\  B_{R^*}(0),\\
		v=0,&\ \ \ \ \ \ \ \ \ \ \ \ \ \ \ \ \ \ \ \ \ \ \ \ \ \ \ \ \ \ \ \ \ \ \ \ \ \ \ \ \ \ \ \ \ \ \ \ \ \ \ \ \ \ \ \ \ \ \ \ \text{on}\ \partial B_{R^*}(0),
	\end{cases}
\end{equation}
where  $ \rho_0>0 $ sufficiently small, %, $ \textbf{1}_{D} $ is Heaviside function of some set $ D $, $ C_{\frac{r_*}{\sqrt{|\ln\varepsilon|}}}=\left \{x\mid |x|=\frac{r_*}{\sqrt{|\ln\varepsilon|}}\right \} $ is the circle centered at the origin of radius $ \frac{r_*}{\sqrt{|\ln\varepsilon|}} $, and $ D_{\frac{\delta}{\sqrt{|\ln\varepsilon|}}}\left (C_{\frac{r_*}{\sqrt{|\ln\varepsilon|}}}\right )=\left \{x\mid |x|\in \left (\frac{r_*-\delta}{\sqrt{|\ln\varepsilon|}},\frac{r_*+\delta}{\sqrt{|\ln\varepsilon|}} \right )\right \} $ is the $ \frac{\delta}{\sqrt{|\ln\varepsilon|}} $-neighborhood of $ C_{\frac{r_*}{\sqrt{|\ln\varepsilon|}}} $.  
$ \alpha $ and $ \beta_i $ in \eqref{equa 4} are chosen by
\begin{equation*}
	\alpha=\frac{\kappa}{4\pi h^2}-\frac{\kappa}{4\pi\lambda_{1,*}^2}-\frac{\mu}{\pi(\lambda_{1,*}^2+\lambda_{2,*}^2)},\ \ \ \ \beta_1=\frac{\kappa}{2\pi},\ \ \ \ \beta_2=\frac{\mu}{2\pi}.
\end{equation*}
Denote $ q_i(x)=\frac{\alpha}{2}|x|^2+\beta_i $ for $ i=1,2. $

We choose 
\begin{equation*} 
	\begin{split}
		\Lambda_{\varepsilon, 4}=\bigg\{Z=(  z_1,  z_2, z_3, z_4) \mid &z_1\in B_{\frac{\rho_0}{2\sqrt{|\ln\varepsilon|}}}\left (\frac{(\lambda_{1,*},0)}{\sqrt{|\ln\varepsilon|}}\right ), z_2\in B_{\frac{\rho_0}{2\sqrt{|\ln\varepsilon|}}}\left (\frac{(0,\lambda_{2,*})}{\sqrt{|\ln\varepsilon|}}\right ),\\
		&z_3=-z_1,z_4=-z_2\bigg\}.
	\end{split}
\end{equation*}
%We will find critical points of $ K_\varepsilon(Z) $ over $ \Lambda_{\varepsilon, 2} $. By rotational-invariance of $ K_\varepsilon(z_1,z_2) $, we assume that $ z_1 $ is on $ x_1 $-axis. Define $ \overline{z_2} $ as the symmetry point of $ z_2 $ with respect to the $ x_1 $-axis. Since $ K_\varepsilon(z_1,z_2)=K_\varepsilon(z_1,\overline{z_2}) $, by  principle of symmetric criticality (see \cite{W}), we  can assume that $ z_2 $ is on $ x_1 $-axis. %Our goal is to construct solutions of \eqref{equa 1} being of the form $ V_{\varepsilon,Z}+\omega_{\varepsilon,Z} $, where $ Z\in \Lambda_{\varepsilon, N,1} $.
%Similar to the proof of Theorem \ref{thm0}, we find critical points of $ K_\varepsilon(Z) $ over $ \Lambda_{\varepsilon, N+1} $. Since $ K_\varepsilon(z_0, z_1,\cdots, z_N)=K_\varepsilon(-z_0, z_1,\cdots, z_N) $, $\nabla K_\varepsilon(0, z_1,\cdots, z_N)=0$ for any $ (z_1,\cdots,z_N) $.  Since $ z_j=e^{\textbf{i}\frac{2\pi(j-1)}{N}}z_1 $,   $ K_\varepsilon(0,z_1,\cdots,z_N) $ is a   function of $ z_1 $ %, i,e, $ K_\varepsilon(Z)=  \hat{K}_\varepsilon(z_1)   $ for some $ \hat{K}_\varepsilon $. Moreover, $ \hat{K}_\varepsilon $ 
%and is rotational-invariant, i.e., $ K_\varepsilon(0,z_1,\cdots,z_N)=\tilde{K}_\varepsilon(|z_1|) $ for some $ \tilde{K}_\varepsilon $. 
From the symmetry of $ K_\varepsilon $, we can assume $ z_1=\left( \frac{\lambda_1}{\sqrt{|\ln\varepsilon|}},0\right)  $ and $ z_2=\left( 0,\frac{\lambda_2}{\sqrt{|\ln\varepsilon|}}\right)  $ for $\lambda_i\in \left( \lambda_{i,*}-\frac{\rho_0}{2}, \lambda_{i,*}+\frac{\rho_0}{2}\right)   $. %Using  Proposition \ref{order of main term} and 
Similar to \eqref{602}, we have
\begin{equation} \label{704}
	\begin{split}
		K_\varepsilon&\left (z_1,z_2,z_3,z_4\right )\\
		=&2\pi\varepsilon^2|\ln\varepsilon|\sum_{i=1}^2\left( q_i^2\sqrt{\det K_H}\right) (0) -\pi\varepsilon^2\left( \left( q_1^2+4q_1q_2+q_2^2 \right) \det K_H \right)(0)\ln|\ln\varepsilon|\\
		&+\pi\varepsilon^2\bigg(    H_4(\lambda_1,\lambda_2)+\frac{(p-1)}{2}   \sum_{i=1}^2\left( q_i^2\sqrt{\det K_H}\right) (0)      -4\pi  \sum_{i=1}^2 \left( q_i^2  \det K_H\right)(0) \bar{S}_K(0,0)\bigg) \\ 
		&+O\left( \frac{\varepsilon^2\ln|\ln\varepsilon|}{|\ln\varepsilon|^{\frac{1}{2}}}\right).
	\end{split}
\end{equation}
where
\begin{equation*} 
\begin{split}
H_4(\lambda_1,\lambda_2)=&\sum_{i=1}^2 \left(q_i^2\sqrt{\det K_H} \right)^{''}(0 )\lambda_i^2+2\sum_{i=1}^2\left( q_i^2 \det K_H\right)(0)\ln 2\lambda_i\\
& +4\left( q_1q_2  \det K_H\right)(0)\ln (\lambda_1^2+\lambda_2^2) .
\end{split}
\end{equation*}
$$$$
It is not hard to verify that 	$ H_4(\lambda_1,\lambda_2) $
	has the unqiue maximizer $ \lambda_1=\lambda_{1,*}, \lambda_2=\lambda_{2,*} $ in $  \left( \lambda_{1,*}-\frac{\rho_0}{2}, \lambda_{1,*}+\frac{\rho_0}{2}\right)\times \left( \lambda_{2,*}-\frac{\rho_0}{2}, \lambda_{2,*}+\frac{\rho_0}{2}\right) $.   
%
%Note that from \eqref{701},
%\begin{equation*}  
%\begin{split}
%K_\varepsilon\left (Z\right )
%=&N\pi\varepsilon^2|\ln\varepsilon|\left( q^2\sqrt{\det K_H}\right)(0)-\frac{N(N-1)\pi\varepsilon^2}{2}\left( q^2\sqrt{\det K_H}\right)(0)\ln|\ln\varepsilon|\\
%&+\pi\varepsilon^2\bigg(  H_1(r)+ \frac{N(p-1)}{4}   \left( q^2\sqrt{\det K_H}\right) (0)      -2N\pi   \left( q^2  \det K_H\right)(0) \bar{S}_K(0,0)\\
%&+\sum_{1\leq j\neq k\leq N}\left( q^2  \det K_H\right)(0)\ln| e^{\textbf{i}\frac{2\pi(j-1)}{N}} - e^{\textbf{i}\frac{2\pi(k-1)}{N}} |\bigg)  
%+\tilde{N}_\varepsilon(z_1).
%\end{split}
%\end{equation*}
%where $ \tilde{N}_\varepsilon(z_1) $ is a $O\left( \frac{\varepsilon^2\ln|\ln\varepsilon|}{|\ln\varepsilon|^{\frac{1}{2}}}\right) -$perturbation term which is invariant under rotation. %Since $ r_* $ is a strict local maximizer of $ H_1 $, 
Thus from \eqref{704}, there exists $ z_{2i-1,\varepsilon}=\left( \frac{(-1)^{i-1}\lambda_{1,\varepsilon}}{\sqrt{|\ln\varepsilon|}},0\right)  $ near $ \left( \frac{(-1)^{i-1}\lambda_{1,*}}{\sqrt{|\ln\varepsilon|}},0\right) $ and $ z_{2i,\varepsilon}=\left( 0,\frac{(-1)^{i-1}\lambda_{2,\varepsilon}}{\sqrt{|\ln\varepsilon|}}\right)  $ near $ \left( 0,\frac{(-1)^{i-1}\lambda_{2,*}}{\sqrt{|\ln\varepsilon|}}\right) $ for $ i=1,2 $ being a maximizer of $ K_\varepsilon  $, which yields a solution $ v_\varepsilon $ of \eqref{equa 4}.   Similar to Theorem \ref{thm03}, we get Theorem \ref{thm04}.






\textbf{Proof of Theorem \ref{thm05}:}
Let $   R^*>0, h>0, N=5, $ $ \kappa_0,\kappa,\mu>0 $ and $ \lambda_{1,*},\lambda_{2,*}>0 $ are constants satisfying the compatibility condition
\begin{equation*} 
\frac{\kappa}{4\pi h^2}-\frac{\kappa_0}{2\pi\lambda_{1,*}^2}-\frac{\kappa}{4\pi\lambda_{1,*}^2}-\frac{\mu}{\pi(\lambda_{1,*}^2+\lambda_{2,*}^2)}=\frac{\mu}{4\pi h^2}-\frac{\kappa_0}{2\pi\lambda_{2,*}^2}-\frac{\mu}{4\pi\lambda_{2,*}^2}-\frac{\kappa}{\pi(\lambda_{1,*}^2+\lambda_{2,*}^2)}.
\end{equation*}. 
Consider  the following problem
\begin{equation}\label{equa 5}
\begin{cases}
-\varepsilon^2\text{div}(K_H(x)\nabla v)=&\sum_{i=1}^2\left (v-\left(\frac{\alpha}{2}|x|^2+\beta_1 \right)|\ln\varepsilon| \right )^{p}_+\textbf{1}_{B_{\frac{\rho_0}{\sqrt{|\ln\varepsilon|}}}\left (\frac{((-1)^{i-1}\lambda_{1,*},0)}{\sqrt{|\ln\varepsilon|}}\right )}\\
&+\sum_{i=1}^2\left (v-\left(\frac{\alpha}{2}|x|^2+\beta_2 \right)|\ln\varepsilon| \right )^{p}_+\textbf{1}_{B_{\frac{\rho_0}{\sqrt{|\ln\varepsilon|}}}\left (\frac{(0,(-1)^{i-1}\lambda_{2,*})}{\sqrt{|\ln\varepsilon|}}\right )}\\
&+\left (v-\left(\frac{\alpha}{2}|x|^2+\beta_0 \right)|\ln\varepsilon| \right )^{p}_+\textbf{1}_{B_{\frac{\rho_0}{\sqrt{|\ln\varepsilon|}}}\left (0\right )},\ \ \ \ \ \ \ \ \ \text{in}\  B_{R^*}(0),\\
v=0,&\ \ \ \ \ \ \ \ \ \ \ \ \ \ \ \ \ \ \ \ \ \ \ \ \ \ \ \ \ \ \ \ \ \ \ \ \ \ \ \ \ \ \ \ \ \ \ \ \ \ \ \ \ \ \ \ \ \ \ \ \text{on}\ \partial B_{R^*}(0),
\end{cases}
\end{equation}
where  $ \rho_0>0 $ sufficiently small, %, $ \textbf{1}_{D} $ is Heaviside function of some set $ D $, $ C_{\frac{r_*}{\sqrt{|\ln\varepsilon|}}}=\left \{x\mid |x|=\frac{r_*}{\sqrt{|\ln\varepsilon|}}\right \} $ is the circle centered at the origin of radius $ \frac{r_*}{\sqrt{|\ln\varepsilon|}} $, and $ D_{\frac{\delta}{\sqrt{|\ln\varepsilon|}}}\left (C_{\frac{r_*}{\sqrt{|\ln\varepsilon|}}}\right )=\left \{x\mid |x|\in \left (\frac{r_*-\delta}{\sqrt{|\ln\varepsilon|}},\frac{r_*+\delta}{\sqrt{|\ln\varepsilon|}} \right )\right \} $ is the $ \frac{\delta}{\sqrt{|\ln\varepsilon|}} $-neighborhood of $ C_{\frac{r_*}{\sqrt{|\ln\varepsilon|}}} $.  
$ \alpha $ and $ \beta_i $ in \eqref{equa 5} are chosen by
\begin{equation*}
\alpha=\frac{\kappa}{4\pi h^2}-\frac{\kappa_0}{2\pi\lambda_{1,*}^2}-\frac{\kappa}{4\pi\lambda_{1,*}^2}-\frac{\mu}{\pi(\lambda_{1,*}^2+\lambda_{2,*}^2)},\ \ \ \ \beta_0=\frac{\kappa_0}{2\pi},\ \ \ \ \beta_1=\frac{\kappa}{2\pi},\ \ \ \ \beta_2=\frac{\mu}{2\pi}.
\end{equation*}
Denote $ q_i(x)=\frac{\alpha}{2}|x|^2+\beta_i $ for $ i=0,1,2. $

We choose 
\begin{equation*} 
\begin{split}
\Lambda_{\varepsilon, 5}=\bigg\{Z=(z_0,  z_1,  z_2, z_3, z_4) \mid &z_1\in B_{\frac{\rho_0}{2\sqrt{|\ln\varepsilon|}}}\left (\frac{(\lambda_{1,*},0)}{\sqrt{|\ln\varepsilon|}}\right ), z_2\in B_{\frac{\rho_0}{2\sqrt{|\ln\varepsilon|}}}\left (\frac{(0,\lambda_{2,*})}{\sqrt{|\ln\varepsilon|}}\right ),\\
&z_3=-z_1,z_4=-z_2, z_0\in B_{\frac{\rho_0}{2\sqrt{|\ln\varepsilon|}}}\left (0\right )\bigg\}.
\end{split}
\end{equation*}
%We will find critical points of $ K_\varepsilon(Z) $ over $ \Lambda_{\varepsilon, 2} $. By rotational-invariance of $ K_\varepsilon(z_1,z_2) $, we assume that $ z_1 $ is on $ x_1 $-axis. Define $ \overline{z_2} $ as the symmetry point of $ z_2 $ with respect to the $ x_1 $-axis. Since $ K_\varepsilon(z_1,z_2)=K_\varepsilon(z_1,\overline{z_2}) $, by  principle of symmetric criticality (see \cite{W}), we  can assume that $ z_2 $ is on $ x_1 $-axis. %Our goal is to construct solutions of \eqref{equa 1} being of the form $ V_{\varepsilon,Z}+\omega_{\varepsilon,Z} $, where $ Z\in \Lambda_{\varepsilon, N,1} $.
%Similar to the proof of Theorem \ref{thm0}, we find critical points of $ K_\varepsilon(Z) $ over $ \Lambda_{\varepsilon, N+1} $. Since $ K_\varepsilon(z_0, z_1,\cdots, z_N)=K_\varepsilon(-z_0, z_1,\cdots, z_N) $, $\nabla K_\varepsilon(0, z_1,\cdots, z_N)=0$ for any $ (z_1,\cdots,z_N) $.  Since $ z_j=e^{\textbf{i}\frac{2\pi(j-1)}{N}}z_1 $,   $ K_\varepsilon(0,z_1,\cdots,z_N) $ is a   function of $ z_1 $ %, i,e, $ K_\varepsilon(Z)=  \hat{K}_\varepsilon(z_1)   $ for some $ \hat{K}_\varepsilon $. Moreover, $ \hat{K}_\varepsilon $ 
%and is rotational-invariant, i.e., $ K_\varepsilon(0,z_1,\cdots,z_N)=\tilde{K}_\varepsilon(|z_1|) $ for some $ \tilde{K}_\varepsilon $. 
From the symmetry of $ K_\varepsilon $, we can assume $ z_0=0 $, $ z_1=\left( \frac{\lambda_1}{\sqrt{|\ln\varepsilon|}},0\right)  $ and $ z_2=\left( 0,\frac{\lambda_2}{\sqrt{|\ln\varepsilon|}}\right)  $ for $\lambda_i\in \left( \lambda_{i,*}-\frac{\rho_0}{2}, \lambda_{i,*}+\frac{\rho_0}{2}\right)   $. %Using  Proposition \ref{order of main term} and 
Similar to \eqref{602}, we have
\begin{equation} \label{705}
\begin{split}
K_\varepsilon&\left (z_0,z_1,z_2,z_3,z_4\right )\\
=&\pi\varepsilon^2|\ln\varepsilon|\left( 2\sum_{i=1}^2\left( q_i^2\sqrt{\det K_H}\right) (0)+\left( q_0^2\sqrt{\det K_H}\right) (0) \right) \\
&-\pi\varepsilon^2\left( \left( q_1^2+4q_1q_2+q_2^2+2q_1q_0+2q_2q_0 \right) \det K_H \right)(0)\ln|\ln\varepsilon|\\
&+\pi\varepsilon^2\bigg(    H_5(\lambda_1,\lambda_2)+\frac{(p-1)}{2}   \sum_{i=1}^2\left( q_i^2\sqrt{\det K_H}\right) (0)   + \frac{(p-1)}{4}\left( q_0^2\sqrt{\det K_H}\right) (0)  \\
& -4\pi  \sum_{i=1}^2 \left( q_i^2  \det K_H\right)(0) \bar{S}_K(0,0)-2\pi   \left( q_0^2  \det K_H\right)(0) \bar{S}_K(0,0)\bigg) \\ 
&+O\left( \frac{\varepsilon^2\ln|\ln\varepsilon|}{|\ln\varepsilon|^{\frac{1}{2}}}\right).
\end{split}
\end{equation}
where
\begin{equation*} 
\begin{split}
H_5(\lambda_1,\lambda_2)=&\sum_{i=1}^2 \left(q_i^2\sqrt{\det K_H} \right)^{''}(0 )\lambda_i^2+2\sum_{i=1}^2\left( q_i^2 \det K_H\right)(0)\ln 2\lambda_i \\
&+4\left( q_1q_2  \det K_H\right)(0)\ln (\lambda_1^2+\lambda_2^2)+4\sum_{i=1}^2\left( q_0q_i \det K_H\right)(0)\ln \lambda_i .
\end{split}
\end{equation*}
$$$$
Since	$ H_5(\lambda_1,\lambda_2) $
has the unqiue maximizer $ \lambda_1=\lambda_{1,*}, \lambda_2=\lambda_{2,*} $ in $  \left( \lambda_{1,*}-\frac{\rho_0}{2}, \lambda_{1,*}+\frac{\rho_0}{2}\right)\times \left( \lambda_{2,*}-\frac{\rho_0}{2}, \lambda_{2,*}+\frac{\rho_0}{2}\right) $,  
%
%Note that from \eqref{701},
%\begin{equation*}  
%\begin{split}
%K_\varepsilon\left (Z\right )
%=&N\pi\varepsilon^2|\ln\varepsilon|\left( q^2\sqrt{\det K_H}\right)(0)-\frac{N(N-1)\pi\varepsilon^2}{2}\left( q^2\sqrt{\det K_H}\right)(0)\ln|\ln\varepsilon|\\
%&+\pi\varepsilon^2\bigg(  H_1(r)+ \frac{N(p-1)}{4}   \left( q^2\sqrt{\det K_H}\right) (0)      -2N\pi   \left( q^2  \det K_H\right)(0) \bar{S}_K(0,0)\\
%&+\sum_{1\leq j\neq k\leq N}\left( q^2  \det K_H\right)(0)\ln| e^{\textbf{i}\frac{2\pi(j-1)}{N}} - e^{\textbf{i}\frac{2\pi(k-1)}{N}} |\bigg)  
%+\tilde{N}_\varepsilon(z_1).
%\end{split}
%\end{equation*}
%where $ \tilde{N}_\varepsilon(z_1) $ is a $O\left( \frac{\varepsilon^2\ln|\ln\varepsilon|}{|\ln\varepsilon|^{\frac{1}{2}}}\right) -$perturbation term which is invariant under rotation. %Since $ r_* $ is a strict local maximizer of $ H_1 $, 
from \eqref{705}, there exists $ z_0=0 $, $ z_{2i-1,\varepsilon}=\left( \frac{(-1)^{i-1}\lambda_{1,\varepsilon}}{\sqrt{|\ln\varepsilon|}},0\right)  $ near $ \left( \frac{(-1)^{i-1}\lambda_{1,*}}{\sqrt{|\ln\varepsilon|}},0\right) $ and $ z_{2i,\varepsilon}=\left( 0,\frac{(-1)^{i-1}\lambda_{2,\varepsilon}}{\sqrt{|\ln\varepsilon|}}\right)  $ near $ \left( 0,\frac{(-1)^{i-1}\lambda_{2,*}}{\sqrt{|\ln\varepsilon|}}\right) $ for $ i=1,2 $ being a critical point of $ K_\varepsilon  $, which yields a solution $ v_\varepsilon $ of \eqref{equa 5}.   Similar to Theorem \ref{thm03}, we get Theorem \ref{thm05}.








\iffalse

Now we consider the existence of critical points of $ K_\varepsilon $ on $ \Lambda_{\varepsilon, N} $ and show that this critical point precisely corresponds to the vorticity field of Euler equations \eqref{Euler eq} concentrated near the filaments in Theorem \ref{thm1}. Recall that $ K=K_H $ satisfies \eqref{coef matrix} and $ q(x)=\frac{\alpha}{2}|x|^2+\beta $, where $ \alpha $ and $ \beta $ satisfies \eqref{choice of a,b}.

We give some notations of symmetry groups. Denote $ \sigma $ a permutation from $ \{1,\cdots,N\} $ to $ \{1,\cdots,N\} $ and  $ S_N  $ the $ N-$element permutation group.
Denote $ SO(2) $ the group of two-dimensional rotational transformations and $ O(2) $ the group of two-dimensional orthogonal  transformations.  

We observe that $ K_\varepsilon(Z) $ is invariant under the group $ S_N, SO(2) $ and $ O(2) $. 
\begin{lemma}
For any $ \sigma\in S_N $, $ \textbf{Q}_1\in SO(2), \textbf{Q}_2\in O(2) $, there holds
\begin{equation*} 
K_\varepsilon(z_1,z_2,\cdots, z_N)=K_\varepsilon(z_{\sigma(1)},z_{\sigma(2)},\cdots, z_{\sigma(N)}).
\end{equation*}
\end{lemma}








\section{Proof of Theorem \ref{thm2}}
It suffices to consider solutions to the problem
\begin{equation}\label{eq01}
\begin{cases}
-\varepsilon^2\text{div}(K_H(x)\nabla u)=  \left(u-\left( \frac{\alpha|x|^2}{2}+\beta\right)\ln\frac{1}{\varepsilon}\right)^{p}_+,\ \ &x\in B_{R^*}(0),\\
u=0,\ \ &x\in\partial  B_{R^*}(0).
\end{cases}
\end{equation}
%where $ \alpha,\beta  $ are any given constants satisfying $ \min_{x\in B_{R^*}(0)}\frac{\alpha|x|^2}{2}+\beta>0 $.
Let $ v=u\backslash|\ln\varepsilon|  $ and $ \delta=\varepsilon|\ln\varepsilon|^{-\frac{p-1}{2}} $, then
\begin{equation}\label{eq02}
\begin{cases}
-\delta^2\text{div}(K_H(x)\nabla v)=  \left( v-\left( \frac{\alpha|x|^2}{2}+\beta\right) \right)^{p}_+,\ \ &x\in B_{R^*}(0),\\
v=0,\ \ &x\in\partial  B_{R^*}(0).
\end{cases}
\end{equation}
Note that \eqref{eq02} coincides with \eqref{111} with $ q=\frac{\alpha|x|^2}{2}+\beta $, $ K=K_H $ and $ \Omega=B_{R^*}(0) $. However,  since $ q^2\sqrt{det(K_H)} $ is a radial function and the set of  extreme points is rotational-invariant, results of Theorem \ref{thm2} can not be deduced directly from those of Theorem \ref{thm1}. %In this case one can also use the reduction procedure to construct solutions of \eqref{eq02}, by using the rotational symmetry of $ K_H, q $ and the domain $ B_{R^*}(0). $

%Let $ h(r)=h(|x|)=q^2\sqrt{det(K_H)}(x) $ for any $ x\in B_{R^*}(0) $. We call $ z^* $ is a  $ strict $ $ local $ $ maximum $ $ (minimum) $ $ point $ of  $q^2\sqrt{det(K_H)}$  $ up $ $ to $ $ ratation $ in $ B_{R^*}(0) $, if $ |z^*| $ is a strict local maximum (minimum) point of $ h $ in $ (0,{R^*}) $.

Let $ x_0  $ be a strict local maximizer of $ q^2\sqrt{\det K_H} $ up to a rotation. %Define $ \mathcal{N}=B_{\bar{\rho}}(z_1) $. Then we can  construct solutions of \eqref{eq02} being of the form $ v_\delta=PV_{\delta,z,\hat{q}}+\omega_{\delta,z} $, where $ z $ is near $ z_1. $
By Lemma \ref{coercive esti} and Proposition \ref{exist and uniq of w}, for any $ Z\in \Lambda_{\varepsilon, m} $ there exists a unique $ \omega_{\delta,Z}\in E_{\delta,Z} $ such that $ Q_\delta L_\delta \omega_{\delta,Z}=Q_\delta l_\delta+Q_\delta R_\delta(\omega_{\delta,Z}). $ So it remains to prove the existence of  maximizers $ Z=Z_{\delta} $ of $ K_\delta  $  near $ x_0 $.  Indeed,  by the rotational symmetry of $ q $ and  $ \det K_H $, one computes directly that
\begin{equation}\label{501}
\begin{array}{ll}
K_\delta(Z)=\sum_{j=1}^m\frac{\pi\delta^2}{\ln\frac{1}{\varepsilon}}q^2\sqrt{\det K}(z_j)-\sum_{1\leq i\neq j\leq m}\frac{2\pi^2\delta^2q(z_i)q(z_j)\sqrt{\det K(z_i)}\sqrt{\det K(z_j)}}{\left( \ln\frac{1}{\varepsilon}\right) ^2}G_K(z_i,z_j)&\\
\qquad\qquad\,+\tilde{N}_\delta(Z),&
\end{array}
\end{equation}
where $ \tilde{N}_\delta(Z) $ is a $ O\left (\frac{\delta^2}{|\ln\varepsilon|^2}\right ) -$perturbation term which is invariant under a rotation. %In fact, we can choose $ T_x^{-1}=\begin{pmatrix}
%\cos\theta_x & -\sin\theta_x  \\
%\sin\theta_x &\cos\theta_x
%\end{pmatrix}\begin{pmatrix}
%\frac{k}{\sqrt{k^2+|x|^2}} & 0  \\
%0 &1
%\end{pmatrix}\begin{pmatrix}
%\cos\theta_x &  \sin\theta_x  \\
%-\sin\theta_x &\cos\theta_x
%\end{pmatrix} $ in \eqref{T_z choice}, where $ (|x|,\theta_x) $ is the polar coordinate of $ x $. By the rotational symmetry of $ q $ and  the domain $ B_R(0) $, one can  prove that for every $ \theta\in[0,2\pi] $,  if we define $ \bar{z}=\bar{R}_\theta(z) $, then
%\begin{equation*}
%V_{\delta,\bar{z}}(\bar{R}_\theta(x))=V_{\delta,z}(x)\ \text{and} \ \omega_{\delta,\bar{z}}(\bar{R}_\theta(x))=\omega_{\delta,z}(x)    \ \ \text{for any}\  x\in B_{R^*}(0).
%\end{equation*}
%Hence one computes directly  that $ P_\delta(\bar{z})=P_\delta(z) $, i.e., $  P_\delta $ is a radially symmetric function. Note that $ q^2\sqrt{det(K_H)}(x)=\left( \frac{\alpha|x|^2}{2}+\beta\right)^2\cdot \frac{k}{\sqrt{k^2+|x|^2}} $ is also radially symmetric. Thus by  Proposition \ref{pro401} and \ref{order of main term},  we have \eqref{501}.
%Since $ z_1 $ is a strict local maximum (minimum) point of $ q^2\sqrt{det(K_H)} $ up to rotation,
So it is not hard to prove  the existence of $ Z_{\delta} $ near $ x_0 $ being a maximizer of $ K_\delta  $, which yields a solution $ v_\delta $ of \eqref{eq02}. Let $ u_\varepsilon=v_\delta|\ln\varepsilon| $, then $ u_\varepsilon $ is a solution of \eqref{eq01}. Moreover, one has
\begin{equation*}
\lim_{\varepsilon\to 0}\frac{1}{\varepsilon^2}\int_{\Omega}\left( u_\varepsilon-q\ln\frac{1}{\varepsilon}\right)^{p}_+dx=2\pi m q \sqrt{\det(K_H)}(x_0)=\frac{km\pi(\alpha|x_0|^2+2\beta)}{\sqrt{k^2+|x_0|^2}}.
\end{equation*}
\textbf{Proof of  Corollary \ref{cor3}}: We choose $ \alpha $ and $ \beta $ such that $ \alpha<0 $ and $ \min_{x\in B_{R^*}(0)}\left( \frac{\alpha|x|^2}{2}+\beta\right)>0 $ in Theorem \ref{thm2}. Then $ (0,0) $ is the unique strict local maximizer of $ \left( \frac{\alpha|x|^2}{2}+\beta\right)^2\sqrt{\det K_H} $ up to a rotation. Thus by Theorem \ref{thm2}, for any $ m\in \mathbb{N}^* $ there exist clustered helical rotational-invariant vorticity fields $ \textbf{w}_\varepsilon $ to \eqref{Euler eq2}  with angular velocity $ \alpha|\ln\varepsilon| $, whose support sets are $ m $ helical tubes and collapse into $ x_3- $axis as $\varepsilon\to 0$. Moreover, the circulations satisfy as $\varepsilon\to 0$
\begin{equation*}
\int_{B_{R^*}(0)} \omega_\varepsilon dx\to 2\pi m \beta.
\end{equation*}
 
\fi



%%%%%%
\iffalse


To conclude, we have
\begin{theorem}\label{thmA}
Let $ \alpha,\beta $ be two constants satisfying $ \min_{x\in B_{R^*}(0)}\left( \frac{\alpha|x|^2}{2}+\beta\right) >0 $ and $ z_1\in B_{R^*}(0)$ be a strict local maximum (minimum) point of $ \left( \frac{\alpha|x|^2}{2}+\beta\right)^2\cdot \frac{k}{\sqrt{k^2+|x|^2}} $ up to rotation. Then there exists $ \varepsilon_0>0 $, such that for any $ \varepsilon\in(0,\varepsilon_0] $, \eqref{eq01} has a solution $ u_\varepsilon $ satisfying the following properties:
\begin{enumerate}
	\item  Define $ A_\varepsilon=\left\{u_\varepsilon>\left( \frac{\alpha|x|^2}{2}+\beta\right)\ln\frac{1}{\varepsilon}\right\} $. Then $ \lim\limits_{\varepsilon\to 0}dist(A_\varepsilon, z_1)=0 $.
	\item $ \lim\limits_{\varepsilon\to 0}\frac{1}{\varepsilon^2}\int_{A_\varepsilon}\left( u_\varepsilon-\left( \frac{\alpha|x|^2}{2}+\beta\right)\ln\frac{1}{\varepsilon}\right)^p_+dx=\frac{k\pi(\alpha|z_1|^2+2\beta)}{\sqrt{k^2+|z_1|^2}}. $
	\item There exist  $ R_1,R_2>0 $ satisfying
	\begin{equation*}
	R_1\varepsilon\leq \text{diam}(A_\varepsilon)\leq R_2\varepsilon.
	\end{equation*}
	
\end{enumerate}
\end{theorem}


\bigskip



\textbf{Proof of Theorem \ref{thm01}}:
To prove Theorem \ref{thm01}, we define for every $ r_*\in (0,R^*) $, $ c>0 $ and
\begin{equation}\label{502}
\alpha=\frac{c}{4\pi k\sqrt{k^2+r_*^2}},\ \ \beta=\frac{\alpha}{2}(3r_*^2+4k^2).
\end{equation}
One computes directly that $ (r_*,0)$ is a strict minimum point of $ q^2\sqrt{det(K_H)}(x)=\left( \frac{\alpha|x|^2}{2}+\beta\right)^2\cdot \frac{k}{\sqrt{k^2+|x|^2}} $ up to rotation and that
\begin{equation*}
2\pi q\sqrt{det(K_H)}((r_*,0))=\frac{k\pi(\alpha r_*^2+2\beta)}{\sqrt{k^2+r_*^2}}=c.
\end{equation*}
Hence by Theorem \ref{thmA}, for any $ \varepsilon $ small there exists a solution $ u_\varepsilon $ of \eqref{eq01} concentrating near $ (r_*,0) $ with $ \alpha=\frac{c}{4\pi k\sqrt{k^2+r_*^2}} $ and $\frac{1}{\varepsilon^2} \int_{A_\varepsilon}\left( u_\varepsilon-\left( \frac{\alpha|x|^2}{2}+\beta\right)\ln\frac{1}{\varepsilon}\right)^p_+dx $ tending to $ c $. Define for any $ (x_1,x_2,x_3)\in B_{R^*}(0)\times\mathbb{R}, t\in\mathbb{R}  $
\begin{equation*}
\mathbf{w}_\varepsilon(x_1,x_2,x_3,t)= \frac{w_\varepsilon(x_1,x_2,x_3,t)}{k}\overrightarrow{\zeta},
\end{equation*}
where $ w_\varepsilon(x_1,x_2,x_3,t) $ is a helical function satisfying
\begin{equation}\label{503}
w_\varepsilon(x_1,x_2,0,t)=\frac{1}{\varepsilon^2}\left( u_\varepsilon(\bar{R}_{-\alpha|\ln\varepsilon| t}(x_1,x_2))-\left( \frac{\alpha|(x_1,x_2)|^2}{2}+\beta\right)\ln\frac{1}{\varepsilon}\right)^p_+.
\end{equation}
Direct computations show that $ w_\varepsilon(x_1,x_2,0,t) $ satisfies \eqref{vor str eq} and $ \mathbf{w}_\varepsilon $ is a left-handed helical vorticity field of \eqref{Euler eq2}. Moreover, $ w_\varepsilon(x_1,x_2,0,t) $ rotates clockwise around the origin with angular velocity $ \alpha|\ln\varepsilon| $. By \eqref{1001}, the circulation of $ \mathbf{w}_\varepsilon $ satisfies
\begin{equation*}
\iint_{A_\varepsilon}\mathbf{w}_\varepsilon\cdot\mathbf{n}d\sigma=\frac{1}{\varepsilon^2} \int_{A_\varepsilon}\left( u_\varepsilon-\left( \frac{\alpha|x|^2}{2}+\beta\right)\ln\frac{1}{\varepsilon}\right)^p_+dx\to c,\ \ \text{as}\ \varepsilon\to0.
\end{equation*}



It suffices to prove that the vorticity field $ \mathbf{w}_\varepsilon $ tends asymptotically to \eqref{1007} in sense of \eqref{1005}.  Define $ P(\tau) $ the intersection point of the curve  parameterized by \eqref{1007} and the $ x_1Ox_2 $ plane. Note that the helix \eqref{1007} corresponds uniquely to the motion of $ P(\tau) $. Taking $ s=\frac{b_1\tau}{k} $ into \eqref{1007}, one computes directly that  $ P(\tau) $ satisfies a 2D point vortex model
\begin{equation*}
\begin{split}
P(\tau)=\bar{R}_{\alpha'\tau}((r_*,0)),
\end{split}
\end{equation*}
where
\begin{equation*}
\alpha'=\frac{1}{\sqrt{k^2+r_*^2}}\left( a_1+\frac{b_1}{k}\right)=\frac{c}{4\pi k\sqrt{k^2+r_*^2}},
\end{equation*}
which is equal to $ \alpha $ in \eqref{502}. Thus by the construction, we  readily check that the support set of $ w_\varepsilon(x_1,x_2,0,|\ln\varepsilon|^{-1}\tau) $ defined by  \eqref{503} concentrates near $ P(\tau) $ as $ \varepsilon\to 0 $, which implies that, the vorticity field $ \mathbf{w}_\varepsilon $ tends asymptotically to \eqref{1007} in sense of \eqref{1005}.
The proof of Theorem \ref{thm01} is thus complete.

\textbf{Proof of Theorem \ref{thm02}}: The proof of Theorem \ref{thm02} is similar to that of Theorem \ref{thm01}, by constructing multiple concentration solutions $ u_\varepsilon $ of \eqref{eq01} (or equivalently, $ v_\delta $ of \eqref{eq02}) with polygonal symmetry. Indeed, let $ m\geq2 $ be an integer, $ r_*\in (0,R) $, $ c>0 $ and $ \alpha,\beta $ satisfying \eqref{502}. For any $ z_1\in B_{\bar{\rho}}((r_*,0)) $, define $ z_i=\bar{R}_{\frac{2(i-1)\pi}{m}}(z_1) $ for $ i=2,\cdots,m. $ Our goal is to construct solutions of \eqref{eq02} being of the form $  \sum_{i=1}^mPV_{\delta,z_i,\hat{q}_i}+\omega_{\delta} $. Using the symmetry of $ K_H, q $ and $ B_{R^*}(0) $, one can  prove that
\begin{equation*}
P_\delta(z_1,\cdots,z_m)=\frac{m\pi\delta^2}{\ln\frac{1}{\varepsilon}}q^2\sqrt{det(K_H)}(z_1)+\hat{N}_\delta(z_1),
\end{equation*}
where $ \hat{N}_\delta(z_1) $ is a $ O\left( \frac{\delta^2\ln|\ln\varepsilon|}{|\ln\varepsilon|^2}\right)  -$perturbation term which is invariant under rotation. The rest of the proof is exactly the same as
in that of Theorem \ref{thm01} and we omit the details.


\fi
%%%%%%








%%%%%%
 
\setcounter{equation}{0}\renewcommand\theequation{A.\arabic{equation}}
\section{Appendix: Proofs of some results used before}
In this Appendix, we give proofs for some results  which have been repeatedly used before. %We first give some properties of  Green's function of $ \mathcal{L}_K=-\text{div}(K(x)\nabla) $ appeared in the linear problem \eqref{diver form}. Note that for any $ f\in L^p(\Omega) $ for $ p>2 $, there exists a unique solution $ u=\mathcal{G}_Kf $ being a weak solution to \eqref{diver form}.
%Clearly from $ L^p $ theory of elliptic equations, there exists a Green's operator $ \mathcal{G}_K  $ such that for every $ f\in L^{p}(\Omega) $ with $ p>1 $, the function $ u=\mathcal{G}_Kf $ is a weak solution to the problem \eqref{diver form}.
%From theorem 1.3 in \cite{CW1}, $ u $ can be represented as follows
%expansion of Green's function of $ \mathcal{G}_K $ as follows

We first give some explicit helical solutions of  the nearly parallel vortex filament model deduced in \cite{KMD} (see also \cite{GM2}). %Let $ N\geq 2 $ and $ \kappa_j $ be $ N $ fixed real numbers. Consider the motion of $ N $ slender vortex tube in 3D Euler equations whose cross-sections are of radius $ \varepsilon $,  circulations are $ \kappa_j $ and centerlines are filaments parameterized as $ \tilde{X}_{j,\varepsilon}(s, t) $ for $ j=1,\cdots,N $ in the complex plane $ \mathbb{C} $. The wave-length of each vortex tube is of the order $ O(1) $ and the mutual distance of any two vortex tubes is of the order $ O(\frac{1}{\sqrt{|\ln\varepsilon|}}) $. By the deduction of \cite{GM2,KMD}, if we define $   \tau=|\ln\varepsilon|t $ and $ X_j(s, \tau)=\sqrt{|\ln\varepsilon|}\tilde{X}_{j,\varepsilon}(s, t) $, then $ X_j(s, \tau) $ satisfies
%\begin{equation*} 
%\partial_\tau X_j=\frac{1}{4\pi}\left(\mathbf{i}\alpha_j\kappa_j\partial_{ss}X_j+2\mathbf{i}\sum_{k\neq j}\kappa_k\frac{X_j-X_k}{|X_j-X_k|^2} \right),
%\end{equation*}
%where $ \mathbf{i} $ is  the imaginary unit and $ \alpha_j $ is the structure constant of the vortex core, which is generally assumed to be 1. 
%If we consider helical solutions of \eqref{NP vor fi} whose pitch is $ h $ for some constant $ h>0, $ i.e., $ X_j(s, \tau)=Z_j(\tau)e^{\mathbf{i}\frac{s}{h}} $. Then $ Z_j $ satisfies for $ j=1,\cdots,N $
%\begin{equation*} 
%\partial_\tau Z_j=\frac{1}{4\pi}\left(-\frac{\mathbf{i}\kappa_j}{h^2}Z_j+2\mathbf{i}\sum_{k\neq j}\kappa_k\frac{Z_j-Z_k}{|Z_j-Z_k|^2} \right).
%\end{equation*} 
We have  exact helical solutions of  \eqref{NP vor fi} as follows.




\begin{customthm}{A.1}  \label{lemA1}
Let $ h>0 $. The following point vortex configurations are co-rotating helical solutions of \eqref{NP vor fi}.
\begin{enumerate}
	\item[\textbf{Case 1:}]  ``$ N $ vortices'' having identical circulation and placed  at each vertex of a regular $ N $-polygon with radius $ r_* $, i.e., $ N\geq 2 $, $ \kappa_j=\kappa $ for some $ \kappa>0 $ and 
	\begin{equation*} 
	X_{j}(s,\tau)=r_*e^{-\mathbf{i}\alpha\tau}e^{\mathbf{i}\frac{s}{h}}e^{\mathbf{i}\frac{2\pi (j-1)}{N}},\ \ \ \ \ \ \ \ \  j=1,\cdots,N, 
	\end{equation*}
	where $ \alpha=\frac{\kappa}{4\pi}\left( \frac{1}{h^2}-\frac{N-1}{r_*^2}\right)  $.
	\item[\textbf{Case 2:}] ``$ N+1  $ vortices'', $ N $ of them having identical circulation and placed  at each vertex of a regular $ N $-polygon with radius $ r_* $ and the $ N +1 $ vortex of arbitrary strength
	at the center of vorticity, i.e., $ N\geq 2 $, $ \kappa_0=\mu $, $ \kappa_j=\kappa $ for $ j=1,\cdots,N $ for   $ \mu, \kappa>0 $  and
	\begin{equation*} 
	X_0(s,\tau)=0,\ \ \ \ X_{j}(s,\tau)=r_*e^{-\mathbf{i}\alpha\tau}e^{\mathbf{i}\frac{s}{h}}e^{\mathbf{i}\frac{2\pi (j-1)}{N}},\ \ \ \ \ \ \  j=1,\cdots,N,
	\end{equation*}
	where $ \alpha=\frac{\kappa}{4\pi}\left( \frac{1}{h^2}-\frac{N-1}{r_*^2}\right)-\frac{\mu}{2\pi r_*^2}  $.
		\item[\textbf{Case 3:}] ``Co-rotating asymmetric $ 2 $ vortices'' with different circulations and rotation radii. In this case, $ N= 2 $, $ \kappa_1,\kappa_2>0 $ and $ \lambda_{1,*},\lambda_{2,*}>0 $ are constants satisfying the compatibility condition
		\begin{equation*} 
		\frac{\kappa_2}{\kappa_1}=\frac{2\lambda_{1,*}h^2+\lambda_{1,*}\lambda_{2,*}(\lambda_{1,*}+\lambda_{2,*})}{2\lambda_{2,*}h^2+\lambda_{1,*}\lambda_{2,*}(\lambda_{1,*}+\lambda_{2,*})},
		\end{equation*}
		and
	\begin{equation*} 
	X_{1}(s,\tau)=\lambda_{1,*}e^{-\mathbf{i}\alpha\tau}e^{\mathbf{i}\frac{s}{h}},\ \ \ \ 	X_{2}(s,\tau)=-\lambda_{2,*}e^{-\mathbf{i}\alpha\tau}e^{\mathbf{i}\frac{s}{h}},
	\end{equation*}
	where $ \alpha=\frac{\kappa_1}{4\pi h^2}  -\frac{\kappa_2}{2\pi \lambda_{1,*}(\lambda_{1,*}+\lambda_{2,*})}=\frac{\kappa_2}{4\pi h^2}  -\frac{\kappa_1}{2\pi \lambda_{2,*}(\lambda_{1,*}+\lambda_{2,*})}  $.
\item[\textbf{Case 4:}] ``Co-rotating  $  2\times2  $ vortices'' with different circulations and rotation radii. In this case, $ N= 4 $, $ \kappa_1=\kappa_3=\kappa>0 $, $ \kappa_2=\kappa_4=\mu>0$. $ \kappa,\mu, \lambda_{1,*},\lambda_{2,*}  $ are constants satisfying the compatibility condition
\begin{equation*} 
\frac{\kappa}{4\pi h^2} -\frac{\kappa}{4\pi \lambda_{1,*}^2} -\frac{\mu}{\pi \left (\lambda_{1,*}^2+\lambda_{2,*}^2\right )}=\frac{\mu}{4\pi h^2} -\frac{\mu}{4\pi \lambda_{2,*}^2} -\frac{\kappa}{\pi \left (\lambda_{1,*}^2+\lambda_{2,*}^2\right )},
\end{equation*}
and
\begin{equation*} 
X_{2j-1}(s,\tau)=\lambda_{1,*}e^{-\mathbf{i}\alpha\tau}e^{\mathbf{i}\frac{s}{h}}e^{\mathbf{i}\pi(j-1)},\ \ \ \ 	X_{2j}(s,\tau)=\lambda_{2,*}e^{-\mathbf{i}\alpha\tau}e^{\mathbf{i}\frac{s}{h}}e^{\mathbf{i}\frac{\pi(2j-1)}{2}},\ \ j=1,2,
\end{equation*}
where $ \alpha=\frac{\kappa}{4\pi h^2} -\frac{\kappa}{4\pi \lambda_{1,*}^2} -\frac{\mu}{\pi \left (\lambda_{1,*}^2+\lambda_{2,*}^2\right )}=\frac{\mu}{4\pi h^2} -\frac{\mu}{4\pi \lambda_{2,*}^2} -\frac{\kappa}{\pi \left (\lambda_{1,*}^2+\lambda_{2,*}^2\right )}  $.
\item[\textbf{Case 5:}] ``Co-rotating  $  2\times2+1  $ vortices'' with different circulations and rotation radii and the $ 5 $-th vortex of arbitrary strength
at the center of vorticity. In this case, $ N= 5 $, $ \kappa_1=\kappa_3=\kappa>0 $, $ \kappa_2=\kappa_4=\mu>0$ and $ \kappa_0>0 $. $ \kappa_0,\kappa,\mu, \lambda_{1,*},\lambda_{2,*}  $ are constants satisfying the compatibility condition
\begin{equation*} 
\frac{\kappa}{4\pi h^2}-\frac{\kappa_0}{2\pi \lambda_{1,*}^2} -\frac{\kappa}{4\pi \lambda_{1,*}^2} -\frac{\mu}{\pi \left (\lambda_{1,*}^2+\lambda_{2,*}^2\right )}=\frac{\mu}{4\pi h^2}-\frac{\kappa_0}{2\pi \lambda_{2,*}^2} -\frac{\mu}{4\pi \lambda_{2,*}^2} -\frac{\kappa}{\pi \left (\lambda_{1,*}^2+\lambda_{2,*}^2\right )},
\end{equation*}
and $ X_0(s,\tau)=0 $,
\begin{equation*} 
X_{2j-1}(s,\tau)=\lambda_{1,*}e^{-\mathbf{i}\alpha\tau}e^{\mathbf{i}\frac{s}{h}}e^{\mathbf{i}\pi(j-1)},\ \ \ \ 	X_{2j}(s,\tau)=\lambda_{2,*}e^{-\mathbf{i}\alpha\tau}e^{\mathbf{i}\frac{s}{h}}e^{\mathbf{i}\frac{\pi(2j-1)}{2}},\ \ j=1,2,
\end{equation*}
where $ \alpha=\frac{\kappa}{4\pi h^2}-\frac{\kappa_0}{2\pi \lambda_{1,*}^2} -\frac{\kappa}{4\pi \lambda_{1,*}^2} -\frac{\mu}{\pi \left (\lambda_{1,*}^2+\lambda_{2,*}^2\right )}=\frac{\mu}{4\pi h^2}-\frac{\kappa_0}{2\pi \lambda_{2,*}^2} -\frac{\mu}{4\pi \lambda_{2,*}^2} -\frac{\kappa}{\pi \left (\lambda_{1,*}^2+\lambda_{2,*}^2\right )}  $.
\end{enumerate}
 
\end{customthm}
%\begin{remark}
%We give some explanation of the compatibility condition in cases 3,4 and 5. Take case 3 for example. The sufficient and necessary condition of the existence of solutions of \eqref{NP vor fi} in the form of \eqref{fila 3} is that the condition $ 	\frac{\kappa_2}{\kappa_1}=\frac{2\lambda_{1,*}h^2+\lambda_{1,*}\lambda_{2,*}(\lambda_{1,*}+\lambda_{2,*})}{2\lambda_{2,*}h^2+\lambda_{1,*}\lambda_{2,*}(\lambda_{1,*}+\lambda_{2,*})} $ holds. Moreover, it is possible to generalize the cases 4 and 5 to $ 2N $ and $ 2N+1 $ with different circulation and rotation radii. 
%\end{remark}
\begin{proof}
For case 1, taking \eqref{fila 1} into \eqref{NP heli vor}, we have $ \alpha=\frac{\kappa}{4\pi h^2}-\frac{\kappa}{2\pi r_*^2}\sum_{k\neq j}\frac{1-e^{\mathbf{i}\frac{2\pi(j-k)}{N}}}{|1-e^{\mathbf{i}\frac{2\pi(j-k)}{N}}|^2} $.  Note that

\begin{equation}\label{A-01}
\sum_{k\neq j}\frac{1-e^{\mathbf{i}\frac{2\pi(j-k)}{N}}}{|1-e^{\mathbf{i}\frac{2\pi(j-k)}{N}}|^2}=\frac{N-1}{2}.
\end{equation} 
Then $ \alpha=\frac{\kappa}{4\pi}\left( \frac{1}{h^2}-\frac{N-1}{r_*^2}\right). $

For case 2, taking \eqref{fila 2} into \eqref{NP heli vor}, we have $ \alpha=\frac{\kappa}{4\pi h^2}-\frac{\kappa}{2\pi r_*^2}\sum_{k\neq j}\frac{1-e^{\mathbf{i}\frac{2\pi(j-k)}{N}}}{|1-e^{\mathbf{i}\frac{2\pi(j-k)}{N}}|^2}-\frac{\mu}{2\pi r_*^2} $. Using \eqref{A-01}, we have $ \alpha=\frac{\kappa}{4\pi}\left( \frac{1}{h^2}-\frac{N-1}{r_*^2}\right)-\frac{\mu}{2\pi r_*^2}  $.

For case 3, taking \eqref{fila 3} into \eqref{NP heli vor}, we have $ \alpha=\frac{\kappa_1}{4\pi h^2}  -\frac{\kappa_2}{2\pi \lambda_{1,*}(\lambda_{1,*}+\lambda_{2,*})}=\frac{\kappa_2}{4\pi h^2}  -\frac{\kappa_1}{2\pi \lambda_{2,*}(\lambda_{1,*}+\lambda_{2,*})}  $. This implies that $ \kappa_1,\kappa_2,  \lambda_{1,*},\lambda_{2,*}  $  satisfies $ 	\frac{\kappa_2}{\kappa_1}=\frac{2\lambda_{1,*}h^2+\lambda_{1,*}\lambda_{2,*}(\lambda_{1,*}+\lambda_{2,*})}{2\lambda_{2,*}h^2+\lambda_{1,*}\lambda_{2,*}(\lambda_{1,*}+\lambda_{2,*})} $. So if this condition holds, \eqref{NP vor fi} has a solution being of the form \eqref{fila 3}.

For case 4, taking \eqref{fila 4} into \eqref{NP heli vor}, we have $ \alpha=\frac{\kappa}{4\pi h^2} -\frac{\kappa}{4\pi \lambda_{1,*}^2} -\frac{\mu}{\pi \left (\lambda_{1,*}^2+\lambda_{2,*}^2\right )}=\frac{\mu}{4\pi h^2} -\frac{\mu}{4\pi \lambda_{2,*}^2} -\frac{\kappa}{\pi \left (\lambda_{1,*}^2+\lambda_{2,*}^2\right )}  $. This implies that $ \kappa_1,\kappa_2,  \lambda_{1,*},\lambda_{2,*}  $  satisfies $\frac{\kappa}{4\pi h^2} -\frac{\kappa}{4\pi \lambda_{1,*}^2} -\frac{\mu}{\pi \left (\lambda_{1,*}^2+\lambda_{2,*}^2\right )}=\frac{\mu}{4\pi h^2} -\frac{\mu}{4\pi \lambda_{2,*}^2} -\frac{\kappa}{\pi \left (\lambda_{1,*}^2+\lambda_{2,*}^2\right )} $. Under this assumption, \eqref{NP vor fi} has a solution being of the form \eqref{fila 4}. The case 5 is similar to the case 4 and so we omit the proof here.  
\end{proof}



 














We now turn to the proof of Lemma \ref{Green expansion2}. 
\begin{customthm}{A.2}\label{lemA2}
There holds
\begin{equation*}
\bar{S}_K(x,y)=-F_{1,y}(x)-F_{2,y}(x)+\bar{H}_1(x,y) \ \ \ \ \forall\ x,y\in\Omega,
\end{equation*}
where $ F_{1,y}(\cdot) $ is defined by \eqref{exp of F_{1,y}} and $ F_{2,y}(\cdot) $ is defined by \eqref{exp of F_{2,y}},
and $ x\to \bar{H}_1(x,y)\in C^{1,\gamma}(\overline{\Omega})$ for all $ y\in \Omega $, $ \gamma\in (0, 1) $.  Moreover, the function $(x, y)\to \bar{H}_1(x,y)\in C^1(\Omega\times\Omega),$ and in particular the corresponding Robin function $x\to \bar{S}_K(x,x)\in C^1(\Omega)$.
\end{customthm}


\begin{proof}
	Let  $ y\in\Omega $ be fixed. In the following,  we always denote $ T_{ij}=(T_y)_{ij} $ the component of row $ i $, column $ j $ of the matrix $ T_y $ for $ i,j=1,2 $. From Lemma \ref{Green expansion}, the regular part $ \bar{S}_K(x,y) $ satisfies
	\begin{equation}\label{equa of S_K}
	\begin{cases}
	-\text{div}\left (K(x)\nabla \bar{S}_K(x,y)\right )=\text{div}\left (\left (K(x)-K(y)\right )\nabla \left( \sqrt{\det K(y)}^{-1}\Gamma\left (T_y(x-y)\right )\right)  \right )\ &\text{in}\ \Omega,\\
	\bar{S}_K(x,y)=-\sqrt{\det K(y)}^{-1}\Gamma\left (T_y(x-y)\right )\ &\text{on}\ \partial\Omega.
	\end{cases}
	\end{equation}
	This implies that
	\begin{equation}\label{3-1}
	\begin{split}
	-\text{div}\left (K(x)\nabla \bar{S}_K(x,y)\right )=&\sum_{i,j=1}^2\partial_{x_i}K_{ij}(x)\partial_{x_j}\left(\sqrt{\det K(y)}^{-1}\Gamma\left (T_y(x-y)\right ) \right) \\
	&+\sum_{i,j=1}^2\left( \left( K_{ij}(x)-K_{ij}(y)\right) \partial_{x_ix_j}\left(\sqrt{\det K(y)}^{-1}\Gamma\left (T_y(x-y)\right ) \right)\right)\\
	=&:A_1+A_2.
	\end{split}
	\end{equation}
	
	As for $ A_1 $, for $ x\in\mathbb{R}^2 $, we denote $ J(x)=-\frac{1}{8\pi}|x|^2\ln|x| $. Then $ \Delta \left( J(x-y)\right)=\Gamma(x-y)-\frac{1}{2\pi}.  $
	Using transformation of coordinates, one computes directly that
	\begin{equation*}
	\text{div}\left(K(y)\nabla J(T_y(x-y)) \right)= \Gamma(T_y(x-y))-\frac{1}{2\pi},
	\end{equation*}
	from which we deduce,
	\begin{equation*}
	\text{div}\left(K(y)\nabla   \partial_{x_j}\left( J(T_y(x-y))\right)  \right)= \partial_{x_j} \left( \Gamma(T_y(x-y))\right) .
	\end{equation*}
	We define for $ x\in\Omega $
	\begin{equation}\label{F_{1,y}}
	\begin{split}
	F_{1,y}(x)=\sum_{i,j=1}^2\partial_{x_i}K_{ij}(y)\cdot\sqrt{\det K(y)}^{-1} \partial_{x_j}\left( J\left (T_y(x-y)\right )\right) ,
	\end{split}
	\end{equation}
	then one has
	\begin{equation}\label{3-2}
	\text{div}\left(K(y)\nabla F_{1,y}(x)  \right)= \sum_{i,j=1}^2\partial_{x_i}K_{ij}(y)\partial_{x_j}\left(\sqrt{\det K(y)}^{-1}\Gamma\left (T_y(x-y)\right ) \right).
	\end{equation}
	
	As for $ A_2 $, using Taylor's expansion we obtain
	\begin{equation}\label{3-3}
	\begin{split}
	\sum_{i,j=1}^2\left( K_{ij}(x)-K_{ij}(y)\right) \partial_{x_ix_j}&\left(\sqrt{\det K(y)}^{-1}\Gamma\left (T_y(x-y)\right ) \right)\\
	=\sum_{\alpha,i,j=1}^2\sqrt{\det K(y)}^{-1}&\partial_{x_\alpha}K_{ij}(y) (x-y)_\alpha\cdot\partial_{x_ix_j}\left( \Gamma\left (T_y(x-y)\right ) \right)+\phi_y(x),
	\end{split}
	\end{equation}
	where $ \phi_y(\cdot)\in L^p(\Omega) $ for all $ p>1. $
	Since
	\begin{equation*}
	\partial_{x_ix_j}\Gamma(x)=-\frac{1}{2\pi}\left(\frac{\delta_{i,j}}{|x|^2}-\frac{2x_ix_j}{|x|^4} \right),
	\end{equation*}
	where $ \delta_{i,j}=1 $ for $ i=j $ and $ \delta_{i,j}=0 $ for $ i\neq j $, we have
	\begin{equation}\label{3-6}
	\begin{split}
	\partial_{x_ix_j}\left( \Gamma(T_y(x-y))\right) =-\frac{1}{2\pi}\sum_{m,n=1}^2T_{mj}T_{ni}\left(\frac{\delta_{m,n}}{|T_y(x-y)|^2}-\frac{2\left( T_y(x-y)\right)_m\left( T_y(x-y)\right)_n}{|T_y(x-y)|^4} \right).
	\end{split}
	\end{equation}
	%Here $ T_{mj}=T_{mj}(y) $. Define $ z=T_y(x-y) $.
	Taking \eqref{3-6} into \eqref{3-3}, we get
	\begin{equation}\label{3-7}
	\begin{split}
	\sum_{i,j=1}^2\left( K_{ij}(x)-K_{ij}(y)\right) \partial_{x_ix_j}&\left(\sqrt{\det K(y)}^{-1}\Gamma\left (T_y(x-y)\right ) \right)\\
	=\sum_{\alpha,\beta,i,j,m,n=1}^2\sqrt{\det K(y)}^{-1}&\partial_{x_\alpha}K_{ij}(y)T^{-1}_{\alpha\beta} \left (T_y(x-y)\right )_\beta\cdot  \\ -\frac{1}{2\pi}T_{mj}T_{ni}&\left(\frac{\delta_{m,n}}{|T_y(x-y)|^2}-\frac{2\left( T_y(x-y)\right)_m\left( T_y(x-y)\right)_n}{|T_y(x-y)|^4} \right)  +\phi_y(x).
	\end{split}
	\end{equation}
	
	
	
	
	Note that
	\begin{equation}\label{3-4}
	\frac{x^p}{|x|^4}=-\frac{1}{8}\Delta\left( \frac{x^p}{|x|^2}\right)+\frac{1}{8}\frac{\Delta x^p}{|x|^2}\ \ \ \ \text{for}\ |p|=3,
	\end{equation}
	where $p=(p_1,p_2)$ is the multi-index and $ x^p=x_1^{p_1}x_2^{p_2} $. From \eqref{3-4}, it is not hard to check that for $ 1\leq m\neq n\leq 2 $
	\begin{equation*}
	\begin{cases}
	\frac{x_m}{|x|^2}&=\Delta\left( \frac{1}{2}x_m\ln|x|\right),\\
	\frac{x_m^2x_n}{|x|^4}&=\Delta\left( -\frac{1}{8}\frac{x_m^2x_n}{|x|^2}+\frac{1}{8}x_n\ln|x|\right),\\
	\frac{x_m^3}{|x|^4}&=\Delta\left( -\frac{1}{8}\frac{x_m^3}{|x|^2}+\frac{3}{8}x_m\ln|x|\right),
	\end{cases}
	\end{equation*}
	which implies that
	\begin{equation}\label{3-5}
	\begin{cases}
	\frac{\left (T_y\left (x-y\right )\right )_m}{|T_y\left (x-y\right )|^2}&=\text{div}\left(K(y) \nabla\left( \frac{1}{2}\left (T_y\left (x-y\right )\right )_m\ln|T_y\left (x-y\right )|\right)\right) ,\\
	\frac{\left (T_y\left (x-y\right )\right )_m^2\left (T_y\left (x-y\right )\right )_n}{|T_y\left (x-y\right )|^4}&=\text{div}\left(K(y) \nabla \left( -\frac{1}{8}\frac{\left (T_y\left (x-y\right )\right )_m^2\left (T_y\left (x-y\right )\right )_n}{|T_y\left (x-y\right )|^2}+\frac{1}{8}\left (T_y\left (x-y\right )\right )_n\ln|T_y\left (x-y\right )|\right)\right) ,\\
	\frac{\left (T_y\left (x-y\right )\right )_m^3}{|T_y\left (x-y\right )|^4}&=\text{div}\left(K(y) \nabla\left( -\frac{1}{8}\frac{\left (T_y\left (x-y\right )\right )_m^3}{|T_y\left (x-y\right )|^2}+\frac{3}{8}\left (T_y\left (x-y\right )\right )_m\ln|T_y\left (x-y\right )|\right)\right).
	\end{cases}
	\end{equation}
	We define for $ x\in \Omega $
	\begin{equation*}
	\begin{split}
	F_{2,y}(x)=-&\frac{1}{2\pi}\sqrt{\det K(y)}^{-1}\sum_{\alpha,\beta,i,j,m,n=1}^2 \partial_{x_\alpha}K_{ij}(y)T^{-1}_{\alpha\beta} T_{mj}T_{ni}\cdot\frac{1}{2}\left (T_y\left (x-y\right )\right )_\beta\ln|T_y\left (x-y\right )|\delta_{m,n}\\
	+&\frac{1}{\pi}\sqrt{\det K(y)}^{-1} \sum_{\alpha,i,j=1}^2\partial_{x_\alpha}K_{ij}(y)\cdot \\
	&\bigg[T^{-1}_{\alpha1}T_{1j}T_{1i}\left( -\frac{1}{8}\frac{\left( T_y\left (x-y\right )\right) _1^3}{|T_y\left (x-y\right )|^2}+\frac{3}{8}\left( T_y\left (x-y\right )\right)_1\ln|T_y\left (x-y\right )|\right) \\
	+&T^{-1}_{\alpha1}T_{1j}T_{2i}\left( -\frac{1}{8}\frac{\left( T_y\left (x-y\right )\right) _1^2\left( T_y\left (x-y\right )\right) _2}{|T_y\left (x-y\right )|^2}+\frac{1}{8}\left( T_y\left (x-y\right )\right)_2\ln|T_y\left (x-y\right )|\right) \\
	+&T^{-1}_{\alpha1}T_{2j}T_{1i}\left( -\frac{1}{8}\frac{\left( T_y\left (x-y\right )\right) _1^2\left( T_y\left (x-y\right )\right) _2}{|T_y\left (x-y\right )|^2}+\frac{1}{8}\left( T_y\left (x-y\right )\right)_2\ln|T_y\left (x-y\right )|\right)\\
	+&T^{-1}_{\alpha1}T_{2j}T_{2i}\left( -\frac{1}{8}\frac{\left( T_y\left (x-y\right )\right) _1\left( T_y\left (x-y\right )\right) _2^2}{|T_y\left (x-y\right )|^2}+\frac{1}{8}\left( T_y\left (x-y\right )\right)_1\ln|T_y\left (x-y\right )|\right)
	\end{split}
	\end{equation*}
	\begin{equation}\label{F_{2,y}}
	\begin{split}
	+&T^{-1}_{\alpha2}T_{1j}T_{1i}\left( -\frac{1}{8}\frac{\left( T_y\left (x-y\right )\right) _1^2\left( T_y\left (x-y\right )\right) _2}{|T_y\left (x-y\right )|^2}+\frac{1}{8}\left( T_y\left (x-y\right )\right)_2\ln|T_y\left (x-y\right )|\right)\\
	+&T^{-1}_{\alpha2}T_{1j}T_{2i}\left( -\frac{1}{8}\frac{\left( T_y\left (x-y\right )\right) _1\left( T_y\left (x-y\right )\right) _2^2}{|T_y\left (x-y\right )|^2}+\frac{1}{8}\left( T_y\left (x-y\right )\right)_1\ln|T_y\left (x-y\right )|\right)\\
	+&T^{-1}_{\alpha2}T_{2j}T_{1i}\left( -\frac{1}{8}\frac{\left( T_y\left (x-y\right )\right) _1\left( T_y\left (x-y\right )\right) _2^2}{|T_y\left (x-y\right )|^2}+\frac{1}{8}\left( T_y\left (x-y\right )\right)_1\ln|T_y\left (x-y\right )|\right)\\
	+&T^{-1}_{\alpha2}T_{2j}T_{2i}\left( -\frac{1}{8}\frac{\left( T_y\left (x-y\right )\right) _2^3 }{|T_y\left (x-y\right )|^2}+\frac{3}{8}\left( T_y\left (x-y\right )\right)_2\ln|T_y\left (x-y\right )|\right)\bigg].
	\end{split}
	\end{equation}
	Combining \eqref{F_{2,y}} with \eqref{3-7} and \eqref{3-5}, we get
	\begin{equation}\label{3-8}
	\begin{split}
	\text{div}\left(K(y) \nabla F_{2,y}(x)  \right)=&\sum_{\alpha,\beta,i,j,m,n=1}^2\sqrt{\det K(y)}^{-1}\partial_{x_\alpha}K_{ij}(y)T^{-1}_{\alpha\beta} \left (T_y(x-y)\right )_\beta\cdot \\
	&-\frac{1}{2\pi}T_{mj}T_{ni}\left(\frac{\delta_{m,n}}{|T_y(x-y)|^2}-\frac{2\left( T_y(x-y)\right)_m\left( T_y(x-y)\right)_n}{|T_y(x-y)|^4} \right)\\
	=&
	\sum_{i,j=1}^2\left( K_{ij}(x)-K_{ij}(y)\right) \partial_{x_ix_j}\left(\sqrt{\det K(y)}^{-1}\Gamma\left (T_y(x-y)\right ) \right)-\phi_y(x).
	\end{split}
	\end{equation}
	
	Now we define $ \bar{H}_{1,y}(x)=\bar{S}_K(x,y)+F_{1,y}(x)+F_{2,y}(x) $. Taking \eqref{3-2} and \eqref{3-8} into \eqref{3-1}, we obtain
	\begin{equation}\label{3-9}
	\begin{split}
	-\text{div}&\left (K(x)\nabla \bar{H}_{1,y}(x)\right )\\
	=&-\text{div}\left (\left( K(x)-K(y)\right) \nabla\left( F_{1,y}(x)+F_{2,y}(x)\right) \right) \\
	&+\sum_{i,j=1}^2\left( \partial_{x_i}K_{ij}(x)-\partial_{x_i}K_{ij}(y)\right) \partial_{x_j}\left(\sqrt{\det K(y)}^{-1}\Gamma\left (T_y(x-y)\right ) \right)+\phi_y(x).
	\end{split}
	\end{equation}
	We can verify that for all $ y\in\Omega $, the right-hand side of \eqref{3-9} belongs to $ L^p(\Omega) $ for all $ p>1. $ Note also that
	\begin{equation*}
	\bar{H}_{1,y}(x)=-\sqrt{\det K(y)}^{-1}\Gamma\left (T_y(x-y)\right )+F_{1,y}(x)+F_{2,y}(x)\ \ \ \ x\in\partial\Omega.
	\end{equation*}
	For $ x,y\in\Omega $, we define $ \bar{H}_1(x,y)=\bar{H}_{1,y}(x) $. Applying the elliptic theory, we obtain that $ x\to \bar{H}_1(x,y) $ is in $ C^{1,\gamma}(\overline{\Omega}) $, for all $ \gamma\in(0,1) $. Furthermore,
	by the continuity of the right-hand side of \eqref{3-9} and the boundary condition with respect to $ y $ in $ L^p(\Omega) $   and $ C^2(\partial\Omega) $, respectively, we can get $ \bar{H}_1(x,y)=\bar{H}_{1,y}(x)\in C(\Omega, C^{1,\gamma}(\overline{\Omega})) $ and thus $ \nabla_x\bar{H}_1(x,y)\in C(\Omega\times \Omega) $.
	
	Similarly,  taking $ \nabla_y $ to both sides of \eqref{3-9}, we can check that $ \nabla_y\bar{H}_{1,y}(x)\in C(\Omega, C^{0,\gamma}(\overline{\Omega})) $, which implies that $ \nabla_y\bar{H}_1(x,y)\in C(\Omega\times \Omega) $, then $ \bar{H}_1 $ is a $ C^1 $ function over $ \Omega\times \Omega $. From \eqref{F_{1,y}} and \eqref{F_{2,y}}, we can prove that \eqref{exp of F_{1,y}} and \eqref{exp of F_{2,y}} hold. Finally, $ \bar{S}_K(x, x) =  \bar{H}_1(x,x) $ is clearly in $ C^1(\Omega) $.
	
\end{proof}







We now give the proof of Lemma \ref{H estimate}.

\begin{customthm}{A.3}\label{lemA3}
	Define $ \zeta_{\varepsilon, \hat{x}, \hat{q}}(x)=H_{\varepsilon, \hat{x}, \hat{q}}(x)+\frac{2\pi\hat{q}\sqrt{\det K(\hat{x})}\ln\frac{1}{\varepsilon}}{\ln s_\varepsilon}\bar{S}_K(x,\hat{x}) $ for $ x\in\Omega $. Then for any $ p\in(1,2) $, there exists a constant $ C>0 $ independent of $ \varepsilon $ such that
\begin{equation*}
||\zeta_{\varepsilon, \hat{x}, \hat{q}}||_{C^{0,2-\frac{2}{p}}(\Omega)}\leq Cs_\varepsilon^{\frac{2}{p}-1}.
\end{equation*}
\end{customthm}


\begin{proof}
We follow the idea of proof of theorem 1.3 in \cite{CW1}. By the construction of Green's function $ G_K $,  on $ \{x\mid |T_{\hat{x}}(x-\hat{x})|>s_\varepsilon\}$ we have
	\begin{equation}\label{3000}
	\text{div}(K(x)\nabla G_K(x,\hat{x}))=0 ,
	\end{equation}
thus by Lemma \ref{Green expansion} we deduce that
	\begin{equation}\label{3001}
	-\sqrt{\det K(\hat{x})}^{-1}\text{div}(K(x)\nabla\Gamma\left (T_{\hat{x}}(x-\hat{x})\right ))=\text{div}(K(x)\nabla\bar{S}_K(x,\hat{x})).
	\end{equation}
Taking \eqref{eq2},\eqref{eq2-2} and \eqref{3001} into \eqref{eq3}, we have on $ \{x\mid |T_{\hat{x}}(x-\hat{x})|>s_\varepsilon\} $,
	\begin{equation*}
	\begin{split}
	-\varepsilon^2\text{div}\left (K(x)\nabla H_{\varepsilon, \hat{x}, \hat{q}}\right )&=\varepsilon^2\text{div}(K(x)\nabla V_{\varepsilon, \hat{x}, \hat{q}})\\
	&=\frac{\varepsilon^2\hat{q}\ln\frac{1}{\varepsilon}}{\ln s_{\varepsilon}}\text{div}(K(x)\nabla \ln|T_{\hat{x}}(x-\hat{x})|)\\
	&=\frac{2\pi\varepsilon^2\hat{q}\sqrt{\det K(\hat{x})}\ln\frac{1}{\varepsilon}}{\ln s_{\varepsilon}}\text{div}(K(x)\nabla \bar{S}_K(x,\hat{x})),
	\end{split}
	\end{equation*}
	which implies that $ -\text{div}(K(x)\nabla \zeta_{\varepsilon, \hat{x}, \hat{q}})\equiv 0 $ on $ \{x\mid |T_{\hat{x}}(x-\hat{x})|>s_\varepsilon\} $.
	
	On the set $ \{x\mid |T_{\hat{x}}(x-\hat{x})|<s_\varepsilon\} $, we note that
	\begin{equation}\label{3-01}
	-\text{div}(K(x)\nabla \zeta_{\varepsilon, \hat{x}, \hat{q}})(x)=\text{div}((K(x)-K(\hat{x}))\nabla V_{\varepsilon, \hat{x}, \hat{q}})(x)-\frac{2\pi\hat{q}\sqrt{\det K(\hat{x})}\ln\frac{1}{\varepsilon}}{\ln s_\varepsilon}\text{div}(K(x)\nabla\bar{S}_K(x,\hat{x})).
	\end{equation}
	From \eqref{eq2-2} and Lemma \ref{Green expansion2}, we get on
	$ \{x\mid |T_{\hat{x}}(x-\hat{x})|<s_\varepsilon\} $
	\begin{equation*}
	|\text{div}((K(x)-K(\hat{x}))\nabla V_{\varepsilon, \hat{x}, \hat{q}})(x)|\leq \frac{C}{s_\varepsilon}\left( \phi'\left(\frac{|T_{\hat{x}}(x-\hat{x})|}{s_\varepsilon}\right)+\phi''\left(\frac{|T_{\hat{x}}(x-\hat{x})|}{s_\varepsilon}\right)\right) 
	\end{equation*}
and
	\begin{equation*}
	\bigg|\frac{2\pi\hat{q}\sqrt{\det K(\hat{x})}\ln\frac{1}{\varepsilon}}{\ln s_\varepsilon}\text{div}(K(x)\nabla\bar{S}_K(x,\hat{x}))\bigg|\leq C\left ( \ln\frac{1}{|T_{\hat{x}}(x-\hat{x})|}+\frac{1}{|T_{\hat{x}}(x-\hat{x})|}+1\right ).
	\end{equation*}
	Taking these into \eqref{3-01}, for any $ p\in(1,2) $ we have
	\begin{equation*}\label{3-02}
	||\text{div}(K(x)\nabla \zeta_{\varepsilon, \hat{x}, \hat{q}})||_{L^p(\{x\mid |T_{\hat{x}}(x-\hat{x})|<s_\varepsilon\})}\leq \frac{C}{s_\varepsilon}\cdot s_\varepsilon^{\frac{2}{p}}=Cs_\varepsilon^{\frac{2}{p}-1}.
	\end{equation*}
%	On $ \partial \Omega $ we have
%	\begin{equation*}
%	H_{\delta, \hat{x}, \hat{q}}(\cdot)=-V_{\delta, \hat{x}, \hat{q}}(\cdot)=-\frac{\hat{q}}{\ln s_{\delta}}\ln|T_{\hat{x}}(\cdot-\hat{x})|=-\frac{2\pi\hat{q}\sqrt{\det K(\hat{x})}}{\ln s_{\delta}}\bar{S}_K(\cdot,\hat{x}),
%	\end{equation*}
Note also that $ \zeta_{\varepsilon, \hat{x}, \hat{q}}\equiv0 $ on $ \partial \Omega. $  By the $ L^p $ theory for second-order elliptic equations, for any $ p\in(1,2) $
	\begin{equation*}
	||\zeta_{\varepsilon, \hat{x}, \hat{q}}||_{W^{2,p}(\Omega)}\leq C||\text{div}(K(x)\nabla \zeta_{\varepsilon, \hat{x}, \hat{q}})||_{L^p(\{x\mid |T_{\hat{x}}(x-\hat{x})|<s_\varepsilon\})}\leq  C s_\varepsilon^{\frac{2}{p}-1}.
	\end{equation*}
Using the Sobolev embedding theorem, the proof is complete.
\end{proof}


 

\iffalse

The following lemma shows the existence of $ \hat{q}_i $ satisfying \eqref{q_i choice}.
\begin{customthm}{A.2}\label{lemA-3}
For any $ Z\in \mathcal{M} $ and $ \delta $ sufficiently small,  there exist $ \hat{q}_i=\hat{q}_{\delta,i}(Z) $ for $ i=1,\cdots,m $ satisfying
\begin{equation}\label{A301}
\hat{q}_i=q(z_i)+\frac{2\pi\hat{q}_i\sqrt{\det K(z_i)}}{\ln s_{\delta,i}}\bar{S}_K(z_i, z_i)+\sum_{j\neq i}\frac{2\pi\hat{q}_j\sqrt{\det K(z_j)}}{\ln s_{\delta,j}}G_K(z_i, z_j),
\end{equation}
where $ s_{\delta,i} $ satisfies
\begin{equation*}
\delta^{\frac{2}{p-1}}s_{\delta,i}^{-\frac{2}{p-1}}\phi'(1)=\hat{q}_i/\ln s_{\delta,i}.
\end{equation*}


\end{customthm}
\begin{proof}
By the Poincar$\acute{\text{e}}$--Miranda Theorem (see \cite{Ku}), it is not hard to prove that there exists the unique $ \hat{q}_i $ satisfying \eqref{A301}. Moreover, we have $ \hat{q}_i=q(z_i)+O(\frac{1}{|\ln\varepsilon|}) $.
\end{proof}

This lemma shows the well-definedness of $ Q_\delta $ defined by \eqref{206}.
\begin{customthm}{A.3}\label{lemA-4}
$ Q_\delta $ is well-defined. Moreover, for any $ q\in[1,+\infty) $, $ u\in L^q(\Omega) $ with $ supp(u)\subset  \cup_{j=1}^mB_{Ls_{\delta,j}}(z_j)$ for some $ L>1 $, there exists a constant $ C>0 $ independent of $ \delta $ satisfying
\begin{equation*}
||Q_\delta u||_{L^q(\Omega)}\leq C||u||_{L^q(\Omega)}.
\end{equation*}
	
\end{customthm}
\begin{proof}
By \eqref{coef of C}, we already know that there exists a unique $ C_{j,h} $ satisfying \eqref{207}. By the assumption, for any $ q\in[1,+\infty) $, $ l=1,\cdots,m, h=1,2, $ one has
\begin{equation*}
\begin{split}
C_{l,h}=&O\left( |\ln\varepsilon|^{p+1}\sum_{i=1}^m\sum_{\hbar=1}^2\int_{\cup_{j=1}^mB_{Ls_{\delta,j}}(z_j)}u  \left( \eta_i\frac{\partial V_{\delta, z_i, \hat{q}_{\delta,i}}}{\partial x_\hbar}\right)   \right) \\
=&O\left( |\ln\varepsilon|^{p+1}\sum_{i=1}^m\sum_{\hbar=1}^2||u||_{L^q(B_{Ls_{\delta,i}}(z_i))}||\frac{\partial V_{\delta, z_i, \hat{q}_{\delta,i}}}{\partial x_\hbar}||_{L^{q'}(B_{Ls_{\delta,i}}(z_i))}\right) \\
=&O\left( |\ln\varepsilon|^{p+1}||u||_{L^q(\Omega)}\sum_{i=1}^m\sum_{\hbar=1}^2\frac{s_{\delta,i}^{\frac{2}{q'}-1}}{|\ln\varepsilon|}\right) \\
=&O\left (|\ln\varepsilon|^{p}\varepsilon^{1-\frac{2}{q}}\right )||u||_{L^q(\Omega)},
\end{split}
\end{equation*}
where $ q' $ denotes the conjugate exponent of $ q $. Hence we get
\begin{equation*}
\begin{split}
\sum_{j=1}^m&\sum_{h=1}^2C_{j,h}\left( -\delta^2\text{div}\left (K(z_j)\nabla   \frac{\partial V_{\delta, z_j, \hat{q}_{\delta,j}}}{\partial x_h}\right )\right)\\
=&p\sum_{j=1}^m\sum_{h=1}^2C_{j,h}(V_{\delta, z_{j}, \hat{q}_{\delta,j}}-\hat{q}_{\delta,j})^{p-1}_+ \frac{\partial V_{\delta, z_{j}, \hat{q}_{\delta,j}}}{\partial x_h} \\
=&O\left( \sum_{j=1}^m\sum_{h=1}^2C_{j,h}|\ln\varepsilon|^{-(p-1)}\cdot \frac{s_{\delta,j}^{\frac{2}{q}-1}}{|\ln\varepsilon|}\right) \\
=&O(1)||u||_{L^q(\Omega)}\ \ \ \ \text{in}\ L^q(\Omega).
\end{split}
\end{equation*}
By the definition of $ Q_\delta $, we have $ ||Q_\delta u||_{L^q(\Omega)}\leq C||u||_{L^q(\Omega)} $.

\end{proof}

\fi







\subsection*{Acknowledgements:}

\par
D. Cao was partially supported by National Key R \& D
Program (2022YFA1005602). J. Wan was supported by NNSF of China (grant No. 12271539 and 12471190).


\subsection*{Conflict of interest statement} On behalf of all authors, the corresponding author states that there is no conflict of interest.

\subsection*{Data availability statement} All data generated or analysed during this study are included in this published article  and its supplementary information files.


\begin{thebibliography}{99}

%\bibitem{AS}
%H. Abidi and S. Sakrani, Global well-posedness of helicoidal Euler equations, \textit{J. Funct. Anal.}, 271(8)(2016), 2177--2214.
%
%\bibitem{ALW}
%W. Ao, Y. Liu and J. Wei,  Clustered travelling vortex rings to the axisymmetric three-dimensional incompressible Euler flows, \textit{Phys. D}, 434(2022), Paper No. 133258.
%
%\bibitem{Ben}
%M. Benvenutti, Nonlinear stability for stationary helical vortices, \textit{NoDEA Nonlinear Differential Equations Appl.}, 27(2020), no. 2, Paper No. 15.
%
%\bibitem{BLN}
%A.C. Bronzi, M.C. Lopes Filho and H.J. Nussenzveig Lopes,  Global existence of a weak solution of the incompressible Euler equations with helical symmetry and $ L^p $ vorticity, \textit{Indiana Univ. Math. J.}, 64(1)(2015), 309--341.
\bibitem{BFM}
V. Banica, E. Faou and E. Miot, Collision of almost parallel vortex filaments, \textit{Comm. Pure Appl. Math.}, 70(2017), no. 2, 378--405.
\bibitem{BM}
V. Banica  and E. Miot, Global existence and collisions for symmetric configurations of nearly parallel vortex filaments, \textit{Ann. Inst. H. Poincar$\acute{\text{e}}$ C Anal. Non Lin$\acute{\text{e}}$aire}, 29(2012), no. 5, 813--832.


%\bibitem{Bur}
%G.R. Burton, Rearrangements of functions, maximization of convex functionals, and
%vortex rings, Math. Ann. 276(2)(1987) 225--253.


%\bibitem{CF}
%L.A. Caffarelli and A. Friedman, Asymptotic estimates for the plasma problem, Duke Math. J. 47 (1980) 705--742.

%\bibitem{CGPY}
%D. Cao, Y. Guo, S. Peng and S. Yan, Local uniqueness for vortex patch problem in incompressible planar steady flow, \textit{J. Math. Pures Appl.}, 131(2019),
%251--289.

\bibitem{CFQW1}
D. Cao, J. Fan, G. Qin and J. Wan, Dynamics and leapfrogging phenomena of multiple helical vortices for  3D incompressible Euler equations, arXiv:2505.12240.

\bibitem{CFQW2}
D. Cao, J. Fan, G. Qin and J. Wan, Dynamics of helical vortex filaments with singular interaction for the 3D Euler equations, preprint.

\bibitem{CLW}
D. Cao, Z. Liu and J. Wei, Regularization of point vortices  for the
Euler equation in dimension two, \textit{Arch. Ration. Mech. Anal.},
212(2014), 179--217.

%\bibitem{CPY1}
%D. Cao, S. Peng and S. Yan, Multiplicity of solutions for the plasma problem
%in two dimensions,  \textit{Adv. Math.}, 225(2010), 2741--2785.
%
\bibitem{CPY}
D. Cao, S. Peng and S. Yan, Planar vortex patch problem in incompressible steady flow,  \textit{Adv. Math.}, 270(2015), 263--301.

\bibitem{CW}
D. Cao and J. Wan, Helical vortices with small cross-section for 3D incompressible Euler equation, \textit{J. Funct. Anal.}, 284(2023), Paper No. 109836.

\bibitem{CW1}
D. Cao and J. Wan, Structure of Green's function of elliptic equations and helical vortex patches for 3D incompressible Euler equations, \textit{Math. Ann.},
 388(2024),  2627--2669.

%\bibitem{CW2}
%D. Cao and J. Wan, Desingularization of 3D steady Euler equations with helical symmetry, \textit{Calc. Var. \& Partial Differential Equations}, to appear, arXiv: 2206.00196.


\bibitem{CW3}
D. Cao and J. Wan, Expansion of Green’s function and regularity of Robin’s function for elliptic operators in divergence form,   \textit{Ann. Sc. Norm. Super. Pisa
Cl. Sci.}, (2024), DOI: https://doi.org/10.2422/2036-2145.202402$ \_ $006.

\bibitem{Cho}
D. Chiron, Vortex helices for the Gross-Pitaevskii equation, \textit{J. Math. Pures Appl.}, 84(9)(2005), no. 11, 1555--1647.

\bibitem{CJ1}
A. Contreras and R.L. Jerrard, Nearly parallel vortex filaments in the 3D Ginzburg-Landau equations, \textit{Geom. Funct. Anal.}, 27(2017), no. 5, 1161--1230.

\bibitem{CGY}
W. Craig, C. Garc$\acute{\text{i}}$a-Azpeitia  and C. Yang, Standing waves in near parallel vortex filaments, \textit{Comm. Math. Phys.}, 350(2017), no. 1, 175--203.


\bibitem{DY}
E.N. Dancer and S. Yan, The Lazer-McKenna conjecture and a free boundary problem
in two dimensions, \textit{J. Lond. Math. Soc.}, 78(2008), 639--662.

\bibitem{DR}
L.S. Da Rios, Sul moto dun liquido indefinito con un filetto vorticoso di forma qualunque, Rendiconti del Circolo Matematico di Palermo (1884-1940), 22(1)(1906), 117--135.

%\bibitem{DR2}
%L.S. Da Rios, Sul moto dei filetti vorticosi di forma qualunque, Rend. R. Acc. Lincei 18(1909)75--79.

	\bibitem{DDMR}
J. D\'avila, M. del~Pino, M. Medina  and R. Rodiac,  Interacting helical
vortex filaments in the 3-dimensional Ginzburg-Landau equation, \textit{J. Eur. Math. Soc. (JEMS)},  24(2022), no. 12, 4143--4199.  %https://doi.org/10.4171/JEMS/1175

\bibitem{DDMR2}
J. D\'avila, M. del~Pino, M. Medina  and R. Rodiac,   Interacting helical traveling waves for the Gross-Pitaevskii equation,
\textit{Ann. Inst. H. Poincar\'e C Anal. Non Lin\'eaire},  39(2022), no. 6, 1319--1367. %https://doi.org/10.4171/AIHPC/32

\bibitem{DDMW}
J. D$\acute{\text{a}}$vila, M. del Pino, M. Musso and J. Wei, Gluing methods for vortex dynamics in Euler flows, \textit{Arch. Ration. Mech. Anal.}, 235(3)(2020), 1467--1530.

\bibitem{DDMW2}
J. D$\acute{\text{a}}$vila, M. del Pino, M. Musso and J. Wei, Travelling helices and the vortex filament conjecture
in the incompressible Euler equations, \textit{Calc. Var. Partial Differential Equations}, 61(2022), art.119.



\bibitem{DLM}
M. Donati, C. Lacave  and E. Miot, Dynamics of helical vortex filaments in non viscous incompressible flows, arXiv: 2403.00389.
%\bibitem{De2}
%J. Dekeyser and J. Van Schaftingen,  Vortex motion for the lake equations, \textit{Comm. Math. Phys.}, 375(2020), 1459--1501.
%\bibitem{De2}
%J. Dekeyser, Asymptotic of steady vortex pair in the lake equation, SIAM J. Math. Anal. 51(2)(2019)1209-1237.

%\bibitem{DV}
%S. de Valeriola and J. Van Schaftingen, Desingularization of vortex rings and shallow water vortices by semilinear elliptic problem, \textit{Arch. Ration. Mech. Anal.}, 210(2)(2013), 409--450.




%\bibitem{Du}
%A. Dutrifoy, Existence globale en temps de solutions h$\acute{\text{e}}$lico$\ddot{\text{i}}$dales des $\acute{\text{e}}$quations d'Euler, \textit{C. R. Acad. Sci. Paris S$\acute{\text{e}}$r. I Math.}, 329(7)(1999), 653--656.



\bibitem{ET}
B. Ettinger and E.S. Titi, Global existence and uniqueness of weak solutions of three-dimensional Euler equations with helical symmetry in the absence of vorticity stretching, \textit{SIAM J. Math. Anal.}, 41(1)(2009), 269--296.


\bibitem{Fr1}
L.E. Fraenkel, On steady vortex rings of small cross-section in an ideal fluid, \textit{Proc. R. Soc. Lond. A.}, 316(1970), 29--62.

\bibitem{FB}
L.E. Fraenkel and M.S. Berger, A global theory of steady vortex rings in an ideal fluid, \textit{Acta Math.}, 132(1974), 13--51.



\bibitem{GT}
D. Gilbarg and N.S. Trudinger,  Elliptic Partial Differential Equations of Second Order, Classics in Mathematics, Springer, Berlin, 2001.

%\bibitem{G}
%W.H. Greub, Linear Algebra, Third edition, Die Grundlehren der Mathematischen Wissenschaften, Band 97 Springer-Verlag New York, Inc., New York 1967, xiii+434 pp.

\bibitem{GM}
I. Guerra and M. Musso, Cluster of vortex helices in the incompressible 3d Euler equations, \textit{Ann. Inst. H. Poincar$\acute{\text{e}}$ C. Anal. Non Lin$\acute{\text{e}}$aire}, 42 (2025), 547--592.

\bibitem{GM2}
I. Guerra and M. Musso, Nearly parallel helical vortex filaments in the three dimensional Euler equations, arXiv: 2502.01470.


\bibitem{GZ}
D. Guo  and L. Zhao, Long time dynamics for helical vortex filament in Euler flows, arXiv: 2403.09071v1.


\bibitem{He}
H. Helmholtz, On integrals of the hydrodynamics equations which express vortex motion, \textit{J. Reine Angew. Math.}, 55(1858), 25--55.


\bibitem{JS}
R.L. Jerrard and C. Seis, On the vortex filament conjecture for Euler flows, \textit{Arch. Ration. Mech. Anal.}, 224(1)(2017), 135--172.

\bibitem{JS2}
R.L. Jerrard and D. Smets, Dynamics of nearly parallel vortex filaments for the Gross-Pitaevskii equation, \textit{Calc. Var. Partial Differential Equations}, 60(2021), art. 127.
%\bibitem{JS2}
%R.L. Jerrard and D. Smets, On the motion of a curve by its binormal curvature, \textit{J. Eur. 	Math. Soc.}, (JEMS) 17(6)(2015), 1487--1515.

%\bibitem{Jiu}
%Q. Jiu, J. Li and D. Niu,  Global existence of weak solutions to the three-dimensional Euler equations with helical symmetry, \textit{J. Differential Equations}, 262(10)(2017), 5179--5205.

\bibitem{KPV}
C. Kenig, G. Ponce and L. Vega, On the interaction of nearly parallel vortex filaments, \textit{Comm. Math. Phys.}, 243(2003), no. 3, 471--483.

\bibitem{KMD}
R. Klein, A. Majda and K. Damodaran, Simplified equations for the interaction of nearly parallel
vortex filaments, \textit{J. Fluid Mech.}, 288(1995), 201--248.


\bibitem{Ku}
W. Kulpa, The Poincar$\acute{\text{e}}$--Miranda theorem, \textit{Amer. Math. Mon.},  104(1997), 545--550.

\bibitem{KVG}
A. Kwiecinski  and  A. Van Gorder, Dynamics of nearly parallel interacting
vortex filaments, \textit{J. Fluid Mech.}, 835(2018),  575--623.
%\bibitem{Lamb}
%H. Lamb, Hydrodynamics Cambridge Mathematical Library, 6th edition. Cambridge University Press, Cambridge, (1932).

\bibitem{LC}
T. Levi-Civita, Sullattrazione esercitata da una linea materiale in punti prossimi alla linea stessa, \textit{Rend. R. Acc. Lincei}, 17(1908), 3--15.

%\bibitem{LC3}
%T. Levi-Civita,  Sulla gravitazione di un tubo sottile con applicazione all'anello di Satumo, Rend. Circ. Mat. Palermo, 33 (1912) 354--374.

%\bibitem{LC2}
%T. Levi-Civita, Attrazione newtoniana dei tubi sottili e vortici filiformi, Annali della Scuola Normale Superiore di Pisa - Classe di Scienze Ser. 2, 1(3) (1932) 229--250.

%\bibitem{LYY}
%G. Li, S. Yan and J. Yang, An elliptic problem related to planar vortex pairs, SIAM J. Math. Anal. 36(2005) 1444--1460.

%\bibitem{Li}
%C.C. Lin, On the motion of vortices in two dimension - I. Existence of the Kirchhoff-Routh function, Proc. Natl. Acad. Sci. USA 27(1941) 570--575.
\bibitem{LM}
P.-L. Lions  and A. Majda, Equilibrium statistical theory for nearly parallel vortex filaments, \textit{Comm. Pure Appl. Math.}, 53(2000), no. 1, 76--142.


\bibitem{MB}
A. Majda and A. Bertozzi, Vorticity and incompressible flow, Cambridge University Press, Cambridge, 2002.

%\bibitem{MR}
%Y. Martel and P. Rapha$\ddot{\text{e}}$l, Strongly interacting blow up bubbles for the mass critical NLS, arXiv: 1512.00900.



%\bibitem{MP}
%C. Marchioro and M. Pulvirenti, Mathematical Theory of Incompressible Nonviscous Fluids, Springer-Verlag, 1994.

%\bibitem{Ric}
%R.L. Ricca, Rediscovery of Da Rios equations, Nature 352 (1991) 561--562.


%\bibitem{Ric2}
%R.L. Ricca, The contributions of Da Rios and Levi-Civita to asymptotic potential theory and vortex filament dynamics, Fluid Dyn. Res. 18(5)(1996) 245--268.



%\bibitem{SV}
%D. Smets and J. Van Schaftingen, Desingularization of vortices for the Euler equation, \textit{Arch. Ration. Mech. Anal.}, 198(3)(2010), 869--925.

\bibitem{WY}
J. Wei and J. Yang, Traveling vortex helices for Schr\"odinger map equations, \textit{Trans. Amer. Math. Soc.}, 368(2016), no. 4, 2589--2622.

\bibitem{WYZ}
J. Wei, D. Ye and F. Zhou, Bubbling solutions for an anisotropic Emden–Fowler equation, \textit{Calc. Var. Partial Differential Equations}, 28(2007), 217--247.

\bibitem{W}
M. Willem, Minimax Theorems, Progr. Nonlinear Differential Equations Appl., 24
Birkhäuser Boston, Inc., Boston, MA, 1996. 
 


%\bibitem{Y}
%V.I. Yudovich, Non-stationary flow of an ideal incompressible fluid,
%\textit{USSR Comp. Math. Math. Phys.}, 3(1963), 1407--1456.



\end{thebibliography}

\end{document}
